\chapter{多元微分学}
在第 I 卷中，我们曾用微分学来深入洞察函数的“精细结构”。在这一过程中，线性近似的思想被证明极为有效。然而，我们至今一直集中于一元函数。

本章主要致力于把微分学推广到多元函数，但我们也将进一步探讨线性近似这一简单的思想。事实上——与一维情形相反——这里的问题本质上更加复杂，因为多维情形下的线性映射所展现的结构比一维情形下的远为丰富。

和以前一样，我们更愿意采用坐标无关的表示。换言之，我们对巴拿赫空间之间的映射建立微分学。这种表示在概念上很简单，并且实际上使许多表达式看起来简单得多。通常坐标表示下的经典导数公式，利用有限维欧几里得空间的积结构，可以很容易地从一般结果推出。

由于巴拿赫空间之间的线性映射是线性近似、因而也是可微性的基础，第一节致力于研究线性算子空间。特别的关注自然落在有限维情形上，我们利用线性代数中关于矩阵的基本法则来展开这一情形。

作为这些结果的一个应用，我们将研究巴拿赫空间的自同态代数中的指数函数，并由此导出关于常微分方程组以及二阶常系数微分方程的基本事实。

第二节建立了另一个核心概念——弗雷歇导数。此外，我们还考虑方向导数，它自然地来自偏导数以及借助雅可比矩阵对 弗雷歇导数的表示。最后，我们考察复变函数的可微性与相应实表示的全可微性之间的联系。我们用柯西–黎曼方程刻画复可微性。

在第三节中，我们汇总求导法则，并通过一元情形的简单推广导出重要的中值定理。

在转向高阶导数之前，我们必须澄清多重线性映射与取值于线性映射的线性算子之间的联系。第四节中得到的简单结果将为多元情形下高阶导数的清晰表示奠定基础。

在第五节中，我们解释高阶导数的概念。特别地，我们证明基本的泰勒公式——既适用于巴拿赫空间之间的映射，也适用于有限多个变量的函数。将一维情形的判据加以推广，我们将给出多元函数局部极值存在的充分条件。

第六节起着特殊的作用。在这里，我们将看到在巴拿赫空间上建立微分学的第一份回报。利用它，我们将解释变分法并导出重要的欧拉–拉格朗日微分方程。在这一过程中，我们希望你能体会到抽象方法的威力：在抽象方法中，函数被看作无穷维巴拿赫空间中的点。“函数的函数”这一简单的几何思想被证明具有深远的后果。在变分法中，我们寻求使某个函数的函数极小化，其形式通常由某个物理原理（如“最小作用量”）确定。若极小值存在，它们出现在临界点处，而这一判据引出了欧拉–拉格朗日微分方程（组），其对数学和物理的重要性怎么强调都不为过。

在这次对“高等分析”的涉猎之后，我们在第七节证明微分学中也许最重要的定理，即反函数定理。与该结果等价的是隐函数定理，我们将在随后的各节中导出它。无须额外费力，我们也能对巴拿赫空间之间的映射证明这一定理。而且坐标无关的表示再次在清晰性与优雅性上带来显著的收益。

在第八节中，我们对非线性常微分方程的理论作一瞥。利用隐函数定理，我们讨论一阶标量微分方程。我们还证明皮卡–林德勒夫定理，它是常微分方程基本的存在唯一性定理。

在本章的剩余部分，我们将阐明隐函数定理在有限维情形中的重要性。借助它，我们将刻画 $ \mathbb{R}^n $ 的子流形；它们是 $ \mathbb{R}^n $ 的子集，对 $ m < n $ 在局部上“看起来像” $ \mathbb{R}^m $。通过曲线、曲面以及高维子流形的众多例子，我们将看到，这种相似性可以通过它们的切空间得到更好的可视化，并得到更精确的描述。这里，我们将集中于欧几里得向量空间的子流形，因为它们在概念上比抽象流形（后者不作为 $ \mathbb{R}^n $ 的子集来定义）更简单。然而，我们将为一般流形上的分析奠定基础，这方面的内容将在第 III 卷中处理。作为这些几何思想的一个实际应用，我们将研究约束极小化问题及其通过拉格朗日乘子法的求解。在导出相关法则之后，我们将处理两个非平凡的例。

\section{连续线性映射}

正如我们刚才所说，多元情形的微分依赖于用仿射函数作局部近似这一思想。当然，由命题 I.12.8，向量空间之间的仿射函数中非平凡的部分正是其线性映射。因此，下面我们集中研究赋范向量空间之间的线性映射。在导出它们若干基本性质之后，我们将用它们来研究指数映射以及常系数线性微分方程理论。

以下，设

• $ E = (E, \|\cdot\|) $ 与 $ F = (F, \|\cdot\|) $ 是域 $ \mathbb{K} $ 上的赋范向量空间\footnote{必要时，我们用相应的指标区分 $ E $ 与 $ F $ 中的范数。}。

\needspace{3\baselineskip}\subsection*{$ \mathcal{L}(E, F) $ 的完备性}

由第 VI.2 节，我们已经知道 $ \mathcal{L}(E, F) $，即从 $ E $ 到 $ F $ 的所有有界线性映射的空间，是一个赋范向量空间。现在我们来研究这个空间的完备性。

\begin{theorem}

若 $ F $ 是巴拿赫空间，则 $ \mathcal{L}(E, F) $ 也是巴拿赫空间。

证明

\begin{subitems}
\item

设 $ (A_{n}) $ 是 $ \mathcal{L}(E, F) $ 中的柯西序列。则对每个 $ \varepsilon > 0 $，存在 $ N(\varepsilon) \in \mathbb{N} $，使得当 $ m, n \ge N(\varepsilon) $ 时 $ \|A_{n} - A_{m}\| < \varepsilon $。特别地，对每个 $ x \in E $，有

\begin{equation}
\|A_{n}x - A_{m}x\| \le \|A_{n} - A_{m}\|\,\|x\| \le \varepsilon\,\|x\| \qquad \text{对 } m, n \ge N(\varepsilon) .\tag{1.1}
\end{equation}

因此 $ (A_{n}x) $ 是 $ F $ 中的柯西序列。由于 $ F $ 完备，存在唯一的 $ y \in F $，使得 $ \lim A_{n}x = y $。于是我们可以通过 $ x \mapsto \lim A_{n}x $ 定义一个映射 $ A : E \to F $。由极限的线性性立即可得 $ A $ 是线性的。

\item

由于 $ (A_{n}) $ 是柯西序列，故它在 $ \mathcal{L}(E, F) $ 中有界（参见命题 II.6.3）。因此存在 $ \alpha \ge 0 $，使得对一切 $ n \in \mathbb{N} $ 有 $ \|A_{n}\| \le \alpha $。由此可得

\[
\|A_{n}x\| \le \|A_{n}\|\,\|x\| \le \alpha\,\|x\| \qquad \text{对 } n \in \mathbb{N} \text{ 且 } x \in E .
\]

舍去中间项，令 $ n \to \infty $ 即得对每个 $ x \in E $ 的估计 $ \|Ax\| \le \alpha\,\|x\| $。因此 $ A $ 属于 $ \mathcal{L}(E, F) $。

\item

最后我们证明 $ \mathcal{L}(E, F) $ 中的序列 $ (A_{n}) $ 收敛到 $ A $。由 (1.1)，可得

\[
\|A_{n}x - A_{m}x\| \le \varepsilon \qquad \text{对 } m, n \ge N(\varepsilon) \text{ 且 } x \in \mathbb{B}_{E} .
\]

令 $ m \to \infty $，这给出

\[
\|A_{n}x - Ax\| \le \varepsilon \qquad \text{对 } n \ge N(\varepsilon) \text{ 且 } x \in \mathbb{B}_{E} ,
\]

并且，正如所断言的那样，我们得到不等式

\[
\|A_{n} - A\| = \sup_{\|x\| \le 1}\|A_{n}x - Ax\| \le \varepsilon \qquad \text{对 } n \ge N(\varepsilon)
\]

其中上确界在 $ \mathbb{B}_{E} $ 上取。\nobreak\hspace{0.35em}■

\end{subitems}

\end{theorem}

\begin{corollary}

\begin{subitems}
\item

$ \mathcal{L}(E, \mathbb{K}) $ 是巴拿赫空间。

\item

若 $ E $ 是巴拿赫空间，则 $ \mathcal{L}(E) $ 是带单位元的巴拿赫代数。

\proofhead
(i) 显然，(ii) 由注 VI.2.4(g) 可得。\nobreak\hspace{0.35em}■
\end{subitems}

\end{corollary}

\needspace{3\baselineskip}\subsection*{有限维巴拿赫空间}

赋范向量空间 $ E $ 与 $ F $ 称为（\textbf{拓扑}）\textbf{同构}的，若存在从 $ E $ 到 $ F $ 的连续线性同构 $ A $，使得 $ A^{-1} $ 也连续，即若 $ A $ 属于 $ \mathcal{L}(F, E) $。此时 $ A $ 是从 $ E $ 到 $ F $ 的\textbf{拓扑同构}。我们用

\[
\mathcal{L}\mathrm{is}(E, F)
\]

表示从 $ E $ 到 $ F $ 的所有拓扑同构的集合，且当 $ \mathcal{L}\mathrm{is}(E, F) $ 非空\footnote{注意，在赋范向量空间中 $ \cong $ 总是指“拓扑同构”，而不仅仅是指存在一个向量空间同构。}时，我们写 $ E \cong F $。此外，我们令

\[
\mathcal{L}\mathrm{aut}(E) := \mathcal{L}\mathrm{is}(E, E) .
\]

因此 $ \mathcal{L}\mathrm{aut}(E) $ 是 $ E $ 的所有\textbf{拓扑自同构}的集合。

\begin{remark}

\begin{items}
\item

空间 $ \mathcal{L}(\mathbb{K}, F) $ 与 $ F $ 是等距同构的\footnote{参见例 III.1.3(o)。}。更确切地说，利用等距同构

\begin{equation}
\mathcal{L}(\mathbb{K}, F) \to F , \quad A \mapsto A1 ,\tag{1.2}
\end{equation}

我们可以规范地把 $ \mathcal{L}(\mathbb{K}, F) $ 与 $ F $ 等同为 $ \mathcal{L}(\mathbb{K}, F) = F $。

\proofhead
显然映射 (1.2) 是线性的且为单射。对 $ v \in F $，考虑 $ A_{v} \in \mathcal{L}(\mathbb{K}, F) $，其中 $ A_{v}x := xv $。则 $ A_{v}1 = v $。因此 $ A \mapsto A1 $ 是向量空间同构。此外，
\[
\|Ax\|_{F} = |x|\,\|A1\|_{F} \le \|A1\|_{F} \qquad \text{对 } x \in \mathbb{B}_{\mathbb{K}} \text{ 且 } A \in \mathcal{L}(\mathbb{K}, F) .
\]

这给出 $ \|A\|_{\mathcal{L}(\mathbb{K}, F)} = \|A1\|_{F} $。因此 (1.2) 是等距。\nobreak\hspace{0.35em}■

\item

$ \mathcal{L}\mathrm{aut}(E) $ 是一个群，即 $ E $ 的\textbf{拓扑自同构群}，其中的乘法由两个线性映射的复合来定义。

\proofhead
这由注 VI.2.4(g) 可得。\nobreak\hspace{0.35em}■
\item

若 $ E $ 与 $ F $ 同构，则 $ E $ 是巴拿赫空间当且仅当 $ F $ 是巴拿赫空间。

\proofhead
容易看出，拓扑同构把柯西序列映为柯西序列，把收敛序列映为收敛序列。\nobreak\hspace{0.35em}■
\item

设 $ E $ 与 $ F $ 是巴拿赫空间，且 $ A \in \mathcal{L}(E, F) $ 是双射。则 $ A $ 是拓扑同构，即 $ A \in \mathcal{L}\mathrm{is}(E, F) $。

\proofhead
这个定理是泛函分析的标准内容（参见巴拿赫同态定理，也称为开映射定理）。\nobreak\hspace{0.35em}■
\end{items}

\end{remark}

\begin{theorem}

设 $ \{b_{1}, \ldots, b_{n}\} $ 是 $ E $ 的一个基。则

\begin{equation}
T : E \to \mathbb{K}^{n} , \quad e = \sum_{j=1}^{n} x_{j}b_{j} \to (x_{1}, \ldots, x_{n})\tag{1.3}
\end{equation}

是拓扑同构，也就是说，每个有限维赋范向量空间都拓扑同构于一个欧几里得向量空间。

\proofhead
显然 $ T $ 良定义、线性且为双射，并且
\[
T^{-1}x = \sum_{j=1}^{n} x^{j}b_{j} \quad \text{对 } x = (x^{1}, \ldots, x^{n}) \in \mathbb{K}^{n} ,
\]

（参见注 I.12.5）。由推论 II.3.9 的柯西-施瓦茨不等式，有

\[
\|T^{-1}x\| \le \sum_{j=1}^{n} |x^{j}|\,\|b_{j}\| \le \left( \sum_{j=1}^{n} |x^{j}|^{2} \right)^{1/2}\left( \sum_{j=1}^{n} \|b_{j}\|^{2} \right)^{1/2} = \beta\,|x|
\]

其中 $ \beta := \left( \sum_{j=1}^{n}\|b_{j}\|^{2} \right)^{1/2} $。因此 $ T^{-1} $ 属于 $ \mathcal{L}(\mathbb{K}^{n}, E) $。

对 $ x \in \mathbb{K}^{n} $，我们令 $ |x|_{\bullet} := \|T^{-1}x\| $。不难看出 $ |\cdot|_{\bullet} $ 是 $ \mathbb{K}^{n} $ 上的一个范数（参见习题 II.3.1）。在例 III.3.9(a) 中我们已经看到，$ \mathbb{K}^{n} $ 上的所有范数都等价。因此存在 $ \alpha > 0 $，使得对 $ x \in \mathbb{K}^{n} $ 有 $ |x| \le \alpha|x|_{\bullet} $。于是我们有

\[
|Te| \le \alpha|Te|_{\bullet} = \alpha\|e\| \qquad \text{对 } e \in E ,
\]

即 $ T \in \mathcal{L}(E, \mathbb{K}^{n}) $，证毕。\nobreak\hspace{0.35em}■

\end{theorem}

\begin{corollary}

若 $ E $ 是有限维的，则下列陈述成立：

\begin{subitems}
\item

$ E $ 上的所有范数都等价。

\item

$ E $ 完备，因而是巴拿赫空间。

\proofhead
令 $ A := T^{-1} $，其中 $ T $ 是 (1.3) 中的映射。
\end{subitems}

(i) 对 $ j = 1, 2 $，设 $ \|\cdot\|_{j} $ 是 $ E $ 上的范数。则 $ x \mapsto |x|_{(j)} := \|Ax\|_{j} $ 是 $ \mathbb{K}^{n} $ 上的范数。因此，由例 III.3.9(a)，存在 $ \alpha \ge 1 $ 使得

\[
\alpha^{-1}|x|_{(1)} \le |x|_{(2)} \le \alpha|x|_{(1)} \qquad \text{对 } x \in \mathbb{K}^{n} .
\]

因为 $ \|e\|_{j} = |A^{-1}e|_{(j)} $，由此可得

\[
\alpha^{-1}\|e\|_{1} \le \|e\|_{2} \le \alpha\|e\|_{1} \qquad \text{对 } e \in E .
\]

因此 $ \|\cdot\|_{1} $ 与 $ \|\cdot\|_{2} $ 是 $ E $ 上的等价范数。

(ii) 这由注 1.3(c) 以及定理 1.4 与 II.6.5 可得。\nobreak\hspace{0.35em}■

\end{corollary}

\begin{theorem}

设 $ E $ 是有限维的。则 $ \operatorname{Hom}(E, F) = \mathcal{L}(E, F) $。换言之，有限维赋范向量空间上的每个线性算子都是连续的。

\proofhead
通过 (1.3) 定义 $ T \in \mathcal{L}\mathrm{is}(E, \mathbb{K}^{n}) $。则存在 $ \tau > 0 $，使得对 $ e \in E $ 有 $ |Te| \le \tau\|e\| $。

现在设 $ A \in \operatorname{Hom}(E, F) $。则 $ Ae = \sum_{j=1}^{n} x^{j}Ab_{j} $。因此，柯西-施瓦茨不等式（取 $ x_{e} := Te $）给出
\[
\|Ae\| \le \left( \sum_{j=1}^{n} \|Ab_{j}\|^{2} \right)^{1/2}|x_{e}| = \alpha|x_{e}| \qquad \text{对 } e \in E ,
\]

其中我们令 $ \alpha := \left( \sum_{j}\|Ab_{j}\|^{2} \right)^{1/2} $。于是我们得到

\[
\|Ae\| \le \alpha|x_{e}| = \alpha|Te| \le \alpha\tau\|e\| \qquad \text{对 } e \in E .
\]

最后，根据定理 VI.2.5，$ A $ 连续。\nobreak\hspace{0.35em}■

\end{theorem}

\begin{remark}

\begin{items}
\item

推论 1.5 与定理 1.6 的陈述对无限维赋范向量空间不成立。
证明

\begin{subitems}
\item

我们令 $ E := BC^{1}((-1, 1), \mathbb{R}) $。在赋范向量空间 $ (E, \|\cdot\|_{\infty}) $ 中，考虑序列 $ (u_{n}) $，其中对 $ t \in [-1, 1] $ 与 $ n \in \mathbb{N}^{\times} $，$ u_{n}(t) := \sqrt{t^{2} + 1/n} $。
容易看出 $ (u_{n}) $ 是 $ (E, \|\cdot\|_{\infty}) $ 中的柯西序列。现在假设 $ (E, \|\cdot\|_{\infty}) $ 完备。则存在 $ u \in E $，使得 $ \lim \|u_{n} - u\|_{\infty} = 0 $。特别地，$ (u_{n}) $ 逐点收敛到 $ u $。
显然，对 $ t \in [-1, 1] $，当 $ n \to \infty $ 时 $ u_{n}(t) = \sqrt{t^{2} + 1/n} \to |t| $。因此，由（逐点收敛情形下）极限的唯一性可得

\[
(t \mapsto |t|) = u \in BC^{1}((-1, 1), \mathbb{R}) ,
\]

这是错误的。这个矛盾表明 $ (E, \|\cdot\|_{\infty}) $ 不完备。
最后，由习题 V.2.10 我们知道，带范数 $ \|u\|_{\bullet} := \|u\|_{\infty} + \|u'\|_{\infty} $ 的 $ (BC^{1}((-1, 1), \mathbb{R}), \|\cdot\|_{\bullet}) $ 是巴拿赫空间。因此，根据注 II.6.7(a)，$ \|\cdot\| $ 与 $ \|\cdot\|_{\bullet} $ 不可能等价。

\item

我们令

\[
E := (C^{1}([0, 1], \mathbb{R}), \|\cdot\|_{\infty}) \quad \text{且} \quad F := (C([0, 1], \mathbb{R}), \|\cdot\|_{\infty})
\]

并考虑 $ A : E \to F $，$ u \mapsto u' $。则 $ E $ 与 $ F $ 是赋范向量空间，且 $ A $ 是线性的。假设 $ A $ 有界，从而连续。因为 $ (u_{n}) $，其中

\[
u_{n}(t) := (1/n)\sin(n^{2}t) \qquad \text{对 } n \in \mathbb{N}^{\times} \text{ 且 } t \in [0, 1]
\]

在 $ E $ 中收敛到零，故在 $ F $ 中有 $ Au_{n} \to 0 $。
由 $ (Au_{n})(t) = n\cos(n^{2}t) $，有 $ (Au_{n})(0) = n $。但由于 $ (Au_{n}) $ 在 $ F $ 中收敛到 0，且一致收敛蕴含逐点收敛，这是不可能的。因此 $ A $ 不属于 $ \mathcal{L}(E, F) $，即 $ \operatorname{Hom}(E, F)\setminus\mathcal{L}(E, F) \ne \varnothing $。\nobreak\hspace{0.35em}■

\end{subitems}

\item

每个有限维内积空间都是希尔伯特空间。

\item

设 $ E $ 是有限维巴拿赫空间。则在 $ E $ 上存在一个等价的\textbf{希尔伯特范数}，即由内积诱导的范数。换言之，每个有限维巴拿赫空间都可以改赋范数成为希尔伯特空间\footnote{此处有限维的假设是重要的。}。

\proofhead
设 $ n := \dim E $ 且 $ T \in \mathcal{L}\mathrm{is}(E, \mathbb{K}^{n}) $。则
\[
(x|y)_{E} := (Tx|Ty) \qquad \text{对 } x, y \in E
\]

定义了 $ E $ 上的一个内积。结论由推论 1.5(i) 可得。\nobreak\hspace{0.35em}■

\end{items}

\end{remark}

\needspace{3\baselineskip}\subsection*{矩阵表示}

设 $ m, n \in \mathbb{N}^{\times} $。我们用 $ \mathbb{K}^{m \times n} $ 表示所有 $ (m \times n) $ 矩阵的集合

\begin{equation*}
[a_{k}^{j}] = \begin{bmatrix} a_{1}^{1} & \cdots & a_{n}^{1} \\ \vdots & & \vdots \\ a_{1}^{m} & \cdots & a_{n}^{m} \end{bmatrix}
\end{equation*}

其中元素 $ a_{k}^{j} $ 属于 $ \mathbb{K} $。这里上指标是行号，下指标是列号\footnote{除 $ a_{k}^{j} $ 外，我们偶尔也写 $ a_{jk} $ 或 $ a^{jk} $，其中第一个指标总是行号，第二个指标总是列号。}。为清晰起见需要时，我们把 $ [a_{k}^{j}] $ 写成

\begin{equation*}
\begin{aligned}
[a_{k}^{j}]_{\substack{1 \le j \le m \\ 1 \le k \le n}} .
\end{aligned}
\end{equation*}

最后，我们令

\begin{equation}
\|M\| := \sqrt{\sum_{j=1}^{m}\sum_{k=1}^{n} |a_{k}^{j}|^{2}} \quad \text{且} \quad M := [a_{k}^{j}] \in \mathbb{K}^{m \times n} .\tag{1.4}
\end{equation}

我们假设你熟悉线性代数中矩阵的基本概念，特别地，你知道 $ \mathbb{K}^{m \times n} $ 关于通常的矩阵加法与数乘是维数为 $ mn $ 的 $ \mathbb{K} $-向量空间。

\qyxffntextdef{1}{另见习题 II.3.14。}
\begin{proposition}

\begin{subitems}
\item

通过 (1.4)，我们在 $ \mathbb{K}^{m \times n} $ 上定义一个范数 $ \|\cdot\| $，称为\textbf{希尔伯特-施密特范数}。因此

\[
\mathbb{K}^{m \times n} := (\mathbb{K}^{m \times n}, \|\cdot\|)
\]

是巴拿赫空间。

\item

映射

\[
\mathbb{K}^{m \times n} \to \mathbb{K}^{mn} , \quad [a_{k}^{j}] \mapsto (a_{1}^{1}, \ldots, a_{n}^{1}, a_{1}^{2}, \ldots, a_{n}^{2}, \ldots, a_{1}^{m}, \ldots, a_{n}^{m})
\]

是等距同构。

\proofhead
我们把这个简单的证明留给你\qyxffnmark{1}。\nobreak\hspace{0.35em}■
\end{subitems}

\end{proposition}

以下，我们总是默认 $ \mathbb{K}^{m \times n} $ 赋有希尔伯特-施密特范数。

设 $ E $ 与 $ F $ 是有限维的，并设

\[
\mathcal{E} = \{e_{1}, \ldots, e_{n}\} \quad \text{且} \quad \mathcal{F} = \{f_{1}, \ldots, f_{m}\}
\]

分别是 $ E $ 与 $ F $ 的（有序）基。根据定理 1.6，$ \operatorname{Hom}(E, F) = \mathcal{L}(E, F) $。我们将用基 $ \mathcal{E} $ 与 $ \mathcal{F} $ 表示 $ A \in \mathcal{L}(E, F) $ 在 $ E $ 中向量上的作用。首先，对 $ k = 1, \ldots, n $，存在唯一的数 $ a_{k}^{1}, \ldots, a_{k}^{m} $，使得 $ Ae_{k} = \sum_{j=1}^{m} a_{k}^{j}f_{j} $。因此由 $ A $ 的线性性，对每个 $ x = \sum_{k=1}^{n} x^{k}e_{k} \in E $，有

\[
Ax = \sum_{k=1}^{n} x^{k}Ae_{k} = \sum_{j=1}^{m}\left( \sum_{k=1}^{n} a_{k}^{j}x^{k} \right)f_{j} = \sum_{j=1}^{m}(Ax)^{j}f_{j}
\]

其中

\[
(Ax)^{j} := \sum_{k=1}^{n} a_{k}^{j}x^{k} \qquad \text{对 } j = 1, \ldots, m .
\]

我们令

\[
[A]_{\mathcal{E},\mathcal{F}} := [a_{k}^{j}] \in \mathbb{K}^{m \times n}
\]

并称 $ [A]_{\mathcal{E},\mathcal{F}} $ 为 $ A $ 关于 $ E $ 的基 $ \mathcal{E} $ 与 $ F $ 的基 $ \mathcal{F} $ 的\textbf{表示矩阵}。（若 $ E = F $，我们简写 $ [A]_{\mathcal{E}} $ 表示 $ [A]_{\mathcal{E},\mathcal{E}} $。）

现在设 $ [a_{k}^{j}] \in \mathbb{K}^{m \times n} $。对 $ x = \sum_{k=1}^{n} x^{k}e_{k} \in E $，我们令

\[
Ax := \sum_{j=1}^{m}\left( \sum_{k=1}^{n} a_{k}^{j}x^{k} \right)f_{j} .
\]

则 $ A := (x \mapsto Ax) $ 是从 $ E $ 到 $ F $ 的线性映射，其表示矩阵为 $ [A]_{\mathcal{E},\mathcal{F}} = [a_{k}^{j}] $。

下一个定理总结了有限维向量空间之间线性映射的表示矩阵的最重要性质。

\qyxffntextdef{1}{对于两个复合的线性映射，我们通常写 $ AB $ 而不写 $ A\circ B $。}
\begin{theorem}

\begin{subitems}
\item

我们令 $ n := \dim E $ 与 $ m := \dim F $，并设 $ \mathcal{E} $ 与 $ \mathcal{F} $ 分别是 $ E $ 与 $ F $ 的基。则\textbf{矩阵表示}

\begin{equation}
\mathcal{L}(E, F) \to \mathbb{K}^{m \times n} , \quad A \mapsto [A]_{\mathcal{E},\mathcal{F}}\tag{1.5}
\end{equation}

是拓扑同构。

\item

设 $ G $ 是有限维赋范向量空间，$ \mathcal{G} $ 是 $ G $ 的基。则，\qyxffnmark{1}

\[
[AB]_{\mathcal{E},\mathcal{G}} = [A]_{\mathcal{F},\mathcal{G}}[B]_{\mathcal{E},\mathcal{F}} \qquad \text{对 } A \in \mathcal{L}(F, G) \text{ 且 } B \in \mathcal{L}(E, F) .
\]

\proofhead
(i) 映射 (1.5) 显然是线性的。上面的讨论表明它还是双射（另见 [Gab96, Chapter D.5.5]）。特别地，$ \dim \mathcal{L}(E, F) = \dim \mathbb{K}^{m \times n} $。因为 $ \mathbb{K}^{m \times n} \cong \mathbb{K}^{mn} $，空间 $ \mathcal{L}(E, F) $ 的维数为 $ mn $。因此我们可以应用定理 1.6，得到
\[
(A \mapsto [A]_{\mathcal{E},\mathcal{F}}) \in \mathcal{L}(\mathcal{L}(E, F), \mathbb{K}^{m \times n}) \quad \text{且} \quad ([A]_{\mathcal{E},\mathcal{F}} \mapsto A) \in \mathcal{L}(\mathbb{K}^{m \times n}, \mathcal{L}(E, F)) .
\]

因此 $ A \mapsto [A]_{\mathcal{E},\mathcal{F}} $ 是拓扑同构。

\end{subitems}

(ii) 这个事实在线性代数中是熟知的（例如 [Gab96, Chapter D.5.5]）。\nobreak\hspace{0.35em}■

\end{theorem}

在分析中，我们要考虑度量空间到两个巴拿赫空间 $ E $ 与 $ F $ 之间连续线性映射空间的映射。若这些空间是有限维的，则根据定理 1.9，我们可以把这些映射视为矩阵。此时，由上面的结果容易证明这些映射连续。正如下一个重要推论所示，只需验证表示矩阵的元素连续即可。

\begin{corollary}

设 $ X $ 是度量空间，$ E $ 与 $ F $ 分别是维数为 $ n $ 与 $ m $ 的有限维巴拿赫空间，并带有序基 $ \mathcal{E} $ 与 $ \mathcal{F} $。设 $ A(\cdot) : X \to \mathcal{L}(E, F) $，并对 $ x \in X $ 令 $ [a_{k}^{j}(x)] $ 为 $ A(x) $ 的表示矩阵 $ [A(x)]_{\mathcal{E},\mathcal{F}} $。则

\[
A(\cdot) \in C(X, \mathcal{L}(E, F)) \iff \left( a_{k}^{j}(\cdot) \in C(X, \mathbb{K}),\ 1 \le j \le m,\ 1 \le k \le n \right) .
\]

\proofhead
这直接由定理 1.9、命题 1.8(ii) 与命题 III.1.10 可得。\nobreak\hspace{0.35em}■
\end{corollary}

约定 除非另有说明，在 $ \mathbb{K}^{n} $ 中我们使用例 I.12.4(a) 中定义的标准基 $ \{e_{1}, \ldots, e_{n}\} $。此外，我们用 $ [A] $ 表示 $ A \in \mathcal{L}(\mathbb{K}^{n}, \mathbb{K}^{m}) $ 的表示矩阵。

\needspace{3\baselineskip}\subsection*{指数映射}

设 $ E $ 是巴拿赫空间，$ A \in \mathcal{L}(E) $。由推论 1.2(ii)，我们知道 $ \mathcal{L}(E) $ 是巴拿赫代数。因此对 $ k \in \mathbb{N} $，$ A^{k} $ 属于 $ \mathcal{L}(E) $，并且

\[
\|A^{k}\| \le \|A\|^{k} \qquad \text{对 } k \in \mathbb{N} .
\]

因此，对每个 $ \alpha \ge \|A\| $，指数级数 $ \sum_{k} \alpha^{k}/k! $ 是\textbf{指数级数}

\begin{equation}
\sum_{k} A^{k}/k!\tag{1.6}
\end{equation}

在 $ \mathcal{L}(E) $ 中的优级数。因此由优级数判别法（定理 II.8.3），级数 (1.6) 在 $ \mathcal{L}(E) $ 中绝对收敛。我们称其和

\[
e^{A} := \sum_{k=0}^{\infty} \frac{A^{k}}{k!}
\]

为 $ A $ 的\textbf{指数}，而

\[
\mathcal{L}(E) \to \mathcal{L}(E) , \quad A \mapsto e^{A}
\]

是 $ \mathcal{L}(E) $ 中的\textbf{指数映射}。

若 $ A $ 与 $ t \in \mathbb{K} $，则 $ tA $ 也属于 $ \mathcal{L}(E) $。因此映射

\begin{equation}
U := U_{A} : \mathbb{K} \to \mathcal{L}(E) , \quad t \mapsto e^{tA}\tag{1.7}
\end{equation}

是良定义的。在下一个定理中，我们汇总函数 $ U $ 的最重要性质。

\begin{theorem}

设 $ E $ 是巴拿赫空间，$ A \in \mathcal{L}(E) $。则

\begin{subitems}
\item

$ U \in C^{\infty}(\mathbb{K}, \mathcal{L}(E)) $ 且 $ \dot{U} = AU $；

\item

$ U $ 是从加法群 $ (\mathbb{K}, +) $ 到乘法群 $ \mathcal{L}\mathrm{aut}(E) $ 的群同态。

\proofhead
(i) 只需证明对每个 $ r \ge 1 $，有
\[
U \in C^{1}(r\mathbb{B}_{\mathbb{C}}, \mathcal{L}(E)) \quad \text{且} \quad \dot{U}(t) = AU(t) \text{ 对 } |t| < r .
\]

结论随之由归纳法得出。
对 $ n \in \mathbb{N} $ 与 $ |t| < r $，定义 $ f_{n}(t) := (tA)^{n}/n! $。则 $ f_{n} \in C^{1}(r\mathbb{B}_{\mathbb{C}}, \mathcal{L}(E)) $，且 $ \sum f_{n} $ 逐点收敛到 $ U|r\mathbb{B}_{\mathbb{C}} $。又 $ \dot{f}_{n} = Af_{n-1} $，从而

\[
\|\dot{f}_{n}(t)\| \le \|A\|\,\|f_{n-1}(t)\| \le (r\|A\|)^{n}/(n-1) \qquad \text{对 } t \in r\mathbb{B}_{\mathbb{C}} \text{ 且 } n \in \mathbb{N}^{\times} .
\]

因此一个标量指数级数是 $ \sum \dot{f}_{n} $ 的优级数。于是 $ \sum \dot{f}_{n} $ 在 $ r\mathbb{B}_{\mathbb{C}} $ 上正规收敛到 $ AU|r\mathbb{B}_{\mathbb{C}} $，定理随之由推论 V.2.9 得出。

\end{subitems}

(ii) 对 $ s, t \in \mathbb{K} $，$ sA $ 与 $ tA $ 在 $ \mathcal{L}(E) $ 中交换。二项式定理（定理 I.8.4）给出

\[
(sA + tA)^{n} = \sum_{k=0}^{n}\binom{n}{k}(sA)^{k}(tA)^{n-k} .
\]

利用指数级数的绝对收敛性，如同例 II.8.12(a) 的证明那样，我们得到

\[
U(s + t) = e^{(s+t)A} = e^{sA+tA} = e^{sA}e^{tA} = U(s)U(t) \qquad \text{对 } s, t \in \mathbb{K} .\quad\blacksquare
\]
\end{theorem}

\begin{remark}

设 $ A \in \mathcal{L}(E) $，$ s, t \in \mathbb{K} $。

\begin{items}
\item

指数 $ e^{sA} $ 与 $ e^{tA} $ 交换：$ e^{sA}e^{tA} = e^{tA}e^{sA} $。

\item

$ \|e^{tA}\| \le e^{|t|\|A\|} $。

\proofhead
这由注 II.8.2(c) 可得。\nobreak\hspace{0.35em}■
\item

$ (e^{tA})^{-1} = e^{-tA} $。

\item

对 $ n \in \mathbb{N} $，$ \partial_{t}^{n}e^{tA} = A^{n}e^{tA} = e^{tA}A^{n} $。

\proofhead
由归纳法，这直接由定理 1.11(i) 得出。\nobreak\hspace{0.35em}■
\end{items}

\end{remark}

\begin{example}

\begin{items}
\item

若 $ E = \mathbb{K} $，则指数映射与指数函数一致。

\item

设 $ N \in \mathcal{L}(E) $ 是\textbf{幂零的}，即存在 $ m \in \mathbb{N} $，使得 $ N^{m+1} = 0 $。则

\[
e^{N} = \sum_{k=0}^{m} \frac{N^{k}}{k!} = 1 + N + \frac{N^{2}}{2!} + \cdots + \frac{N^{m}}{m!} ,
\]

其中 1 表示 $ \mathcal{L}(E) $ 的单位元。

\item

设 $ E $ 与 $ F $ 是巴拿赫空间，$ A \in \mathcal{L}(E) $，$ B \in \mathcal{L}(F) $。又设

\[
A \otimes B : E \times F \to E \times F , \quad (x, y) \mapsto (Ax, By) .
\]

则

\[
e^{A \otimes B} = e^{A} \otimes e^{B} .
\]

\proofhead
这容易由 $ (A \otimes B)^{n} = A^{n} \otimes B^{n} $ 得出。\nobreak\hspace{0.35em}■
\item

对 $ A \in \mathcal{L}(\mathbb{K}^{m}) $ 与 $ B \in \mathcal{L}(\mathbb{K}^{n}) $，有关系式

\begin{equation*}
[e^{A \otimes B}] = \begin{bmatrix} [e^{A}] & 0 \\ 0 & [e^{B}] \end{bmatrix} .
\end{equation*}

\proofhead
这由 (c) 可得。\nobreak\hspace{0.35em}■
\item

设 $ \lambda_{1}, \ldots, \lambda_{m} \in \mathbb{K} $，并通过下式定义 $ \operatorname{diag}(\lambda_{1}, \ldots, \lambda_{m}) \in \mathcal{L}(\mathbb{K}^{m}) $：

\begin{equation*}
[\operatorname{diag}(\lambda_{1}, \ldots, \lambda_{m})] := \begin{bmatrix} \lambda_{1} & & 0 \\ & \ddots & \\ 0 & & \lambda_{m} \end{bmatrix} =: \operatorname{diag}[\lambda_{1}, \ldots, \lambda_{m}] .
\end{equation*}

则

\[
e^{\operatorname{diag}(\lambda_{1}, \ldots, \lambda_{m})} = \operatorname{diag}(e^{\lambda_{1}}, \ldots, e^{\lambda_{m}}) .
\]

\proofhead
这通过利用 (a) 与 (d) 归纳得出。\nobreak\hspace{0.35em}■
\item

对满足下式的 $ A \in \mathcal{L}(\mathbb{R}^{2}) $：

\begin{equation*}
[A] = \begin{bmatrix} 0 & -1 \\ 1 & 0 \end{bmatrix} ,
\end{equation*}

我们有

\begin{equation*}
[e^{tA}] = \begin{bmatrix} \cos t & -\sin t \\ \sin t & \cos t \end{bmatrix} \qquad \text{对 } t \in \mathbb{R} .
\end{equation*}

\proofhead
习题 I.11.10 容易给出
\[
(tA)^{2n+1} = (-1)^{n}t^{2n+1}A , \quad (tA)^{2n} = (-1)^{n}t^{2n}1_{\mathbb{R}^{2}} \qquad \text{对 } t \in \mathbb{R} \text{ 且 } n \in \mathbb{N} .
\]

因为绝对收敛级数可以任意重排，我们得到

\begin{equation*}
\begin{aligned}
\sum_{k=0}^{\infty} \frac{(tA)^{k}}{k!} &= \sum_{n=0}^{\infty} \frac{(tA)^{2n}}{(2n)!} + \sum_{n=0}^{\infty} \frac{(tA)^{2n+1}}{(2n+1)!} \\ &= \left( \sum_{n=0}^{\infty} (-1)^{n}\frac{t^{2n}}{(2n)!} \right)1_{\mathbb{R}^{2}} + \left( \sum_{n=0}^{\infty} (-1)^{n}\frac{t^{2n+1}}{(2n+1)!} \right)A ,
\end{aligned}
\end{equation*}

这就证明了该断言。\nobreak\hspace{0.35em}■

\end{items}

\end{example}

\needspace{3\baselineskip}\subsection*{线性微分方程}

对 $ a \in E $，映射 $ u(\cdot, a) : \mathbb{R} \to E $，$ t \mapsto e^{tA}a $ 是光滑的，且满足“微分方程” $ \dot{x} = Ax $，这立即由定理 1.11(i) 得出。现在我们将更深入地研究这类微分方程。

设 $ A \in \mathcal{L}(E) $，$ f \in C(\mathbb{R}, E) $。则

\begin{equation}
\dot{x} = Ax + f(t) \qquad \text{对 } t \in \mathbb{R} ,\tag{1.8}
\end{equation}

称为 $ E $ 中的\textbf{一阶线性微分方程}。若 $ f = 0 $，则 (1.8) 是\textbf{齐次的}；否则它是\textbf{非齐次的}。若 $ A $ 与“时间” $ t $ 无关，我们称 $ \dot{x} = Ax $ 是\textbf{常系数}线性齐次微分方程。所谓 (1.8) 的\textbf{解}，我们指的是逐点满足 (1.8) 的函数 $ x \in C^{1}(\mathbb{R}, E) $，即对 $ t \in \mathbb{R} $ 有 $ \dot{x}(t) = A(x(t)) + f(t) $。

设 $ a \in E $。则

\begin{equation}
\dot{x} = Ax + f(t) \text{ 对 } t \in \mathbb{R} , \quad \text{满足} \quad x(0) = a\tag*{(1.9)${}_{a}$}
\end{equation}

是微分方程 (1.8) 的\textbf{初值问题}，$ a $ 是\textbf{初值}。这个初值问题的\textbf{解}是指逐点满足 $ (1.9)_{a} $ 的函数 $ x \in C^{1}(\mathbb{R}, E) $。

\begin{remark}

\begin{items}
\item

通过在 (1.8) 与 $ (1.9)_{a} $ 中标出 $ f $ 的自变量 $ t $（但不标出 $ x $ 的），我们采用的是一种不太精确的、符号化的记法，我们这样做是因为它强调 $ f $ 通常依赖于 $ t $。为了完整描述 (1.8) 或 $ (1.9)_{a} $，必须始终指明这类微分方程的解的含义。

\item

为简单起见，这里我们只考虑 $ f $ 在整个 $ \mathbb{R} $ 上有定义的情形。当 $ f $ 只在一个完全子区间上有定义且连续时，所需的明显修改留给你。\nobreak\hspace{0.35em}■

\item

因为 $ C^{1}(\mathbb{R}, E) \subset C(\mathbb{R}, E) $ 且导数算子 $ \partial $ 是线性的（参见定理 IV.1.12），映射

\[
\partial - A : C^{1}(\mathbb{R}, E) \to C(\mathbb{R}, E) , \quad u \mapsto \dot{u} - Au
\]

是线性的。定义

\[
\ker(\partial - A) = \left\{ x \in C^{1}(\mathbb{R}, E) ;\ x \text{ 是 } \dot{x} = Ax \text{ 的解} \right\} =: V
\]

表明齐次微分方程 $ \dot{x} = Ax $ 的\textbf{解空间} $ V $ 是 $ C^{1}(\mathbb{R}, E) $ 的向量子空间。

\item

(1.8) 的全部解构成 $ C^{1}(\mathbb{R}, E) $ 的一个仿射子空间 $ u + V $，其中 $ u $ 是 (1.8) 的（任一）解。

\proofhead
对 $ w := u + v \in u + V $，有 $ (\partial - A)w = (\partial - A)u + (\partial - A)v = f $。因此 $ w $ 是 (1.8) 的解。反之，若 $ w $ 是 (1.8) 的解，则 $ v := u - w $ 属于 $ V $，因为 $ (\partial - A)v = (\partial - A)u - (\partial - A)w = f - f = 0 $。\nobreak\hspace{0.35em}■
\item

对 $ a \in E $ 与 $ x \in C(\mathbb{R}, E) $，令

\[
\Phi_{a}(x)(t) := a + \int_{0}^{t} \left( Ax(s) + f(s) \right)ds .
\]

显然，$ \Phi_{a} $ 是从 $ C(\mathbb{R}, E) $ 到 $ C^{1}(\mathbb{R}, E) $ 的仿射映射，且 $ x $ 是 $ (1.9)_{a} $ 的解当且仅当 $ x $ 是 $ \Phi_{a} $ 在\footnote{当然，我们把 $ \Phi_{a} $ 等同于 $ i \circ \Phi_{a} $，其中 $ i $ 表示 $ C^{1}(\mathbb{R}, E) $ 到 $ C(\mathbb{R}, E) $ 的包含映射。} $ C(\mathbb{R}, E) $ 中的不动点，即当 $ x $ 解线性\textbf{积分方程}

\[
x(t) = a + \int_{0}^{t} \left( Ax(s) + f(s) \right)ds \qquad \text{对 } t \in \mathbb{R} .\quad\blacksquare
\]
\end{items}

\end{remark}

\needspace{3\baselineskip}\subsection*{格朗沃尔引理}

研究微分方程的一个重要工具是下面这个估计，称为\textbf{格朗沃尔引理}。

\begin{lemma}

（格朗沃尔）设 $ J $ 是区间，$ t_{0} \in J $，$ \alpha, \beta \in [0, \infty) $。进一步设 $ y : J \to [0, \infty) $ 连续，且满足

\begin{equation}
y(t) \le \alpha + \beta\left| \int_{t_{0}}^{t} y(\tau)\,d\tau \right| \qquad \text{对 } t \in J .\tag{1.10}
\end{equation}

则

\[
y(t) \le \alpha e^{\beta|t - t_{0}|} \qquad \text{对 } t \in J .
\]

证明

\begin{subitems}
\item

首先考虑 $ t \ge t_{0} $ 的情形，并令

\[
h : [t_{0}, t] \to \mathbb{R}^{+} , \quad s \mapsto \beta e^{\beta(t_{0} - s)}\int_{t_{0}}^{s} y(\tau)\,d\tau \qquad \text{对 } t \in J .
\]

由定理 VI.4.12 与 (1.10)，可得

\[
h'(s) = -\beta h(s) + \beta e^{\beta(t_{0} - s)}y(s) \le \alpha\beta e^{\beta(t_{0} - s)} = \frac{d}{ds}\left( -\alpha e^{\beta(t_{0} - s)} \right) \qquad \text{对 } s \in [t_{0}, t] .
\]

将该不等式从到 $ t_{0} $ 到 $ t $ 积分，得

\[
h(t) = \beta e^{\beta(t_{0} - t)}\int_{t_{0}}^{t} y(\tau)\,d\tau \le \alpha - \alpha e^{\beta(t_{0} - t)} ,
\]

其中我们用了 $ h(t_{0}) = 0 $。于是有

\[
\beta\int_{t_{0}}^{t} y(\tau)\,d\tau \le \alpha e^{\beta(t - t_{0})} - \alpha ,
\]

而 (1.10) 完成了该情形的证明。

\item

当 $ t \le t_{0} $ 时，我们令

\[
h(s) := -\beta e^{\beta(s - t_{0})}\int_{s}^{t_{0}} y(\tau)\,d\tau \qquad \text{对 } s \in [t, t_{0}] ,
\]

并如同 (i) 那样进行。\nobreak\hspace{0.35em}■

\end{subitems}

\end{lemma}

利用格朗沃尔引理，我们可以轻松证明关于方程 (1.8) 的解空间的下一个陈述。

\begin{proposition}

\begin{subitems}
\item

设 $ x \in C^{1}(\mathbb{R}, E) $ 解 $ \dot{x} = Ax $。则

\[
x = 0 \iff x(t_{0}) = 0 \text{ 对某个 } t_{0} \in \mathbb{R} .
\]
\item

初值问题 $ (1.9)_{a} $ 至多有一个解。

\item

对每个 $ a \in E $，初值问题 $ \dot{x} = Ax $，$ x(0) = a $ 有唯一解 $ t \mapsto e^{tA}a $。

\item

$ \dot{x} = Ax $ 的解空间 $ V \subset C^{1}(\mathbb{R}, E) $ 是同构于 $ E $ 的向量空间。一个同构由下式给出

\[
E \to V , \quad a \mapsto (t \mapsto e^{tA}a) .
\]
\item

设 $ x_{1}, \ldots, x_{m} $ 解 $ \dot{x} = Ax $。则下列陈述等价：
(α) $ x_{1}, \ldots, x_{m} $ 在 $ C^{1}(\mathbb{R}, E) $ 中线性无关。
(β) 对每个 $ t \in \mathbb{R} $，$ x_{1}(t), \ldots, x_{m}(t) $ 在 $ E $ 中线性无关。
(γ) 存在 $ t_{0} \in \mathbb{R} $，使得 $ x_{1}(t_{0}), \ldots, x_{m}(t_{0}) $ 在 $ E $ 中线性无关。

\proofhead
(i) 只需验证蕴含“$ \Longleftarrow $”。于是设 $ t_{0} \in \mathbb{R} $ 满足 $ x(t_{0}) = 0 $。则
\[
x(t) = A\int_{t_{0}}^{t} x(\tau)\,d\tau \qquad \text{对 } t \in \mathbb{R}
\]

（参见注 1.14(e) 与习题 VI.3.3）。则 $ y := \|x\| $ 满足不等式

\[
y(t) \le \|A\|\left\| \int_{t_{0}}^{t} x(\tau)\,d\tau \right\| \le \|A\|\int_{t_{0}}^{t} \|x(\tau)\|\,d\tau = \|A\|\int_{t_{0}}^{t} y(\tau)\,d\tau \qquad \text{对 } t \in \mathbb{R} ,
\]

结论由格朗沃尔引理可得。

\end{subitems}

\begin{subitems}
\item[（ii）]

若 $ u $ 与 $ v $ 解 $ (1.9)_{a} $，则 $ w = u - v $ 也解 $ \dot{x} = Ax $，且 $ w(0) = 0 $。由 (i)，得 $ w = 0 $。

\item[（iii）]

这由定理 1.11 与 (ii) 可得。

\item[（iv）]

这是 (iii) 的推论。

\item[（v）]

由于 (iv) 与定理 1.11(ii)，只需证明蕴含“(α)$ \Longrightarrow $(β)”。因此设对某个 $ t_{0} \in \mathbb{R} $，$ \lambda_{j} \in \mathbb{K} $ 满足 $ \sum_{j=1}^{m} \lambda_{j}x_{j}(t_{0}) = 0 $（$ j = 1, \ldots, m $）。因为 $ \sum_{j=1}^{m} \lambda_{j}x_{j} $ 解 $ \dot{x} = Ax $，由 (i) 得 $ \sum_{j=1}^{m} \lambda_{j}x_{j} = 0 $。由假设，$ x_{1}, \ldots, x_{m} $ 在 $ C^{1}(\mathbb{R}, E) $ 中线性无关。因此，由 $ x_{1}(t_{0}), \ldots, x_{m}(t_{0}) $ 在 $ E $ 中的线性无关性，得 $ \lambda_{1} = \cdots = \lambda_{m} = 0 $。\nobreak\hspace{0.35em}■

\end{subitems}

\end{proposition}

\needspace{3\baselineskip}\subsection*{常数变易公式}

由命题 1.16(iii)，齐次方程 $ \dot{x} = Ax $ 的初值问题对每一个初值都有唯一解。为解非齐次方程 (1.8)，我们参照注 1.14(d)：若齐次方程 $ \dot{x} = Ax $ 的解空间 $ V $ 已知，则 (1.8) 的任一解事实上就给出全部解。

为构造这样一个\textbf{特解}，设 $ x \in C^{1}(\mathbb{R}, E) $ 解 (1.8)，且

\begin{equation}
T \in C^{1}(\mathbb{R}, \mathcal{L}(E)) \quad \text{使得 } T(t) \in \mathcal{L}_{\mathrm{aut}}(E)\tag{1.11}
\end{equation}

对 $ t \in \mathbb{R} $。进一步定义

\[
y(t) := T^{-1}(t)x(t) , \qquad \text{其中 } T^{-1}(t) := [T(t)]^{-1} \text{ 且 } t \in \mathbb{R} .
\]

选择“变换” $ T $ 的自由度可用来寻找 $ y $ 的“尽可能简单”的微分方程。首先，乘积法则\footnote{容易验证，定理 IV.1.6(ii) 中乘积法则的证明在当前这个更一般的情形下仍然成立。另见例 4.8(b)。}以及 $ x $ 解 (1.8) 给出

\begin{equation}
\begin{aligned}
\dot{y}(t) &= (T^{-1})^{\cdot}(t)x(t) + T^{-1}(t)\dot{x}(t) \\ &= (T^{-1})^{\cdot}(t)x(t) + T^{-1}(t)Ax(t) + T^{-1}(t)f(t) \\ &= (T^{-1})^{\cdot}(t)T(t)y(t) + T^{-1}(t)AT(t)y(t) + T^{-1}(t)f(t) .
\end{aligned}\tag{1.12}
\end{equation}

接着，对恒等式

\[
T^{-1}(t)T(t) = \mathrm{id}_{E} \quad \text{对 } t \in \mathbb{R} ,
\]

（利用乘积法则）求导，我们得到

\[
(T^{-1})^{\cdot}(t)T(t) + T^{-1}(t)\dot{T}(t) = 0 \quad \text{对 } t \in \mathbb{R} ,
\]

从而

\[
(T^{-1})^{\cdot}(t) = -T^{-1}(t)\dot{T}(t)T^{-1}(t) \quad \text{对 } t \in \mathbb{R} .
\]

与 (1.12) 结合，它变成

\begin{equation}
\dot{y}(t) = T^{-1}(t)\left[AT(t) - \dot{T}(t)\right]y(t) + T^{-1}(t)f(t) \quad \text{对 } t \in \mathbb{R} .\tag{1.13}
\end{equation}

根据定理 1.11，$ t \mapsto T(t) := e^{tA} $ 满足 $ AT - \dot{T} = 0 $ 和 (1.11)。对 $ T $ 的这一选择，由 (1.13) 可得

\[
\dot{y}(t) = T^{-1}(t)f(t) \quad \text{对 } t \in \mathbb{R} .
\]

因此

\[
y(t) = y(0) + \int_{0}^{t}e^{-\tau A}f(\tau)\,d\tau \quad \text{对 } t \in \mathbb{R} .
\]

因为 $ y(0) = x(0) $，我们有

\begin{equation}
x(t) = e^{tA}x(0) + \int_{0}^{t}e^{(t-\tau)A}f(\tau)\,d\tau \quad \text{对 } t \in \mathbb{R} .\tag{1.14}
\end{equation}

现在我们可以容易地证明 $ (1.9)_{a} $ 的存在唯一性结果。

\begin{theorem}

对每个 $ a \in E $ 和 $ f \in C(\mathbb{R}, E) $，存在唯一解 $ u(\cdot\,;a) $ 满足

\[
\dot{x} = Ax + f(t) \quad \text{且} \quad x(0) = a
\]

它由\textbf{常数变易公式}

\begin{equation}
u(t;a) = e^{tA}a + \int_{0}^{t}e^{(t-\tau)A}f(\tau)\,d\tau \quad \text{对 } t \in \mathbb{R} .\tag{1.15}
\end{equation}

给出。

\proofhead
由命题 1.16(ii)，我们只需证明 (1.15) 解 $ (1.9)_{a} $。而这可由定理 1.11 与 VI.4.12 立即推出。\nobreak\hspace{0.35em}■
\end{theorem}

\begin{remark}

为阐明“常数变易公式”这一名称，我们考虑最简单的情形，即 $ \mathbb{R} $ 值方程

\begin{equation}
\dot{x} = ax + f(t) \quad \text{对 } t \in \mathbb{R} .\tag{1.16}
\end{equation}

若 $ a \in \mathbb{R} $，齐次方程 $ \dot{x} = ax $ 有解 $ v(t) := e^{ta} $。因此齐次方程的\textbf{通解}由 $ cv $ 给出，其中 $ c \in \mathbb{R} $ 是常数。也就是说，$ \{cv \;;\; c \in \mathbb{R}\} $ 是齐次方程的解空间。

为指定 (1.16) 的一个特解，我们“变动常数”，即做试设 $ x(t) := c(t)v(t) $，其中 $ c $ 是待定函数。因为 $ x $ 必须满足方程 (1.16)，我们有

\[
\dot{c}v + c\dot{v} = acv + f ,
\]

从而，由 $ c\dot{v} = c(av) $，我们有 $ \dot{c} = f/v $。于是由积分即得 (1.14)。\nobreak\hspace{0.35em}■

\end{remark}

\needspace{3\baselineskip}\subsection*{行列式与特征值}

接下来我们汇集线性代数中的一些重要结果。

\begin{remark}

\begin{items}
\item

$ [a_{k}^{j}] \in \mathbb{K}^{m\times m} $ 的\textbf{行列式} $ \det[a_{k}^{j}] $ 通过符号公式定义

\[
\det[a_{k}^{j}] := \sum_{\sigma \in S_{m}}\operatorname{sign}(\sigma)a_{\sigma(1)}^{1} \cdots\cdots a_{\sigma(m)}^{m}
\]

（见 [Gab96, Chapter A3]），其中 $ \operatorname{sign} $ 是习题 I.9.6 中的符号函数。\footnote{$ \operatorname{sign}(\sigma) $ 也称为置换 $ \sigma $ 的符号。}

\item

设 $ E $ 是有限维巴拿赫空间。\textbf{行列式函数}

\[
\det : \mathcal{L}(E) \to \mathbb{K}
\]

对 $ A \in \mathcal{L}(E) $ 可表示为 $ \det(A) := \det([A]_{\mathcal{E}}) $，其中 $ \mathcal{E} $ 是 $ E $ 的一个基。于是 $ \det $ 是良定义的，即不依赖于特定的基 $ \mathcal{E} $，且我们有

\begin{equation}
\det(1_{E}) = 1 \quad \text{且} \quad \det(AB) = \det(A)\det(B) \quad \text{对 } A, B \in \mathcal{L}(E) ,\tag{1.17}
\end{equation}

且

\begin{equation}
\det(A) \ne 0 \iff A \in \mathcal{L}_{\mathrm{aut}}(E) .\tag{1.18}
\end{equation}

令 $ \lambda - A := \lambda 1_{E} - A $ 且 $ m := \dim(E) $，我们有

\begin{equation}
\det(\lambda - A) = \sum_{k=0}^{m}(-1)^{k}\alpha_{m-k}\lambda^{m-k} \quad \text{对 } \lambda \in \mathbb{C} ,\tag{1.19}
\end{equation}

其中 $ \alpha_{k} \in \mathbb{C} $ 且

\begin{equation}
\alpha_{m} = 1 , \quad \alpha_{m-1} = \operatorname{tr}(A) , \quad \alpha_{0} = \det(A) .\tag{1.20}
\end{equation}

\textbf{特征多项式} (1.19) 的零点是 $ A $ 的\textbf{特征值}，而所有特征值之集就是 $ A $ 的\textbf{谱} $ \sigma(A) $。
数 $ \lambda \in \mathbb{K} $ 是 $ A $ 的特征值当且仅当\textbf{特征值方程}

\begin{equation}
Ax = \lambda x\tag{1.21}
\end{equation}

有非平凡解 $ x \in E \setminus \{0\} $，即若 $ \ker(\lambda - A) \ne \{0\} $。对\footnote{注意，若 $ E $ 是实巴拿赫空间，则 $ A $ 的谱也是 $ \mathbb{C} $ 的子集。若 $ \mathbb{K} = \mathbb{R} $ 且 $ \lambda $ 属于 $ \sigma(A)\setminus\mathbb{R} $，则方程 (1.21) 无意义。此时我们必须允许 $ E $ 与 $ A $ 的复化，从而使 $ \lambda $ 可以作为特征值方程的解（例如见 [Ama95, § 12]）。于是在每种情形下 $ \lambda \in \sigma(A) $ 的几何重数都有定义。} $ \lambda \in \sigma(A) \cap \mathbb{K} $，$ \lambda $ 的\textbf{几何重数}由 $ \dim(\ker(\lambda - A)) $ 给出，且每个 $ x \in \ker(\lambda - A)\setminus\{0\} $ 是 $ A $ 的以 $ \lambda $ 为\textbf{特征值}的\textbf{特征向量}。$ \lambda $ 作为特征多项式的零点出现的次数就是 $ \lambda $ 的\textbf{代数重数}。几何重数不大于代数重数。最后，若 $ \lambda $ 的代数重数为 1，从而 $ \lambda $ 是特征多项式的单零点，则 $ \lambda $ 是 $ A $ 的\textbf{单特征值}。若几何重数与代数重数相等，则 $ \lambda $ 是\textbf{半单的}。
每个 $ A \in \mathcal{L}(E) $ 恰有 $ m $ 个特征值 $ \lambda_{1},\ldots,\lambda_{m} $（若按（代数）重数计数）。$ A $ 的特征多项式的系数 $ \alpha_{k} $ 是 $ \lambda_{1},\ldots,\lambda_{m} $ 的初等对称函数。特别地，

\begin{equation}
\operatorname{tr}(A) = \lambda_{1} + \cdots + \lambda_{m} \quad \text{且} \quad \det(A) = \lambda_{1}\cdots\cdots\lambda_{m} .\tag{1.22}
\end{equation}

若 $ \mathcal{E} $ 是 $ E $ 的一个基，使得 $ [A]_{\mathcal{E}} $ 是（上）三角矩阵，

\begin{equation}
[A]_{\mathcal{E}} = \begin{bmatrix} a_{1}^{1} & a_{2}^{1} & \cdots & a_{m}^{1} \\ & a_{2}^{2} & \cdots & a_{m}^{2} \\ & & \ddots & \vdots \\ 0 & & & a_{m}^{m} \end{bmatrix} ,\tag{1.23}
\end{equation}

即 $ a_{k}^{j} = 0 $（$ k < j $），则 $ A $ 的特征值连同其代数重数沿对角线出现。

\proofhead
方阵行列式的性质 (1.17)–(1.20) 是线性代数中的标准结果，我们假定读者已了解。
若 $ \mathcal{F} $ 是 $ E $ 的另一个基，则存在可逆矩阵 $ T \in \mathbb{K}^{m\times m} $ 使得
\[
[A]_{\mathcal{F}} = T[A]_{\mathcal{E}}T^{-1}
\]

（例如 [Koe83, Chapter 9.3.2]）。因此由矩阵法则 (1.17) 可得 $ \det([A]_{\mathcal{F}}) = \det([A]_{\mathcal{E}}) $。这表明 $ \det(A) $ 对 $ A \in \mathcal{L}(E) $ 是良定义的。现在显然 (1.17)–(1.20) 对 $ A \in \mathcal{L}(E) $ 成立。
$ A $ 的特征值的存在性，以及 $ A $ 按代数重数计数时恰有 $ m $ 个特征值的断言，都由代数基本定理（见例 III.3.9(b)）推出。\nobreak\hspace{0.35em}■

\item

设 $ p \in K[X] $ 是域 $ K $ 上的 $ n \ge 1 $ 次多项式。若存在 $ k, \nu_{1},\ldots,\nu_{k} \in \mathbb{N}^{\times} $ 以及 $ a, \lambda_{1},\ldots,\lambda_{k} \in K $ 使得

\[
p = a\prod_{j=1}^{k}(X - \lambda_{j})^{\nu_{j}} ,
\]

则 $ p $ 在 $ K $ 中\textbf{分裂}。代数基本定理蕴含每个非常数多项式 $ p \in \mathbb{C}[X] $ 都在 $ \mathbb{C} $ 中分裂（见例 III.3.9(b)）。一般地，多项式未必分裂。例如，若 $ K $ 是有序域，则 $ X^{2} + 1 $ 在 $ K $ 中不分裂。

\item

设 $ \lambda \in \mathbb{K} $ 且 $ p \in \mathbb{N}^{\times} $。则由

\begin{equation*}
J(\lambda,1) := [\lambda] \quad \text{且} \quad J(\lambda,p) := \begin{bmatrix} \lambda & 1 & & & 0 \\ & \ddots & \ddots & & \\ & & \ddots & \ddots & \\ & & & \ddots & 1 \\ 0 & & & & \lambda \end{bmatrix}
\end{equation*}

定义的 $ J(\lambda,p) \in \mathbb{K}^{p\times p} $（$ p \ge 2 $）称为 $ \lambda $ 的 $ p $ 阶\textbf{初等若尔当矩阵}。
对满足 $ \alpha \in \mathbb{R} $ 且 $ \omega > 0 $ 的 $ \mu = \alpha + i\omega \in \mathbb{C} $，令

\begin{equation*}
A_{2}(\mu) := \begin{bmatrix} \alpha & -\omega \\ \omega & \alpha \end{bmatrix} \in \mathbb{R}^{2\times 2} ,
\end{equation*}

并令 $ 1_{2} $ 表示 $ \mathbb{R}^{2\times 2} $ 中的单位元。则我们把由

\begin{equation*}
\widetilde{J}(\mu,1) := A_{2}(\mu) \quad \text{且} \quad \widetilde{J}(\mu,p) := \begin{bmatrix} A_{2}(\mu) & 1_{2} & & & 0 \\ & \ddots & \ddots & & \\ & & \ddots & \ddots & \\ & & & \ddots & 1_{2} \\ 0 & & & & A_{2}(\mu) \end{bmatrix}
\end{equation*}

定义的 $ \widetilde{J}(\mu,p) \in \mathbb{R}^{2p\times 2p} $（$ p \ge 2 $）称为 $ \mu $ 的 $ 2p $ 阶\textbf{扩展若尔当矩阵}。

\item

设 $ E $ 是 $ \mathbb{K} $ 上的有限维巴拿赫空间且 $ A \in \mathcal{L}(E) $。又设 $ A $ 的特征多项式在 $ \mathbb{K} $ 中分裂。在线性代数中可证，此时存在 $ E $ 的一个基 $ \mathcal{E} $ 以及 $ p_{1},\ldots,p_{r} \in \mathbb{N}^{\times} $ 使得

\begin{equation}
[A]_{\mathcal{E}} = \begin{bmatrix} J(\lambda_{1},p_{1}) & & 0 \\ & \ddots & \\ 0 & & J(\lambda_{r},p_{r}) \end{bmatrix} \in \mathbb{K}^{n\times n} ,\tag{1.24}
\end{equation}

其中 $ \{\lambda_{1},\ldots,\lambda_{r}\} = \sigma(A) $。(1.24) 称为 $ A $ 的\textbf{若尔当标准形}。它在分块 $ J(\lambda_{j},p_{j}) $ 的排列顺序意义上唯一。基 $ \mathcal{E} $ 称为 $ E $ 关于 $ A $ 的\textbf{若尔当基}。若 $ A $ 只有半单特征值，则 $ r = \dim(E) $ 且 $ p_{j} = 1 $（$ j = 1,\ldots,r $）。这样的 $ A $ 称为\textbf{可对角化的}。

\item

矩阵

\begin{equation*}
A = \begin{bmatrix} 0 & 1 & 0 \\ -1 & 0 & 0 \\ 0 & 0 & 1 \end{bmatrix} \in \mathbb{R}^{3\times 3}
\end{equation*}

的特征多项式为

\[
p = X^{3} - X^{2} + X - 1 = (X - 1)(X^{2} + 1) .
\]

因此 $ p $ 在 $ \mathbb{R} $ 中不分裂，且 $ \sigma(A) = \{1, i, -i\} $。
此时 $ A $ 不能化为 (1.24) 的形式。但它可以化为扩展若尔当标准形。

\item

设 $ E $ 是有限维实巴拿赫空间，且 $ A \in \mathcal{L}(E) $ 的特征多项式 $ p $ 在 $ \mathbb{R} $ 中不分裂。则 $ A $ 有一个非实特征值 $ \mu $。因为特征多项式 $ p = \sum_{k=0}^{n}a_{k}X^{k} $ 只有实系数，由

\[
0 = \bar{0} = \overline{p(\mu)} = \sum_{k=0}^{n}a_{k}\mu^{k} = \sum_{k=0}^{n}a_{k}\bar{\mu}^{k} = p(\bar{\mu})
\]

可得 $ \bar{\mu} $ 也是 $ A $ 的特征值。
现在设

\[
\sigma := \{\mu_{1},\bar{\mu}_{1},\ldots,\mu_{\ell},\bar{\mu}_{\ell}\} \subset \sigma(A) \quad \text{对 } \ell \ge 1 ,
\]

是 $ A $ 的全部非实特征值之集，其中，不失一般性，我们可设 $ \mu_{j} = \alpha_{j} + i\omega_{j} $，其中 $ \alpha_{j} \in \mathbb{R} $ 且 $ \omega_{j} > 0 $。若 $ A $ 还有实特征值，我们用 $ \lambda_{1},\ldots,\lambda_{k} $ 记它们。
则存在 $ E $ 的一个基 $ \mathcal{E} $ 以及 $ p_{1},\ldots,p_{r},q_{1},\ldots,q_{s} \in \mathbb{N}^{\times} $ 使得

\begin{equation}
[A]_{\mathcal{E}} = \begin{bmatrix} J(\lambda_{1},p_{1}) & & & & & 0 \\ & \ddots & & & & \\ & & J(\lambda_{r},p_{r}) & & & \\ & & & \widetilde{J}(\mu_{1},q_{1}) & & \\ & & & & \ddots & \\ 0 & & & & & \widetilde{J}(\mu_{s},q_{s}) \end{bmatrix} \in \mathbb{R}^{n\times n} .\tag{1.25}
\end{equation}

这里，$ \lambda_{j} \in \sigma(A)\setminus\sigma $（$ j = 1,\ldots,r $）且 $ \mu_{k} \in \sigma $（$ k = 1,\ldots,s $）。我们把 (1.25) 称为 $ A $ 的\textbf{扩展若尔当标准形}。它在各组成矩阵的排列顺序意义上唯一。（见 [Gab96, Chapter A5]。）\nobreak\hspace{0.35em}■

\end{items}

\end{remark}

\needspace{3\baselineskip}\subsection*{基本解矩阵}

设 $ E $ 是 $ \mathbb{K} $ 上的 $ n $ 维向量空间且 $ A \in \mathcal{L}(E) $。又用 $ V \subset C^{1}(\mathbb{R}, E) $ 表示齐次线性微分方程 $ \dot{x} = Ax $ 的解空间。由命题 1.16(iv)，$ V $ 与 $ E $ 同构，特别地有相同维数。我们把 $ V $ 的每个基称为 $ \dot{x} = Ax $ 的\textbf{基本解组}。

\begin{example}

\begin{items}
\item

令

\begin{equation*}
A := \begin{bmatrix} a & b \\ 0 & c \end{bmatrix} \in \mathbb{K}^{2\times 2} .
\end{equation*}

为求 $ \dot{x} = Ax $ 的一个基本解组，我们先解初值问题

\begin{equation}
\begin{aligned}
\dot{x}^{1} &= ax^{1} + bx^{2} , & x^{1}(0) &= 1 , \\ \dot{x}^{2} &= cx^{2} , & x^{2}(0) &= 0 .
\end{aligned}\tag{1.26}
\end{equation}

由应用 IV.3.9(b) 可得 $ x^{2} = 0 $ 且 $ x^{1}(t) = e^{at} $（$ t \in \mathbb{R} $）。因此 $ x_{1} := (t \mapsto (e^{at},0)) $ 是 (1.26) 的解。我们通过解

\begin{equation*}
\begin{aligned}
\dot{x}^{1} &= ax^{1} + bx^{2} , & x^{1}(0) &= 0 , \\ \dot{x}^{2} &= cx^{2} , & x^{2}(0) &= 1 .
\end{aligned}
\end{equation*}

得到与 $ x_{1} $ 线性无关的第二个解 $ x_{2} $。此时我们有 $ x_{2}^{2}(t) = e^{ct} $（$ t \in \mathbb{R} $），而常数变易公式给出

\begin{equation*}
\begin{aligned}
x_{2}^{1}(t) &= \int_{0}^{t}e^{a(t-\tau)}be^{c\tau}\,d\tau = be^{at}\int_{0}^{t}e^{\tau(c-a)}\,d\tau \\ &= \begin{cases} b(e^{ct} - e^{at})/(c - a) , & a \ne c , \\ bte^{ct} , & a = c . \end{cases}
\end{aligned}
\end{equation*}

因此

\begin{equation*}
x_{1}(t) = (e^{at},0) , \quad x_{2}(t) = \begin{cases} \left( b(e^{ct} - e^{at})/(c - a), e^{ct} \right) , & a \ne c , \\ (bte^{at}, e^{at}) , & a = c , \end{cases}
\end{equation*}

（其中 $ t \in \mathbb{R} $）是 $ \dot{x} = Ax $ 的一个基本解组。

\item

设 $ \omega > 0 $ 且

\begin{equation*}
A := \begin{bmatrix} 0 & -\omega \\ \omega & 0 \end{bmatrix} \in \mathbb{K}^{2\times 2} .
\end{equation*}

由

\[
\dot{x}^{1} = -\omega x^{2} \quad \text{且} \quad \dot{x}^{2} = \omega x^{1}
\]

可得 $ \ddot{x}^{j} + \omega^{2}x^{j} = 0 $（$ j = 0,1 $）。习题 IV.3.2(a) 表明

\[
x_{1}(t) = (\cos(\omega t),\sin(\omega t)) \quad \text{且} \quad x_{2}(t) = (-\sin(\omega t),\cos(\omega t)) \quad \text{对 } t \in \mathbb{R} ,
\]

是 $ \dot{x} = Ax $ 的一个基本解组。\nobreak\hspace{0.35em}■

\end{items}

若 $ A \in \mathbb{K}^{n\times n} $ 且 $ \{x_{1},\ldots,x_{n}\} $ 是 $ \dot{x} = Ax $ 的一个基本解组，则我们把映射

\begin{equation*}
X : \mathbb{R} \to \mathbb{K}^{n\times n} , \quad t \mapsto [x_{1}(t),\ldots,x_{n}(t)] = \begin{bmatrix} x_{1}^{1}(t) & \cdots & x_{n}^{1}(t) \\ \vdots & & \vdots \\ x_{1}^{n}(t) & \cdots & x_{n}^{n}(t) \end{bmatrix}
\end{equation*}

称为 $ \dot{x} = Ax $ 的\textbf{基本解矩阵}。若 $ X(0) = 1_{n} $，则我们把 $ X $ 称为\textbf{主基本解矩阵}。

\end{example}

\begin{remark}

设 $ A \in \mathbb{K}^{n\times n} $。

\begin{items}
\item

映射

\[
\mathbb{K}^{n\times n} \to \mathbb{K}^{n\times n} , \quad X \mapsto AX
\]

是线性的。因此，除 $ \mathbb{K}^{n} $ 中的线性微分方程 $ \dot{x} = Ax $ 之外，我们还可考虑 $ \mathbb{K}^{n\times n} $ 中的线性微分方程 $ \dot{X} = AX $，它可由 $ \dot{x} = Ax $ 的任一基本解矩阵来解。

\item

设 $ X $ 是 $ \dot{x} = Ax $ 的一个基本解矩阵且 $ t \in \mathbb{R} $。则

\begin{subitems}
\item

$ X(t) \in \mathcal{L}_{\mathrm{aut}}(\mathbb{K}^{n}) $；

\item

$ X(t) = e^{tA}X(0) $。

\proofhead
(i) 是命题 1.16(v) 的推论。(ii) 由 (a)、定理 1.11 与命题 1.16(iii) 推出。\nobreak\hspace{0.35em}■
\end{subitems}

\item

映射 $ t \mapsto e^{tA} $ 是 $ \dot{x} = Ax $ 的唯一主基本解矩阵。

\proofhead
这由 (b) 推出。\nobreak\hspace{0.35em}■
\item

对 $ s, t \in \mathbb{R} $，我们有

\[
e^{(t-s)A} = X(t)X^{-1}(s) .
\]

\proofhead
这也由 (b) 推出。\nobreak\hspace{0.35em}■
\item

利用若尔当标准形和例 1.13(b)–(e)，我们可以计算 $ e^{tA} $。我们就 $ n = 2 $ 的情形加以演示。
矩阵

\begin{equation*}
A := \begin{bmatrix} a & b \\ c & d \end{bmatrix} \in \mathbb{R}^{2\times 2}
\end{equation*}

的特征多项式给出 $ \det(\lambda - A) = \lambda^{2} - \lambda(a + d) + ad - bc $，且其零点为

\[
\lambda_{1,2} = \left( a + d \pm \sqrt{D} \right)/2 , \quad \text{其中} \quad D := (a - d)^{2} + 4bc .
\]

\textbf{情形 1：} $ \mathbb{K} = \mathbb{C} $，$ D < 0 $。此时 $ A $ 有两个单的、互为共轭的复特征值 $ \lambda := \lambda_{1} $ 和 $ \bar{\lambda} = \lambda_{2} $。根据注 1.19(e)，存在 $ \mathbb{C}^{2} $ 的一个基 $ \mathcal{B} $ 使得 $ [A]_{\mathcal{B}} = \mathrm{diag}[\lambda,\bar{\lambda}] $，且由例 1.13(e) 可得

\[
e^{t[A]_{\mathcal{B}}} = \mathrm{diag}[e^{\lambda t}, e^{\bar{\lambda} t}] \quad \text{对 } t \in \mathbb{R} .
\]

用 $ T \in \mathcal{L}_{\mathrm{aut}}(\mathbb{C}^{2}) $ 表示从标准基到 $ \mathcal{B} $ 的基变换。则利用定理 1.9，我们有

\[
[A]_{\mathcal{B}} = [TAT^{-1}] = [T]A[T]^{-1} .
\]

因此，由习题 11，我们得到

\[
e^{t[A]_{\mathcal{B}}} = [T]e^{tA}[T]^{-1} ,
\]

并求得

\[
e^{tA} = [T]^{-1}\mathrm{diag}[e^{\lambda t}, e^{\bar{\lambda} t}][T] \quad \text{对 } t \in \mathbb{R} .
\]

\textbf{情形 2：} $ \mathbb{K} = \mathbb{C} $ 或 $ \mathbb{K} = \mathbb{R} $，$ D > 0 $。则 $ A $ 有两个单的实特征值 $ \lambda_{1} $ 和 $ \lambda_{2} $，同情形 1 一样可得

\[
e^{tA} = [T]^{-1}\mathrm{diag}[e^{\lambda_{1}t}, e^{\lambda_{2}t}][T] \quad \text{对 } t \in \mathbb{R} .
\]

\textbf{情形 3：} $ \mathbb{K} = \mathbb{C} $ 或 $ \mathbb{K} = \mathbb{R} $，$ D = 0 $。此时 $ A $ 有代数重数为 2 的特征值 $ \lambda := (a + d)/2 $。因此，根据注 1.19(e)，存在 $ \mathbb{K}^{2} $ 的一个基 $ \mathcal{B} $ 使得

\begin{equation*}
[A]_{\mathcal{B}} = \begin{bmatrix} \lambda & 1 \\ 0 & \lambda \end{bmatrix} .
\end{equation*}

因为矩阵

\begin{equation*}
\begin{bmatrix} \lambda & 0 \\ 0 & \lambda \end{bmatrix} \quad \text{且} \quad \begin{bmatrix} 0 & 1 \\ 0 & 0 \end{bmatrix}
\end{equation*}

可交换，由习题 11 和例 1.13(e) 可得

\begin{equation*}
e^{t[A]_{\mathcal{B}}} = e^{\lambda t}\exp\left( t\begin{bmatrix} 0 & 1 \\ 0 & 0 \end{bmatrix} \right) = e^{\lambda t}\left( 1_{2} + t\begin{bmatrix} 0 & 1 \\ 0 & 0 \end{bmatrix} \right) = e^{\lambda t}\begin{bmatrix} 1 & t \\ 0 & 1 \end{bmatrix} .
\end{equation*}

于是我们得到

\begin{equation*}
e^{tA} = e^{\lambda t}[T]^{-1}\begin{bmatrix} 1 & t \\ 0 & 1 \end{bmatrix}[T] \quad \text{对 } t \in \mathbb{R} .
\end{equation*}

\textbf{情形 4：} $ \mathbb{K} = \mathbb{R} $，$ D < 0 $。$ A $ 有两个互为共轭的复特征值 $ \lambda := \alpha + i\omega $ 和 $ \bar{\lambda} $，其中 $ \alpha := (a + d)/2 $ 且 $ \omega := \sqrt{-D}/2 $。由注 1.19(g)，存在 $ \mathbb{R}^{2} $ 的一个基 $ \mathcal{B} $ 使得

\begin{equation*}
[A]_{\mathcal{B}} = \begin{bmatrix} \alpha & -\omega \\ \omega & \alpha \end{bmatrix} .
\end{equation*}

矩阵

\begin{equation*}
\begin{bmatrix} \alpha & 0 \\ 0 & \alpha \end{bmatrix} \quad \text{且} \quad \begin{bmatrix} 0 & -\omega \\ \omega & 0 \end{bmatrix}
\end{equation*}

可交换。于是由例 1.13(f) 我们有

\begin{equation*}
\exp\left( t\begin{bmatrix} \alpha & -\omega \\ \omega & \alpha \end{bmatrix} \right) = e^{t\alpha}\exp\left( t\omega\begin{bmatrix} 0 & -1 \\ 1 & 0 \end{bmatrix} \right) = e^{t\alpha}\begin{bmatrix} \cos(\omega t) & -\sin(\omega t) \\ \sin(\omega t) & \cos(\omega t) \end{bmatrix} .
\end{equation*}

由此我们推出

\begin{equation*}
e^{tA} = e^{\alpha t}[T]^{-1}\begin{bmatrix} \cos(\omega t) & -\sin(\omega t) \\ \sin(\omega t) & \cos(\omega t) \end{bmatrix}[T] \quad \text{对 } t \in \mathbb{R} ,
\end{equation*}

其中 $ T $ 是从标准基到 $ \mathcal{B} $ 的基变换（见例 1.13(f)）。\nobreak\hspace{0.35em}■

\end{items}

\end{remark}

\needspace{3\baselineskip}\subsection*{二阶线性微分方程}

迄今为止，我们只考虑了一阶线性微分方程，其中待求函数与其导数之间是线性关系。然而许多应用需要高阶微分方程。二阶的情形经常出现，我们现在就集中考虑这种情形。\footnote{对更高阶的微分方程，可查阅常微分方程的文献，例如 [Ama95]。}

下面，设

• $ b, c \in \mathbb{R} $ 且 $ g \in C(\mathbb{R}, \mathbb{K}) $。

对由下式给出的\textbf{常系数二阶线性微分方程}

\begin{equation}
\ddot{u} + b\dot{u} + cu = g(t) \quad \text{对 } t \in \mathbb{R} ,\tag{1.27}
\end{equation}

我们说 $ u \in C^{2}(\mathbb{R}, \mathbb{K}) $ 是一个\textbf{解}，若它逐点满足 (1.27)。下面的结果表明，(1.27) 等价于 $ (u,\dot{u}) $ 平面（称为\textbf{相平面}）中形如 (1.8) 的一阶微分方程。

\begin{lemma}

\begin{subitems}
\item

若 $ u \in C^{2}(\mathbb{R}, \mathbb{K}) $ 解 (1.27)，则 $ (u,\dot{u}) \in C^{1}(\mathbb{R}, \mathbb{K}^{2}) $ 解方程

\begin{equation}
\dot{x} = Ax + f(t)\tag{1.28}
\end{equation}

（在 $ \mathbb{K}^{2} $ 中），其中

\begin{equation}
A := \begin{bmatrix} 0 & 1 \\ -c & -b \end{bmatrix} \quad \text{且} \quad f := (0,g) .\tag{1.29}
\end{equation}
\item

若 $ x \in C^{1}(\mathbb{R}, \mathbb{K}^{2}) $ 解 (1.28)，则 $ u := \mathrm{pr}_{1}x $ 解 (1.27)。

\proofhead
(i) 令 $ x := (u,\dot{u}) $，由 (1.27) 得
\begin{equation*}
\dot{x} = \begin{bmatrix} \dot{u} \\ \ddot{u} \end{bmatrix} = \begin{bmatrix} \dot{u} \\ -cu - b\dot{u} + g(t) \end{bmatrix} = Ax + f(t) .
\end{equation*}
\end{subitems}

(ii) 设 $ x = (u,v) \in C^{1}(\mathbb{R}, \mathbb{K}^{2}) $ 解 (1.28)。则

\begin{equation}
\dot{u} = v , \quad \dot{v} = -cu - bv + g(t) .\tag{1.30}
\end{equation}

第一个方程表明 $ u $ 属于 $ C^{2}(\mathbb{R}, \mathbb{K}) $，而这两个方程合在一起蕴含 $ u $ 解微分方程 (1.27)。\nobreak\hspace{0.35em}■

\end{lemma}

引理 1.22 使我们可以把关于一阶微分方程的知识应用到 (1.27) 上。我们需要 $ A $ 的特征值。由 $ \det(\lambda - A) = \lambda(\lambda + b) + c $，我们发现 $ A $ 的特征值 $ \lambda_{1} $ 和 $ \lambda_{2} $ 是多项式 $ X^{2} + bX + c $（称为 $ \ddot{u} + b\dot{u} + cu = 0 $ 的\textbf{特征多项式}）的零点，且我们得到

\[
\lambda_{1,2} = \left( -b \pm \sqrt{D} \right)/2 , \quad \text{其中} \quad D = b^{2} - 4c .
\]

现在我们考虑初值问题

\begin{equation}
\ddot{u} + b\dot{u} + cu = g(t) \text{ 对 } t \in \mathbb{R} , \quad \text{满足} \quad u(0) = a_{1} , \quad \dot{u}(0) = a_{2} ,\tag{1.31}
\end{equation}

其中 $ (a_{1},a_{2}) \in \mathbb{K}^{2} $。由以上，我们可以容易地证明一个基本定理。

\begin{theorem}

\begin{subitems}
\item

对每个 $ (a_{1},a_{2}) \in \mathbb{K}^{2} $，初值问题 (1.31) 存在唯一解 $ u \in C^{2}(\mathbb{R}, \mathbb{K}) $。

\item

$ \ddot{u} + b\dot{u} + cu = 0 $ 的解集通过

\begin{equation*}
\begin{cases} \{e^{\lambda_{1}t}, e^{\lambda_{2}t}\} & \text{若 } D > 0 \text{ 或 } (D < 0 \text{ 且 } \mathbb{K} = \mathbb{C}) , \\ \{e^{\alpha t}, te^{\alpha t}\} & \text{若 } D = 0 , \\ \{e^{\alpha t}\cos(\omega t), e^{\alpha t}\sin(\omega t)\} & \text{若 } D < 0 \text{ 且 } \mathbb{K} = \mathbb{R} , \end{cases}
\end{equation*}

张成 $ C^{2}(\mathbb{R}, \mathbb{K}) $ 的二维向量子空间 $ V $，其中 $ \alpha := -b/2 $ 且 $ \omega := \sqrt{-D}/2 $。

\item

(1.27) 的解集构成 $ C^{2}(\mathbb{R}, \mathbb{K}) $ 的二维仿射子空间 $ v + V $，其中 $ v $ 是 (1.27) 的任一解。

\item

设方程 $ \ddot{u} + b\dot{u} + cu = 0 $ 和 $ w := u_{1}\dot{u}_{2} - \dot{u}_{1}u_{2} $ 有线性无关的解 $ u_{1}, u_{2} \in V $。则

\[
v(t) := \int_{0}^{t}\frac{u_{1}(\tau)g(\tau)}{w(\tau)}\,d\tau\,u_{2}(t) - \int_{0}^{t}\frac{u_{2}(\tau)g(\tau)}{w(\tau)}\,d\tau\,u_{1}(t) \quad \text{对 } t \in \mathbb{R}
\]

解 (1.27)。

\proofhead
(i) 这由引理 1.22 与定理 1.17 立即可得。
\end{subitems}

(ii) 和 (iii) 这些由引理 1.22、命题 1.16(iv) 以及注 1.14(d) 和 1.21(e) 蕴含。

(iv) 设 $ A $ 和 $ f $ 如 (1.29) 所示。则由引理 1.22 和注 1.21(a)，

\begin{equation*}
X := \begin{bmatrix} u_{1} & u_{2} \\ \dot{u}_{1} & \dot{u}_{2} \end{bmatrix}
\end{equation*}

是 $ \dot{x} = Ax $ 的一个基本解矩阵。因为

\[
e^{(t-\tau)A} = X(t)X^{-1}(\tau) \quad \text{对 } t, \tau \in \mathbb{R}
\]

（见注 1.21(d)），由定理 1.17 的常数变易公式可得

\begin{equation}
y(t) = X(t)\int_{0}^{t}X^{-1}(\tau)f(\tau)\,d\tau \quad \text{对 } t \in \mathbb{R} ,\tag{1.32}
\end{equation}

解 $ \dot{x} = Ax + f(t) $。因为 $ \det(X) = w $，我们有

\begin{equation*}
X^{-1} = \frac{1}{w}\begin{bmatrix} \dot{u}_{2} & -u_{2} \\ -\dot{u}_{1} & u_{1} \end{bmatrix} ,
\end{equation*}

其中函数 $ w $ 无零点，这是由命题 1.16(v) 得到的。因此我们断定 $ X^{-1}f = (-u_{2}g/w, u_{1}g/w) $。于是由引理 1.22 和 (1.32) 即得本定理。\nobreak\hspace{0.35em}■

\end{theorem}

\begin{example}

\begin{items}
\item

方程

\[
\ddot{u} - 2\dot{u} + 10u = 0
\]

的特征多项式 $ X^{2} - 2X + 10 $ 有零点 $ 1 \pm 3i $。因此 $ \{e^{t}\cos(3t), e^{t}\sin(3t)\} $ 是一个基本解组。\footnote{与一阶线性微分方程一样，我们把解空间的任一基称为基本解组。}因此通解为

\[
u(t) = e^{t}\left( a_{1}\cos(3t) + a_{2}\sin(3t) \right) \quad \text{对 } t \in \mathbb{R} ,
\]

其中 $ a_{1}, a_{2} \in \mathbb{R} $。

\end{items}

\bookfigure[0.66]{figures/666f6df421144cfc.png}

\begin{items}
\item[（b）]

为解初值问题

\begin{equation}
\ddot{u} - 2\dot{u} + u = e^{t} , \quad u(0) = 0 , \quad \dot{u}(0) = 1 ,\tag{1.33}
\end{equation}

我们先求齐次方程 $ \ddot{u} - 2\dot{u} + u = 0 $ 的一个基本解组。相应的特征多项式为

\[
X^{2} - 2X + 1 = (X - 1)^{2} .
\]

则由定理 1.23(iii)，$ u_{1}(t) := e^{t} $ 和 $ u_{1}(t) := te^{t} $（$ t \in \mathbb{R} $）构成一个基本解组。因为

\[
w(\tau) = u_{1}(\tau)\dot{u}_{2}(\tau) - \dot{u}_{1}(\tau)u_{2}(\tau) = e^{2\tau} \quad \text{对 } \tau \in \mathbb{R} ,
\]

方程

\[
v(t) = \int_{0}^{t}\frac{u_{1}g}{w}\,d\tau\,u_{2}(t) - \int_{0}^{t}\frac{u_{2}g}{w}\,d\tau\,u_{1}(t) = \frac{t^{2}}{2}e^{t} \quad \text{对 } t \in \mathbb{R} ,
\]

是 (1.33) 的一个特解。因此通解为

\[
x(t) = a_{1}e^{t} + a_{2}te^{t} + \frac{t^{2}}{2}e^{t} \quad \text{对 } t \in \mathbb{R} \text{ 且 } a_{1}, a_{2} \in \mathbb{R} .
\]

因为 $ x(0) = a_{1} $ 且 $ \dot{x}(0) = a_{1} + a_{2} $，我们最终求得 (1.33) 的唯一解为

\[
u(t) = (t + t^{2}/2)e^{t} \quad \text{对 } t \in \mathbb{R} .
\]
\item[（c）]

\textbf{谐波振子}（或\textbf{无阻尼振子}）的微分方程为

\[
\ddot{u} + \omega_{0}^{2}u = 0 ,
\]

其中 $ \omega_{0} > 0 $ 是（角）\textbf{频率}。由定理 1.23，通解具有形式

\[
u(t) = a_{1}\cos(\omega_{0}t) + a_{2}\sin(\omega_{0}t) \quad \text{对 } t \in \mathbb{R} ,
\]

其中 $ a_{1}, a_{2} \in \mathbb{R} $。显然，所有解都以 $ 2\pi/\omega_{0} $ 为周期，且它们绕相平面的中心旋转。

\end{items}

\bookfigure[0.74]{figures/23699141e0463b22.png}

\begin{items}
\item[（d）]

\textbf{阻尼振子}的微分方程为

\[
\ddot{u} + 2\alpha\dot{u} + \omega_{0}^{2}u = 0 .
\]

这里，$ \alpha > 0 $ 是\textbf{阻尼常数}，$ \omega_{0} > 0 $ 是无阻尼振子的频率。特征多项式的零点为

\[
\lambda_{1,2} = -\alpha \pm \sqrt{\alpha^{2} - \omega_{0}^{2}} .
\]
\end{items}

(i) （强阻尼：$ \alpha > \omega_{0} $）此时有两个负特征值 $ \lambda_{1} < \lambda_{2} < 0 $，通解为

\[
u(t) = a_{1}e^{\lambda_{1}t} + a_{2}e^{\lambda_{2}t} \quad \text{对 } t \in \mathbb{R} \text{ 且 } a_{1}, a_{2} \in \mathbb{R} .
\]

因此所有解都指数式衰减。在相平面中，原点处有一个“稳定结点”。

\bookfigure[0.64]{figures/ac28931b3866f214.png}

\begin{subitems}
\item[（ii）]

（弱阻尼：$ \alpha < \omega_{0} $）此时特征多项式有两个互为共轭的复特征值 $ \lambda_{1,2} = -\alpha \pm i\omega $，其中 $ \omega := \sqrt{\omega_{0}^{2} - \alpha^{2}} $。因此根据定理 1.22，通解为

\[
u(t) = e^{-\alpha t}\left( a_{1}\cos(\omega t) + a_{2}\sin(\omega t) \right) \quad \text{对 } t \in \mathbb{R} \text{ 且 } a_{1}, a_{2} \in \mathbb{R} .
\]

此时解也指数式阻尼衰减，且相平面的原点是“稳定涡旋”。

\end{subitems}

\bookfigure[0.71]{figures/20e0a4885532e136.png}

\begin{subitems}
\item[（iii）]

（临界阻尼：$ \alpha = \omega_{0} $）这里，$ \lambda = -\alpha $ 是特征多项式的代数重数为 2 的特征值。因此通解为

\[
u(t) = e^{-\alpha t}(a_{1} + a_{2}t) \quad \text{对 } t \in \mathbb{R} \text{ 且 } a_{1}, a_{2} \in \mathbb{R} .
\]

相平面的原点现在称为“稳定虚结点”。\nobreak\hspace{0.35em}■

\end{subitems}

\bookfigure[0.74]{figures/19abb1ccde27fd67.png}
\end{example}

\exerciseshead

\begin{exercises}
\item

设 $ E $ 与 $ F_j $ 为巴拿赫空间，且对 $ j = 1, \ldots, m $ 有 $ A_j \in \operatorname{Hom}(E, F_j) $。又令 $ F := \prod_{j=1}^{m} F_j $。对 $ A := (x \mapsto (A_1x, \ldots, A_mx)) \in \operatorname{Hom}(E, F) $，证明

\[
A \in \mathcal{L}(E, F) \iff (A_j \in \mathcal{L}(E, F_j), \ j = 1, \ldots, m) .
\]
\item

对方阵 $ A := [a_{k}^{j}] \in \mathbb{K}^{m \times m} $，其\textbf{迹}定义为

\[
\operatorname{tr}(A) := \sum_{j=1}^{m} a_{j}^{j} .
\]

证明

\begin{subitems}
\item

映射 $ \operatorname{tr} : \mathbb{K}^{m \times m} \to \mathbb{K} $ 是线性的；

\item

对 $ A, B \in \mathbb{K}^{m \times m} $ 有 $ \operatorname{tr}(AB) = \operatorname{tr}(BA) $。

\end{subitems}

\item

若存在可逆矩阵 $ S \in \mathbb{K}^{m \times m} $（即 $ S $ 属于 $ \mathcal{L}\mathrm{aut}(\mathbb{K}^{m}) $）使得 $ A = SBS^{-1} $，则称两个矩阵 $ A, B \in \mathbb{K}^{m \times m} $ \textbf{相似}。证明：若 $ A $ 与 $ B $ 相似，则 $ \operatorname{tr}(A) = \operatorname{tr}(B) $。

\item

设 $ E $ 为有限维赋范向量空间。证明：对 $ A \in \mathcal{L}(E) $，有

\[
\operatorname{tr}([A]_{\mathcal{E}}) = \operatorname{tr}([A]_{\mathcal{F}}) ,
\]

其中 $ \mathcal{E} $ 与 $ \mathcal{F} $ 是 $ E $ 的基。

由此可得，$ A \in \mathcal{L}(E) $ 的\textbf{迹}通过下式良好定义：

\[
\operatorname{tr}(A) := \operatorname{tr}([A]_{\mathcal{E}}) ,
\]

其中 $ \mathcal{E} $ 是 $ E $ 的任一基。再证明

\begin{subitems}
\item

$ \operatorname{tr} \in \mathcal{L}(\mathcal{L}(E), \mathbb{K}) $；

\item

对 $ A, B \in \mathcal{L}(E) $ 有 $ \operatorname{tr}(AB) = \operatorname{tr}(BA) $。

\end{subitems}

\item

设 $ (E, (\cdot|\cdot)_E) $ 与 $ (F, (\cdot|\cdot)_F) $ 为有限维希尔伯特空间。验证：对每个 $ A \in \mathcal{L}(E, F) $，存在唯一的 $ A^{*} \in \mathcal{L}(F, E) $，称为 $ A $ 的\textbf{伴随算子}（或简称伴随）$ A^{*} $，使得

\[
(Ax|y)_F = (x|A^{*}y)_E \quad \text{对 } x \in E \text{ 且 } y \in F .
\]

证明：映射

\[
\mathcal{L}(E, F) \to \mathcal{L}(F, E) , \quad A \mapsto A^{*}
\]

是共轭线性的，且满足

\begin{subitems}
\item

$ (AB)^{*} = B^{*}A^{*} $，

\item

$ (A^{*})^{*} = A $，且

\item

$ (C^{-1})^{*} = (C^{*})^{-1} $
对 $ A, B \in \mathcal{L}(E, F) $ 及 $ C \in \mathcal{L}\mathrm{is}(E, F) $。对 $ A = [a_{k}^{j}] \in \mathbb{K}^{m \times n} $，证明 $ A^{*} = \left[ \overline{a_{j}^{k}} \right] \in \mathbb{K}^{n \times m} $。此矩阵即矩阵 $ A $ 的\textbf{埃尔米特共轭}或\textbf{埃尔米特伴随}。
若 $ A = A^{*} $（从而 $ E = F $），则称 $ A $ 为\textbf{自伴的}或\textbf{对称的}。证明：对 $ A \in \mathcal{L}(E, F) $，$ A^{*}A $ 与 $ AA^{*} $ 是自伴的。

\end{subitems}

\item

设 $ E := (E, (\cdot|\cdot)) $ 为有限维希尔伯特空间。证明：对 $ A \in \mathcal{L}(E) $，

\begin{subitems}
\item

$ \operatorname{tr}(A^{*}) = \overline{\operatorname{tr}(A)} $；

\item

$ \operatorname{tr}(A) = \sum_{j=1}^{n} (A\varphi_j|\varphi_j) $，其中 $ \{\varphi_1, \ldots, \varphi_n\} $ 是 $ E $ 的规范正交基。

\end{subitems}

\item

设 $ E $ 与 $ F $ 为有限维希尔伯特空间。证明：

\begin{subitems}
\item

$ \mathcal{L}(E, F) $ 在内积

\[
(A, B) \mapsto A : B := \operatorname{tr}(B^{*}A) ;
\]

下是希尔伯特空间；

\item

若 $ E := \mathbb{K}^{n} $ 且 $ F := \mathbb{K}^{m} $，则

\[
|A|_{\mathcal{L}(E, F)} \le \sqrt{A : A} = \|A\| \quad \text{对 } A \in \mathcal{L}(E, F) .
\]
\end{subitems}

\item

设 $ [g_{jk}] \in \mathbb{R}^{n \times n} $ 是对称且\textbf{正定的}，即 $ g_{jk} = g_{kj} $，且存在 $ \gamma > 0 $ 使得

\[
\sum_{j,k=1}^{n} g_{jk}\xi^{j}\xi^{k} \ge \gamma|\xi|^{2} \quad \text{对 } \xi \in \mathbb{R}^{n} .
\]

证明

\[
(x|y)^{g} := \sum_{j,k=1}^{n} g_{jk}x^{j}y^{k} \quad \text{对 } x, y \in \mathbb{R}^{n}
\]

定义了 $ \mathbb{R}^{n} $ 上的一个内积。

\item

设 $ E $ 为巴拿赫空间，且 $ A, B \in \mathcal{L}(E) $ 可交换，即 $ AB = BA $。证明 $ Ae^{B} = e^{B}A $ 与 $ e^{A+B} = e^{A}e^{B} $。

\item

求 $ A, B \in \mathbb{K}^{2 \times 2} $ 使得

\begin{subitems}
\item

$ AB \ne BA $ 且 $ e^{A+B} \ne e^{A}e^{B} $；

\item

$ AB \ne BA $ 且 $ e^{A+B} = e^{A}e^{B} $。
（(ii) 的提示：对 $ k \in \mathbb{Z} $，$ e^{2k\pi i} = 1 $。）

\end{subitems}

\item

设 $ E $ 为巴拿赫空间，$ A \in \mathcal{L}(E) $ 且 $ B \in \mathcal{L}\mathrm{aut}(E) $。证明

\[
e^{BAB^{-1}} = Be^{A}B^{-1} .
\]
\item

设 $ E $ 为有限维希尔伯特空间，$ A \in \mathcal{L}(E) $。证明

\begin{subitems}
\item

$ (e^{A})^{*} = e^{A^{*}} $；

\item

若 $ A $ 自伴，则 $ e^{A} $ 也自伴；

\item

若 $ A $ 是\textbf{反自伴的}或\textbf{斜对称的}，即 $ A^{*} = -A $，则 $ [e^{A}]^{*}e^{A} = 1 $。

\end{subitems}

\item

对具有下列表示矩阵的 $ A \in \mathcal{L}(\mathbb{R}^{4}) $ 计算 $ e^{A} $：

\begin{equation*}
\begin{bmatrix} 1 & 1 & 1 & 1 \\ 1 & 1 & 1 & 1 \\ 1 & 1 & 1 & 1 \\ 1 & 1 & 1 & 1 \end{bmatrix} , \quad \begin{bmatrix} 2 & 1 & 0 & 0 \\ 0 & 2 & 1 & 0 \\ 0 & 0 & 2 & 1 \\ 0 & 0 & 0 & 2 \end{bmatrix} , \quad \begin{bmatrix} 2 & 1 & 0 & 0 \\ 0 & 2 & 1 & 0 \\ 0 & 0 & 2 & 0 \\ 0 & 0 & 0 & 2 \end{bmatrix} , \quad \begin{bmatrix} 2 & 1 & 0 & 0 \\ 0 & 2 & 0 & 0 \\ 0 & 0 & 2 & 0 \\ 0 & 0 & 0 & 2 \end{bmatrix} .
\end{equation*}
\item

求使得下式成立的 $ A \in \mathbb{R}^{3 \times 3} $：

\begin{equation*}
e^{tA} = \begin{bmatrix} \cos t & -\sin t & 0 \\ \sin t & \cos t & 0 \\ 0 & 0 & 1 \end{bmatrix} \quad \text{对 } t \in \mathbb{R} .
\end{equation*}
\item

求下列方程组的通解：

\begin{equation*}
\begin{aligned}
\dot{x} &= 3x - 2y + t , \\ \dot{y} &= 4x - y + t^{2} .
\end{aligned}
\end{equation*}
\item

解初值问题

\begin{equation*}
\begin{aligned}
\dot{x} &= z - y , & x(0) &= 0 , \\ \dot{y} &= z + e^{t} , & y(0) &= 3/2 , \\ \dot{z} &= z - x , & z(0) &= 1 .
\end{aligned}
\end{equation*}
\item

设 $ A = [a_{k}^{j}] \in \mathbb{R}^{n \times n} $ 为\textbf{马尔可夫矩阵}，即对所有 $ j \ne k $ 有 $ a_{k}^{j} \ge 0 $，且对 $ j = 1, \ldots, n $ 有 $ \sum_{k=1}^{n} a_{k}^{j} = 1 $。又设

\[
H_c := \left\{ (x^{1}, \ldots, x^{n}) \in \mathbb{R}^{n} \;;\ \sum_{j=1}^{n} x^{j} = c \right\} \quad \text{对 } c \in \mathbb{R} .
\]

证明：初值在 $ H_c $ 中的 $ \dot{x} = Ax $ 的每个解对一切时刻仍停留在 $ H_c $ 中，即对 $ t \in \mathbb{R} $ 有 $ e^{tA}H_c \subset H_c $。

\item

设 $ b, c \in \mathbb{R} $，且 $ z \in C^{2}(\mathbb{R}, \mathbb{C}) $ 满足 $ \ddot{u} + b\dot{u} + cu = 0 $。证明 $ \operatorname{Re} z $ 与 $ \operatorname{Im} z $ 是 $ \ddot{u} + b\dot{u} + cu = 0 $ 的实解。

\item

设 $ b, c \in \mathbb{R} $，且 $ u $ 满足 $ \ddot{u} + b\dot{u} + cu = g(t) $。证明：

\begin{subitems}
\item

$ k \in \mathbb{N} \cup \{\infty\} $ 且 $ g \in C^{k}(\mathbb{R}, \mathbb{K}) $ 蕴含 $ u \in C^{k+2}(\mathbb{R}, \mathbb{K}) $；

\item

若 $ g \in C^{\omega}(\mathbb{R}) $，则 $ u \in C^{\omega}(\mathbb{R}) $。

\end{subitems}

\item

\textbf{阻尼受迫振子}的微分方程为

\[
\ddot{u} + 2\alpha\dot{u} + \omega_0^{2}u = c\sin(\omega t) \quad \text{对 } t \in \mathbb{R} ,
\]

其中 $ \alpha, \omega_0, \omega > 0 $ 且 $ c \in \mathbb{R} $。求其通解。

\item

在 $ (-\pi/2, \pi/2) $ 上求下列方程的通解：

\begin{subitems}
\item

$ \ddot{u} + u = 1/\cos t $；

\item

$ \ddot{u} + 4u = 2\tan t $。

\end{subitems}

\item

设 $ E $ 为有限维巴拿赫空间，$ A \in \mathcal{L}(E) $。证明下列陈述等价：

\begin{subitems}
\item

$ \dot{x} = Ax $ 在 $ E $ 中的每个解 $ u $ 都满足 $ \lim_{t \to \infty} u(t) = 0 $。

\item

对 $ \lambda \in \sigma(A) $ 有 $ \operatorname{Re}\lambda < 0 $。
在下列习题中，$ E $ 为巴拿赫空间，$ A \in \mathcal{L}(E) $。

\end{subitems}

\item

若 $ \lambda \in \mathbb{K} $ 满足 $ \operatorname{Re}\lambda > \|A\| $，则

\[
(\lambda - A)^{-1} = \int_{0}^{\infty} e^{-t(\lambda - A)} \, dt .
\]

（提示：令 $ R(\lambda) := \int_{0}^{\infty} e^{-t(\lambda - A)} \, dt $，并计算 $ (\lambda - A)R(\lambda) $。）

\item

证明指数映射具有表示

\[
e^{tA} = \lim_{n \to \infty} \left( 1 - \frac{t}{n}A \right)^{-n} .
\]

（提示：定理 III.6.23 的证明。）

\item

设 $ C $ 为 $ E $ 的闭凸子集。证明下列陈述等价：

\begin{subitems}
\item

对所有 $ t \in \mathbb{R} $，$ e^{tA}C \subset C $；

\item

对所有满足 $ \operatorname{Re}\lambda > \|A\| $ 的 $ \lambda \in \mathbb{K} $，$ (\lambda - A)^{-1}C \subset C $。
（提示：习题 23 与 24。）

\end{subitems}

\end{exercises}

\section{可微性}

在本节中，我们解释函数\footnote{记号 $ f : X \subset E \to F $ 是 $ X \subset E $ 与 $ f : X \to F $ 的简写。} $ f : X \subset E \to F $ 的“可微性”“导数”与“方向导数”等核心概念。这里 $ E $ 与 $ F $ 是巴拿赫空间，$ X $ 是 $ E $ 的子集。此后，我们将通过对空间 $ E $ 与 $ F $ 作具体假设来展示这些概念的威力与灵活性。

若 $ E = \mathbb{R}^n $，我们解释偏导数的概念，并研究可微性与偏导数的存在性之间的联系。

在应用中，$ E = \mathbb{R}^n $ 且 $ F = \mathbb{R} $ 的情形尤为重要。这里我们引入梯度的概念，并借助里斯表示定理阐明导数与梯度之间的关系。

在 $ E = F = \mathbb{C} $ 的情形，我们导出重要的柯西–黎曼方程。它们把 $ f = u + iv : \mathbb{C} \to \mathbb{C} $ 的复可微性与 $ (u, v) : \mathbb{R}^2 \to \mathbb{R}^2 $ 的可微性联系起来。

最后，我们证明这些新概念与第 IV.1 节中业已确立的 $ E = \mathbb{K} $ 情形相一致。

在下文中，设

• $ E = (E, \|\cdot\|) $ 与 $ F = (F, \|\cdot\|) $ 是域 $ \mathbb{K} $ 上的巴拿赫空间；$ X $ 是 $ E $ 的开子集。

\needspace{3\baselineskip}\subsection*{定义}

函数 $ f : X \to F $ 在 $ x_0 \in X $ 处\textbf{可微}，如果存在 $ A_{x_0} \in \mathcal{L}(E, F) $ 使得

\begin{equation}
\lim_{x \to x_0} \frac{f(x) - f(x_0) - A_{x_0}(x - x_0)}{\|x - x_0\|} = 0 .\tag{2.1}
\end{equation}

下一个定理重新表述了这个定义，并给出可微函数的最初几条性质。

\begin{proposition}

设 $ f : X \to F $ 且 $ x_0 \in X $。

\begin{subitems}
\item

下列陈述等价：
(α) $ f $ 在 $ x_0 $ 处可微。
(β) 存在 $ A_{x_0} \in \mathcal{L}(E, F) $ 与 $ r_{x_0} : X \to F $，其中 $ r_{x_0} $ 在 $ x_0 $ 处连续且满足 $ r_{x_0}(x_0) = 0 $，使得

\[
f(x) = f(x_0) + A_{x_0}(x - x_0) + r_{x_0}(x) \|x - x_0\| \quad \text{对 } x \in X .
\]

(γ) 存在 $ A_{x_0} \in \mathcal{L}(E, F) $ 使得

\begin{equation}
f(x) = f(x_0) + A_{x_0}(x - x_0) + o(\|x - x_0\|) \quad (x \to x_0) .\tag{2.2}
\end{equation}
\item

若 $ f $ 在 $ x_0 $ 处可微，则 $ f $ 在 $ x_0 $ 处连续。

\item

设 $ f $ 在 $ x_0 $ 处可微。则来自 (2.1) 的线性算子 $ A_{x_0} \in \mathcal{L}(E, F) $ 被唯一确定。

\proofhead
(i) 证明与定理 IV.1.1 的证明相同，留给你来完成。
\end{subitems}

\begin{subitems}
\item[（ii）]

若 $ f $ 在 $ x_0 $ 处可微，则它在 $ x_0 $ 处的连续性直接由 (i$ \beta $) 得出。

\item[（iii）]

设 $ B \in \mathcal{L}(E, F) $，且

\begin{equation}
f(x) = f(x_0) + B(x - x_0) + o(\|x - x_0\|) \quad (x \to x_0) .\tag{2.3}
\end{equation}

从 (2.2) 中减去 (2.3)，并把结果除以 $ \|x - x_0\| $，我们得到

\[
\lim_{x \to x_0} (A_{x_0} - B)\left( \frac{x - x_0}{\|x - x_0\|} \right) = 0 .
\]

设 $ y \in E $ 满足 $ \|y\| = 1 $，且对 $ n \in \mathbb{N}^{\times} $ 令 $ x_n := x_0 + y/n $。则因为 $ \lim x_n = x_0 $ 且 $ (x_n - x_0)/\|x_n - x_0\| = y $，我们得到

\[
(A_{x_0} - B)y = \lim_n (A_{x_0} - B)\left( \frac{x_n - x_0}{\|x_n - x_0\|} \right) = 0 .
\]

由于这对每个 $ y \in \partial\mathbb{B}_E $ 都成立，由此得 $ A_{x_0} = B $。\nobreak\hspace{0.35em}■

\end{subitems}

\end{proposition}

\needspace{3\baselineskip}\subsection*{导数}

设 $ f : X \to F $ 在 $ x_0 \in X $ 处可微。则我们用 $ \partial f(x_0) $ 表示由命题 2.1 唯一确定的线性算子 $ A_{x_0} \in \mathcal{L}(E, F) $。它称为 $ f $ 在 $ x_0 $ 处的\textbf{导数}，也写作

\[
Df(x_0) \quad \text{或} \quad f'(x_0) .
\]

因此 $ \partial f(x_0) \in \mathcal{L}(E, F) $，且

\[
\lim_{x \to x_0} \frac{f(x) - f(x_0) - \partial f(x_0)(x - x_0)}{\|x - x_0\|} = 0 .
\]

若 $ f : X \to F $ 在每一点 $ x \in X $ 都可微，我们称 $ f $ \textbf{可微}，并称映射\footnote{为了清楚地区分在一点 $ x_0 $ 处的导数 $ \partial f(x_0) \in \mathcal{L}(E, F) $ 与导数 $ \partial f : X \to \mathcal{L}(E, F) $，我们也可以说导函数 $ \partial f : X \to \mathcal{L}(E, F) $。}

\[
\partial f : X \to \mathcal{L}(E, F) , \quad x \mapsto \partial f(x)
\]

为 $ f $ 的\textbf{导数}。

当 $ \mathcal{L}(E, F) $ 是巴拿赫空间时，我们可以有意义地谈论导数的连续性。若 $ \partial f $ 连续，即 $ \partial f \in C(X, \mathcal{L}(E, F)) $，我们称 $ f $ \textbf{连续可微}。我们令

\[
C^1(X, F) := \{ f : X \to F ; f \text{ 是连续可微的} \} .
\]

\begin{remark}

\begin{items}
\item

下列陈述等价：

\begin{subitems}
\item

$ f : X \to F $ 在 $ x_0 \in X $ 处可微。

\item

存在 $ \partial f(x_0) \in \mathcal{L}(E, F) $ 使得

\[
f(x) = f(x_0) + \partial f(x_0)(x - x_0) + o(\|x - x_0\|) \quad (x \to x_0) .
\]
\item

存在 $ \partial f(x_0) \in \mathcal{L}(E, F) $ 与 $ r_{x_0} : X \to F $，其中 $ r_{x_0} $ 在 $ x_0 $ 处连续且满足 $ r_{x_0}(x_0) = 0 $，使得

\[
f(x) = f(x_0) + \partial f(x_0)(x - x_0) + r_{x_0}(x) \|x - x_0\| \quad \text{对 } x \in X .
\]
\end{subitems}

\item

由 (a) 得，$ f $ 在 $ x_0 $ 处可微，当且仅当 $ f $ 在 $ x_0 $ 处被仿射映射

\[
g : E \to F , \quad x \mapsto f(x_0) + \partial f(x_0)(x - x_0)
\]

逼近，使得误差比 $ x $ 趋于 $ x_0 $ 更快地趋于零，即

\[
\lim_{x \to x_0} \frac{\|f(x) - g(x)\|}{\|x - x_0\|} = 0 .
\]

因此，映射 $ f $ 在 $ x_0 $ 处可微，当且仅当它在 $ x_0 $ 处近似线性（见推论 IV.1.3）。

\item

“可微性”与“导数”的概念不依赖于 $ E $ 与 $ F $ 中等价范数的选取。

\proofhead
例如，这由注 II.3.13(d) 得出。\nobreak\hspace{0.35em}■
\item

有时我们不说“可微”，而说“全可微”或“弗雷歇可微”。

\item

（$ E = \mathbb{K} $ 的情形）在注 1.3(a) 中我们看到，映射 $ \mathcal{L}(\mathbb{K}, F) \to F $，$ A \mapsto A1 $ 是等距同构，借此我们在典范意义下把 $ \mathcal{L}(\mathbb{K}, F) $ 与 $ F $ 等同起来。这个等同表明，若 $ E = \mathbb{K} $，则可微性与导数的新定义与第 IV.1 节中的定义一致。

\item

$ C^1(X, F) \subset C(X, F) $，即每个连续可微函数也是连续的。\nobreak\hspace{0.35em}■

\end{items}

\end{remark}

\begin{example}

\begin{items}
\item

对 $ A \in \mathcal{L}(E, F) $，函数 $ A = (x \mapsto Ax) $ 是连续可微的，且对 $ x \in E $ 有 $ \partial A(x) = A $。

\proofhead
因为 $ Ax = Ax_0 + A(x - x_0) $，结论由注 2.2(a) 得出。\nobreak\hspace{0.35em}■
\item

设 $ y_0 \in F $。则常值映射 $ k_{y_0} : E \to F $，$ x \mapsto y_0 $ 是连续可微的，且 $ \partial k_{y_0} = 0 $。

\proofhead
这是显然的，因为对 $ x, x_0 \in E $ 有 $ k_{y_0}(x) = k_{y_0}(x_0) $。\nobreak\hspace{0.35em}■
\item

设 $ H $ 是希尔伯特空间，定义 $ b : H \to \mathbb{K} $，$ x \mapsto \|x\|^2 $。则 $ b $ 是连续可微的，且

\[
\partial b(x) = 2 \operatorname{Re}(x|\cdot) \quad \text{对 } x \in H .
\]

\proofhead
对任意满足 $ x \ne x_0 $ 的 $ x, x_0 \in H $，我们有
\[
\|x\|^2 - \|x_0\|^2 = \|x - x_0 + x_0\|^2 - \|x_0\|^2 = \|x - x_0\|^2 + 2 \operatorname{Re}(x_0 | x - x_0)
\]

从而

\[
\frac{b(x) - b(x_0) - 2 \operatorname{Re}(x_0 | x - x_0)}{\|x - x_0\|} = \|x - x_0\| .
\]

这蕴含对 $ h \in H $ 有 $ \partial b(x_0)h = 2 \operatorname{Re}(x_0 | h) $。\nobreak\hspace{0.35em}■

\end{items}

\end{example}

\needspace{3\baselineskip}\subsection*{方向导数}

设 $ f : X \to F $，$ x_0 \in X $ 且 $ v \in E \backslash \{0\} $。因为 $ X $ 是开集，存在 $ \varepsilon > 0 $ 使得对 $ |t| < \varepsilon $ 有 $ x_0 + tv \in X $。因此函数

\[
(-\varepsilon, \varepsilon) \to F , \quad t \mapsto f(x_0 + tv)
\]

是有定义的。当这个函数在点 $ 0 $ 处可微时，我们称它的导数为 $ f $ 在点 $ x_0 $ 处沿方向 $ v $ 的\textbf{方向导数}，并记作 $ D_v f(x_0) $。于是

\[
D_v f(x_0) = \lim_{t \to 0} \frac{f(x_0 + tv) - f(x_0)}{t} .
\]

\begin{remark}

函数

\[
f_{v,x_0} : (-\varepsilon, \varepsilon) \to F , \quad t \mapsto f(x_0 + tv)
\]

可看作 $ E \times F $ 中的一条“曲线”，它位于“$ f $ 的图像上”。则当 $ \|v\| = 1 $ 时，$ D_v f(x_0) $ 表示这条曲线在点 $ (x_0, f(x_0)) $ 处切线的斜率（见注 IV.1.4(a)）。\nobreak\hspace{0.35em}■

\bookfigure[0.37]{figures/e9c2b26c3eb8c187.png}

下一个结果说的是，若 $ f $ 在 $ x_0 $ 处可微，则 $ f $ 在 $ x_0 $ 处沿每个方向都有方向导数。

\end{remark}

\begin{proposition}

设 $ f : X \to F $ 在 $ x_0 \in X $ 处可微。则对每个 $ v \in E \backslash \{0\} $，$ D_v f(x_0) $ 存在，且 $ D_v f(x_0) = \partial f(x_0)v $。

\proofhead
对 $ v \in E \backslash \{0\} $，由注 2.2(a) 得
\[
f(x_0 + tv) = f(x_0) + \partial f(x_0)(tv) + o(\|tv\|) = f(x_0) + t \partial f(x_0)v + o(|t|)
\]

（当 $ t \to 0 $ 时）。定理现在由注 IV.3.1(c) 得出。\nobreak\hspace{0.35em}■

\end{proposition}

\begin{remark}

命题 2.5 的逆命题不成立，也就是说，在每一方向都有方向导数的函数未必可微。

\proofhead
我们考虑由下式定义的函数 $ f : \mathbb{R}^2 \to \mathbb{R} $：
\begin{equation*}
f(x,y) := \begin{cases} \frac{x^2 y}{x^2 + y^2} , & (x,y) \ne (0,0) , \\ 0 , & (x,y) = (0,0) . \end{cases}
\end{equation*}

对每个 $ v = (\xi, \eta) \in \mathbb{R}^2 \backslash \{(0,0)\} $，我们有

\[
f(tv) = \frac{t^3 \xi^2 \eta}{t^2(\xi^2 + \eta^2)} = t f(v) .
\]

于是

\[
D_v f(0) = \lim_{t \to 0} f(tv)/t = f(v) .
\]

若 $ f $ 在 $ 0 $ 处可微，则命题 2.5 将蕴含对每个 $ v \in \mathbb{R}^2 \backslash \{(0,0)\} $ 有 $ \partial f(0)v = D_v f(0) = f(v) $。因为 $ v \mapsto \partial f(0)v $ 是线性的，而 $ f $ 不是，这是不可能的。因此 $ f $ 在 $ 0 $ 处不可微。\nobreak\hspace{0.35em}■

\end{remark}

\needspace{3\baselineskip}\subsection*{偏导数}

若 $ E = \mathbb{R}^n $，则出于实用与历史原因，沿坐标轴方向的方向导数尤为重要。为它们引入专门的记号是方便的。因此我们用 $ \partial_k $ 或 $ \partial/\partial x^k $ 表示沿标准基向量\footnote{偶尔我们也把 $ \partial/\partial x^k $ 写成 $ \partial_{x^k} $。} $ e_k $（$ k = 1, \ldots, n $）方向的方向导数。于是

\[
\partial_k f(x_0) := \frac{\partial f}{\partial x^k}(x_0) := D_{e_k} f(x_0) = \lim_{t \to 0} \frac{f(x_0 + te_k) - f(x_0)}{t} \quad \text{对 } 1 \le k \le n ,
\]

且 $ \partial_k f(x_0) $ 称为 $ f $ 在 $ x_0 $ 处关于 $ x^k $ 的\textbf{偏导数}。函数 $ f $ 称为在 $ x_0 $ 处偏可微的，如果所有 $ \partial_1 f(x_0), \ldots, \partial_n f(x_0) $ 都存在；函数称为\textbf{[连续]偏可微的}，如果它在 $ X $ 的每一点偏可微[如果 $ \partial_k f : X \to F $ 对 $ 1 \le k \le n $ 连续]。

\begin{remark}

\begin{items}
\item

当第 $ k $ 个偏导数在 $ x_0 \in X $ 处存在时，我们有

\[
\partial_k f(x_0) = \lim_{h \to 0} \frac{f(x_0^1, \ldots, x_0^{k-1}, x_0^k + h, x_0^{k+1}, \ldots, x_0^n) - f(x_0)}{h} \quad \text{对 } 1 \le k \le n .
\]

因此 $ f $ 在 $ x_0 $ 处关于第 $ k $ 个坐标偏可微，当且仅当单实变量函数 $ t \mapsto f(x_0^1, \ldots, x_0^{k-1}, t, x_0^{k+1}, \ldots, x_0^n) $ 在 $ x_0^k $ 处可微。

\item

若 $ x_0^k $ 是集合

\[
\{ \xi \in \mathbb{R} ; (x_0^1, \ldots, x_0^{k-1}, \xi, x_0^{k+1}, \ldots, x_0^n) \in X \} .
\]

的聚点，则偏导数 $ \partial_k f(x_0) $ 可以定义。特别地，$ X $ 必须不是开集。

\item

若 $ f $ 在 $ x_0 $ 处可微，则 $ f $ 在 $ x_0 $ 处偏可微。

\proofhead
这是命题 2.5 的特殊情形。\nobreak\hspace{0.35em}■
\item

若 $ f $ 在 $ x_0 $ 处偏可微，并不能推出 $ f $ 在 $ x_0 $ 处可微。

\proofhead
见注 2.6。\nobreak\hspace{0.35em}■
\item

$ f $ 在 $ x_0 $ 处偏可微并不意味着它在该点连续。

\proofhead
我们考虑 $ f : \mathbb{R}^2 \to \mathbb{R} $，其中
\begin{equation*}
f(x,y) := \begin{cases} \frac{xy}{(x^2 + y^2)^2} , & (x,y) \ne (0,0) , \\ 0 , & (x,y) = (0,0) . \end{cases}
\end{equation*}

则对一切 $ h \in \mathbb{R} $ 有 $ f(h,0) = f(0,h) = 0 $。因此 $ \partial_1 f(0,0) = \partial_2 f(0,0) = 0 $。于是 $ f $ 在 $ (0,0) $ 处偏可微。
因为对 $ n \in \mathbb{N}^{\times} $ 有 $ f(0,0) = 0 $ 与 $ f(1/n, 1/n) = n^2/4 $，我们看到 $ f $ 在 $ (0,0) $ 处不连续。\nobreak\hspace{0.35em}■

\item

为了清楚地区分“可微性”与“偏可微性”，我们有时也说 $ f $ \textbf{全可微}[在 $ x_0 $ 处]，如果 $ f $ 可微[在 $ x_0 $ 处]。\nobreak\hspace{0.35em}■

\end{items}

接下来，我们寻求当 $ E = \mathbb{R}^n $ 或 $ F = \mathbb{R}^m $ 时 $ \partial f(x_0) $ 的具体表示。这些表示将使我们能够显式地计算偏导数。

\end{remark}

\begin{theorem}

\begin{subitems}
\item

设 $ E = \mathbb{R}^n $ 且 $ f : X \to F $ 在 $ x_0 $ 处可微。则

\[
\partial f(x_0)h = \sum_{k=1}^{n} \partial_k f(x_0) h^k \quad \text{对 } h = (h^1, \ldots, h^n) \in \mathbb{R}^n .
\]
\item

设 $ E $ 是巴拿赫空间且 $ f = (f^1, \ldots, f^m) : X \to \mathbb{K}^m $。则 $ f $ 在 $ x_0 $ 处可微，当且仅当所有坐标函数 $ f^j $（$ 1 \le j \le m $）都在 $ x_0 $ 处可微。此时

\[
\partial f(x_0) = (\partial f^1(x_0), \ldots, \partial f^m(x_0)) ,
\]

即向量被逐分量地求导。

\proofhead
(i) 因为对 $ h = (h^1, \ldots, h^n) \in \mathbb{R}^n $ 有 $ h = \sum_k h^k e_k $，由 $ \partial f(x_0) $ 的线性性与命题 2.5 得
\[
\partial f(x_0)h = \sum_{k=1}^{n} h^k \partial f(x_0)e_k = \sum_{k=1}^{n} \partial_k f(x_0) h^k .
\]
\end{subitems}

(ii) 对 $ A = (A^1, \ldots, A^m) \in \operatorname{Hom}(E, \mathbb{K}^m) $，我们有

\[
A \in \mathcal{L}(E, \mathbb{K}^m) \iff A^j \in \mathcal{L}(E, \mathbb{K}) \quad \text{对 } 1 \le j \le m ,
\]

（见习题 1.1）。现在由命题 II.3.14 与命题 2.1(iii) 得

\[
\lim_{x \to x_0} \frac{f(x) - f(x_0) - \partial f(x_0)(x - x_0)}{\|x - x_0\|} = 0
\]

等价于

\[
\lim_{x \to x_0} \frac{f^j(x) - f^j(x_0) - \partial f^j(x_0)(x - x_0)}{\|x - x_0\|} = 0 \quad \text{对 } 1 \le j \le m ,
\]

正如定理所要求。\nobreak\hspace{0.35em}■

\end{theorem}

\needspace{3\baselineskip}\subsection*{雅可比矩阵}

设 $ X $ 是 $ \mathbb{R}^n $ 中的开集，$ f = (f^1, \ldots, f^m) : X \to \mathbb{R}^m $ 在 $ x_0 $ 处偏可微。则我们称

\begin{equation*}
[\partial_k f^j(x_0)] = \begin{bmatrix} \partial_1 f^1(x_0) & \cdots & \partial_n f^1(x_0) \\ \vdots & & \vdots \\ \partial_1 f^m(x_0) & \cdots & \partial_n f^m(x_0) \end{bmatrix}
\end{equation*}

为 $ f $ 在 $ x_0 $ 处的\textbf{雅可比矩阵}。

\begin{corollary}

若 $ f $ 在 $ x_0 $ 处可微，则每个坐标函数 $ f^j $ 都在 $ x_0 $ 处偏可微，且

\begin{equation*}
[\partial f(x_0)] = [\partial_k f^j(x_0)] = \begin{bmatrix} \partial_1 f^1(x_0) & \cdots & \partial_n f^1(x_0) \\ \vdots & & \vdots \\ \partial_1 f^m(x_0) & \cdots & \partial_n f^m(x_0) \end{bmatrix} ,
\end{equation*}

即 $ f $ 的导数的表示矩阵（在标准基下）就是 $ f $ 的雅可比矩阵。

\proofhead
对 $ k = 1, \ldots, n $，我们有 $ \partial f(x_0)e_k = \sum_{j=1}^{m} a_k^j e_j $，其中 $ a_k^j \in \mathbb{R} $ 唯一确定。由 $ \partial f(x_0) $ 的线性性与命题 2.8 得
\[
\partial f(x_0)e_k = (\partial f^1(x_0)e_k, \ldots, \partial f^m(x_0)e_k) = (\partial_k f^1(x_0), \ldots, \partial_k f^m(x_0)) = \sum_{j=1}^{m} \partial_k f^j(x_0) e_j .
\]

因此 $ a_k^j = \partial_k f^j $。\nobreak\hspace{0.35em}■

\end{corollary}

\needspace{3\baselineskip}\subsection*{一个可微性判别准则}

注 2.6 表明，一个函数所有偏导数的存在性并不蕴含其全可微性。为保证可微性，需要额外的假设。其中一个假设就是偏导数连续。正如下一个定理——基本可微性判别准则——所示，这意味着该映射实际上是连续可微的。

\begin{theorem}

设 $ X $ 是 $ \mathbb{R}^n $ 中的开集，$ F $ 是巴拿赫空间。则 $ f : X \to F $ 连续可微，当且仅当 $ f $ 有连续的偏导数。

\proofhead
“$ \Rightarrow $” 由命题 2.5 很容易得出。

“$ \Leftarrow $” 设 $ x \in X $。我们通过下式定义线性映射 $ A(x) : \mathbb{R}^n \to F $：
\[
h := (h^1, \ldots, h^n) \mapsto A(x)h := \sum_{k=1}^{n} \partial_k f(x) h^k .
\]

定理 1.6 表明 $ A(x) $ 属于 $ \mathcal{L}(\mathbb{R}^n, F) $。我们的目标是证明 $ \partial f(x) = A(x) $，即

\[
\lim_{h \to 0} \frac{f(x + h) - f(x) - A(x)h}{|h|} = 0 .
\]

我们选取 $ \varepsilon > 0 $ 使得 $ \mathbb{B}(x, \varepsilon) \subset X $，并对 $ 1 \le k \le n $ 与 $ h = (h^1, \ldots, h^n) \in \mathbb{B}(x, \varepsilon) $ 令 $ x_0 := x $ 与 $ x_k := x_0 + \sum_{j=1}^{k} h^j e_j $。则我们得到

\[
f(x + h) - f(x) = \sum_{k=1}^{n} (f(x_k) - f(x_{k-1})) ,
\]

而微积分基本定理蕴含

\[
f(x + h) - f(x) = \sum_{k=1}^{n} h^k \int_0^1 \partial_k f(x_{k-1} + th^k e_k) \, dt .
\]

由此我们得到表示

\[
f(x + h) - f(x) - A(x)h = \sum_{k=1}^{n} h^k \int_0^1 (\partial_k f(x_{k-1} + th^k e_k) - \partial_k f(x)) \, dt ,
\]

它给出估计

\[
\|f(x + h) - f(x) - A(x)h\|_F \le |h|_\infty \sum_{k=1}^{n} \sup_{|y - x|_\infty \le |h|_\infty} \|\partial_k f(y) - \partial_k f(x)\|_F .
\]

$ \partial_k f $ 的连续性蕴含

\[
\lim_{h \to 0} \left( \sum_{k=1}^{n} \sup_{|y - x|_\infty \le |h|_\infty} \|\partial_k f(y) - \partial_k f(x)\|_F \right) = 0 .
\]

因此

\[
f(x + h) - f(x) - A(x)h = o(|h|_\infty) \quad (h \to 0) .
\]

于是，由注 2.2(a) 与 2.2(c)，$ f $ 在 $ x $ 处可微，且 $ \partial f(x) = A(x) $。利用

\[
\|(\partial f(x) - \partial f(y))h\|_F \le \sum_{k=1}^{n} \|\partial_k f(x) - \partial_k f(y)\|_F |h^k| \le \left( \sum_{k=1}^{n} \|\partial_k f(x) - \partial_k f(y)\|_F \right) |h|_\infty
\]

以及 $ \mathbb{R}^n $ 的范数的等价性，$ \partial f $ 的连续性由 $ \partial_k f $（$ 1 \le k \le n $）的连续性推出。\nobreak\hspace{0.35em}■

\end{theorem}

\begin{corollary}

设 $ X $ 是 $ \mathbb{R}^n $ 中的开集。则 $ f : X \to \mathbb{R}^m $ 连续可微，当且仅当每个坐标函数 $ f^j : X \to \mathbb{R} $ 都有连续的偏导数，此时

\[
[\partial f(x)] = [\partial_k f^j(x)] \in \mathbb{R}^{m \times n} .
\]
\end{corollary}

\begin{example}

函数

\[
f : \mathbb{R}^3 \to \mathbb{R}^2 , \quad (x,y,z) \mapsto (e^x \cos y, \sin(xz))
\]

是连续可微的，且

\begin{equation*}
[\partial f(x,y,z)] = \begin{bmatrix} e^x \cos y & -e^x \sin y & 0 \\ z \cos(xz) & 0 & x \cos(xz) \end{bmatrix}
\end{equation*}

对 $ (x,y,z) \in \mathbb{R}^3 $。\nobreak\hspace{0.35em}■

多实变量的实值函数的导数有一个重要的几何解释，我们接下来为之做准备。

\end{example}

\needspace{3\baselineskip}\subsection*{里斯表示定理}

设 $ E $ 是 $ \mathbb{K} $ 上的赋范向量空间。我们称 $ E' := \mathcal{L}(E, \mathbb{K}) $ 为 $ E $ 的（连续）\textbf{对偶空间}。$ E' $ 的元素是 $ E $ 上的\textbf{连续线性形式}。

\begin{remark}

\begin{items}
\item

按约定，我们给 $ E' $ 以 $ \mathcal{L}(E, \mathbb{K}) $ 上的算子范数，即

\[
\|f\| := \sup_{\|x\| \le 1} |f(x)| \quad \text{对 } f \in E' .
\]

则推论 1.2 保证 $ E' := (E', \|\cdot\|) $ 是巴拿赫空间。我们也说 $ E' $ 的范数是 $ E $ 的范数的\textbf{对偶}（或简称为对偶范数）。

\item

设 $ (E, (\cdot|\cdot)) $ 是内积空间。对 $ y \in E $，令

\[
f_y(x) := (x|y) \quad \text{对 } x \in E .
\]

则 $ f_y $ 是 $ E $ 上的连续线性形式，且 $ \|f_y\|_{E'} = \|y\|_E $。

\proofhead
因为内积关于第一个变量是线性的，我们有 $ f_y \in \operatorname{Hom}(E, \mathbb{K}) $。因为 $ f_0 = 0 $，只需考虑 $ y \in E \backslash \{0\} $ 的情形。
由柯西–施瓦茨不等式，我们得到
\[
|f_y(x)| = |(x|y)| \le \|x\| \, \|y\| \quad \text{对 } x \in E .
\]

因此 $ f_y \in \mathcal{L}(E, \mathbb{K}) = E' $ 且 $ \|f_y\| \le \|y\| $。由 $ f_y(y/\|y\|) = (y/\|y\| | y) = \|y\| $ 得

\[
\|f_y\| = \sup_{\|x\| \le 1} |f_y(x)| \ge |f_y(y/\|y\|)| = \|y\| .
\]

总起来，我们得到 $ \|f_y\| = \|y\| $。\nobreak\hspace{0.35em}■

\item

设 $ (E, (\cdot|\cdot)) $ 是内积空间。则映射

\[
T : E \to E' , \quad y \mapsto f_y := (\cdot|y)
\]

是共轭线性的，并且是等距。

\proofhead
这由 $ T(y + \lambda z) = (\cdot | y + \lambda z) = (\cdot | y) + \bar{\lambda}(\cdot | z) = Ty + \bar{\lambda}Tz $ 以及 (b) 得出。\nobreak\hspace{0.35em}■
\end{items}

根据注 2.13(b)，每个 $ y \in E $ 都对应着 $ E $ 上的一个连续线性形式 $ f_y $。下一个定理表明，对于有限维希尔伯特空间，不存在其他连续线性形式。\footnote{由定理 1.6 我们知道，有限维希尔伯特空间上的每个线性形式都是连续的。此外，推论 1.5 保证每个有限维内积空间都是希尔伯特空间。}换句话说，每个 $ f \in E' $ 都确定唯一的 $ y \in E $ 使得 $ f = f_y $。

\end{remark}

\begin{theorem}

（里斯表示定理）设 $ (E, (\cdot|\cdot)) $ 是有限维希尔伯特空间。则对每个 $ f \in E' $，存在唯一的 $ y \in E $ 使得对所有 $ x \in E $ 有 $ f(x) = (x|y) $。因此

\[
T : E \to E' , \quad y \mapsto (\cdot|y)
\]

是双射的共轭线性等距。

\proofhead
由注 2.13(c)，我们知道 $ T : E \to E' $ 是共轭线性等距。特别地，$ T $ 是单射。令 $ n := \dim(E) $，由定理 1.9 与命题 1.8(ii) 得
\[
E' = \mathcal{L}(E, \mathbb{K}) \cong \mathbb{K}^{1 \times n} \cong \mathbb{K}^n ,
\]

从而 $ \dim(E') = \dim(E) $。线性代数的\textbf{秩公式}，即

\begin{equation}
\dim(\ker(T)) + \dim(\operatorname{im}(T)) = \dim E\tag{2.4}
\end{equation}

（例如见 [Gab96, ????? D.5.4]），现在蕴含 $ T $ 的满射性。\nobreak\hspace{0.35em}■

\end{theorem}

\begin{remark}

\begin{items}
\item

采用定理 2.14 的记号，令

\[
[\cdot|\cdot] : E' \times E' \to \mathbb{K} , \quad [f|g] := (T^{-1}g | T^{-1}f) .
\]

则 $ (E', [\cdot|\cdot]) $ 是希尔伯特空间。

\proofhead
我们把简单的验证留给你。\nobreak\hspace{0.35em}■
\item

设 $ (E, (\cdot|\cdot)) $ 是实有限维希尔伯特空间。则映射 $ E \to E' $，$ y \mapsto (\cdot|y) $ 是等距同构。特别地，$ (\mathbb{R}^m)' $ 通过\textbf{典范同构}等距同构于 $ \mathbb{R}^m $，其中 $ (\cdot|\cdot) $ 表示 $ \mathbb{R}^m $ 上的欧几里得内积。人们经常利用这个同构等同空间 $ \mathbb{R}^m $ 与 $ (\mathbb{R}^m)' $。

\item

里斯表示定理对无限维希尔伯特空间也成立。对于完全一般性的证明，我们建议你查阅泛函分析的文献，或鼓励你选修这门课程。\nobreak\hspace{0.35em}■

\end{items}

\end{remark}

\needspace{3\baselineskip}\subsection*{梯度}

设 $ X $ 是 $ \mathbb{R}^n $ 中的开集，$ f : X \to \mathbb{R} $ 在 $ x_0 \in X $ 处可微。则我们也把导数 $ \partial f(x_0) $ 称为 $ f $ 在 $ x_0 $ 处的\textbf{微分}，并把它写成\footnote{在第 VIII.3 节中，我们将阐明这个记号与注 VI.5.2 中关于诱导微分的同一记号之间的关系。} $ df(x_0) $。因此 $ f $ 在 $ x_0 $ 处的微分是 $ \mathbb{R}^n $ 上的一个（连续）线性形式。利用里斯表示定理，存在唯一的 $ y \in \mathbb{R}^n $ 使得

\[
df(x_0)h = (h|y) = (y|h) \quad \text{对 } h \in \mathbb{R}^n .
\]

这个由 $ f $ 与 $ x_0 $ 唯一确定的 $ y \in \mathbb{R}^n $ 称为 $ f $ 在 $ x_0 $ 处的\textbf{梯度}。我们把它记作 $ \nabla f(x_0) $ 或 $ \operatorname{grad} f(x_0) $。因此 $ f $ 在 $ x_0 $ 处的微分与梯度由基本关系式联系起来

\[
df(x_0)h = (\nabla f(x_0)|h) \quad \text{对 } h \in \mathbb{R}^n .
\]

注意，微分 $ df(x_0) $ 是 $ \mathbb{R}^n $ 上的线性形式，而梯度 $ \nabla f(x_0) $ 是 $ \mathbb{R}^n $ 中的向量。

\begin{proposition}

我们有

\[
\nabla f(x_0) = (\partial_1 f(x_0), \ldots, \partial_n f(x_0)) \in \mathbb{R}^n .
\]

\proofhead
由命题 2.5，对 $ k = 1, \ldots, n $ 有 $ df(x_0)e_k = \partial_k f(x_0) $。现在由此得
\[
(\nabla f(x_0)|h) = df(x_0)h = df(x_0) \sum_{k=1}^{n} h^k e_k = \sum_{k=1}^{n} \partial_k f(x_0) h^k = (y|h)
\]

对 $ h = (h^1, \ldots, h^n) \in \mathbb{R}^n $，其中 $ y := (\partial_1 f(x_0), \ldots, \partial_n f(x_0)) \in \mathbb{R}^n $。因为这对每个 $ h \in \mathbb{R}^n $ 都成立，故定理得证。\nobreak\hspace{0.35em}■

\end{proposition}

\begin{remark}

\begin{items}
\item

若 $ \nabla f(x_0) = 0 $，我们称点 $ x_0 $ 为 $ f $ 的\textbf{临界点}。

\item

设 $ x_0 $ 不是 $ f $ 的临界点。令 $ h_0 := \nabla f(x_0)/|\nabla f(x_0)| $。注 2.13(b) 的证明则表明

\begin{equation*}
\begin{aligned}
df(x_0)h_0 = |\nabla f(x_0)| = \max_{\substack{h \in \mathbb{R}^n \\ |h| \le 1}} df(x_0)h .
\end{aligned}
\end{equation*}

因为 $ df(x_0)h $ 是 $ f $ 在 $ x_0 $ 处沿方向 $ h $ 的方向导数，我们有

\[
f(x_0 + th_0) = f(x_0) + t\,df(x_0)h_0 + o(t) = f(x_0) + t \max_{|h| = 1} df(x_0)h + o(t)
\]

（当 $ t \to 0 $ 时），即向量 $ \nabla f(x_0) $ 指向 $ f $ 有最大方向导数 $ f $ 的方向，即 $ f $ 的\textbf{最速上升}方向。
若 $ n = 2 $，则 $ f $ 的图像

\[
M := \{ (x, f(x)) ; x \in X \} \subset \mathbb{R}^2 \times \mathbb{R} \cong \mathbb{R}^3 ,
\]

可以解释为 $ \mathbb{R}^3 $ 中的一个曲面。\footnote{见推论 8.9。}我们现在按如下方式设想在 $ M $ 中的一条路径上取一点 $ m $：我们从一点 $ (x_0, f(x_0)) \in M $ 出发，其中 $ x_0 $ 不是 $ f $ 的临界点。然后移动 $ m $，使其“速度向量”的投影始终沿 $ f $ 的梯度方向。$ m $ 以这种方式移动所走的路径称为\textbf{最速上升曲线}。如果我们让 $ m $ 沿梯度的反方向移动，它的路径就称为最速下降曲线。

\end{items}

\bookfigure[0.25]{figures/9a8e3373f87dc5ba.png}
\bookfigure[0.26]{figures/7ccba2062c98265c.png}

\begin{items}
\item[（c）]

上面的注表明，$ \nabla f(x_0) $ 获得了不依赖于坐标或内积的特殊选取的几何解释。然而，$ \nabla f(x_0) $ 的表示确实依赖于内积的选取。特别地，命题 2.16 的表示只在 $ \mathbb{R}^n $ 被赋予欧几里得内积时才成立。

\proofhead
设 $ [g_{jk}] \in \mathbb{R}^{n \times n} $ 是对称正定矩阵，即对 $ 1 \le j, k \le n $ 有 $ g_{jk} = g_{kj} $，且存在 $ \gamma > 0 $ 使得
\[
\sum_{j,k=1}^{n} g_{jk} \xi^j \xi^k \ge \gamma |\xi|^2 \quad \text{对 } \xi \in \mathbb{R}^n .
\]

则

\[
(x|y)^g := \sum_{j,k=1}^{n} g_{jk} x^j y^k \quad \text{对 } x, y \in \mathbb{R}^n
\]

在 $ \mathbb{R}^n $ 上定义了一个内积（见习题 1.8）。\footnote{其逆命题也成立：$ \mathbb{R}^n $ 上的每个内积都具有 $ (\cdot|\cdot)^g $ 的形式（例如见 [Art93, Section VII.1]）。}于是定理 2.14 表明，存在唯一的 $ y \in \mathbb{R}^n $ 使得对 $ h \in \mathbb{R}^n $ 有 $ df(x_0)h = (y|h)^g $。我们称 $ \nabla^g f(x_0) := y $ 为 $ f $ 在 $ x_0 $ 处关于内积 $ (\cdot|\cdot)^g $ 的梯度。为确定它的分量，我们考虑

\[
\sum_{j=1}^{n} \partial_j f(x_0) h^j = df(x_0)h = (y|h)^g = \sum_{j,k=1}^{n} g_{jk} y^k h^j \quad \text{对 } h \in \mathbb{R}^n .
\]

因此，

\begin{equation}
\sum_{k=1}^{n} g_{jk} y^k = \partial_j f(x_0) \quad \text{对 } j = 1, \ldots, n .\tag{2.5}
\end{equation}

由假设，矩阵 $ g $ 可逆，且其逆也是对称正定的。\footnote{见习题 5.18。}我们用 $ g^{jk} $ 表示逆矩阵 $ [g_{jk}] $ 的元素，即 $ [g^{jk}] = [g_{jk}]^{-1} $。于是由 (2.5) 得

\[
y^k = \sum_{j=1}^{n} g^{kj} \partial_j f(x_0) \quad \text{对 } k = 1, \ldots, n .
\]

这意味着

\begin{equation}
\nabla^g f(x_0) = \left( \sum_{k=1}^{n} g^{1k} \partial_k f(x_0), \ldots, \sum_{k=1}^{n} g^{nk} \partial_k f(x_0) \right)\tag{2.6}
\end{equation}

是 $ f $ 关于由 $ g = [g_{jk}] $ 诱导的内积的梯度。\nobreak\hspace{0.35em}■

\end{items}

\end{remark}

\needspace{3\baselineskip}\subsection*{复可微性}

借助等同 $ \mathcal{L}(\mathbb{C}) = \mathbb{C} $ 与 $ \mathbb{C}^{1 \times 1} = \mathbb{C} $，我们可以把 $ A \in \mathcal{L}(\mathbb{C}) $ 与其矩阵表示等同起来。因此，对 $ z \in \mathbb{C} $ 与某个 $ a \in \mathbb{C} $ 有 $ Az = a \cdot z $。照例，我们把 $ \mathbb{C} = \mathbb{R} + i\mathbb{R} $ 与 $ \mathbb{R}^2 $ 等同起来：

\[
\mathbb{C} = \mathbb{R} + i\mathbb{R} \ni z = x + iy \leftrightarrow (x,y) \in \mathbb{R}^2 , \quad \text{其中 } x = \operatorname{Re} z \text{ 且 } y = \operatorname{Im} z .
\]

至于 $ A $ 在这个等同下的作用，通过等同 $ a = \alpha + i\beta \leftrightarrow (\alpha, \beta) $，我们得到

\[
Az = a \cdot z = (\alpha + i\beta)(x + iy) = (\alpha x - \beta y) + i(\beta x + \alpha y) ,
\]

因此

\begin{equation}
Az \leftrightarrow \begin{bmatrix} \alpha & -\beta \\ \beta & \alpha \end{bmatrix} \begin{bmatrix} x \\ y \end{bmatrix} .\tag{2.7}
\end{equation}

设 $ X $ 是 $ \mathbb{C} $ 中的开集。对 $ f : X \to \mathbb{C} $，令 $ u := \operatorname{Re} f $ 与 $ v := \operatorname{Im} f $。则

\[
\mathbb{C}^X \ni f = u + iv \leftrightarrow (u,v) =: F \in (\mathbb{R}^2)^X ,
\]

其中在最后一项中，$ X $ 被理解为 $ \mathbb{R}^2 $ 的子集。在这个记号下，下面的基本定理给出了复\footnote{回忆，$ f $ 在一点 $ z_0 $ 处复可微，当且仅当 $ h $ 在 $ \mathbb{C} $ 中趋于 $ 0 $ 时 $ (f(z_0 + h) - f(z_0))/h $ 在 $ \mathbb{C} $ 中有极限。}可微性与全可微性之间的联系。

\qyxffntextdef{1}{对于两个或三个实变量的函数 $ f $，常用的记号是 $ f_x := \partial_1 f $、$ f_y := \partial_2 f $ 与 $ f_z := \partial_3 f $。}
\begin{theorem}

函数 $ f $ 在 $ z_0 = x_0 + iy_0 $ 处复可微，当且仅当 $ F := (u,v) $ 在 $ (x_0, y_0) $ 处全可微，且满足\textbf{柯西–黎曼方程}\qyxffnmark{1}

\[
u_x = v_y , \quad u_y = -v_x ,
\]

在 $ (x_0, y_0) $ 处。此时

\[
f'(z_0) = u_x(x_0, y_0) + i v_x(x_0, y_0) .
\]

证明

\begin{subitems}
\item

设 $ f $ 在 $ z_0 $ 处复可微。令

\begin{equation*}
A := \begin{bmatrix} \alpha & -\beta \\ \beta & \alpha \end{bmatrix} ,
\end{equation*}

其中 $ \alpha := \operatorname{Re} f'(z_0) $ 与 $ \beta := \operatorname{Im} f'(z_0) $。则对 $ h = \xi + i\eta \leftrightarrow (\xi, \eta) $，我们有

\begin{equation*}
\begin{aligned}
\lim_{(\xi,\eta) \to (0,0)} \frac{|F(x_0 + \xi, y_0 + \eta) - F(x_0, y_0) - A(\xi, \eta)|}{|(\xi, \eta)|} &= \lim_{h \to 0} \left| \frac{f(z_0 + h) - f(z_0) - f'(z_0)h}{h} \right| \\ &= 0 .
\end{aligned}
\end{equation*}

因此 $ F $ 在 $ (x_0, y_0) $ 处全可微（见注 II.2.1(a)），且

\begin{equation*}
[\partial F(x_0, y_0)] = \begin{bmatrix} \partial_1 u(x_0, y_0) & \partial_2 u(x_0, y_0) \\ \partial_1 v(x_0, y_0) & \partial_2 v(x_0, y_0) \end{bmatrix} = \begin{bmatrix} \alpha & -\beta \\ \beta & \alpha \end{bmatrix} .
\end{equation*}

于是我们得到柯西–黎曼方程

\begin{equation}
\partial_1 u(x_0, y_0) = \partial_2 v(x_0, y_0) , \quad \partial_2 u(x_0, y_0) = -\partial_1 v(x_0, y_0) .\tag{2.8}
\end{equation}
\item

若 $ F $ 在 $ (x_0, y_0) $ 处全可微且 (2.8) 成立，我们可以令

\[
a := \partial_1 u(x_0, y_0) + i \partial_1 v(x_0, y_0) .
\]

则由于 (2.7)，我们有

\begin{equation*}
\begin{aligned}
\lim_{h \to 0} \left| \frac{f(z_0 + h) - f(z_0) - ah}{h} \right| &= \lim_{(\xi,\eta) \to (0,0)} \frac{|F(x_0 + \xi, y_0 + \eta) - F(x_0, y_0) - \partial F(x_0, y_0)(\xi, \eta)|}{|(\xi, \eta)|} \\ &= 0 .
\end{aligned}
\end{equation*}

从而 $ f $ 在 $ z_0 $ 处复可微，且 $ f'(z_0) = a $。\nobreak\hspace{0.35em}■

\end{subitems}

\end{theorem}

\begin{example}

\begin{items}
\item

根据例 IV.1.13(b)，$ f : \mathbb{C} \to \mathbb{C} $，$ z \mapsto z^2 $ 处处复可微，且 $ f'(z) = 2z $。因为

\[
f(x + iy) = (x + iy)^2 = x^2 - y^2 + i(2xy) \leftrightarrow (u(x,y), v(x,y)) ,
\]

我们有

\[
u_x = v_y = 2x \quad \text{且} \quad u_y = -v_x = -2y .
\]

因此柯西–黎曼方程被满足（正如它们必须被满足），且 $ f'(z) = f'(x + iy) = 2x + i2y = 2(x + iy) = 2z $。

\item

映射 $ f : \mathbb{C} \to \mathbb{C} $，$ z \mapsto \bar{z} $ 处处复不可微，\footnote{见习题 IV.1.4。}因为，其中

\[
f(x + iy) = x - iy \leftrightarrow (u(x,y), v(x,y)) ,
\]

柯西–黎曼方程从不满足：

\[
u_x = 1 , \quad u_y = 0 , \quad v_x = 0 , \quad v_y = -1 .
\]

上面的映射给出了一个单复变量的连续但处处不可微的复值函数的一个非常简单的例子。不过要注意，“同一个”函数

\[
F = (u,v) : \mathbb{R}^2 \to \mathbb{R}^2 , \quad (x,y) \mapsto (x,-y)
\]

处处全可微，且有常值导数

\begin{equation*}
[\partial F(x_0, y_0)] = \begin{bmatrix} 1 & 0 \\ 0 & -1 \end{bmatrix}
\end{equation*}

对 $ (x_0, y_0) \in \mathbb{R}^2 $。因此 $ F $ 实际上是连续可微的。\nobreak\hspace{0.35em}■

\end{items}

\end{example}

\exerciseshead

\begin{exercises}
\item

对 $ f : (0,\infty)^{2} \to \mathbb{R} $，其中

\[
f(x,y) := \left(\left(\left(x^{x}\right)^{x}\right)^{x}\right)^{y} + \log(x)\arctan\left(\arctan\left(\arctan\left(\sin\left(\cos(xy) - \log(x + y)\right)\right)\right)\right) .
\]

计算 $ \partial_{2} f(1,y) $。

\item

设 $ f : \mathbb{R}^{2} \to \mathbb{R} $ 定义为

\begin{equation*}
f(x,y) := \begin{cases} \sqrt{x^{2} + y^{2}} , & y > 0 , \\ x , & y = 0 , \\ -\sqrt{x^{2} + y^{2}} , & y < 0 . \end{cases}
\end{equation*}

证明

\begin{items}
\item

$ f $ 在 $ (0,0) $ 处不可微；

\item

$ f $ 在 $ (0,0) $ 处的每个方向导数都存在。

\end{items}

\item

映射

\begin{equation*}
\mathbb{R}^{2} \to \mathbb{R} , \quad (x,y) \mapsto \begin{cases} x^{3}/\sqrt{x^{2} + y^{2}} , & (x,y) \ne (0,0) , \\ 0 , & (x,y) = (0,0) , \end{cases}
\end{equation*}

在哪些点可微？

\item

设 $ f : \mathbb{R}^{2} \to \mathbb{R} $ 定义为

\begin{equation*}
f(x,y) := \begin{cases} xy/\sqrt{x^{2} + y^{2}} , & (x,y) \ne (0,0) , \\ 0 , & (x,y) = (0,0) . \end{cases}
\end{equation*}

证明：

\begin{items}
\item

$ \partial_{1} f(0,0) $ 与 $ \partial_{2} f(0,0) $ 存在。

\item

对 $ v \in \mathbb{R}^{2}\setminus\{e_{1},e_{2}\} $，$ D_{v} f(0,0) $ 不存在。

\item

$ f $ 在 $ (0,0) $ 处不可微。

\end{items}

\item

计算下列映射的雅可比矩阵：

\begin{equation*}
\begin{aligned}
\mathbb{R}^{3} \to \mathbb{R} , \quad (x,y,z) &\mapsto 3x^{2}y + e^{xz^{2}} + 4z^{3} ; \\ \mathbb{R}^{2} \to \mathbb{R}^{3} , \quad (x,y) &\mapsto \left(xy, \cosh(xy), \log(1 + x^{2})\right) ; \\ \mathbb{R}^{3} \to \mathbb{R}^{3} , \quad (x,y,z) &\mapsto \left(\log(1 + x^{2} + z^{2}), z^{2} + y^{2} - x^{2}, \sin(xz)\right) ; \\ \mathbb{R}^{3} \to \mathbb{R}^{3} , \quad (x,y,z) &\mapsto \left(x\sin y\cos z, x\sin y\sin z, x\cos y\right) .
\end{aligned}
\end{equation*}
\item

设 $ X $ 是 $ \mathbb{R}^{n} $ 中的开集，$ F $ 是巴拿赫空间，$ f : X \to F $。又设 $ x_{0} \in X $，并取 $ \varepsilon > 0 $ 使得 $ \mathbb{B}(x_{0},\varepsilon) \subset X $。最后，令

\[
x_{k}(h) := x_{0} + h^{1}e_{1} + \cdots + h^{k}e_{k} \quad \text{对 } k = 1,\ldots,n \text{ 且 } h \in \mathbb{B}(x_{0},\varepsilon) .
\]

证明 $ f $ 在 $ x_{0} $ 处可微当且仅当：对每个满足 $ h^{k} \ne 0 $（$ 1 \le k \le n $）的 $ h \in \mathbb{B}(x_{0},\varepsilon) $，极限

\begin{equation*}
\begin{aligned}
\lim_{\substack{h^{k} \to 0 \\ h^{k} \ne 0}} \frac{f\left(x_{k}(h)\right) - f\left(x_{k-1}(h)\right)}{h^{k}} \quad \text{对 } 1 \le k \le n
\end{aligned}
\end{equation*}

在 $ F $ 中存在。

\item

证明注 2.15(a)。（提示：定理 2.14。）

\item

确定下列函数的梯度：

\begin{equation*}
\begin{aligned}
\mathbb{R}^{m} \to \mathbb{R} , \quad x &\mapsto (x_{0}|x) ; \\ \mathbb{R}^{m}\setminus\{0\} \to \mathbb{R} , \quad x &\mapsto |x| ; \\ \mathbb{R}^{m} \to \mathbb{R} , \quad x &\mapsto \left|(x_{0}|x)\right|^{2} ; \\ \mathbb{R}^{m}\setminus\{0\} \to \mathbb{R} , \quad x &\mapsto 1/|x| .
\end{aligned}
\end{equation*}
\item

映射 $ \mathbb{C} \to \mathbb{C} $，$ z \mapsto z|z| $ 在哪些点可微？如果可能，求出其导数。

\item

映射 $ \mathbb{R}^{m} \to \mathbb{R}^{m} $，$ x \mapsto x|x|^{k} $（$ k \in \mathbb{N} $）在哪些点可微？在导数存在的每一点处求出其导数。

\item

求所有满足 $ f(\mathbb{C}) \subset \mathbb{R} $ 的可微函数 $ f : \mathbb{C} \to \mathbb{C} $。

\item

对 $ p \in [1,\infty] $，令 $ f_{p} : \mathbb{R}^{m} \to \mathbb{R} $，$ x \mapsto |x|_{p} $。$ f_{p} $ 在哪些点可微？如果可能，求出 $ \nabla f_{p}(x) $。

\item

设 $ X $ 是 $ \mathbb{C} $ 中的开集，$ f : X \to \mathbb{C} $ 可微。再对 $ z \in X^{*} := \{\, \overline{z} \in \mathbb{C} ; z \in X \,\} $ 定义 $ f^{*}(z) := \overline{f(\overline{z})} $。证明 $ f^{*} : X^{*} \to \mathbb{C} $ 可微。

\end{exercises}

\section{多元微分法则}

我们在此汇集多元微分最重要的法则。照例，我们假设

• $ E $ 与 $ F $ 是域 $ \mathbb{K} $ 上的巴拿赫空间；$ X $ 是 $ E $ 的开子集。

\needspace{3\baselineskip}\subsection*{线性}

下一个定理表明，与一元情形一样，微分是线性的。

\begin{proposition}

设 $ f, g : X \to F $ 在 $ x_0 $ 处可微，且 $ \alpha \in \mathbb{K} $。则 $ f + \alpha g $ 也在 $ x_0 $ 处可微，并且

\[
\partial(f + \alpha g)(x_0) = \partial f(x_0) + \alpha\partial g(x_0) .
\]

\proofhead
由假设，我们可以写出
\begin{equation*}
\begin{aligned}
f(x) &= f(x_0) + \partial f(x_0)(x - x_0) + r_{x_0}(x)\|x - x_0\| , \\ g(x) &= g(x_0) + \partial g(x_0)(x - x_0) + s_{x_0}(x)\|x - x_0\| ,
\end{aligned}
\end{equation*}

其中函数 $ r_{x_0}, s_{x_0} : X \to F $ 连续且在 $ x_0 $ 处取值为零。于是，

\[
(f + \alpha g)(x) = (f + \alpha g)(x_0) + [\partial f(x_0) + \alpha\partial g(x_0)](x - x_0) + t_{x_0}(x)\|x - x_0\|
\]

其中 $ t_{x_0} := r_{x_0} + \alpha s_{x_0} $。命题 2.1 于是推出该定理。\nobreak\hspace{0.35em}■

\end{proposition}

\begin{corollary}

$ C^1(X, F) $ 是 $ C(X, F) $ 的向量子空间，并且

\[
\partial : C^1(X, F) \to C(X, \mathcal{L}(E, F)) , \quad f \mapsto \partial f
\]

是线性的。

\end{corollary}

\needspace{3\baselineskip}\subsection*{链式法则}

在第 IV 章中，我们看到了链式法则对一元函数的巨大重要性，而正如我们在此将看到的，多元情形也是如此。

\qyxffntextdef{1}{参见定理 IV.1.7 的证明。}
\begin{theorem}

（链式法则）设 $ Y $ 是 $ F $ 中的开集，$ G $ 是巴拿赫空间。又设 $ f : X \to F $ 在 $ x_0 $ 处可微，$ g : Y \to G $ 在 $ y_0 := f(x_0) $ 处可微，并且 $ f(X) \subset Y $。则 $ g \circ f : X \to G $ 在 $ x_0 $ 处可微，其导数由下式给出

\[
\partial(g \circ f)(x_0) = \partial g(f(x_0))\partial f(x_0) .
\]

证明\qyxffnmark{1} 若 $ A := \partial g(f(x_0))\partial f(x_0) $，则 $ A \in \mathcal{L}(E, G) $。命题 2.1 蕴含

\begin{equation}
\begin{aligned}
f(x) &= f(x_0) + \partial f(x_0)(x - x_0) + r(x)\|x - x_0\| \qquad \text{对 } x \in X , \\ g(y) &= g(y_0) + \partial g(y_0)(y - y_0) + s(y)\|y - y_0\| \qquad \text{对 } y \in Y ,
\end{aligned}\tag{3.1}
\end{equation}

其中 $ r : X \to F $ 与 $ s : Y \to G $ 分别在 $ x_0 $ 处与 $ y_0 $ 处连续，并且 $ r(x_0) = 0 $ 且 $ s(y_0) = 0 $。我们通过 $ t(x_0) := 0 $ 以及下式定义 $ t : X \to G $：

\[
t(x) := \partial g(f(x_0))r(x) + s(f(x))\left\|\partial f(x_0)\frac{x - x_0}{\|x - x_0\|} + r(x)\right\| \qquad \text{对 } x \ne x_0 .
\]

则 $ t $ 在 $ x_0 $ 处连续。由 (3.1)，并取 $ y := f(x) $，我们推出关系式

\begin{equation*}
\begin{aligned}
(g \circ f)(x) &= g(f(x_0)) + A(x - x_0) + \partial g(f(x_0))r(x)\|x - x_0\| \\ &\quad + s(f(x))\left\|\partial f(x_0)(x - x_0) + r(x)\|x - x_0\|\right\| \\ &= (g \circ f)(x_0) + A(x - x_0) + t(x)\|x - x_0\| .
\end{aligned}
\end{equation*}

于是该定理由命题 2.1 推出。\nobreak\hspace{0.35em}■

\end{theorem}

\begin{corollary}

（坐标形式下的链式法则）设 $ X $ 是 $ \mathbb{R}^n $ 中的开集，$ Y $ 是 $ \mathbb{R}^m $ 中的开集。又设 $ f : X \to \mathbb{R}^m $ 在 $ x_0 $ 处可微，$ g : Y \to \mathbb{R}^\ell $ 在 $ y_0 := f(x_0) $ 处可微，并且 $ f(X) \subset Y $。则 $ h := g \circ f : X \to \mathbb{R}^\ell $ 在 $ x_0 $ 处可微，并且

\begin{equation}
[\partial h(x_0)] = [\partial g(f(x_0))][\partial f(x_0)] ,\tag{3.2}
\end{equation}

也就是说，复合映射 $ h = g \circ f $ 的 雅可比矩阵是 $ g $ 与 $ f $ 的 雅可比矩阵之积。

\proofhead
这由定理 3.3、推论 2.9 与定理 1.9(ii) 推出。\nobreak\hspace{0.35em}■
\end{corollary}

\begin{remark}

采用推论 3.4 的记号，$ h $ 的 雅可比矩阵的各个元素为

\[
\frac{\partial h^j(x_0)}{\partial x^k} = \sum_{i=1}^{m} \frac{\partial g^j(f(x_0))}{\partial y^i}\frac{\partial f^i(x_0)}{\partial x^k} \qquad \text{对 } 1 \le j \le \ell \text{ 且 } 1 \le k \le n .
\]

利用这个基本公式，我们可以在许多具体情形中计算复合函数的偏导数。

\proofhead
利用矩阵乘法法则，这立即由 (3.2) 推出。\nobreak\hspace{0.35em}■
\end{remark}

\begin{example}

\begin{items}
\item

我们考虑映射

\begin{equation*}
\begin{aligned}
f : \mathbb{R}^2 \to \mathbb{R}^3 , &\quad (x, y) \mapsto (x^2, xy, xy^2) , \\ g : \mathbb{R}^3 \to \mathbb{R}^2 , &\quad (\xi, \eta, \zeta) \mapsto (\sin \xi, \cos(\xi \eta \zeta)) .
\end{aligned}
\end{equation*}

对于 $ h := g \circ f : \mathbb{R}^2 \to \mathbb{R}^2 $，我们于是有 $ h(x, y) = (\sin x^2, \cos(x^4y^3)) $。因此 $ h $ 连续可微，并且

\begin{equation}
[\partial h(x, y)] = \begin{bmatrix} 2x\cos x^2 & 0 \\ -4x^3y^3\sin(x^4y^3) & -3x^4y^2\sin(x^4y^3) \end{bmatrix} .\tag{3.3}
\end{equation}

至于 $ g $ 与 $ f $ 的 雅可比矩阵，我们分别求得

\begin{equation*}
\begin{bmatrix} \cos \xi & 0 & 0 \\ -\eta\zeta\sin(\xi \eta \zeta) & -\xi\zeta\sin(\xi \eta \zeta) & -\xi\eta\sin(\xi \eta \zeta) \end{bmatrix} \quad \text{与} \quad \begin{bmatrix} 2x & 0 \\ y & x \\ y^2 & 2xy \end{bmatrix} .
\end{equation*}

现在我们可以容易地验证：这两个矩阵的乘积在位置 $ (\xi, \eta, \zeta) = f(x, y) $ 处与 (3.3) 一致。

\item

设 $ X $ 是 $ \mathbb{R}^n $ 中的开集，$ f \in C^1(X, \mathbb{R}) $。又令 $ I $ 是 $ \mathbb{R} $ 中的开区间，并设 $ \varphi \in C^1(I, \mathbb{R}^n) $ 满足 $ \varphi(I) \subset X $。则 $ f \circ \varphi $ 属于 $ C^1(I, \mathbb{R}) $，并且

\[
(f \circ \varphi)'(t) = (\nabla f(\varphi(t)) \mid \dot\varphi(t)) \qquad \text{对 } t \in I .
\]

\proofhead
由链式法则，我们得到
\[
(f \circ \varphi)'(t) = df(\varphi(t))\dot\varphi(t) = (\nabla f(\varphi(t)) \mid \dot\varphi(t)) \qquad \text{对 } t \in I .\quad\blacksquare
\]
\item

我们沿用上述记号，考虑一条道路 $ \varphi : I \to X $，它保持在 $ f $ 的一个“水平集”之内，即存在 $ y \in \mathrm{im}(f) $，使得对每个 $ t \in I $ 有 $ f(\varphi(t)) = y $。由 (b) 于是可得

\[
(\nabla f(\varphi(t)) \mid \dot\varphi(t)) = 0
\]

对 $ t \in I $ 成立。这表明梯度 $ \nabla f(x) $ 在点 $ x = \varphi(t) $ 处与道路 $ \varphi $ 正交，从而也与 $ \varphi $ 过 $ (t, \varphi(t)) $ 的切线正交（参见注 IV.1.4(a)）。稍微不严格地说，我们说“梯度正交于水平集”。\nobreak\hspace{0.35em}■

\end{items}

\bookfigure[0.26]{figures/a4b3ecf4e3ad49b8.png}

函数 $ (x, y, z) \mapsto x^2 + y^2 - z^2 $ 的水平集

\end{example}

\needspace{3\baselineskip}\subsection*{乘积法则}

实值函数的乘积法则是链式法则的又一应用。\footnote{另见例 4.8(b)。}

\begin{proposition}

设 $ f, g \in C^1(X, \mathbb{R}) $。则 $ fg $ 也属于 $ C^1(X, \mathbb{R}) $，且乘积法则给出

\begin{equation}
\partial(fg) = g\partial f + f\partial g .\tag{3.4}
\end{equation}

\proofhead
对于
\[
m : \mathbb{R}^2 \to \mathbb{R} , \quad (\alpha, \beta) \mapsto \alpha\beta
\]

我们有 $ m \in C^1(\mathbb{R}^2, \mathbb{R}) $ 且 $ \nabla m(\alpha, \beta) = (\beta, \alpha) $。令 $ F := m \circ (f, g) $，我们得到对 $ x \in X $ 有 $ F(x) = f(x)g(x) $。由链式法则可得 $ F \in C^1(X, \mathbb{R}) $，并且

\[
\partial F(x) = \partial m(f(x), g(x)) \circ (\partial f(x), \partial g(x)) = g(x)\partial f(x) + f(x)\partial g(x) .\quad\blacksquare
\]
\end{proposition}

\begin{corollary}

在 $ E = \mathbb{R}^n $ 的情形下，我们有

\[
d(fg) = g\,df + f\,dg \quad \text{与} \quad \nabla(fg) = g\nabla f + f\nabla g .
\]

\proofhead
第一个公式是 (3.4) 的另一种写法。因为
\begin{equation*}
\begin{aligned}
(\nabla(fg)(x) \mid h) &= d(fg)(x)h = f(x)dg(x)h + g(x)df(x)h \\ &= (f(x)\nabla g(x) + g(x)\nabla f(x) \mid h)
\end{aligned}
\end{equation*}

对 $ x \in X $ 与 $ h \in \mathbb{R}^n $，第二个公式也成立。\nobreak\hspace{0.35em}■

\end{corollary}

\needspace{3\baselineskip}\subsection*{中值定理}

与一元情形一样，多元情形也有中值定理，这意味着我们可以利用导数来估计函数值之差。

在下文中，我们用第 III.4 节引入的记号 $ \llbracket x, y \rrbracket $ 表示点 $ x, y \in E $ 之间的直线道路 $ \{x + t(y - x) ; t \in [0, 1]\} $。

\begin{theorem}

（中值定理）设 $ f : X \to F $ 可微。则

\begin{equation}
\|f(x) - f(y)\| \le \sup_{0 \le t \le 1}\|\partial f(x + t(y - x))\|\|y - x\|\tag{3.5}
\end{equation}

对满足 $ \llbracket x, y \rrbracket \subset X $ 的所有 $ x, y \in X $ 成立。

\proofhead
对 $ t \in [0, 1] $ 我们令 $ \varphi(t) := f(x + t(y - x)) $。因为 $ \llbracket x, y \rrbracket \subset X $，我们知道 $ \varphi $ 有定义。链式法则表明 $ \varphi $ 可微，并且
\[
\dot\varphi(t) = \partial f(x + t(y - x))(y - x) .
\]

由定理 IV.2.18（一元向量值函数的中值定理）可得

\[
\|f(y) - f(x)\| = \|\varphi(1) - \varphi(0)\| \le \sup_{0 \le t \le 1}\|\dot\varphi(t)\| .
\]

我们再考虑下式即可完成定理的证明：

\[
\|\dot\varphi(t)\| \le \|\partial f(x + t(y - x))\|\|y - x\| \qquad \text{对 } t \in [0, 1] .\quad\blacksquare
\]
\end{theorem}

\begin{theorem}

（积分形式的中值定理）设 $ f \in C^1(X, F) $。则我们有

\begin{equation}
f(y) - f(x) = \int_0^1 \partial f(x + t(y - x))(y - x)\,dt\tag{3.6}
\end{equation}

对满足 $ \llbracket x, y \rrbracket \subset X $ 的 $ x, y \in X $ 成立。

\proofhead
前一证明中的辅助函数 $ \varphi $ 现在是连续可微的。因此我们可以应用微积分基本定理（推论 VI.4.14），它给出
\[
f(y) - f(x) = \varphi(1) - \varphi(0) = \int_0^1 \dot\varphi(t)\,dt = \int_0^1 \partial f(x + t(y - x))(y - x)\,dt .\quad\blacksquare
\]
\end{theorem}

\begin{remark}

设 $ f : X \to F $ 可微。

\begin{items}
\item

若 $ \partial f $ 连续，则表示式 (3.6) 给出估计 (3.5)。

\proofhead
应用命题 VI.4.3 与算子范数的定义即可。\nobreak\hspace{0.35em}■
\item

若 $ X $ 是凸集且 $ \partial f : X \to \mathcal{L}(E, F) $ 有界，则 $ f $ 利普希茨连续。

\proofhead
令 $ \alpha := \sup_{x \in X}\|\partial f(x)\| $。则由定理 3.9，我们得到
\begin{equation}
\|f(y) - f(x)\| \le \alpha\|y - x\| \qquad \text{对 } x, y \in X .\tag{3.7}
\end{equation}
\item

若 $ X $ 连通且 $ \partial f = 0 $，则 $ f $ 是常值函数。

\proofhead
设 $ x_0 \in X $ 且 $ r > 0 $ 满足 $ \mathbb{B}(x_0, r) \subset X $。令 $ y_0 := f(x_0) $，则由 (3.7) 可得对一切 $ x \in \mathbb{B}(x_0, r) $ 有 $ f(x) = y_0 $。因为 $ x_0 $ 是任意的，$ f $ 局部常值。因此 $ f^{-1}(y_0) $ 非空且在 $ X $ 中开。另外，由命题 2.1(ii) 与例 III.2.22(a)，$ f^{-1}(y_0) $ 在 $ X $ 中闭。于是该断言由注 III.4.3 推出。\nobreak\hspace{0.35em}■
\end{items}

\end{remark}

\needspace{3\baselineskip}\subsection*{函数序列的极限的可微性}

函数序列的极限的可微性定理（定理 V.2.8）现在可以容易地推广到一般情形。

\begin{theorem}

设对 $ k \in \mathbb{N} $ 有 $ f_k \in C^1(X, F) $。又假设存在函数 $ f \in F^X $ 与 $ g : X \to \mathcal{L}(E, F) $，使得

\begin{subitems}
\item

$ (f_k) $ 逐点收敛于 $ f $；

\item

$ (\partial f_k) $ 局部一致收敛于 $ g $。
则 $ f $ 属于 $ C^1(X, F) $，且我们有 $ \partial f = g $。

\proofhead
在应用定理 3.9 之下，定理 V.2.8 的证明逐字逐句仍然成立。\nobreak\hspace{0.35em}■
\end{subitems}

\end{theorem}

\needspace{3\baselineskip}\subsection*{局部极值的必要条件}

在重要的定理 IV.2.1 中，我们对一元实变量函数的局部极值的存在性给出了一个必要判据。利用这里建立的方法，我们现在可以处理多个变量。

\begin{theorem}

设映射 $ f : X \to \mathbb{R} $ 在 $ x_0 $ 处有局部极值点，且其所有方向导数在该点都存在。则

\[
D_vf(x_0) = 0 \qquad \text{对 } v \in E\setminus\{0\} .
\]

\proofhead
设 $ v \in E\setminus\{0\} $。我们选取 $ r > 0 $，使得对 $ t \in (-r, r) $ 有 $ x_0 + tv \in X $，并考虑一个一元实变量函数：
\[
\varphi : (-r, r) \to \mathbb{R} , \quad t \mapsto f(x_0 + tv) .
\]

则 $ \varphi $ 在 $ 0 $ 处可微且在 $ 0 $ 处有局部极值。由定理 IV.2.1，我们有 $ \dot\varphi(0) = D_vf(x_0) = 0 $。\nobreak\hspace{0.35em}■

\end{theorem}

\begin{remark}

\begin{items}
\item

设 $ f : X \to \mathbb{R} $ 在 $ x_0 $ 处可微。则若 $ df(x_0) = 0 $，称 $ x_0 $ 为 $ f $ 的\textbf{临界点}。若 $ f $ 在 $ x_0 $ 处有局部极值，则 $ x_0 $ 是临界点。当 $ E = \mathbb{R}^n $ 时，此定义与注 2.17(a) 的定义一致。

\proofhead
这由命题 2.5 与定理 3.13 推出。\nobreak\hspace{0.35em}■
\item

临界点 $ x_0 $ \emph{未必}是局部极值点。\footnote{另请与注 IV.2.2(c) 比较。}

\proofhead
考虑
\[
f : \mathbb{R}^2 \to \mathbb{R} , \quad (x, y) \mapsto x^2 - y^2 .
\]

则 $ \nabla f(0, 0) = 0 $，但 $ (0, 0) $ 不是 $ f $ 的极值点。它反而是“鞍点”。\nobreak\hspace{0.35em}■

\end{items}

\bookfigure[0.23]{figures/64750456f97b94d9.png}
\end{remark}

\exerciseshead

\begin{exercises}
\item

我们称函数 $ f : E \to F $ 为 $ \alpha \in \mathbb{R} $ 次\textbf{正齐次}的，如果

\[
f(tx) = t^{\alpha} f(x) \quad \text{对 } t > 0 \text{ 且 } x \in E \setminus \{0\} .
\]

证明：若 $ f \in C^{1}(E, F) $ 是 1 次正齐次的，则 $ f \in \mathcal{L}(E, F) $。

\item

设 $ f : \mathbb{R}^{m} \to \mathbb{R} $ 在 $ \mathbb{R}^{m} \setminus \{0\} $ 上可微。证明 $ f $ 是 $ \alpha $ 次正齐次的，如果它满足\textbf{欧拉齐次性关系}

\[
(\nabla f(x)|x) = \alpha f(x) \quad \text{对 } x \in \mathbb{R}^{m} \setminus \{0\} .
\]
\item

设 $ X $ 是 $ \mathbb{R}^{m} $ 中的开集，$ f \in C^{1}(X, \mathbb{R}^{n}) $。证明

\[
g : X \to \mathbb{R} , \quad x \mapsto \sin(|f(x)|^{2})
\]

是连续可微的，并确定 $ \nabla g $。

\item

对于 $ f \in C^{1}(\mathbb{R}^{n}, \mathbb{R}) $ 与 $ A \in \mathcal{L}(\mathbb{R}^{n}) $，证明

\[
\nabla(f \circ A)(x) = A^{*} \nabla f(Ax) \quad \text{对 } x \in \mathbb{R}^{n} .
\]
\item

设 $ X $ 是 $ \mathbb{R}^{k} $ 中的开集，$ Y $ 是 $ \mathbb{R}^{n} $ 中的有界开集。进一步设 $ f \in C(\overline{X \times Y}, \mathbb{R}) $ 在 $ X \times Y $ 中可微，且存在 $ \xi \in C^{1}(X, \mathbb{R}^{n}) $ 使得 $ \operatorname{im}(\xi) \subset Y $ 且

\[
m(x) := \min_{y \in \overline{Y}} f(x, y) = f(x, \xi(x)) \quad \text{对 } x \in X .
\]

计算 $ m : X \to \mathbb{R} $ 的梯度。

\item

对于 $ g \in C^{1}(\mathbb{R}^{m}, \mathbb{R}) $ 与 $ f_{j} \in C^{1}(\mathbb{R}, \mathbb{R}) $，其中 $ j = 1, \ldots, m $，计算以下映射的梯度：

\[
\mathbb{R}^{m} \to \mathbb{R} , \quad (x_{1}, \ldots, x_{m}) \mapsto g(f_{1}(x_{1}), \ldots, f_{m}(x_{m})) .
\]
\item

函数 $ f \in C^{1}(\mathbb{R}^{2}, \mathbb{R}) $ 满足 $ \partial_{1} f = \partial_{2} f $ 且 $ f(0, 0) = 0 $。证明存在 $ g \in C(\mathbb{R}^{2}, \mathbb{R}) $ 使得 $ f(x, y) = g(x, y)(x + y) $ 对 $ (x, y) \in \mathbb{R}^{2} $ 成立。

\item

对指数映射验证

\[
\exp \in C^{1}(\mathcal{L}(E), \mathcal{L}(E)) \quad \text{且} \quad \partial \exp(0) = \operatorname{id}_{\mathcal{L}(E)} .
\]
\end{exercises}

\section{多重线性映射}

设 $ E $ 与 $ F $ 为巴拿赫空间，$ X $ 是 $ E $ 的开子集。对 $ f \in C^1(X, F) $，有 $ \partial f \in C(X, \mathcal{L}(E, F)) $。令 $ g := \partial f $ 且 $ \mathbb{F} := \mathcal{L}(E, F) $。依定理 1.1，$ \mathbb{F} $ 是巴拿赫空间，因此我们可以研究映射 $ g \in C(X, \mathbb{F}) $ 的可微性。若 $ g $ [连续]可微，则称 $ f $ 二次[连续]可微，而 $ \partial^2 f := \partial g $ 是 $ f $ 的二阶导数。于是

\[
\partial^2 f(x) \in \mathcal{L}(E, \mathcal{L}(E, F)) \qquad \text{对 } x \in X .
\]

定理 1.1 表明 $ \mathcal{L}(E, \mathcal{L}(E, F)) $ 也是巴拿赫空间。因此我们可以研究三阶导数 $ \partial^3 f := \partial(\partial^2 f) $（若其存在），并同样求得

\[
\partial^3 f(x) \in \mathcal{L}(E, \mathcal{L}(E, \mathcal{L}(E, F))) \qquad \text{对 } x \in X .
\]

在这种记号下，随着 $ n $ 增大，$ \partial^n f $ 所占据的空间显然变得越来越复杂。然而，情况并不像看上去那样糟糕。我们将证明，这个看似复杂的空间 $ \mathcal{L}(E, \mathcal{L}(E, \ldots, \mathcal{L}(E, F) \cdots)) $（其中 $ E $ 出现 $ n $ 次）等距同构于从 $ E \times \cdots \times E $（同样有 $ n $ 个 $ E $）到 $ F $ 的所有连续多重线性映射的空间。因此，多重线性映射是理解高阶导数的恰当框架。

\needspace{3\baselineskip}\subsection*{连续的多重线性映射}

在下文中，对 $ m \ge 2 $ 的 $ E_1, \ldots, E_m $、$ E $ 与 $ F $ 都是域 $ \mathbb{K} $ 上的巴拿赫空间。映射 $ \varphi : E_1 \times \cdots \times E_m \to F $ 是\textbf{多重线性}的，或等价地，是 $ m $ 线性的\footnote{当 $ m = 2 $ [或 $ m = 3 $] 时，称其为双线性 [或三线性] 映射。}，若对每个 $ k \in \{1, \ldots, m\} $ 以及满足 $ j = 1, \ldots, m $ 且 $ j \ne k $ 的任意选取 $ x_j \in E_j $，映射

\[
\varphi(x_1, \ldots, x_{k-1}, \cdot, x_{k+1}, \ldots, x_m) : E_k \to F
\]

是线性的，也就是说，$ \varphi $ 是多重线性的，若它在每个变量上都是线性的。

首先我们证明，多重线性映射连续当且仅当它有界。\footnote{与定理 VI.2.5 比较。}

\begin{proposition}

对 $ m $ 线性映射 $ \varphi : E_1 \times \cdots \times E_m \to F $，下列陈述等价：

\begin{subitems}
\item

$ \varphi $ 连续。

\item

$ \varphi $ 在 $ 0 $ 处连续。

\item

$ \varphi $ 在有界集上有界。

\item

存在 $ \alpha \ge 0 $ 使得

\[
\|\varphi(x_1, \ldots, x_m)\| \le \alpha\|x_1\| \cdots\cdots \|x_m\| \qquad \text{对 } x_j \in E_j ,\ 1 \le j \le m .
\]

\proofhead
“$ \text{(i)}\Rightarrow\text{(ii)} $”这一蕴含是显然的。
“$ \text{(ii)}\Rightarrow\text{(iii)} $” 假设 $ B \subset E_1 \times \cdots \times E_m $ 有界，则依例 II.3.3(c) 与注 II.3.2(a)，存在 $ \beta > 0 $，使得对 $ (x_1, \ldots, x_m) \in B $ 及 $ 1 \le j \le m $ 有 $ \|x_j\| \le \beta $。因为 $ \varphi $ 在 $ 0 $ 处连续，存在 $ \delta > 0 $，使得
\[
\|\varphi(y_1, \ldots, y_m)\| \le 1 \qquad \text{对 } y_j \in E_j ,\ \|y_j\| \le \delta ,\ 1 \le j \le m .
\]

对 $ (x_1, \ldots, x_m) \in B $ 及 $ 1 \le j \le m $，令 $ y_j := \delta x_j/\beta $。则 $ \|y_j\| \le \delta $，并且

\[
(\delta/\beta)^m\|\varphi(x_1, \ldots, x_m)\| = \|\varphi(y_1, \ldots, y_m)\| \le 1 .
\]

因此 $ \varphi(B) $ 有界。
“$ \text{(iii)}\Rightarrow\text{(iv)} $” 由假设，存在 $ \alpha > 0 $，使得

\[
\|\varphi(x_1, \ldots, x_m)\| \le \alpha \qquad \text{对 } (x_1, \ldots, x_m) \in \overline{\mathbb{B}} .
\]

对 $ y_j \in E_j\setminus\{0\} $，令 $ x_j := y_j/\|y_j\| $。则 $ (x_1, \ldots, x_m) $ 属于 $ \overline{\mathbb{B}} $，并且我们求得

\[
\frac{1}{\|y_1\| \cdots\cdots \|y_m\|}\|\varphi(y_1, \ldots, y_m)\| = \|\varphi(x_1, \ldots, x_m)\| \le \alpha ,
\]

这就证明了该断言。
“$ \text{(iv)}\Rightarrow\text{(i)} $” 设 $ y = (y_1, \ldots, y_m) $ 是 $ E_1 \times \cdots \times E_m $ 中的点，$ (x^n) $ 是 $ E_1 \times \cdots \times E_m $ 中满足 $ \lim_n x^n = y $ 的序列。令 $ (x_1^n, \ldots, x_m^n) := x^n $，则由假设我们有

\begin{equation*}
\begin{aligned}
\|\varphi(y_1, \ldots, y_m) - \varphi(x_1^n, \ldots, x_m^n)\| &\le \|\varphi(y_1 - x_1^n, y_2, \ldots, y_m)\| + \|\varphi(x_1^n, y_2 - x_2^n, y_3, \ldots, y_m)\| \\ &\quad + \cdots + \|\varphi(x_1^n, x_2^n, \ldots, y_m - x_m^n)\| \\ &\le \alpha\Big( \|y_1 - x_1^n\|\|y_2\| \cdots\cdots \|y_m\| + \|x_1^n\|\|y_2 - x_2^n\| \cdots\cdots \|y_m\| \\ &\quad + \cdots + \|x_1^n\|\|x_2^n\| \cdots\cdots \|y_m - x_m^n\| \Big) .
\end{aligned}
\end{equation*}

因为序列 $ (x^n) $ 在 $ E_1 \times \cdots \times E_m $ 中有界，这一估计连同（命题 II.3.14 的）类似物表明，$ \varphi(x^n) $ 收敛于 $ \varphi(y) $。\nobreak\hspace{0.35em}■

\end{subitems}

\end{proposition}

引入一个缩写记号是有用的。我们用

\[
\mathcal{L}(E_1, \ldots, E_m; F)
\]

表示从 $ E_1 \times \cdots \times E_m $ 到 $ F $ 的所有连续多重线性映射之集。显然 $ \mathcal{L}(E_1, \ldots, E_m; F) $ 是 $ C(E_1 \times \cdots \times E_m, F) $ 的向量子空间。常常，所有 $ E_j $ 都相同，此时我们采用记号

\[
\mathcal{L}^m(E, F) := \mathcal{L}(E, \ldots, E; F) .
\]

另外，我们令 $ \mathcal{L}^1(E, F) := \mathcal{L}(E, F) $ 以及 $ \mathcal{L}^0(E, F) := F $。最后，令

\begin{equation}
\|\varphi\| := \inf\{ \alpha \ge 0 ; \|\varphi(x_1, \ldots, x_m)\| \le \alpha\|x_1\| \cdots\cdots \|x_m\|, x_j \in E_j \}\tag{4.1}
\end{equation}

对 $ \varphi \in \mathcal{L}(E_1, \ldots, E_m; F) $。

\begin{theorem}

\begin{subitems}
\item

对 $ \varphi \in \mathcal{L}(E_1, \ldots, E_m; F) $，我们有

\[
\|\varphi\| = \sup\{ \|\varphi(x_1, \ldots, x_m)\| ; \|x_j\| \le 1,\ 1 \le j \le m \}
\]

并且

\[
\|\varphi(x_1, \ldots, x_m)\| \le \|\varphi\|\|x_1\| \cdots\cdots \|x_m\|
\]

对 $ (x_1, \ldots, x_m) \in E_1 \times \cdots \times E_m $。

\item

关系式 (4.1) 定义了一个范数，并且

\[
\mathcal{L}(E_1, \ldots, E_m; F) := (\mathcal{L}(E_1, \ldots, E_m; F), \|\cdot\|)
\]

是巴拿赫空间。

\item

当对 $ 1 \le j \le m $ 有 $ \dim E_j < \infty $ 时，每个 $ m $ 线性映射都连续。

\proofhead
对 $ m = 1 $，这些由命题 VI.2.3、结论 VI.2.4(e) 以及定理 1.1 与 1.6 推出。一般情形可通过对这些证明作显而易见的修改来验证，我们将其作为习题留给读者。\nobreak\hspace{0.35em}■
\end{subitems}

\end{theorem}

\needspace{3\baselineskip}\subsection*{典范同构}

$ \mathcal{L}(E_1, \ldots, E_m; F) $ 上的范数是 $ \mathcal{L}(E, F) $ 上算子范数的自然推广。下一个定理表明，这一范数在另一种意义上也是自然的。

\begin{theorem}

空间 $ \mathcal{L}(E_1, \ldots, E_m; F) $ 与 $ \mathcal{L}(E_1, \mathcal{L}(E_2, \ldots, \mathcal{L}(E_m, F) \cdots)) $ 等距同构。

\proofhead
我们对 $ m = 2 $ 验证该断言。一般情形通过简单的归纳论证得到。

\begin{subitems}
\item

对 $ T \in \mathcal{L}(E_1, \mathcal{L}(E_2, F)) $ 我们令

\[
\varphi_T(x_1, x_2) := (Tx_1)x_2 \qquad \text{对 } (x_1, x_2) \in E_1 \times E_2 .
\]

则 $ \varphi_T : E_1 \times E_2 \to F $ 是双线性的，并且

\[
\|\varphi_T(x_1, x_2)\| \le \|T\|\|x_1\|\|x_2\| \qquad \text{对 } (x_1, x_2) \in E_1 \times E_2 .
\]

因此 $ \varphi_T $ 属于 $ \mathcal{L}(E_1, E_2; F) $，且 $ \|\varphi_T\| \le \|T\| $。

\item

设 $ \varphi \in \mathcal{L}(E_1, E_2; F) $。则我们令

\[
T_\varphi(x_1)x_2 := \varphi(x_1, x_2) \qquad \text{对 } (x_1, x_2) \in E_1 \times E_2 .
\]

因为

\[
\|T_\varphi(x_1)x_2\| = \|\varphi(x_1, x_2)\| \le \|\varphi\|\|x_1\|\|x_2\| \qquad \text{对 } (x_1, x_2) \in E_1 \times E_2 ,
\]

我们得到

\[
T_\varphi(x_1) \in \mathcal{L}(E_2, F) \qquad \text{对 } \|T_\varphi(x_1)\| \le \|\varphi\|\|x_1\|
\]

对每个 $ x_1 \in E_1 $。因此

\[
T_\varphi := [x_1 \mapsto T_\varphi(x_1)] \in \mathcal{L}(E_1, \mathcal{L}(E_2, F)) \quad \text{且} \quad \|T_\varphi\| \le \|\varphi\| .
\]
\item

综上所述，我们已经证明了映射

\[
T \mapsto \varphi_T : \mathcal{L}(E_1, \mathcal{L}(E_2, F)) \to \mathcal{L}(E_1, E_2; F)
\]

与

\begin{equation}
\varphi \mapsto T_\varphi : \mathcal{L}(E_1, E_2; F) \to \mathcal{L}(E_1, \mathcal{L}(E_2, F))\tag{4.2}
\end{equation}

是线性的、连续的，且互为逆映射。因此它们是拓扑同构。特别地，$ T_{\varphi_T} = T $，并且我们得到

\[
\|T\| = \|T_{\varphi_T}\| \le \|\varphi_T\| \le \|T\| .
\]

因此映射 $ T \mapsto \varphi_T $ 是等距。\nobreak\hspace{0.35em}■

\end{subitems}

\end{theorem}

约定 用等距同构 (4.2) 将 $ \mathcal{L}(E_1, \ldots, E_m; F) $ 与 $ \mathcal{L}(E_1, \mathcal{L}(E_2, \ldots, \mathcal{L}(E_m, F) \cdots)) $ 等同起来。

\begin{conclusion}

\begin{items}
\item

对 $ m \in \mathbb{N} $，我们有

\[
\mathcal{L}(E, \mathcal{L}^m(E, F)) = \mathcal{L}^{m+1}(E, F) .
\]

\proofhead
这由定理 4.3 与上述约定立即推出。\nobreak\hspace{0.35em}■
\item

$ \mathcal{L}^m(E, F) $ 是巴拿赫空间。

\proofhead
这由定理 4.3 与注 1.3(c) 推出。\nobreak\hspace{0.35em}■
\end{items}

\end{conclusion}

\needspace{3\baselineskip}\subsection*{对称多重线性映射}

设 $ m \ge 2 $ 且 $ \varphi : E^m \to F $ 是 $ m $ 线性的。我们称 $ \varphi $ 是\textbf{对称}的，若

\[
\varphi(x_{\sigma(1)}, \ldots, x_{\sigma(m)}) = \varphi(x_1, \ldots, x_m)
\]

对每个 $ (x_1, \ldots, x_m) $ 以及 $ \{1, \ldots, m\} $ 的每个置换 $ \sigma $ 成立。我们令

\[
\mathcal{L}_{\mathrm{sym}}^m(E, F) := \{ \varphi \in \mathcal{L}^m(E, F) ; \varphi \text{ 是对称的} \} .
\]

\qyxffntextdef{1}{回顾 $ \mathrm{S}_n $ 是集合 $ \{1, \ldots, n\} $ 的置换群（见第 I.7 节末尾）。}
\begin{proposition}

$ \mathcal{L}_{\mathrm{sym}}^m(E, F) $ 是 $ \mathcal{L}^m(E, F) $ 的闭向量子空间，因此它本身也是巴拿赫空间。

\proofhead
设 $ (\varphi_k) $ 是 $ \mathcal{L}_{\mathrm{sym}}^m(E, F) $ 中的序列，在 $ \mathcal{L}^m(E, F) $ 中收敛于 $ \varphi $。对每个 $ (x_1, \ldots, x_m) \in E^m $ 以及每个置换\qyxffnmark{1} $ \sigma \in S_m $，我们于是有
\begin{equation*}
\begin{aligned}
\varphi(x_{\sigma(1)}, \ldots, x_{\sigma(m)}) &= \lim_k \varphi_k(x_{\sigma(1)}, \ldots, x_{\sigma(m)}) = \lim_k \varphi_k(x_1, \ldots, x_m) \\ &= \varphi(x_1, \ldots, x_m) .
\end{aligned}
\end{equation*}

因此 $ \varphi $ 是对称的。\nobreak\hspace{0.35em}■

\end{proposition}

\needspace{3\baselineskip}\subsection*{多重线性映射的导数}

下一个定理表明，$ m $ 线性映射实际上是连续可微的，并且其导数是一些 $ (m-1) $ 线性函数的和。

\begin{proposition}

$ \mathcal{L}(E_1, \ldots, E_m; F) $ 是 $ C^1(E_1 \times \cdots \times E_m, F) $ 的向量子空间。并且，对 $ \varphi \in \mathcal{L}(E_1, \ldots, E_m; F) $ 及 $ (x_1, \ldots, x_m) \in E_1 \times \cdots \times E_m $，我们有

\[
\partial\varphi(x_1, \ldots, x_m)(h_1, \ldots, h_m) = \sum_{j=1}^{m}\varphi(x_1, \ldots, x_{j-1}, h_j, x_{j+1}, \ldots, x_m)
\]

对 $ (h_1, \ldots, h_m) \in E_1 \times \cdots \times E_m $。

\proofhead
我们用 $ A_x $ 表示映射
\[
(h_1, \ldots, h_m) \mapsto \sum_{j=1}^{m}\varphi(x_1, \ldots, x_{j-1}, h_j, x_{j+1}, \ldots, x_m)
\]

从 $ E_1 \times \cdots \times E_m $ 到 $ F $。则不难验证

\[
(x \mapsto A_x) \in C(E_1 \times \cdots \times E_m, \mathcal{L}(E_1 \times \cdots \times E_m, F)) .
\]

令 $ y_k := x_k + h_k $，则等式

\[
\varphi(y_1, \ldots, y_m) = \varphi(x_1, \ldots, x_m) + \sum_{k=1}^{m}\varphi(x_1, \ldots, x_{k-1}, h_k, y_{k+1}, \ldots, y_m)
\]

由 $ \varphi $ 的多重线性性推出。因为每个映射

\[
E_{k+1} \times \cdots \times E_m \to F , \quad (z_{k+1}, \ldots, z_m) \mapsto \varphi(x_1, \ldots, x_{k-1}, h_k, z_{k+1}, \ldots, z_m)
\]

是 $ (m-k) $ 线性的，我们同样得到

\begin{equation*}
\begin{aligned}
\varphi(x_1, \ldots, x_{k-1}, h_k, y_{k+1}, \ldots, y_m) &= \varphi(x_1, \ldots, x_{k-1}, h_k, x_{k+1}, \ldots, x_m) \\ &\quad + \sum_{j=1}^{m-k}\varphi(x_1, \ldots, x_{k-1}, h_k, x_{k+1}, \ldots, x_{k+j-1}, h_{k+j}, y_{k+j+1}, \ldots, y_m) .
\end{aligned}
\end{equation*}

因此，我们求得

\[
\varphi(x_1 + h_1, \ldots, x_m + h_m) - \varphi(x_1, \ldots, x_m) - A_x h = r(x, h) ,
\]

其中 $ r(x, h) $ 是若干函数值多重线性映射的和，其中每个被加项对适当的 $ i \in \{1, \ldots, m\} $ 都至少含有两个不同的 $ h_i $。于是定理 4.2(i) 蕴含当 $ h \to 0 $ 时 $ r(x, h) = o(\|h\|) $。\nobreak\hspace{0.35em}■

\end{proposition}

\begin{corollary}

设 $ X $ 是 $ \mathbb{K} $ 中的开集，$ \varphi \in \mathcal{L}(E_1, \ldots, E_m; F) $ 且对 $ 1 \le j \le m $ 有 $ f_j \in C^1(X, E_j) $。则映射

\[
\varphi(f_1, \ldots, f_m) : X \to F , \quad x \mapsto \varphi(f_1(x), \ldots, f_m(x))
\]

属于 $ C^1(X, F) $，并且我们有

\[
\partial(\varphi(f_1, \ldots, f_m)) = \sum_{j=1}^{m}\varphi(f_1, \ldots, f_{j-1}, f_j', f_{j+1}, \ldots, f_m) .
\]

\proofhead
这由 $ f_j'(x) \in E_j $（命题 4.6）以及链式法则推出。\nobreak\hspace{0.35em}■
\end{corollary}

符号公式

\[
\det[a_k^j] = \sum_{\sigma \in S_m}\operatorname{sign}(\sigma)a_{\sigma(1)}^1 \cdots\cdots a_{\sigma(m)}^m
\]

表明 $ \det[a_k^j] $ 是行向量 $ a_{j\bullet} := (a_j^1, \ldots, a_j^m) $（或者，等价地，列向量 $ a_{\bullet k} := (a_k^1, \ldots, a_k^m) $）的 $ m $ 线性函数。

对满足 $ a_k = (a_k^1, \ldots, a_k^m) $ 的 $ a_1, \ldots, a_m \in \mathbb{K}^m $，我们令

\[
[a_1, \ldots, a_m] := [a_k^j] \in \mathbb{K}^{m \times m} .
\]

换句话说，$ [a_1, \ldots, a_m] $ 是以列向量 $ a_{\bullet k} := a_k $（$ 1 \le k \le m $）为列的方阵。因此，行列式函数

\[
\det : \underbrace{\mathbb{K}^m \times \cdots \times \mathbb{K}^m}_{m} \to \mathbb{K} , \quad (a_1, \ldots, a_m) \mapsto \det[a_1, \ldots, a_m]
\]

是一个良定义的 $ m $ 线性映射。

\begin{example}

设 $ X $ 是 $ \mathbb{K} $ 中的开集。

\begin{items}
\item

对 $ a_1, \ldots, a_m \in C^1(X, \mathbb{K}^m) $，我们有 $ \det[a_1, \ldots, a_m] \in C^1(X, \mathbb{K}) $，并且

\[
(\det[a_1, \ldots, a_m])' = \sum_{j=1}^{m}\det[a_1, \ldots, a_{j-1}, a_j', a_{j+1}, \ldots, a_m] .
\]
\item

设 $ \varphi \in \mathcal{L}(E_1, E_2; F) $ 且 $ (f, g) \in C^1(X, E_1 \times E_2) $。则\textbf{（广义）乘积法则}为

\[
\partial\varphi(f, g) = \varphi(\partial f, g) + \varphi(f, \partial g) .
\]
\end{items}

\end{example}

\exerciseshead

\begin{exercises}
\item

设 $ H := (H,(\cdot \mid \cdot)) $ 是有限维希尔伯特空间。证明这些事实：

\begin{items}
\item

对每个 $ a \in \mathcal{L}^{2}(H,\mathbb{K}) $，恰好存在一个 $ A \in \mathcal{L}(H) $，即由 $ a $ 诱导的\textbf{线性算子}，满足

\[
a(x,y) = (Ax \mid y) \quad \text{对 } x, y \in H .
\]

由此诱导的映射

\begin{equation}
\mathcal{L}^{2}(H,\mathbb{K}) \to \mathcal{L}(H) , \quad a \mapsto A\tag{4.3}
\end{equation}

是等距同构。

\item

在实的情形中，有 $ a \in \mathcal{L}^{2}_{\mathrm{sym}}(H,\mathbb{R}) \iff A = A^{*} $。

\item

设 $ H = \mathbb{R}^{m} $ 且 $ (\cdot \mid \cdot) $ 是欧几里得内积。证明

\[
a(x,y) = \sum_{j,k=1}^{m} a_{jk}x^{j}y^{k} \quad \text{对 } x = (x^{1},\ldots,x^{m}) , \quad y = (y^{1},\ldots,y^{m}) , \quad x, y \in \mathbb{K}^{m} ,
\]

其中 $ [a_{jk}] $ 是由 $ a $ 诱导的线性算子的矩阵表示。
（(a) 的提示：里斯表示定理。\footnote{因为里斯表示定理对无限维希尔伯特空间也成立，所以 (4.3) 在这样的空间中也成立。}）

\end{items}

\item

设 $ X $ 在 $ E $ 中开且 $ f \in C^{1}(X,\mathcal{L}^{m}(E,F)) $ 满足 $ f(X) \subset \mathcal{L}^{m}_{\mathrm{sym}}(E,F) $。则证明对 $ x \in X $ 与 $ h \in E $ 有 $ \partial f(x)h \in \mathcal{L}^{m}_{\mathrm{sym}}(E,F) $。

\item

对 $ A \in \mathbb{K}^{n \times n} $，证明 $ \det e^{A} = e^{\operatorname{tr}(A)} $。（提示：令 $ h(t) := \det e^{tA} - e^{t\operatorname{tr}(A)} $，其中 $ t \in \mathbb{R} $。借助例 4.8(a)，推出 $ h' = 0 $。）

\item

对 $ T_{1}, T_{2} \in \mathcal{L}(F,E) $ 与 $ S \in \mathcal{L}(E,F) $，令 $ g(T_{1},T_{2})(S) := T_{1}ST_{2} $。验证

\begin{items}
\item

$ g(T_{1},T_{2}) \in \mathcal{L}(\mathcal{L}(E,F),\mathcal{L}(F,E)) $；

\item

$ g \in \mathcal{L}^{2}(\mathcal{L}(F,E);\mathcal{L}(\mathcal{L}(E,F),\mathcal{L}(F,E))) $。

\end{items}

\item

设 $ H $ 是希尔伯特空间 $ A \in \mathcal{L}(H) $ 且 $ a(x) := (Ax \mid x) $，其中 $ x \in H $。证明

\begin{items}
\item

对 $ x, y \in H $，$ \partial a(x)y = (Ax \mid y) + (Ay \mid x) $；

\item

若 $ H = \mathbb{R}^{n} $，则对 $ x \in \mathbb{R}^{n} $ 有 $ \nabla a(x) = (A + A^{*})x $。

\end{items}

\end{exercises}

\section{高阶导数}

现在我们可以定义高阶导数。我们将看到，如果一个函数连续可微，那么它的高阶导数就是对称多重线性映射。

本节主要把泰勒展开推广到多于一个变量的函数。作为它的一个推论，它将像一元情形那样给出局部极值存在性的一个充分判据。

在下文中，假设

• $ E $ 与 $ F $ 是巴拿赫空间，$ X $ 是 $ E $ 的开子集。

\needspace{3\baselineskip}\subsection*{定义}

如前所述，我们将像一元情形那样归纳地定义高阶导数。

设 $ f : X \to F $ 且 $ x_{0} \in X $。我们令 $ \partial^{0}f := f $。因此 $ \partial^{0}f(x_{0}) $ 属于 $ F = \mathcal{L}^{0}(E, F) $。现在设 $ m \in \mathbb{N}^{\times} $ 且 $ \partial^{m-1}f : X \to \mathcal{L}^{m-1}(E, F) $ 已有定义。若

\[
\partial^{m}f(x_{0}) := \partial(\partial^{m-1}f)(x_{0}) \in \mathcal{L}\left( E, \mathcal{L}^{m-1}(E, F) \right) = \mathcal{L}^{m}(E, F)
\]

存在，我们称 $ \partial^{m}f(x_{0}) $ 为 $ f $ 在 $ x_{0} $ 处的 $ m $ 阶导数，并称 $ f $ 在 $ x_{0} $ 处 $ m $ 次可微。当 $ \partial^{m}f(x) $ 对每个 $ x \in X $ 都存在时，我们称

\[
\partial^{m}f : X \to \mathcal{L}^{m}(E, F)
\]

为 $ f $ 的 $ m $ 阶导数，并说 $ f $ 是 $ m $ 次可微的。如果 $ \partial^{m}f $ 还连续，我们说 $ f $ 是 $ m $ 次连续可微的。我们令

\[
C^{m}(X, F) := \{ f : X \to F ;\ f \text{ 是 } m\text{ 次连续可微的} \}
\]

以及

\[
C^{\infty}(X, F) := \bigcap_{m \in \mathbb{N}}C^{m}(X, F) .
\]

显然 $ C^{0}(X, F) = C(X, F) $，并且若 $ f $ 属于 $ C^{\infty}(X, F) $，则 $ f $ 是\textbf{光滑的}或\textbf{无穷次连续可微的}。

除 $ \partial^{m}f $ 之外，我们有时也写 $ D^{m}f $ 或 $ f^{(m)} $；我们也令 $ f' := f^{(1)} $、$ f'' := f^{(2)} $、$ f''' := f^{(3)} $ 等等。

\begin{remark}

设 $ m \in \mathbb{N} $。

\begin{items}
\item

$ m $ 阶导数是线性的，即

\[
\partial^{m}(f + \alpha g)(x_{0}) = \partial^{m}f(x_{0}) + \alpha\partial^{m}g(x_{0})
\]

对 $ \alpha \in \mathbb{K} $ 以及满足 $ f $ 与 $ g $ 在 $ x_{0} $ 处 $ m $ 次可微的 $ f, g : X \to F $ 成立。

\proofhead
这由命题 3.1 用归纳法推出。\nobreak\hspace{0.35em}■
\item

$ C^{m}(X, F) = \left\{ f : X \to F ;\ \partial^{j}f \in C\left( X, \mathcal{L}^{j}(E, F) \right),\ 0 \le j \le m \right\} $。

\proofhead
这由命题 2.1(ii) 与 $ m $ 阶导数的定义推出。\nobreak\hspace{0.35em}■
\item

现在容易验证 $ C^{m+1}(X, F) $ 是 $ C^{m}(X, F) $ 的向量子空间。$ C^{\infty}(X, F) $ 也是 $ C^{m}(X, F) $ 的向量子空间。特别地，我们有包含链

\[
C^{\infty}(X, F) \subset \cdots \subset C^{m+1}(X, F) \subset C^{m}(X, F) \subset \cdots \subset C(X, F) .
\]

此外，对 $ k \in \mathbb{N} $，映射

\[
\partial^{k} : C^{m+k}(X, F) \to C^{m}(X, F)
\]

有定义且是线性的。

\proofhead
我们把简单的验证留给你。\nobreak\hspace{0.35em}■
\end{items}

\end{remark}

\qyxffntextdef{1}{为了清楚起见，我们把多重线性映射的各个自变量用方括号括了起来。}
\begin{theorem}

设 $ f \in C^{2}(X, F) $。则

\[
\partial^{2}f(x) \in \mathcal{L}^{2}_{\mathrm{sym}}(E, F) \quad \text{对 } x \in X ,
\]

也就是说，\qyxffnmark{1}

\[
\partial^{2}f(x)[h, k] = \partial^{2}f(x)[k, h] \quad \text{对 } x \in X \text{ 且 } h, k \in E .
\]

证明

\begin{subitems}
\item

设 $ x \in X $ 且 $ r > 0 $ 满足 $ \mathbb{B}(x, 2r) \subset X $，并设 $ h, k \in r\mathbb{B} $。我们令 $ g(y) := f(y + h) - f(y) $ 以及

\[
r(h, k) := \int_{0}^{1}\left[ \int_{0}^{1}\left\{ \partial^{2}f(x + sh + tk) - \partial^{2}f(x) \right\}h\, ds \right]k\, dt .
\]

由中值定理（定理 3.10）与导数的线性性，我们得到

\begin{equation*}
\begin{aligned}
f(x + h + k) - f(x + k) - f(x + h) + f(x) &= g(x + k) - g(x) \\ &= \int_{0}^{1}\partial g(x + tk)k\, dt \\ &= \int_{0}^{1}\left( \partial f(x + h + tk) - \partial f(x + tk) \right)k\, dt \\ &= \int_{0}^{1}\left( \int_{0}^{1}\partial^{2}f(x + sh + tk)h\, ds \right)k\, dt \\ &= \partial^{2}f(x)[h, k] + r(h, k) ,
\end{aligned}
\end{equation*}

因为对 $ s, t \in [0, 1] $ 有 $ x + sh + tk $ 属于 $ \mathbb{B}(x, 2r) $。

\item

令 $ \tilde{g}(y) := f(y + k) - f(y) $ 以及

\[
\tilde{r}(k, h) := \int_{0}^{1}\left( \int_{0}^{1}\left\{ \partial^{2}f(x + sk + th) - \partial^{2}f(x) \right\}k\, ds \right)h\, dt ,
\]

类似地可得

\[
f(x + h + k) - f(x + h) - f(x + k) + f(x) = \partial^{2}f(x)[k, h] + \tilde{r}(k, h) .
\]

于是我们得到

\[
\partial^{2}f(x)[h, k] - \partial^{2}f(x)[k, h] = \tilde{r}(k, h) - r(h, k) .
\]

再注意到估计

\[
\|\tilde{r}(k, h)\| \vee \|r(h, k)\| \le \sup_{0 \le s,t \le 1}\|\partial^{2}f(x + sh + tk) - \partial^{2}f(x)\|\,\|h\|\,\|k\| ,
\]

我们得到

\[
\|\partial^{2}f(x)[h, k] - \partial^{2}f(x)[k, h]\| \le 2\sup_{0 \le s,t \le 1}\|\partial^{2}f(x + sh + tk) - \partial^{2}f(x)\|\,\|h\|\,\|k\| .
\]

在这个不等式中，我们可以用 $ \tau h $ 与 $ \tau k $（其中 $ \tau \in (0, 1] $）替换 $ h $ 与 $ k $。于是我们有

\[
\|\partial^{2}f(x)[h, k] - \partial^{2}f(x)[k, h]\| \le 2\sup_{0 \le s,t \le 1}\left\|\partial^{2}f\left( x + \tau(sh + tk) \right) - \partial^{2}f(x)\right\|\,\|h\|\,\|k\| .
\]

最后，$ \partial^{2}f $ 的连续性蕴含

\[
\sup_{0 \le s,t \le 1}\left\|\partial^{2}f\left( x + \tau(sh + tk) \right) - \partial^{2}f(x)\right\| \to 0 \quad (\tau \to 0) .
\]

因此对 $ h, k \in r\mathbb{B} $ 有 $ \partial^{2}f(x)[h, k] = \partial^{2}f(x)[k, h] $，证明完成。\nobreak\hspace{0.35em}■

\end{subitems}

\end{theorem}

\begin{corollary}

对满足 $ m \ge 2 $ 的 $ f \in C^{m}(X, F) $，我们有

\[
\partial^{m}f(x) \in \mathcal{L}^{m}_{\mathrm{sym}}(E, F) \quad \text{对 } x \in X .
\]

\proofhead
仿照定理 5.2，我们现在构造一个关于 $ m $ 的归纳证明。归纳步骤 $ m \to m + 1 $ 如下。因为 $ \mathcal{L}^{m}_{\mathrm{sym}}(E, F) $ 是闭向量子空间或 $ \mathcal{L}(E, F) $，由归纳假设得
\[
\partial^{m+1}f(x)h_{1} \in \mathcal{L}^{m}_{\mathrm{sym}}(E, F) \quad \text{对 } h_{1} \in E ,
\]

（参见习题 4.2）。特别地，我们得到

\[
\partial^{m+1}f(x)[h_{1}, h_{\sigma(2)}, h_{\sigma(3)}, \ldots, h_{\sigma(m+1)}] = \partial^{m+1}f(x)[h_{1}, h_{2}, h_{3}, \ldots, h_{m+1}]
\]

对每个 $ (h_{1}, \ldots, h_{m+1}) \in E^{m+1} $ 与 $ \{2, \ldots, m+1\} $ 的每个置换 $ \sigma $ 成立。因为每个 $ \tau \in \mathrm{S}_{m+1} $ 都可表示成对换的复合，只需验证

\[
\partial^{m+1}f(x)[h_{1}, h_{2}, h_{3}, \ldots, h_{m+1}] = \partial^{m+1}f(x)[h_{2}, h_{1}, h_{3}, \ldots, h_{m+1}] .
\]

因为 $ \partial^{m+1}f(x) = \partial^{2}(\partial^{m-1}f)(x) $，这由定理 5.2 推出。\nobreak\hspace{0.35em}■

\end{corollary}

\needspace{3\baselineskip}\subsection*{高阶偏导数}

我们首先考虑 $ E = \mathbb{R}^{n} $ 且 $ n \ge 2 $ 的情形。

对 $ q \in \mathbb{N}^{\times} $ 与指标 $ j_{1}, \ldots, j_{q} \in \{1, \ldots, n\} $，我们称

\[
\frac{\partial^{q}f(x)}{\partial x^{j_{1}}\partial x^{j_{2}}\cdots\partial x^{j_{q}}} := \partial_{j_{1}}\partial_{j_{2}}\cdots\partial_{j_{q}}f(x) \quad \text{对 } x \in X ,
\]

为 $ f : X \to F $ 在 $ x $ 处的 $ q $ 阶偏导数\footnote{在一个显然的记号下，我们把对同一变量的多个相继导数简写为，例如 $ \frac{\partial^{3}f}{\partial x^{1}\partial x^{4}\partial x^{4}} = \frac{\partial^{3}f}{\partial x^{1}(\partial x^{4})^{2}} $、$ \frac{\partial^{4}f}{\partial x^{2}\partial x^{3}\partial x^{3}\partial x^{2}} = \frac{\partial^{4}f}{\partial x^{2}(\partial x^{3})^{2}\partial x^{2}} $ 等等。不过要注意，一般仍应保持不同变量之间的求导次序，因为一般有 $ \frac{\partial^{2}f}{\partial x^{1}\partial x^{2}} \ne \frac{\partial^{2}f}{\partial x^{2}\partial x^{1}} $（参见注 5.6）。}。映射 $ f $ 称为 $ q $ 次[连续]偏可微的，如果它直到 $ q $ 阶（含 $ q $ 阶）的导数存在[且连续]。

\qyxffnnumdef{1}
\begin{theorem}

设 $ X $ 是 $ \mathbb{R}^{n} $ 中的开集，$ f : X \to F $，且 $ m \in \mathbb{N}^{\times} $。则下列陈述成立：

\begin{subitems}
\item

$ f $ 属于 $ C^{m}(X, F) $，当且仅当 $ f $ 是 $ m $ 次连续偏可微的。

\item

对 $ f \in C^{m}(X, F) $，我们有

\[
\frac{\partial^{q}f}{\partial x^{j_{1}}\cdots\partial x^{j_{q}}} = \frac{\partial^{q}f}{\partial x^{j_{\sigma(1)}}\cdots\partial x^{j_{\sigma(q)}}} \quad \text{对 } 1 \le q \le m ,
\]

对每个置换 $ \sigma \in \mathrm{S}_{q} $ 成立，也就是说，偏导数与求导次序无关。\qyxffnmark{1}

\proofhead
首先，我们注意到，对满足 $ q \le m $ 的 $ h_{1}, \ldots, h_{q} \in \mathbb{R}^{n} $，我们有
\[
\partial^{q}f(x)[h_{1}, \ldots, h_{q}] = \partial(\cdots\partial(\partial f(x)h_{1})h_{2}\cdots)h_{q} \quad \text{对 } x \in X .
\]

特别地，由此得

\begin{equation}
\partial^{q}f(x)[e_{j_{1}}, \ldots, e_{j_{q}}] = \partial_{j_{q}}\partial_{j_{q-1}}\cdots\partial_{j_{1}}f(x) \quad \text{对 } x \in X .\tag{5.1}
\end{equation}

于是，我们得到

\begin{equation}
\partial^{q}f(x)[h_{1}, \ldots, h_{q}] = \sum_{j_{1},\ldots,j_{q}=1}^{n}\partial_{j_{q}}\partial_{j_{q-1}}\cdots\partial_{j_{1}}f(x)h_{1}^{j_{1}}\cdots h_{q}^{j_{q}}\tag{5.2}
\end{equation}

对 $ x \in X $ 以及满足 $ 1 \le i \le q $ 的 $ h_{i} = (h_{i}^{1}, \ldots, h_{i}^{n}) \in \mathbb{R}^{n} $ 成立。

\end{subitems}

(i) 公式 (5.1) 立即蕴含 $ C^{m}(X, F) $ 的每个元素都是 $ m $ 次连续偏可微的。

反之，设 $ f $ 是 $ m $ 次连续偏可微的。在 $ m = 1 $ 的情形，由定理 2.10 我们知道 $ f $ 属于 $ C^{1}(X, F) $。我们假设该陈述对 $ m-1 \ge 1 $ 成立。由该假设与 (5.2)，我们看到 $ \partial_{j}(\partial^{m-1}f) $ 对 $ j \in \{1, \ldots, n\} $ 存在且连续。于是定理 2.10 说

\[
\partial(\partial^{m-1}f) : X \to \mathcal{L}\left( \mathbb{R}^{n}, \mathcal{L}^{m-1}(\mathbb{R}^{n}, F) \right) = \mathcal{L}^{m}(\mathbb{R}^{n}, F)
\]

存在且连续。因此 $ f \in C^{m}(X, F) $。

(ii) 这立即由推论 5.3 与 (5.1) 推出。\nobreak\hspace{0.35em}■

\end{theorem}
\qyxffntextof{1}{可以证明，当 $ f : X \subset \mathbb{R}^{n} \to F $ 是 $ m $ 次全可微时，偏导数已经与该次序无关（参见习题 17）。一般地，当 $ f $ 只是 $ m $ 次偏可微时，偏导数的先后次序不能交换（参见注 5.6）。}

\begin{corollary}

设 $ f : X \subset \mathbb{R}^{n} \to \mathbb{K} $。则有：

\begin{subitems}
\item

$ f $ 属于 $ C^{2}(X, \mathbb{K}) $，当且仅当

\[
f, \partial_{j}f, \partial_{j}\partial_{k}f \in C(X, \mathbb{K}) \quad \text{对 } 1 \le j, k \le n .
\]
\item

（施瓦茨定理）若 $ f $ 属于 $ C^{2}(X, \mathbb{K}) $，则

\[
\partial_{j}\partial_{k}f(x) = \partial_{k}\partial_{j}f(x) \quad \text{对 } x \in X , \quad 1 \le j, k \le n .
\]
\end{subitems}

\end{corollary}

\begin{remark}

当 $ f $ 只是两次偏可微时，施瓦茨定理不成立。

\proofhead
我们考虑 $ f : \mathbb{R}^{2} \to \mathbb{R} $，其中
\begin{equation*}
f(x, y) := \begin{cases} \frac{xy(x^{2} - y^{2})}{x^{2} + y^{2}} , & (x, y) \ne (0, 0) , \\ 0 , & (x, y) = (0, 0) . \end{cases}
\end{equation*}

于是我们有

\begin{equation*}
\partial_{1}f(x, y) = \begin{cases} \frac{y(x^{4} + 4x^{2}y^{2} - y^{4})}{(x^{2} + y^{2})^{2}} , & (x, y) \ne (0, 0) , \\ 0 , & (x, y) = (0, 0) , \end{cases}
\end{equation*}

并且由此得

\[
\partial_{2}\partial_{1}f(0, 0) = \lim_{h \to 0}\frac{\partial_{1}f(0, h) - \partial_{1}f(0, 0)}{h} = -1 .
\]

利用 $ f(y, x) = -f(x, y) $，我们得到

\[
\partial_{2}f(y, x) = \lim_{h \to 0}\frac{f(y, x + h) - f(y, x)}{h} = -\lim_{h \to 0}\frac{f(x + h, y) - f(x, y)}{h} = -\partial_{1}f(x, y) .
\]

因此对 $ (x, y) \in \mathbb{R}^{2} $ 有 $ \partial_{1}\partial_{2}f(y, x) = -\partial_{2}\partial_{1}f(x, y) $。于是 $ \partial_{2}\partial_{1}f(0, 0) \ne 0 $ 给出 $ \partial_{1}\partial_{2}f(0, 0) \ne -\partial_{2}\partial_{1}f(0, 0) $。\nobreak\hspace{0.35em}■

\end{remark}

\needspace{3\baselineskip}\subsection*{链式法则}

链式法则对高阶导数也很重要，尽管复合函数的高阶导数一般没有简单的表示。

\begin{theorem}

（链式法则）设 $ Y $ 是 $ F $ 中的开集，$ G $ 是巴拿赫空间。又设 $ m \in \mathbb{N}^{\times} $ 与 $ f \in C^{m}(X, F) $ 满足 $ f(X) \subset Y $，并设 $ g \in C^{m}(Y, G) $。则我们有 $ g \circ f \in C^{m}(X, G) $。

\proofhead
利用定理 3.3，该陈述由关于 $ m $ 的归纳法推出。我们把细节留给你作为练习。\nobreak\hspace{0.35em}■
\end{theorem}

\needspace{3\baselineskip}\subsection*{泰勒公式}

我们现在把泰勒定理推广到多元函数。为此，我们使用下面的记号：

\begin{equation*}
\partial^{k}f(x)[h]^{k} := \begin{cases} \partial^{k}f(x)[\underbrace{h, \ldots, h}_{k\text{ 次}}] , & 1 \le k \le q , \\ f(x) , & k = 0 , \end{cases}
\end{equation*}

对 $ x \in X $、$ h \in E $ 与 $ f \in C^{q}(X, F) $ 成立。

\qyxffnnumdef{1}
\begin{theorem}

（泰勒公式）设 $ X $ 是 $ E $ 中的开集，$ q \in \mathbb{N}^{\times} $，且 $ f $ 属于 $ C^{q}(X, F) $。则

\[
f(x + h) = \sum_{k=0}^{q}\frac{1}{k!}\partial^{k}f(x)[h]^{k} + R_{q}(f, x; h)
\]

对满足 $ \llbracket x, x + h \rrbracket \subset X $ 的 $ x \in X $ 与 $ h \in E $ 成立。这里

\[
R_{q}(f, x; h) := \int_{0}^{1}\frac{(1 - t)^{q-1}}{(q - 1)!}\left[ \partial^{q}f(x + th) - \partial^{q}f(x) \right][h]^{q}\, dt \in F
\]

是 $ f $ 在点 $ x $ 处的 $ q $ 阶余项。

\proofhead
当 $ q = 1 $ 时，该公式化为积分形式的中值定理（定理 3.10）：
\begin{equation*}
\begin{aligned}
f(x + h) &= f(x) + \int_{0}^{1}\partial f(x + th)h\, dt \\ &= f(x) + \partial f(x)h + \int_{0}^{1}\left[ \partial f(x + th) - \partial f(x) \right]h\, dt .
\end{aligned}
\end{equation*}

现在设 $ q = 2 $。令 $ u(t) := \partial f(x + th)h $ 与 $ v(t) := t - 1 $（$ t \in [0, 1] $），我们有 $ u'(t) = \partial^{2}f(x + th)[h]^{2} $ 且 $ v' = 1 $。因此，由广义乘积法则（例 4.8(b)）并对上式积分，我们得到\qyxffnmark{1}

\begin{equation*}
\begin{aligned}
f(x + h) &= f(x) + \partial f(x)h + \int_{0}^{1}(1 - t)\partial^{2}f(x + th)[h]^{2}\, dt \\ &= \sum_{k=0}^{2}\frac{1}{k!}\partial^{k}f(x)[h]^{k} + R_{2}(f, x; h) .
\end{aligned}
\end{equation*}

对 $ q = 3 $，我们令 $ u(t) := \partial^{2}f(x + th)[h]^{2} $ 与 $ v(t) := -(1 - t)^{2}/2 $。于是类似地得到

\[
f(x + h) = f(x) + \partial f(x)h + \frac{1}{2}\partial^{2}f(x)[h]^{2} + \int_{0}^{1}\frac{(1 - t)^{2}}{2}\partial^{3}f(x + th)[h]^{3}\, dt .
\]

在一般情形 $ q > 3 $，该断言用简单的归纳论证推出。\nobreak\hspace{0.35em}■

\end{theorem}
\qyxffntextof{1}{参见命题 VI.5.4 的证明。}

\begin{remark}

对 $ f \in C^{q}(X, F) $ 以及满足 $ \llbracket x, x + h \rrbracket \subset X $ 的 $ x, h \in E $，我们有

\[
\|R_{q}(f, x; h)\| \le \frac{1}{q!}\max_{0 \le t \le 1}\|\partial^{q}f(x + th) - \partial^{q}f(x)\|\,\|h\|^{q} .
\]

特别地，我们得到

\[
R_{q}(f, x; h) = o(\|h\|^{q}) \quad (h \to 0) .
\]

\proofhead
这由命题 VI.4.3 与 $ \partial^{q}f $ 的连续性推出。\nobreak\hspace{0.35em}■
\end{remark}

\qyxffntextdef{1}{对充分小的 $ h $，$ \llbracket x, x + h \rrbracket $ 落在 $ X $ 中。}
\begin{corollary}

设 $ f \in C^{q}(X, F) $，其中 $ q \in \mathbb{N}^{\times} $，且 $ x \in X $。则我们有\qyxffnmark{1}

\[
f(x + h) = \sum_{k=0}^{q}\frac{1}{k!}\partial^{k}f(x)[h]^{k} + o(\|h\|^{q}) \quad (h \to 0) .
\]
\end{corollary}

\needspace{3\baselineskip}\subsection*{$ m $ 个变量的函数}

在本节余下部分，我们考虑 $ m $ 个实变量的函数，即令 $ E = \mathbb{R}^{m} $。在这种情形，用多重指标表示偏导数是方便的。

设 $ \alpha = (\alpha_{1}, \ldots, \alpha_{m}) \in \mathbb{N}^{m} $ 是长度为\footnote{参见第 I.8 节。} $ |\alpha| = \sum_{j=1}^{m}\alpha_{j} $ 的多重指标，并设 $ f \in C^{|\alpha|}(X, F) $。则我们写

\[
\partial^{\alpha}f := \partial_{1}^{\alpha_{1}}\partial_{2}^{\alpha_{2}}\cdots\partial_{m}^{\alpha_{m}}f = \frac{\partial^{|\alpha|}f}{(\partial x^{1})^{\alpha_{1}}(\partial x^{2})^{\alpha_{2}}\cdots(\partial x^{m})^{\alpha_{m}}} \quad \text{对 } \alpha \ne 0
\]

以及 $ \partial^{0}f := f $。

\begin{theorem}

（泰勒定理）设 $ f \in C^{q}(X, F) $，其中 $ q \in \mathbb{N}^{\times} $，且 $ x \in X $。则

\[
f(y) = \sum_{|\alpha| \le q}\frac{\partial^{\alpha}f(x)}{\alpha!}(y - x)^{\alpha} + o(|x - y|^{q}) \quad (x \to y) .
\]

\proofhead
我们令 $ h := y - x $ 并写 $ h = \sum_{j=1}^{m}h_{j}e_{j} $。则
\begin{equation}
\begin{aligned}
\partial^{k}f(x)[h]^{k} &= \partial^{k}f(x)\left[ \sum_{j_{1}=1}^{m}h_{j_{1}}e_{j_{1}}, \ldots, \sum_{j_{k}=1}^{m}h_{j_{k}}e_{j_{k}} \right] \\ &= \sum_{j_{1}=1}^{m}\cdots\sum_{j_{k}=1}^{m}\partial^{k}f(x)[e_{j_{1}}, \ldots, e_{j_{k}}]h_{j_{1}}\cdots h_{j_{k}}
\end{aligned}\tag{5.3}
\end{equation}

对 $ 1 \le k \le q $ 成立。使得所有数 $ \ell \in \{1, \ldots, m\} $ 都恰好出现 $ \alpha_{\ell} $ 次的 $ k $ 元组 $ (j_{1}, \ldots, j_{k}) $（其中 $ 1 \le j_{i} \le m $）的个数等于

\begin{equation}
\frac{k!}{(\alpha_{1})!(\alpha_{2})!\cdots(\alpha_{m})!} = \frac{k!}{\alpha!} .\tag{5.4}
\end{equation}

由 (5.1) 与推论 5.3，该断言由

\begin{equation}
\partial^{k}f(x)[e_{j_{1}}, e_{j_{2}}, \ldots, e_{j_{k}}] = \partial_{m}^{\alpha_{m}}\partial_{m-1}^{\alpha_{m-1}}\cdots\partial_{1}^{\alpha_{1}}f(x) = \partial^{\alpha}f(x) .\tag{5.5}
\end{equation}

推出。综合起来，由 (5.3)–(5.5) 我们得到

\[
\partial^{k}f(x)[h]^{k} = \sum_{|\alpha| = k}\frac{k!}{\alpha!}\partial^{\alpha}f(x)h^{\alpha} \quad \text{对 } |\alpha| = k ,
\]

于是该定理由推论 5.10 推出。\nobreak\hspace{0.35em}■

\end{theorem}

\begin{remark}

设 $ f \in C^{q}(X, \mathbb{R}) $，其中 $ q \in \mathbb{N}^{\times} $。

\begin{items}
\item

若 $ q = 1 $，我们有

\[
f(y) = f(x) + df(x)(y - x) + o(|y - x|) = f(x) + (\nabla f(x)|y - x) + o(|y - x|)
\]

当 $ y \to x $ 时。

\item

若 $ q = 2 $，我们有

\begin{equation*}
\begin{aligned}
f(y) &= f(x) + df(x)(y - x) + \frac{1}{2}\partial^{2}f(x)[y - x]^{2} + o(|y - x|^{2}) \\ &= f(x) + (\nabla f(x)|y - x) + \frac{1}{2}(H_{f}(x)(y - x)|y - x) + o(|y - x|^{2})
\end{aligned}
\end{equation*}

当 $ y \to x $ 时。这里，\footnote{我们用 $ \mathbb{R}^{m\times m}_{\mathrm{sym}} $ 表示 $ \mathbb{R}^{m\times m} $ 中由所有对称 $ (m \times m) $ 矩阵组成的向量子空间。}

\[
H_{f}(x) := [\partial_{j}\partial_{k}f(x)] \in \mathbb{R}^{m\times m}_{\mathrm{sym}}
\]

表示 $ f $ 在 $ x $ 处的\textbf{海森矩阵}。换句话说，$ H_{f}(x) $ 是由双线性型 $ \partial^{2}f(x)[\,\cdot\,\cdot\,] $ 在 $ \mathbb{R}^{m} $ 上诱导的线性算子的表示矩阵（参见习题 4.1）。\nobreak\hspace{0.35em}■

\end{items}

\end{remark}

\needspace{3\baselineskip}\subsection*{局部极值的充分判据}

在下面的注中，我们汇集线性代数的几个结果。

\begin{remark}

\begin{items}
\item

设 $ H := (H, (\cdot|\cdot)) $ 是有限维实希尔伯特空间。双线性型 $ b \in \mathcal{L}^{2}(H, \mathbb{R}) $ 称为\textbf{[半]正定}的，如果它对称且

\[
b(x, x) > 0 \quad [b(x, x) \ge 0] \quad \text{对 } x \in H\setminus\{0\} .
\]

它称为\textbf{[半]负定}的，如果 $ -b $ 是\textbf{[半]正定}的。如果 $ b $ 既不是半正定的也不是半负定的，则 $ b $ 是\textbf{不定}的。型 $ b $ 诱导出线性算子\footnote{参见习题 4.1。} $ B \in \mathcal{L}(H) $，满足对 $ x, y \in H $ 有 $ (Bx|y) = b(x, y) $；当 $ b $ 具有相应性质时，称该算子为\textbf{正定的}或\textbf{[半]负定}的。
设 $ A \in \mathcal{L}(H) $。则下列陈述等价：

\begin{subitems}
\item

$ A $ 是[半]正定的。

\item

$ A = A^{*} $，且对 $ x \in H\setminus\{0\} $ 有 $ (Ax|x) > 0 $ [$ \ge 0 $]。

\item

$ A = A^{*} $，且存在 $ \alpha > 0 $ [$ \alpha \ge 0 $]，使得对 $ x \in H $ 有 $ (Ax|x) \ge \alpha|x|^{2} $。

\proofhead
这由习题 4.1 与 III.4.11 推出。\nobreak\hspace{0.35em}■
\end{subitems}

\item

设 $ H = (H, (\cdot|\cdot)) $ 是维数为 $ m $ 的希尔伯特空间，且 $ A \in \mathcal{L}(H) $ 是自伴的，即 $ A = A^{*} $。则 $ A $ 的所有特征值都是实的且半单的。此外，存在特征向量 $ h_{1}, \ldots, h_{m} $ 的一个标准正交基，它给出 $ A $ 的\textbf{谱表示}

\begin{equation}
A = \sum_{j=1}^{m}\lambda_{j}(\cdot|h_{j})h_{j} ,\tag{5.6}
\end{equation}

其中，在 $ \lambda_{1}, \ldots, \lambda_{m} $ 中，$ A $ 的特征值按其重数出现，且对 $ 1 \le j \le m $ 有 $ Ah_{j} = \lambda_{j}h_{j} $。这特别意味着 $ A $ 在此基下的矩阵是对角的，即 $ [A] = \operatorname{diag}(\lambda_{1}, \ldots, \lambda_{m}) $。
最后，$ A $ 是正定的当且仅当它的所有特征值都是正的；$ A $ 是半正定的当且仅当它的特征值中没有负的。

\proofhead
$ \sigma(A) \subset \mathbb{R} $ 以及谱表示存在这两个陈述在线性代数中是标准的（例如，[Art93, Section VII.5]）。\footnote{另参见例 10.17(b)。}
由 (5.6)，我们读出
\[
(Ax|x) = \sum_{j=1}^{m}\lambda_{j}|(x|h_{j})|^{2} \quad \text{对 } x \in H .
\]

现在，所断言的 $ A $ 的刻画是所假设的特征值的[半]正定性的一个容易的推论。\nobreak\hspace{0.35em}■

\end{items}

我们现在可以导出——作为应用 IV.3.9 的结果的一个（部分）推广——多个实变量的实值函数的极大值与极小值的一个充分判据。

\end{remark}

\begin{theorem}

设 $ X $ 是 $ \mathbb{R}^{m} $ 中的开集，且 $ x_{0} \in X $ 是 $ f \in C^{2}(X, \mathbb{R}) $ 的临界点。则

\begin{subitems}
\item

若 $ \partial^{2}f(x_{0}) $ 正定，则 $ f $ 在 $ x_{0} $ 处有孤立局部极小值；

\item

若 $ \partial^{2}f(x_{0}) $ 负定，则 $ f $ 在 $ x_{0} $ 处有孤立局部极大值；

\item

若 $ \partial^{2}f(x_{0}) $ 不定，则 $ f $ 在 $ x_{0} $ 处没有局部极值。

\proofhead
因为 $ x_{0} $ 是 $ f $ 的临界点，根据注 5.12(b)，我们有
\[
f(x_{0} + \xi) = f(x_{0}) + \frac{1}{2}\partial^{2}f(x_{0})[\xi]^{2} + o(|\xi|^{2}) \quad (\xi \to 0) .
\]
\end{subitems}

(i) 因为 $ \partial^{2}f(x_{0}) $ 正定，存在 $ \alpha > 0 $ 使得

\[
\partial^{2}f(x_{0})[\xi]^{2} \ge \alpha|\xi|^{2} \quad \text{对 } \xi \in \mathbb{R}^{m} .
\]

我们取定一个 $ \delta > 0 $，使得对 $ \xi \in \bar{\mathbb{B}}^{m}(0, \delta) $ 有 $ |o(|\xi|^{2})| \le \alpha|\xi|^{2}/4 $，然后得到估计

\[
f(x_{0} + \xi) \ge f(x_{0}) + \frac{\alpha}{2}|\xi|^{2} - \frac{\alpha}{4}|\xi|^{2} = f(x_{0}) + \frac{\alpha}{4}|\xi|^{2} \quad \text{对 } |\xi| \le \delta .
\]

因此 $ f $ 在 $ x_{0} $ 处有局部极小值。

\begin{subitems}
\item[（ii）]

把 (i) 应用于 $ -f $ 即得该陈述。

\item[（iii）]

因为 $ \partial^{2}f(x_{0}) $ 不定，存在 $ \xi_{1}, \xi_{2} \in \mathbb{R}^{m}\setminus\{0\} $ 使得

\[
\alpha := \partial^{2}f(x_{0})[\xi_{1}]^{2} > 0 \quad \text{且} \quad \beta := \partial^{2}f(x_{0})[\xi_{2}]^{2} < 0 .
\]

此外，我们找到 $ t_{j} > 0 $ 使得 $ \llbracket x_{0}, x_{0} + t_{j}\xi_{j} \rrbracket \subset X $ 且

\[
\frac{\alpha}{2} + \frac{o(t^{2}|\xi_{1}|^{2})}{t^{2}|\xi_{1}|^{2}}|\xi_{1}|^{2} > 0 \quad \text{且} \quad \frac{\beta}{2} + \frac{o(t^{2}|\xi_{2}|^{2})}{t^{2}|\xi_{2}|^{2}}|\xi_{2}|^{2} < 0
\]

分别对 $ 0 < t < t_{1} $ 与 $ 0 < t < t_{2} $ 成立。于是，

\[
f(x_{0} + t\xi_{1}) = f(x_{0}) + t^{2}\left( \frac{\alpha}{2} + \frac{o(t^{2}|\xi_{1}|^{2})}{t^{2}|\xi_{1}|^{2}}|\xi_{1}|^{2} \right) > f(x_{0}) \quad \text{对 } 0 < t < t_{1} ,
\]

且

\[
f(x_{0} + t\xi_{2}) = f(x_{0}) + t^{2}\left( \frac{\beta}{2} + \frac{o(t^{2}|\xi_{2}|^{2})}{t^{2}|\xi_{2}|^{2}}|\xi_{2}|^{2} \right) < f(x_{0}) \quad \text{对 } 0 < t < t_{2} .
\]

因此 $ x_{0} $ 不是 $ f $ 的局部极值点。\nobreak\hspace{0.35em}■

\end{subitems}

\end{theorem}

\begin{example}

我们考虑

\[
f : \mathbb{R}^{2} \to \mathbb{R} , \quad (x, y) \mapsto c + \delta x^{2} + \varepsilon y^{2}
\]

其中 $ c \in \mathbb{R} $ 且 $ \delta, \varepsilon \in \{-1, 1\} $。则

\begin{equation*}
\nabla f(x, y) = 2(\delta x, \varepsilon y) \quad \text{且} \quad H_{f}(x, y) = 2\begin{bmatrix} \delta & 0 \\ 0 & \varepsilon \end{bmatrix} .
\end{equation*}

于是 $ (0, 0) $ 是 $ f $ 的唯一临界点。

\begin{items}
\item

取 $ f(x, y) = c + x^{2} + y^{2} $ 使得 $ \delta = \varepsilon = 1 $，我们得到 $ H_{f}(0, 0) $ 正定。于是 $ f $ 在 $ (0, 0) $ 处有孤立（绝对）极小值。

\item

取 $ f(x, y) = c - x^{2} - y^{2} $，我们得到 $ H_{f}(0, 0) $ 负定，且 $ f $ 在 $ (0, 0) $ 处有孤立（绝对）极大值。

\item

取 $ f(x, y) = c + x^{2} - y^{2} $ 使得 $ H_{f}(0, 0) $ 不定。因此 $ f $ 在 $ (0, 0) $ 处没有极值；更确切地说，$ (0, 0) $ 是一个“鞍点”。

\end{items}

\bookfigure[0.91]{figures/a67583ce17a1f40c.png}

\begin{items}
\item[（d）]

当 $ \partial^{2}f $ 在 $ f $ 的临界点处半定时，研究二阶导数不能给我们提供更多信息。例如，考虑映射 $ f_{j} : \mathbb{R}^{2} \to \mathbb{R} $（$ j = 1, 2, 3 $），其中

\[
f_{1}(x, y) := x^{2} + y^{4} , \quad f_{2}(x, y) := x^{2} , \quad f_{3}(x, y) := x^{2} + y^{3} .
\]

在每种情形，$ (0, 0) $ 都是 $ f_{j} $ 的临界点，且海森矩阵 $ H_{f_{j}}(0, 0) $ 是半正定的。

\end{items}

\bookfigure[0.85]{figures/15807451fd01c254.png}

然而，我们容易看出，$ (0, 0) $ 是 $ f_{1} $ 的孤立极小值点，不是 $ f_{2} $ 的孤立极小值点，也不是 $ f_{3} $ 的局部极值点。\nobreak\hspace{0.35em}■

\end{example}

\exerciseshead

\begin{exercises}
\item

若 $ A \in \mathcal{L}(E, F) $，证明 $ A \in C^{\infty}(E, F) $ 且 $ \partial^{2}A = 0 $。

\item

若 $ \varphi \in \mathcal{L}(E_{1}, \ldots, E_{m}; F) $，证明 $ \varphi \in C^{\infty}(E_{1} \times \cdots \times E_{m}, F) $ 且 $ \partial^{m+1}\varphi = 0 $。

\item

设 $ X $ 是 $ \mathbb{R}^{m} $ 中的开集。（$ m $ 维）\textbf{拉普拉斯算子} $ \Delta $ 定义为

\[
\Delta : C^{2}(X, \mathbb{R}) \to C(X, \mathbb{R}) , \quad u \mapsto \Delta u := \sum_{j=1}^{m}\partial_{j}^{2}u .
\]

若 $ \Delta u = 0 $，则称函数 $ u \in C^{2}(X, \mathbb{R}) $ 在 $ X $ 中\textbf{调和}。因此，调和函数构成线性映射 $ \Delta $ 的核。验证由下式定义的 $ g_{m} : \mathbb{R}^{m} \setminus \{0\} \to \mathbb{R} $

\begin{equation*}
g_{m}(x) := \begin{cases} \log|x| , & m = 2 , \\ |x|^{2-m} , & m > 2 , \end{cases}
\end{equation*}

在 $ \mathbb{R}^{m} \setminus \{0\} $ 中调和。

\item

对 $ f, g \in C^{2}(X, \mathbb{R}) $，证明

\[
\Delta(fg) = g\Delta f + 2(\nabla f|\nabla g) + f\Delta g .
\]
\item

设 $ g : [0, \infty) \times \mathbb{R} \to \mathbb{R}^{2} $，$ (r, \varphi) \mapsto (r\cos\varphi, r\sin\varphi) $。验证下列各条：

\begin{items}
\item

$ g \mid (0, \infty) \times (-\pi, \pi) \to \mathbb{R}^{2} \setminus H $（其中 $ H := \{(x, y) \in \mathbb{R}^{2} ; \ x \le 0, \ y = 0\} = (-\mathbb{R}^{+}) \times \{0\} $）是拓扑的。

\item

$ \operatorname{im}(g) = \mathbb{R}^{2} $。

\item

若 $ X $ 是 $ \mathbb{R}^{2} \setminus H $ 中的开集且 $ f \in C^{2}(X, \mathbb{R}) $，则

\[
(\Delta f) \circ g = \frac{\partial^{2}(f \circ g)}{\partial r^{2}} + \frac{1}{r}\frac{\partial(f \circ g)}{\partial r} + \frac{1}{r^{2}}\frac{\partial^{2}(f \circ g)}{\partial \varphi^{2}}
\]

在 $ g^{-1}(X) $ 上成立。

\end{items}

\item

令 $ X := \mathbb{R}^{n} \times (0, \infty) $，且对 $ (x, y) \in X $ 令 $ p(x, y) := y/(|x|^{2} + y^{2}) $。计算 $ \Delta p $。

\item

验证

\[
(\Delta u) \circ A = \Delta(u \circ A) \quad \text{对 } u \in C^{2}(\mathbb{R}^{n}, \mathbb{R})
\]

若 $ A \in \mathcal{L}(\mathbb{R}^{n}) $ 是\textbf{正交的}，即若 $ A^{*}A = 1 $。

\item

设 $ X $ 是 $ \mathbb{R}^{m} $ 中的开集，$ E $ 是巴拿赫空间。对 $ k \in \mathbb{N} $，令

\[
BC^{k}(X, E) := \left( \left\{ u \in BC(X, E) ; \ \partial^{\alpha}u \in BC(X, E), \ |\alpha| \le k \right\}, \ \|\cdot\|_{k,\infty} \right) ,
\]

其中

\[
\|u\|_{k,\infty} := \max_{|\alpha| \le k}\|\partial^{\alpha}u\|_{\infty} .
\]

验证下列各条：

\begin{items}
\item

$ BC^{k}(X, E) $ 是巴拿赫空间；

\item

$ BUC^{k}(X, E) := \{u \in BC(X, E) ; \ \partial^{\alpha}u \in BUC(X, E), \ |\alpha| \le k\} $ 是 $ BC^{k}(X, E) $ 的闭向量子空间，因此自身也是巴拿赫空间。

\end{items}

\item

令 $ q \in \mathbb{N} $，且对满足 $ |\alpha| \le q $ 的 $ \alpha \in \mathbb{N}^{m} $ 令 $ a_{\alpha} \in \mathbb{K} $。又设对某个满足 $ |\alpha| = q $ 的 $ \alpha \in \mathbb{N}^{m} $ 有 $ a_{\alpha} \ne 0 $。则称

\[
\mathcal{A}(\partial) : C^{q}(X, \mathbb{K}) \to C(X, \mathbb{K}) , \quad u \mapsto \mathcal{A}(\partial)u := \sum_{|\alpha| \le q}a_{\alpha}\partial^{\alpha}u
\]

为一个 $ q $ 阶常系数\textbf{线性微分算子}。对 $ k \in \mathbb{N} $，证明

\begin{items}
\item

$ \mathcal{A}(\partial) \in \mathcal{L}\left( BC^{k+q}(X, \mathbb{K}), BC^{k}(X, \mathbb{K}) \right) $；

\item

$ \mathcal{A}(\partial) \in \mathcal{L}\left( BUC^{k+q}(X, \mathbb{K}), BUC^{k}(X, \mathbb{K}) \right) $。

\end{items}

\item

对 $ u \in C^{2}(\mathbb{R} \times \mathbb{R}^{m}, \mathbb{R}) $ 与 $ (t, x) \in \mathbb{R} \times \mathbb{R}^{m} $，令

\[
\square u := \partial_{t}^{2}u - \Delta_{x}u
\]

以及

\[
(\partial_{t} - \Delta)u := \partial_{t}u - \Delta_{x}u ,
\]

其中 $ \Delta_{x} $ 是关于 $ x \in \mathbb{R}^{m} $ 的拉普拉斯算子。称 $ \square $ 为\textbf{波动算子}（或\textbf{达朗贝尔算子}）。称 $ \partial_{t} - \Delta $ 为\textbf{热算子}。

\begin{items}
\item

计算 $ (\partial_{t} - \Delta)k $ 在 $ (0, \infty) \times \mathbb{R}^{m} $ 上的值，其中

\[
k(t, x) := t^{-m/2}\exp\left( -|x|^{2}/4t \right) \quad \text{对 } (t, x) \in (0, \infty) \times \mathbb{R}^{m} .
\]
\item

设 $ g \in C^{2}(\mathbb{R}, \mathbb{R}) $、$ c > 0 $ 且 $ v \in S^{m-1} $。计算 $ \square w $，其中

\[
w(t, x) := g(v \cdot x - tc) \quad \text{对 } (t, x) \in \mathbb{R} \times \mathbb{R}^{m} .
\]
\end{items}

\item

设 $ X $ 是巴拿赫空间 $ E $ 的开凸子集，$ f \in C^{2}(X, \mathbb{R}) $。证明下列陈述等价：

\begin{subitems}
\item

$ f $ 是凸的；

\item

对 $ a, x \in X $，$ f(x) \ge f(a) + \partial f(a)(x - a) $；

\item

对 $ a \in X $ 与 $ h \in E $，$ \partial^{2}f(a)[h]^{2} \ge 0 $。
若 $ E = \mathbb{R}^{m} $，则这些陈述还等价于

\item

对 $ a \in X $，$ H_{f}(a) $ 是半正定的。

\end{subitems}

\item

对所有 $ \alpha \in \mathbb{R} $，把函数

\[
f_{\alpha} : \mathbb{R}^{2} \to \mathbb{R} , \quad (x, y) \mapsto x^{3} - y^{3} + 3\alpha xy ,
\]

的临界点分类为极大值、极小值或鞍点。

\item

设 $ f : \mathbb{R}^{2} \to \mathbb{R} $，$ (x, y) \mapsto (y - x^{2})(y - 2x^{2}) $。证明下列各条：

\begin{items}
\item

$ (0, 0) $ 不是 $ f $ 的局部极小值点。

\item

对每个 $ (x_{0}, y_{0}) \in \mathbb{R}^{2} \setminus \{(0, 0)\} $，$ \mathbb{R} \to \mathbb{R} \quad t \mapsto f(tx_{0}, ty_{0}) $ 在 $ 0 $ 处有孤立局部极小值。

\end{items}

\item

确定（到含二阶为止）函数

\[
(0, \infty) \times (0, \infty) \to \mathbb{R} , \quad (x, y) \mapsto (x - y)/(x + y)
\]

在点 $ (1, 1) $ 处的泰勒展开。

\item

设 $ X $ 是 $ \mathbb{R}^{m} $ 中的开集。对 $ f, g \in C^{1}(X, \mathbb{R}^{m}) $，由下式定义 $ [f, g] \in C(X, \mathbb{R}^{m}) $：

\[
[f, g](x) := \partial f(x)g(x) - \partial g(x)f(x) \quad \text{对 } x \in X .
\]

称 $ [f, g] $ 为 $ f $ 与 $ g $ 的\textbf{李括号}。验证下列陈述：

\begin{subitems}
\item

$ [f, g] = -[g, f] $；

\item

对 $ \alpha, \beta \in \mathbb{R} $ 与 $ h \in C^{1}(X, \mathbb{R}^{m}) $，$ [\alpha f + \beta g, h] = \alpha[f, h] + \beta[g, h] $；

\item

对 $ \varphi, \psi \in C^{1}(X, \mathbb{R}) $，$ [\varphi f, \psi g] = \varphi\psi[f, g] + (\nabla\varphi|f)\psi f - (\nabla\psi|g)\varphi g $；

\item

\textbf{（雅可比恒等式）} 若 $ f, g, h \in C^{2}(X, \mathbb{R}^{m}) $，则 $ [[f, g], h] + [[g, h], f] + [[h, f], g] = 0 $。

\end{subitems}

\item

证明

\begin{equation*}
\mathbb{R}^{n} \to \mathbb{R} , \quad x \mapsto \begin{cases} \exp\left( 1/(|x|^{2} - 1) \right) , & |x| < 1 , \\ 0 , & |x| \ge 1 , \end{cases}
\end{equation*}

是光滑的。

\item

设 $ X $ 是 $ \mathbb{R}^{m} $ 中的开集，$ f : X \to F $ 是 $ m $ 次可微的。证明

\[
\partial^{q}f(x) \in \mathcal{L}_{\mathrm{sym}}^{q}(X, F) \quad \text{对 } x \in X \text{ 且 } 1 \le q \le m .
\]

特别地，对每个置换 $ \sigma \in S_{q} $，证明

\[
\frac{\partial^{q}f(x)}{\partial x^{j_{1}} \cdot \ldots \cdot \partial x^{j_{q}}} = \frac{\partial^{q}f(x)}{\partial x^{j_{\sigma(1)}} \cdot \ldots \cdot \partial x^{j_{\sigma(q)}}} \quad \text{对 } x \in X .
\]

（提示：只需考虑 $ q = m = 2 $ 的情形。对

\[
[0, 1] \to F , \quad t \mapsto f(x + tsh_{1} + tsh_{2}) - f(x + tsh_{1}) ,
\]

应用中值定理，并验证

\begin{equation*}
\begin{aligned}
\lim_{\substack{s \to 0 \\ s > 0}}\frac{f(x + sh_{1} + sh_{2}) - f(x + sh_{1}) - f(x + sh_{2}) + f(x)}{s^{2}} = \partial^{2}f(x)[h_{1}, h_{2}] .)
\end{aligned}
\end{equation*}
\item

设 $ H $ 是有限维希尔伯特空间，$ A \in \mathcal{L}(H) $ 正定 [或负定]。证明 $ A $ 属于 $ \mathcal{L}\mathrm{aut}(H) $，且 $ A^{-1} $ 正定 [或负定]。

\item

设 $ H $ 是有限维希尔伯特空间，$ A \in C^{1}([0, T], \mathcal{L}(H)) $。又设对每个 $ t \in [0, T] $，$ A(t) $ 对称。证明下列各条：

\begin{items}
\item

对 $ t \in [0, T] $，$ \partial A(t) $ 对称。

\item

若对每个 $ t \in [0, T] $，$ \partial A(t) $ 正定，则对 $ 0 \le \sigma < \tau \le T $，有

\end{items}

\begin{subitems}
\item

$ A(\tau) - A(\sigma) $ 正定；

\item

若对每个 $ t \in [0, T] $，$ A(t) $ 是自同构，则 $ A^{-1}(\tau) - A^{-1}(\sigma) $ 负定。
（(ii) 的提示：对 $ t \mapsto A^{-1}(t)A(t) $ 求导。）

\end{subitems}

\item

设 $ H $ 是有限维希尔伯特空间，且令 $ A, B \in \mathcal{L}(H) $。证明：若 $ A $、$ B $ 与 $ A - B $ 正定，则 $ A^{-1} - B^{-1} $ 负定。

（提示：考虑对 $ t \in [0, 1] $ 的 $ t \mapsto B + t(A - B) $，并研究习题 19。）

\item

设 $ X $ 是 $ \mathbb{R}^{m} $ 中的开集，$ f, g \in C^{q}(X, \mathbb{K}) $。证明 $ fg $ 属于 $ C^{q}(X, \mathbb{K}) $，并证明\textbf{莱布尼茨法则}：

\[
\partial^{\alpha}(fg) = \sum_{\beta \le \alpha}\binom{\alpha}{\beta}\partial^{\beta}f\partial^{\alpha - \beta}g \quad \text{对 } \alpha \in \mathbb{N}^{m} \text{ 且 } |\alpha| \le q .
\]

这里 $ \beta \le \alpha $ 表示对 $ 1 \le j \le m $ 有 $ \beta_{j} \le \alpha_{j} $，且

\[
\binom{\alpha}{\beta} := \binom{\alpha_{1}}{\beta_{1}} \cdot \ldots \cdot \binom{\alpha_{m}}{\beta_{m}}
\]

对满足 $ \alpha = (\alpha_{1}, \ldots, \alpha_{m}) $ 与 $ \beta = (\beta_{1}, \ldots, \beta_{m}) $ 的 $ \alpha, \beta \in \mathbb{N}^{m} $。（提示：定理 IV.1.12(ii) 与归纳法。）

\end{exercises}

\section{涅梅茨基算子与变分法}

尽管我们为一般的巴拿赫空间建立了微分学，但迄今为止我们几乎只用有限维的例子来演示它。在本节中，我们通过让你领略一般理论的范围来弥补这一点。为此，我们首先考虑函数空间之间最简单的非线性映射，即“覆盖算子”，它用一个固定的非线性映射把两个函数空间连接起来。我们限制在连续函数巴拿赫空间中的覆盖算子，并研究它们的连续性与可微性性质。

作为应用，我们将研究变分法的基本问题，即刻画并求出“无穷多个变量”的实值函数的极值的问题。特别地，我们导出欧拉–拉格朗日方程，满足该方程是变分问题具有局部极值的必要条件。

\needspace{3\baselineskip}\subsection*{涅梅茨基算子}

设 $ T $、$ X $ 和 $ Y $ 是非空集合，$ \varphi $ 是从 $ T \times X $ 到 $ Y $ 的映射。那么我们通过下式定义\textbf{诱导涅梅茨基算子}或\textbf{覆盖算子} $ \varphi^{\sharp} $

\[
\varphi^{\sharp} : X^{T} \to Y^{T} , \quad u \mapsto \varphi(\cdot, u(\cdot)) .
\]

这意味着 $ \varphi^{\sharp} $ 是把每个函数 $ u : T \to X $ 对应到函数 $ \varphi^{\sharp}(u) : T \to Y $，$ t \mapsto \varphi(t, u(t)) $ 的映射。

在下面假设

• $ T $ 是紧致度量空间；
  $ E $ 和 $ F $ 是巴拿赫空间；
  $ X $ 在 $ E $ 中开。

当不会引起混淆时，我们把 $ E $ 和 $ F $ 上的范数都记为 $ \left|\cdot\right| $。

\begin{lemma}

$ C(T, X) $ 在巴拿赫空间 $ C(T, E) $ 中是开的。

\proofhead
由定理 V.2.6，我们知道 $ C(T, E) $ 是巴拿赫空间。

只需考虑 $ X \ne E $ 的情形。于是设 $ u \in C(T, X) $。那么，由于 $ u(T) $ 是紧致度量空间的连续像，它是 $ X $ 的紧致子集（见定理 III.3.6）。于是根据例 III.3.9(c)，存在 $ r > 0 $ 使得 $ d(u(T), X^{c}) \ge 2r $。于是，对于 $ v \in u + r\mathbb{B}_{C(T,E)} $，由 $ |u(t) - v(t)| < r $ 对 $ t \in T $ 成立可知 $ v(T) $ 落在 $ E $ 中 $ u(T) $ 的 $ r $ 邻域内，从而落在 $ X $ 中。因此 $ v $ 属于 $ C(T, X) $。 \nobreak\hspace{0.35em}■
\end{lemma}

\needspace{3\baselineskip}\subsection*{涅梅茨基算子的连续性}

我们首先证明连续映射在连续函数空间中诱导出连续的涅梅茨基算子。

\begin{theorem}

设 $ \varphi \in C(T \times X, F) $。那么

\begin{subitems}
\item

$ \varphi^{\sharp} \in C(C(T, X), C(T, F)) $；

\item

如果 $ \varphi $ 在有界集上有界，那么 $ \varphi^{\sharp} $ 也是如此。

\proofhead
对于 $ u \in C(T, X) $，设 $ \tilde{u}(t) := (t, u(t)) $ 对 $ t \in T $。那么 $ \tilde{u} $ 显然属于 $ C(T, T \times X) $（见命题 III.1.10）。于是定理 III.1.8 蕴含 $ \varphi^{\sharp}(u) = \varphi \circ \tilde{u} $ 属于 $ C(T, F) $。
\end{subitems}

(i) 设 $ (u_{j}) $ 是 $ C(T, X) $ 中的序列，且在 $ C(T, X) $ 中 $ u_{j} \to u_{0} $。那么根据例 III.3.1(a)，集合 $ M := \{ u_{j} ;\ j \in \mathbb{N}^{\times} \} \cup \{u_{0}\} $ 在 $ C(T, X) $ 中是紧致的。

我们令 $ M(T) := \bigcup\{ m(T) ;\ m \in M \} $，并要证明 $ M(T) $ 在 $ X $ 中是紧致的。为此，设 $ (x_{j}) $ 是 $ M(T) $ 中的序列。那么存在 $ t_{j} \in T $ 和 $ m_{j} \in M $ 使得 $ x_{j} = m_{j}(t_{j}) $。由于 $ T $ 和 $ M $ 是序列紧致的，并且根据习题 III.3.1，我们找到 $ (t, m) \in T \times M $ 以及一个子序列 $ ((t_{j_{k}}, m_{j_{k}}))_{k \in \mathbb{N}} $，它在 $ T \times M $ 中收敛到 $ (t, m) $。由此我们得到

\begin{equation*}
\begin{aligned}
|x_{j_{k}} - m(t)| &\le |m_{j_{k}}(t_{j_{k}}) - m(t_{j_{k}})| + |m(t_{j_{k}}) - m(t)| \\ &\le \|m_{j_{k}} - m\|_{\infty} + |m(t_{j_{k}}) - m(t)| \to 0
\end{aligned}
\end{equation*}

当 $ k \to \infty $ 时成立，因为 $ m $ 连续。于是 $ (x_{j_{k}}) $ 收敛到 $ M(T) $ 中的元素 $ m(t) $，这就证明了 $ M(T) $ 的序列紧致性，从而证明了它的紧致性。

于是，根据习题 III.3.1，$ T \times M(T) $ 是紧致的。现在定理 III.3.13 表明 $ \varphi $ 在 $ T \times M(T) $ 上的限制是一致连续的。由此我们得到

\[
\|\varphi^{\sharp}(u_{j}) - \varphi^{\sharp}(u_{0})\|_{\infty} = \max_{t \in T}\left|\varphi(t, u_{j}(t)) - \varphi(t, u_{0}(t))\right| \to 0
\]

当 $ j \to \infty $ 时成立，因为 $ \|u_{j} - u_{0}\|_{\infty} \to 0 $。我们还发现，对于 $ t \in T $ 和 $ j \in \mathbb{N} $，$ (t, u_{j}(t)) $ 和 $ (t, u_{0}(t)) $ 属于 $ T \times M(T) $。因此 $ \varphi^{\sharp} $ 是序列连续的，从而根据定理 III.1.4 是连续的。

(ii) 设 $ B \subset C(T, X) $，且存在 $ R > 0 $ 使得对 $ u \in B $ 有 $ \|u\|_{\infty} \le R $。那么

\[
B(T) := \bigcup\{ u(T) ;\ u \in B \} \subset X ,
\]

并且对 $ x \in B(T) $ 有 $ |x| \le R $。因此 $ B(T) $ 在 $ X $ 中有界。于是 $ T \times B(T) $ 在 $ T \times X $ 中有界。由于 $ \varphi $ 在有界集上有界，存在 $ r > 0 $ 使得对 $ (t, x) \in T \times B(T) $ 有 $ |\varphi(t, x)| \le r $。这蕴含对 $ u \in B $ 有 $ \|\varphi^{\sharp}(u)\|_{\infty} \le r $。因此 $ \varphi^{\sharp} $ 在有界集上有界，因为 $ \varphi $ 具有这一性质。 \nobreak\hspace{0.35em}■

\end{theorem}

\begin{remark}

当 $ E $ 是有限维时，每个 $ \varphi \in C(T \times X, F) $ 在形如 $ T \times B $ 的集合上都有界，其中 $ B $ 在 $ E $ 中有界，且 $ B $ 在 $ E $ 中的闭包 $ \overline{B} $ 落在 $ X $ 中。

\proofhead
这由定理 III.3.5、推论 III.3.7 和定理 1.4 得到。 \nobreak\hspace{0.35em}■
\end{remark}

\needspace{3\baselineskip}\subsection*{涅梅茨基算子的可微性}

设 $ p \in \mathbb{N} \cup \{\infty\} $。那么我们写

\[
\varphi \in C^{0,p}(T \times X, F)
\]

如果对每个 $ t \in T $，函数 $ \varphi(t, \cdot) : X \to F $ 是 $ p $ 次可微的，并且它的导数（我们记为 $ \partial_{2}^{q}\varphi $）满足

\[
\partial_{2}^{q}\varphi \in C(T \times X, \mathcal{L}^{q}(E, F)) \qquad \text{对 } q \in \mathbb{N} \text{ 且 } q \le p .
\]

因而 $ C^{0,0}(T \times X, F) = C(T \times X, F) $。

\qyxffntextdef{1}{这里以及在类似情形中，所有关于导数的陈述都适用于自然的情形 $ p \ge 1 $。}
\qyxffntextdef{2}{$ C(T, X) \times C(T, E) $ 与 $ C(T, X \times E) $ 按显然的方式等同。}
\begin{theorem}

对于 $ \varphi \in C^{0,p}(T \times X, F) $，$ \varphi^{\sharp} $ 属于 $ C^{p}(C(T, X), C(T, F)) $，并且\qyxffnmark{1}

\[
[\partial\varphi^{\sharp}(u)h](t) = \partial_{2}\varphi(t, u(t))h(t) \quad \text{对 } t \in T ,
\]

$ u \in C(T, X) $，$ h \in C(T, E) $。

\proofhead
根据定理 6.2，我们可以假设 $ p \ge 1 $。首先我们令
\[
G : C(T, X) \times C(T, E) \to F^{T} , \quad (u, h) \mapsto \partial_{2}\varphi(\cdot, u(\cdot))h(\cdot)
\]

为映射\qyxffnmark{2}

\[
[(t, (x, \xi)) \mapsto \partial_{2}\varphi(t, x)\xi] \in C(T \times (X \times E), F) .
\]

所诱导的涅梅茨基算子。那么定理 6.2 蕴含 $ G(u, h) $ 属于 $ C(T, F) $。显然

\begin{equation}
\|G(u, h)\|_{C(T,F)} \le \max_{t \in T}\|\partial_{2}\varphi(t, u(t))\|_{\mathcal{L}(E,F)}\|h\|_{C(T,E)} ,\tag{6.1}
\end{equation}

且 $ G(u, \cdot) $ 是线性的。于是我们看到

\[
A(u) := G(u, \cdot) \in \mathcal{L}(C(T, E), C(T, F)) \quad \text{对 } u \in C(T, X) .
\]

函数 $ \partial_{2}\varphi \in C(T \times X, \mathcal{L}(E, F)) $ 诱导出一个涅梅茨基算子 $ \widetilde{A} $。那么，由定理 6.2 和 1.1，

\begin{equation}
\widetilde{A} \in C(C(T, X), C(T, \mathcal{L}(E, F))) .\tag{6.2}
\end{equation}

现在设 $ u \in C(T, X) $。我们取 $ \varepsilon > 0 $ 使得 $ u(T) + \varepsilon\mathbb{B}_{E} \subset X $。那么对于每个满足 $ \|h\|_{\infty} < \varepsilon $ 的 $ h \in C(T, E) $，$ u + h $ 属于 $ C(T, X) $，并且积分形式的中值定理（定理 3.10）蕴含

\begin{equation*}
\begin{aligned}
&\left|\varphi^{\sharp}(u + h)(t) - \varphi^{\sharp}(u)(t) - (A(u)h)(t)\right| \\ &\quad = \left|\varphi(t, u(t) + h(t)) - \varphi(t, u(t)) - \partial_{2}\varphi(t, u(t))h(t)\right| \\ &\quad = \left|\int_{0}^{1}[\partial_{2}\varphi(t, u(t) + \tau h(t)) - \partial_{2}\varphi(t, u(t))]h(t) d\tau\right| \\ &\quad \le \|h\|_{C(T,E)}\int_{0}^{1}\|\widetilde{A}(u + \tau h) - \widetilde{A}(u)\|_{C(T,\mathcal{L}(E,F))} d\tau .
\end{aligned}
\end{equation*}

于是，由 (6.2) 可得

\[
\varphi^{\sharp}(u + h) - \varphi^{\sharp}(u) - A(u)h = o(\|h\|_{C(T,E)}) \quad (h \to 0) .
\]

于是，$ \partial\varphi^{\sharp}(u) $ 存在且等于 $ A(u) $。与 (6.1) 类似，我们考虑

\[
\|\partial\varphi^{\sharp}(u) - \partial\varphi^{\sharp}(v)\|_{\mathcal{L}(C(T,E),C(T,F))} \le \|\widetilde{A}(u) - \widetilde{A}(v)\|_{C(T,\mathcal{L}(E,F))} ,
\]

由 (6.2)，这蕴含

\[
\partial\varphi^{\sharp} \in C(C(T, X), \mathcal{L}(C(T, E), C(T, F))) .
\]

这证明了 $ p = 1 $ 时的定理。

一般情形现在简单地由归纳法得到，我们把它留给你。 \nobreak\hspace{0.35em}■

\end{theorem}

\begin{corollary}

如果 $ \partial_{2}\varphi $ 在有界集上有界，那么 $ \partial\varphi^{\sharp} $ 也是如此。

\proofhead
这由 (6.1) 和定理 6.2(ii) 得到。 \nobreak\hspace{0.35em}■
\end{corollary}

\begin{example}

\begin{items}
\item

设 $ \varphi(\xi) := \sin\xi $ 对 $ \xi \in \mathbb{R} $，且 $ T := [0, 1] $。那么涅梅茨基算子

\[
\varphi^{\sharp} : C(T) \to C(T) , \quad u \mapsto \sin u(\cdot)
\]

由 $ \varphi $ 诱导的是 $ C^{\infty} $ 映射，

\[
(\partial\varphi^{\sharp}(u)h)(t) = [\cos u(t)]h(t) \quad \text{对 } t \in T
\]

且 $ u, h \in C(T) $。这意味着在这种情况下，线性映射 $ \partial\varphi^{\sharp}(u) \in \mathcal{L}(C(T)) $ 就是被理解为\textbf{乘法算子}的函数 $ \cos u(\cdot) $，即

\[
C(T) \to C(T) , \quad h \mapsto [\cos u(\cdot)]h .
\]
\item

设 $ -\infty < \alpha < \beta < \infty $ 且 $ \varphi \in C^{0,p}([\alpha, \beta] \times X, F) $。再令

\begin{equation}
f(u) := \int_{\alpha}^{\beta}\varphi(t, u(t)) dt \quad \text{对 } u \in C([\alpha, \beta], X) .\tag{6.3}
\end{equation}

那么 $ f \in C^{p}(C([\alpha, \beta], X), F) $ 且

\[
\partial f(u)h = \int_{\alpha}^{\beta}\partial_{2}\varphi(t, u(t))h(t) dt \quad \text{对 } u \in C([\alpha, \beta], X) , \quad h \in C([\alpha, \beta], E) .
\]

\proofhead
令 $ T := [\alpha, \beta] $，定理 6.4 蕴含 $ \varphi^{\sharp} $ 属于 $ C^{p}(C(T, X), C(T, F)) $，且 $ \partial\varphi^{\sharp}(u) $ 是乘法算子
\[
C(T, E) \to C(T, F) , \quad h \mapsto \partial_{2}\varphi(\cdot, u(\cdot))h .
\]

由第 VI.3 节，我们知道 $ \int_{\alpha}^{\beta} $ 属于 $ \mathcal{L}(C(T, F), F) $。因此，由链式法则（定理 3.3）和例 2.3(a) 可得

\[
f = \int_{\alpha}^{\beta} \circ\, \varphi^{\sharp} \in C^{p}(C(T, X), F)
\]

且

\[
\partial f(u) = \int_{\alpha}^{\beta} \circ\, \partial\varphi^{\sharp}(u) \quad \text{对 } u \in C(T, X) .\quad\blacksquare
\]
\item

设 $ -\infty < \alpha < \beta < \infty $ 且 $ \varphi \in C^{0,p}([\alpha, \beta] \times (X \times E), F) $。那么 $ C^{1}([\alpha, \beta], X) $ 在 $ C^{1}([\alpha, \beta], E) $ 中开，并且对于由下式定义的映射

\[
\Phi(u)(t) := \varphi(t, u(t), \dot{u}(t)) \quad \text{对 } u \in C^{1}([\alpha, \beta], X) \text{ 且 } \alpha \le t \le \beta ,
\]

我们有

\[
\Phi \in C^{p}(C^{1}([\alpha, \beta], X), C([\alpha, \beta], F))
\]

且

\[
[\partial\Phi(u)h](t) = \partial_{2}\varphi(t, u(t), \dot{u}(t))h(t) + \partial_{3}\varphi(t, u(t), \dot{u}(t))\dot{h}(t)
\]

对 $ t \in [\alpha, \beta] $，$ u \in C^{1}([\alpha, \beta], X) $，$ h \in C^{1}([\alpha, \beta], E) $。

\proofhead
典范包含
\begin{equation}
i : C^{1}([\alpha, \beta], E) \to C([\alpha, \beta], E) , \quad u \mapsto u\tag{6.4}
\end{equation}

是线性且连续的。这是因为 $ C([\alpha, \beta], E) $ 在 $ C^{1}([\alpha, \beta], E) $ 上诱导的最大范数比 $ C^{1} $ 范数弱。因此，根据引理 6.1 和定理 III.2.20，

\[
C^{1}([\alpha, \beta], X) = i^{-1}(C([\alpha, \beta], X))
\]

在 $ C^{1}([\alpha, \beta], E) $ 中是开的。
由 (6.4) 可得

\[
\Psi := [u \mapsto (u, \dot{u})] \in \mathcal{L}(C^{1}(T, X), C(T, X) \times C(T, E))
\]

对 $ T := [\alpha, \beta] $。由于 $ \Phi = \varphi^{\sharp} \circ \Psi $，我们用定理 6.4 和链式法则得出结论。 \nobreak\hspace{0.35em}■

\item

设 (c) 的假设满足，且

\[
f(u) := \int_{\alpha}^{\beta}\varphi(t, u(t), \dot{u}(t)) dt \quad \text{对 } u \in C^{1}([\alpha, \beta], X) .
\]

那么

\[
f \in C^{p}(C^{1}([\alpha, \beta], X), F)
\]

且

\[
\partial f(u)h = \int_{\alpha}^{\beta}\{\partial_{2}\varphi(t, u(t), \dot{u}(t))h(t) + \partial_{3}\varphi(t, u(t), \dot{u}(t))\dot{h}(t)\} dt
\]

对 $ u \in C^{1}([\alpha, \beta], X) $ 和 $ h \in C^{1}([\alpha, \beta], E) $。

\proofhead
由于 $ f = \int_{\alpha}^{\beta} \circ\, \Phi $，这由 (c) 和链式法则得到。 \nobreak\hspace{0.35em}■
\end{items}

\end{example}

\needspace{3\baselineskip}\subsection*{含参积分的可微性}

作为上述例子的简单推论，我们将研究\textbf{含参积分}

\[
\int_{\alpha}^{\beta}\varphi(t, x) dt .
\]

在第 III 卷中，我们将在勒贝格积分理论的框架下证明下面关于含参积分的连续性与可微性的定理的一般版本。

\qyxffntextdef{1}{我们也建议你构造一个不利用例 6.6 的直接证明。}
\qyxffntextdef{2}{这里以及在下面的陈述中，我们理解关于导数的陈述适用于 $ p \ge 1 $ 的情形。}
\begin{proposition}

设 $ -\infty < \alpha < \beta < \infty $ 且 $ p \in \mathbb{N} \cup \{\infty\} $。再假设 $ \varphi $ 属于 $ C^{0,p}([\alpha, \beta] \times X, F) $，并令

\[
\Phi(x) := \int_{\alpha}^{\beta}\varphi(t, x) dt \quad \text{对 } x \in X .
\]

那么我们有 $ \Phi \in C^{p}(X, F) $ 且

\begin{equation}
\partial\Phi = \int_{\alpha}^{\beta}\partial_{2}\varphi(t, \cdot) dt .\tag{6.5}
\end{equation}

证明\qyxffnmark{1} 对于 $ x \in X $，设 $ u_{x} $ 是常值映射 $ [\alpha, \beta] \to E $，$ t \mapsto x $。那么对于由 $ x \mapsto u_{x} $ 定义的函数 $ \psi $，我们显然有

\[
\psi \in C^{\infty}(X, C([\alpha, \beta], X)) ,
\]

其中 $ \partial\psi(x)\xi = u_{\xi} $ 对 $ \xi \in E $。利用 (6.3)，我们得到 $ \Phi(x) = f(u_{x}) = f \circ \psi(x) $。因此，由例 6.6(b) 和链式法则，可得 $ \Phi \in C^{p}(X, F) $ 且\qyxffnmark{2}

\[
\partial\Phi(x)\xi = \partial f(\psi(x))\partial\psi(x)\xi = \partial f(u_{x})u_{\xi} = \int_{\alpha}^{\beta}\partial_{2}\varphi(t, x)\xi dt
\]

对 $ x \in X $ 和 $ \xi \in E $。由于 $ \partial_{2}\varphi(\cdot, x) \in C([\alpha, \beta], \mathcal{L}(E, F)) $ 对 $ x \in X $，我们有

\[
\int_{\alpha}^{\beta}\partial_{2}\varphi(t, x) dt \in \mathcal{L}(E, F) ,
\]

从而

\[
\int_{\alpha}^{\beta}\partial_{2}\varphi(t, x) dt\,\xi = \int_{\alpha}^{\beta}\partial_{2}\varphi(t, x)\xi dt \quad \text{对 } \xi \in E
\]

（见习题 VI.3.3）。由此我们得到

\[
\partial\Phi(x)\xi = \int_{\alpha}^{\beta}\partial\varphi(t, x) dt\,\xi \quad \text{对 } x \in X \text{ 且 } \xi \in E
\]

从而得到 (6.5)。 \nobreak\hspace{0.35em}■

\end{proposition}

\begin{corollary}

设命题 6.7 的假设在 $ p \le 1 $ 和 $ E := \mathbb{R}^{m} $ 下满足，并且

\[
\Psi(x) := \int_{a(x)}^{b(x)}\varphi(t, x) dt \quad \text{对 } x \in X
\]

且 $ a, b \in C(X, (\alpha, \beta)) $。那么 $ \Psi $ 属于 $ C^{p}(X, F) $，并且对 $ x \in X $ 我们有

\[
\partial\Psi(x) = \int_{a(x)}^{b(x)}\partial_{2}\varphi(t, x) dt + \varphi(b(x), x)\partial b(x) - \varphi(a(x), x)\partial a(x) .
\]

\proofhead
我们令
\[
\Phi(x, y, z) := \int_{y}^{z}\varphi(t, x) dt \quad \text{对 } (x, y, z) \in X \times (\alpha, \beta) \times (\alpha, \beta) .
\]

命题 6.7 保证 $ \Phi(\cdot, y, z) \in C^{p}(X, F) $ 使得

\[
\partial_{1}\Phi(x, y, z) := \int_{y}^{z}\partial_{2}\varphi(t, x) dt
\]

对每个 $ (x, y, z) \in X \times (\alpha, \beta) \times (\alpha, \beta) $ 成立，若 $ p = 1 $。

设 $ (x, y) \in X \times (\alpha, \beta) $。那么由定理 VI.4.12 可得 $ \Phi(x, y, \cdot) $ 属于 $ C^{p}((\alpha, \beta), F) $ 且

\[
\partial_{3}\Phi(x, y, z) = \varphi(z, x) \quad \text{对 } z \in (\alpha, \beta) .
\]

类似地，注意到由 (VI.4.3) 有 $ (x, z) \in X \times (\alpha, \beta) $，我们发现 $ \Phi(x, \cdot, z) $ 属于 $ C^{p}((\alpha, \beta), F) $ 且

\[
\partial_{2}\Phi(x, y, z) = -\varphi(y, x) \quad \text{对 } y \in (\alpha, \beta) .
\]

由此我们容易导出

\[
\partial_{k}\Phi \in C(X \times (\alpha, \beta) \times (\alpha, \beta), F) \quad \text{对 } 1 \le k \le 3
\]

并由定理 2.10 得到 $ \Phi \in C^{p}(X \times (\alpha, \beta) \times (\alpha, \beta), F) $。现在链式法则蕴含

\[
\Psi = \Phi(\cdot, a(\cdot), b(\cdot)) \in C^{p}(X, F)
\]

且

\begin{equation*}
\begin{aligned}
\partial\Psi(x) &= \partial_{1}\Phi(x, a(x), b(x)) + \partial_{2}\Phi(x, a(x), b(x))\partial a(x) \\ &\quad + \partial_{3}\Phi(x, a(x), b(x))\partial b(x) \\ &= \int_{a(x)}^{b(x)}\partial_{2}\varphi(t, x) dt + \varphi(b(x), x)\partial b(x) - \varphi(a(x), x)\partial a(x)
\end{aligned}
\end{equation*}

对 $ x \in X $。 \nobreak\hspace{0.35em}■

\end{corollary}

由 $ \partial\Psi $ 的公式，简单的归纳法表明，若 $ \varphi $ 属于 $ C^{p-1,p}([0, 1] \times X, F) $，则 $ \Psi \in C^{p}(X, F) $ 对 $ p \ge 2 $ 成立，其中 $ C^{p-1,p} $ 按显然的方式定义。

\needspace{3\baselineskip}\subsection*{变分问题}

在本节剩余部分，我们假设\footnote{在大多数应用中，$ L $ 只定义在形如 $ [\alpha, \beta] \times X \times Y $ 的集合上，其中 $ Y $ 在 $ \mathbb{R}^{m} $ 中开（见习题 10）。我们把纳入这种情形所需的简单修改留给你。}

• $ -\infty < \alpha < \beta < \infty $；
  $ X $ 在 $ \mathbb{R}^{m} $ 中开；
  $ a, b \in X $；
  $ L \in C^{0,1}([\alpha, \beta] \times (X \times \mathbb{R}^{m}), \mathbb{R}) $。

那么积分

\begin{equation}
f(u) := \int_{\alpha}^{\beta}L(t, u(t), \dot{u}(t)) dt\tag{6.6}
\end{equation}

对每个 $ u \in C^{1}([\alpha, \beta], X) $ 存在。当我们试图在 $ X $ 中所有的 $ C^{1} $ 函数 $ u $ 上最小化该积分时，定义 (6.6) 产生一个\textbf{自由边界条件的变分问题}。我们把该问题写为

\begin{equation}
\int_{\alpha}^{\beta}L(t, u, \dot{u}) dt \Rightarrow \mathrm{Min} \quad \text{对 } u \in C^{1}([\alpha, \beta], X) .\tag{6.7}
\end{equation}

于是，问题 (6.7) 变成在集合 $ U := C^{1}([\alpha, \beta], X) $ 上最小化 $ f $ 的问题。

当我们在一个约束集上最小化函数 $ f $ 时，定义 (6.6) 也可以是一个\textbf{固定边界条件的变分问题}的一部分

\[
U_{a,b} := \{ u \in C^{1}([\alpha, \beta], X) ;\ u(\alpha) = a ,\ u(\beta) = b \} .
\]

我们把该问题写为

\begin{equation}
\int_{\alpha}^{\beta}L(t, u, \dot{u}) dt \Rightarrow \mathrm{Min} \quad \text{对 } u \in C^{1}([\alpha, \beta], X) , \quad u(\alpha) = a , \quad u(\beta) = b .\tag{6.8}
\end{equation}

（6.7）和（6.8）的每个解分别称为变分问题（6.7）和（6.8）的\textbf{极值曲线}\footnote{在变分法中，我们也用极值曲线这一术语来描述欧拉–拉格朗日方程的解。}。

为了有效地处理这个问题，我们需要如下内容：

\begin{lemma}

令 $ E := C^{1}([\alpha, \beta], \mathbb{R}^{m}) $ 且

\[
E_{0} := \{ u \in E ;\ u(\alpha) = u(\beta) = 0 \} =: C_{0}^{1}([\alpha, \beta], \mathbb{R}^{m}) .
\]

那么 $ E_{0} $ 是 $ E $ 的闭向量子空间，因此它本身是巴拿赫空间。对于 $ \overline{u} \in U_{a,b} $，$ U_{a,b} - \overline{u} $ 在 $ E_{0} $ 中开。

\proofhead
显然 $ E_{0} $ 是 $ E $ 的闭向量子空间。由于根据例 6.6(c)，$ U $ 在 $ E $ 中开，并且由于 $ U_{a,b} = U \cap (\overline{u} + E_{0}) $，我们知道 $ U_{a,b} $ 在仿射空间 $ \overline{u} + E_{0} $ 中开。现在显然 $ U_{a,b} - \overline{u} $ 在 $ E_{0} $ 中开。 \nobreak\hspace{0.35em}■
\end{lemma}

\begin{lemma}

\begin{subitems}
\item

$ f \in C^{1}(U, \mathbb{R}) $ 且

\[
\partial f(u)h = \int_{\alpha}^{\beta}\{\partial_{2}L(t, u(t), \dot{u}(t))h(t) + \partial_{3}L(t, u(t), \dot{u}(t))\dot{h}(t)\} dt
\]

对 $ u \in U $ 和 $ h \in E $。

\item

设 $ \overline{u} \in U_{a,b} $，并令 $ g(v) := f(\overline{u} + v) $ 对 $ v \in V := U_{a,b} - \overline{u} $。那么 $ g $ 属于 $ C^{1}(V, \mathbb{R}) $ 且

\begin{equation}
\partial g(v)h = \partial f(\overline{u} + v)h \quad \text{对 } v \in V \text{ 且 } h \in E_{0} .\tag{6.9}
\end{equation}

\proofhead
(i) 是例 6.6(d) 的推论，(ii) 由 (i) 得到。 \nobreak\hspace{0.35em}■
\end{subitems}

\end{lemma}

\begin{lemma}

对于 $ u \in C([\alpha, \beta], \mathbb{R}^{m}) $，设

\begin{equation}
\int_{\alpha}^{\beta}(u(t)|v(t)) dt = 0 \quad \text{对 } v \in C_{0}^{1}([\alpha, \beta], \mathbb{R}^{m}) .\tag{6.10}
\end{equation}

那么 $ u = 0 $。

\proofhead
设 $ u \ne 0 $。那么存在 $ k \in \{1, \ldots, m\} $ 和 $ t_{0} \in (\alpha, \beta) $ 使得 $ u_{k}(t_{0}) \ne 0 $。只需考虑 $ u_{k}(t_{0}) > 0 $ 的情形。由 $ u_{k} $ 的连续性，存在 $ \alpha < \alpha' < \beta' < \beta $ 使得对 $ t \in (\alpha', \beta') $ 有 $ u_{k}(t) > 0 $（见例 III.1.3(d)）。

取 $ v \in C_{0}^{1}([\alpha, \beta], \mathbb{R}^{m}) $ 使得对 $ j \ne k $ 有 $ v_{j} = 0 $，由 (6.10) 可得
\begin{equation}
\int_{\alpha}^{\beta}u_{k}(t)w(t) dt = 0 \quad \text{对 } w \in C_{0}^{1}([\alpha, \beta], \mathbb{R}) .\tag{6.11}
\end{equation}

现在设 $ w \in C_{0}^{1}([\alpha, \beta], \mathbb{R}) $ 使得对 $ t \in (\alpha', \beta') $ 有 $ w(t) > 0 $，并且对区间 $ (\alpha', \beta') $ 之外的 $ t $ 有 $ w(t) = 0 $（见习题 7）。那么我们有

\[
\int_{\alpha}^{\beta}u_{k}(t)w(t) dt = \int_{\alpha'}^{\beta'}u_{k}(t)w(t) dt > 0 ,
\]

这与 (6.11) 矛盾。 \nobreak\hspace{0.35em}■

\end{lemma}

\needspace{3\baselineskip}\subsection*{欧拉–拉格朗日方程}

现在我们可以证明变分法的基本结果：

\begin{theorem}

设 $ u $ 是自由边界条件的变分问题 (6.7) 或（备选地）固定边界条件的变分问题 (6.8) 的极值曲线。再设

\begin{equation}
[t \mapsto \partial_{3}L(t, u(t), \dot{u}(t))] \in C^{1}([\alpha, \beta], \mathbb{R}) .\tag{6.12}
\end{equation}

那么 $ u $ 满足\textbf{欧拉–拉格朗日方程}

\begin{equation}
\left[\partial_{3}L(\cdot, u, \dot{u})\right]^{\cdot} = \partial_{2}L(\cdot, u, \dot{u}) .\tag{6.13}
\end{equation}

在自由边界条件的问题中，$ u $ 还满足\textbf{自然边界条件}

\begin{equation}
\partial_{3}L(\alpha, u(\alpha), \dot{u}(\alpha)) = \partial_{3}L(\beta, u(\beta), \dot{u}(\beta)) = 0 .\tag{6.14}
\end{equation}

对于固定边界条件，$ u $ 满足

\begin{equation}
u(\alpha) = a \text{ 且 } u(\beta) = b .\tag{6.15}
\end{equation}

证明

\begin{subitems}
\item

我们首先考虑问题 (6.7)。由假设，定义在 (6.6) 中的函数 $ f $ 在 $ u \in U $ 处取最小值。于是由引理 6.10(i)、定理 3.13 和命题 2.5 可得

\begin{equation}
\int_{\alpha}^{\beta}\{\partial_{2}L(t, u(t), \dot{u}(t))h(t) + \partial_{3}L(t, u(t), \dot{u}(t))\dot{h}(t)\} dt = 0\tag{6.16}
\end{equation}

对 $ h \in E = C^{1}([\alpha, \beta], \mathbb{R}^{m}) $。利用假设 (6.12)，我们可以分部积分得到

\[
\int_{\alpha}^{\beta}\partial_{3}L(\cdot, u, \dot{u})\dot{h} dt = \partial_{3}L(\cdot, u, \dot{u})h\big|_{\alpha}^{\beta} - \int_{\alpha}^{\beta}\left[\partial_{3}L(\cdot, u, \dot{u})\right]^{\cdot}h dt .
\]

然后由 (6.16) 我们得到

\begin{equation}
\int_{\alpha}^{\beta}\left\{\partial_{2}L(\cdot, u, \dot{u}) - \left[\partial_{3}L(\cdot, u, \dot{u})\right]^{\cdot}\right\}h dt + \partial_{3}L(\cdot, u, \dot{u})h\big|_{\alpha}^{\beta} = 0\tag{6.17}
\end{equation}

对 $ h \in E $。特别地，我们有

\begin{equation}
\int_{\alpha}^{\beta}\left\{\partial_{2}L(\cdot, u, \dot{u}) - \left[\partial_{3}L(\cdot, u, \dot{u})\right]^{\cdot}\right\}h dt = 0 \quad \text{对 } h \in E_{0} ,\tag{6.18}
\end{equation}

因为

\begin{equation}
\partial_{3}L(\cdot, u, \dot{u})h\big|_{\alpha}^{\beta} = 0 \quad \text{对 } h \in E_{0} .\tag{6.19}
\end{equation}

现在欧拉–拉格朗日方程 (6.13) 由 (6.18) 和引理 6.11 蕴含。于是，由 (6.17) 可得

\begin{equation}
\partial_{3}L(\cdot, u, \dot{u})h\big|_{\alpha}^{\beta} = 0 \quad \text{对 } h \in E_{0} .\tag{6.20}
\end{equation}

对于 $ \xi, \eta \in \mathbb{R}^{m} $，我们令

\[
h_{\xi,\eta}(t) := \frac{\beta - t}{\beta - \alpha}\xi + \frac{t - \alpha}{\beta - \alpha}\eta \quad \text{对 } \alpha \le t \le \beta .
\]

那么 $ h_{\xi,\eta} $ 属于 $ E $，并且由 (6.20) 对 $ h = h_{\xi,\eta} $ 我们得到

\begin{equation}
\partial_{3}L(\beta, u(\beta), \dot{u}(\beta))\eta - \partial_{3}L(\alpha, u(\alpha), \dot{u}(\alpha))\xi = 0 .\tag{6.21}
\end{equation}

因为 (6.21) 对每个 $ \xi, \eta \in \mathbb{R}^{m} $ 都成立，所以自然边界条件 (6.14) 满足。

\item

现在我们考虑问题 (6.8)。我们固定一个 $ \overline{u} \in U_{a,b} $ 并令 $ v := u - \overline{u} $。那么由引理 6.10(ii) 可得，$ g $ 在巴拿赫空间 $ E_{0} $ 的开集 $ V := U_{a,b} - \overline{u} $ 的点 $ v $ 处取最小值。因此由引理 6.10(ii)、定理 3.13 和命题 2.5 可得 (6.16) 在这种情况下也成立，尽管只对 $ h \in E_{0} $ 成立。另一方面，由于 (6.19)，我们由 (6.17) 得到 (6.18) 的有效性，从而得到欧拉–拉格朗日方程。陈述 (6.15) 是平凡的。 \nobreak\hspace{0.35em}■

\end{subitems}

\end{theorem}

\begin{remark}

\begin{items}
\item

在定理 6.12 中，假设变分问题有解，即存在极值曲线 $ u $。在 (6.12) 中我们作了一个额外的假设，对极值曲线 $ u $ 施加一个隐含的正则性条件。只有在这个额外假设下，欧拉–拉格朗日方程才是存在极值曲线 $ u $ 的必要条件。
\textbf{变分法的间接方法}赋予欧拉–拉格朗日方程核心地位。在其中，我们试图验证存在函数 $ u $ 满足方程 (6.13) 和边界条件 (6.14) 或 (6.15)。由于欧拉–拉格朗日方程是（一般为非线性的）微分方程，这引向微分方程的\textbf{边值问题}这一主题。如果能够用边值问题的方法证明存在解，那么我们就在第二步从中选出那些实际上是极值曲线的解。
与之相对的是\textbf{变分法的直接方法}，在其中人们试图直接证明函数 $ f $ 在点 $ u \in U $（或函数 $ g $ 在点 $ v \in V $）处取最小值。那么定理 3.13 和命题 2.5 表明 $ \partial f(u) = 0 $（或 $ \partial g(v) = 0 $），从而关系式 (6.16) 对所有 $ h \in E_{1} $（或所有 $ h \in E_{0} $）成立。如果辅助的“正则性问题”能够解决——即能够证明极值曲线 $ u $ 满足条件 (6.12)——那么它也满足欧拉–拉格朗日微分方程。
因此定理 6.12 有两种可能的应用。在第一种情形中，我们可以用微分方程边值问题的理论来解欧拉–拉格朗日方程，即最小化积分。这是变分法的经典方法。
然而，一般说来，证明边值问题的存在性结果是非常困难的。因此第二种情形采取另一个观点：我们给定一个边值问题，然后试图找出产生它的变分问题，使得该微分方程可以解释为欧拉–拉格朗日方程。如果这样的变分问题存在，那么我们就试图证明它有一个满足正则性条件 (6.12) 的极值解。这样的证明就蕴含原来的边值问题也有解。这被证明是研究边值问题的一种有力方法。
遗憾的是，细节超出了本书的范围。欲知更多，你应查阅文献或修读变分方法、微分方程和（非线性）泛函分析的课程。

\item

显然，如果 $ u $ 只是 $ f $（或 $ g $）的临界点，即如果 $ \partial f(u) = 0 $（或 $ \partial g(v) = 0 $，其中 $ v := u - \overline{u} $），定理 6.12 也成立。在这种情况下，人们称 $ u $ 是积分 (6.6) 的\textbf{驻值}。

\item

设 $ E $ 是巴拿赫空间，$ Z $ 在 $ E $ 中开，$ F : Z \to \mathbb{R} $。如果 $ F $ 的方向导数对某个 $ h \in E \setminus \{0\} $ 在 $ z_{0} \in Z $ 存在，那么（在变分法中）我们称 $ D_{h}F(z_{0}) $ 为 $ F $ 沿方向 $ h $ 的第一变分，并把它写为 $ \delta F(z_{0}; h) $。如果 $ F $ 的第一变分在每个方向都存在，我们称

\[
\delta F(z_{0}) := \delta F(z_{0}; \cdot) : E \to \mathbb{R}
\]

为 $ F $ 在 $ z_{0} $ 处的\textbf{第一变分}。
如果 $ F $ 在 $ z_{0} \in Z $ 处有局部极值且第一变分在该点存在，那么 $ \delta F(z_{0}) = 0 $。换句话说：第一变分的消失是存在局部极值的必要条件。

\proofhead
这是定理 3.13 的重新表述。 \nobreak\hspace{0.35em}■
\item

在我们的假设下，$ f $（或 $ g $）是连续可微的。因此第一变分存在，并与 $ \partial f $（或 $ \partial g $）一致。然而，在许多情况下，欧拉–拉格朗日方程已经可以从只保证第一变分存在的较弱假设推出。

\item

在变分法中，通常分别把 $ \partial_{2}L $ 写成 $ L_{u} $，把 $ \partial_{3}L $ 写成 $ L_{\dot{u}} $。那么欧拉–拉格朗日方程取缩写形式

\begin{equation}
\frac{d}{dt}(L_{\dot{u}}) = L_{u} ,\tag{6.22}
\end{equation}

即它失去了自变量。如果再假设 $ L $ 和极值曲线 $ u $ 是二次可微的，我们可以把 (6.22) 写为 $ L_{t,\dot{u}} + L_{u,\dot{u}}\dot{u} + L_{\dot{u},\dot{u}}\ddot{u} = L_{u} $，其中 $ L_{t,\dot{u}} = \partial_{1}\partial_{3}L $ 等等。 \nobreak\hspace{0.35em}■

\end{items}

\end{remark}

\needspace{3\baselineskip}\subsection*{经典力学}

变分法在物理学中很重要。为了说明这一点，我们考虑一个具有 $ m $ 个自由度的有质量质点系统，它由\textbf{（广义）坐标} $ q = (q^{1}, \ldots, q^{m}) $ 确定。经典力学的基本问题是，已知系统在时刻 $ t_{0} $ 的初始位置 $ q_{0} $，确定系统在给定时刻 $ t $ 的状态 $ q(t) $。

我们用 $ \dot{q} := dq/dt $ 表示（广义）\textbf{速度坐标}，用 $ T(t, q, \dot{q}) $ 表示\textbf{动能}，用 $ U(t, q) $ 表示\textbf{势能}。然后我们作出\textbf{哈密顿最小作用量原理}的基本假设。它说，在每一对时刻 $ t_{0} $ 和 $ t_{1} $ 之间，系统以最小化

\[
\int_{t_{0}}^{t_{1}}[T(t, q, \dot{q}) - U(t, q)] dt ,
\]

的方式从 $ q_{0} $ 到 $ q_{1} $；也就是说，系统实际选择的路径是在所有从 $ t_{0} $ 时刻的 $ q_{0} $ 出发、在 $ t_{1} $ 时刻到达 $ q_{1} $ 的（虚拟）路径中使这个积分最小化的那条。在此背景下，$ L := T - U $ 称为系统的\textbf{拉格朗日量}。

哈密顿原理意味着系统所走的路径 $ \{ q(t) ;\ t_{0} \le t \le t_{1} \} $ 是固定边界条件变分问题的极值

\[
\int_{t_{0}}^{t_{1}}L(t, q, \dot{q}) dt \Rightarrow \mathrm{Min} \quad \text{对 } q \in C^{1}([t_{0}, t_{1}], \mathbb{R}^{m}) , \quad q(t_{0}) = q_{0} , \quad q(t_{1}) = q_{1} .
\]

然后由定理 6.12 可得 $ q $ 满足欧拉–拉格朗日方程

\[
\frac{d}{dt}(L_{\dot{q}}) = L_{q}
\]

如果适当的正则性条件也满足。

在物理学、工程学以及使用变分方法的其他领域（例如经济学）中，通常预设正则性条件满足，并假设欧拉–拉格朗日方程成立（并且把极值路径的存在性视为物理上显然的）。然后利用欧拉–拉格朗日方程来提取极值路径的形式，进而理解系统的行为。

\begin{example}

\begin{items}
\item

我们考虑一个质量为（正的）$ m $ 的无约束质点在与时间无关的势场 $ U(t, q) = U(q) $ 作用下的三维运动。我们假设动能只依赖于 $ \dot{q} $，且由下式给出

\[
T(\dot{q}) = m\left|\dot{q}\right|^{2}/2 .
\]

欧拉–拉格朗日微分方程取形式

\begin{equation}
m\ddot{q} = -\nabla U(q) .\tag{6.23}
\end{equation}

由于 $ \ddot{q} $ 是系统的\textbf{加速度}，所以 (6.23) 与受\textbf{保守}\footnote{如果一个力场有势，则称其为保守的（见注 VIII.4.10(c)）。}力 $ -\nabla U $ 作用的质点的运动所满足的\textbf{牛顿运动方程}相同。

\proofhead
我们用里斯表示定理把 $ (\mathbb{R}^{3})' = \mathcal{L}(\mathbb{R}^{3}, \mathbb{R}) $ 与 $ \mathbb{R}^{3} $ 等同。然后由 $ L(t, q, \dot{q}) = T(\dot{q}) - U(q) $ 我们得到关系式
\[
\partial_{2}L(t, q, \dot{q}) = -\nabla U(q) \text{ 且 } \partial_{3}L(t, q, \dot{q}) = \nabla T(\dot{q}) = m\dot{q} ,
\]

它们证明了结论。 \nobreak\hspace{0.35em}■

\item

作为 (a) 的推广，我们考虑势场中 $ N $ 个无约束的有质量质点的运动。我们用 $ x = (x_{1}, \ldots, x_{N}) $ 表示坐标，其中 $ x_{j} \in \mathbb{R}^{3} $ 确定第 $ j $ 个质点的位置。再设 $ X $ 在 $ \mathbb{R}^{3N} $ 中开且 $ U \in C^{1}(X, \mathbb{R}) $。令 $ m_{j} $ 为第 $ j $ 个质点的质量，系统的动能为

\[
T(\dot{x}) := \sum_{j=1}^{N}\frac{m_{j}}{2}\left|\dot{x}_{j}\right|^{2} .
\]

那么我们得到一个由 $ N $ 个欧拉–拉格朗日方程构成的方程组：

\[
-m_{j}\ddot{x}_{j} = \nabla_{x_{j}}U(x) \quad \text{对 } 1 \le j \le N .
\]

这里，$ \nabla_{x_{j}} $ 表示 $ U $ 关于变量 $ x_{j} \in \mathbb{R}^{3} $ 的梯度。 \nobreak\hspace{0.35em}■

\end{items}

\end{example}

\exerciseshead

\begin{exercises}
\item

设 $ I = [\alpha, \beta] $，其中 $ -\infty < \alpha < \beta < \infty $，$ k \in C(I \times I, \mathbb{R}) $，$ \varphi \in C^{0,1}(I \times E, F) $。又令

\[
\Phi(u)(t) := \int_{\alpha}^{\beta} k(t, s)\varphi(s, u(s)) \, ds \quad \text{对 } t \in I
\]

且 $ u \in C(I, E) $。证明 $ \Phi \in C^{1}(C(I, E), C(I, F)) $，且

\[
(\partial\Phi(u)h)(t) = \int_{\alpha}^{\beta} k(t, s)\partial_{2}\varphi(s, u(s))h(s) \, ds \quad \text{对 } t \in I
\]

且 $ u, h \in C(I, E) $。

\item

设 $ I $ 与 $ J := [\alpha, \beta] $ 是紧致完全区间，且 $ f \in C(I \times J, E) $。证明

\[
\int_{J} \left( \int_{I} f(s, t) \, ds \right) dt = \int_{I} \left( \int_{J} f(s, t) \, dt \right) ds .
\]

（提示：考虑

\[
J \to E , \quad y \mapsto \int_{I} \left( \int_{\alpha}^{y} f(s, t) \, dt \right) ds
\]

并应用命题 6.7 与推论 VI.4.14。）

\item

对 $ f \in C([\alpha, \beta], E) $，证明

\[
\int_{\alpha}^{s} \left( \int_{\alpha}^{t} f(\tau) \, d\tau \right) dt = \int_{\alpha}^{s} (s - t)f(t) \, dt \quad \text{对 } s \in [\alpha, \beta] .
\]
\item

设 $ I $ 与 $ J $ 是紧致完全区间，且 $ \rho \in C(I \times J, \mathbb{R}) $。验证

\[
u(x, y) := \int_{I} \left( \int_{J} \log\left( (x - s)^{2} + (y - t)^{2} \right)\rho(s, t) \, dt \right) ds \quad \text{对 } (x, y) \in \mathbb{R}^{2}\setminus(I \times J)
\]

是调和的。

\item

定义

\begin{equation*}
h : \mathbb{R} \to [0, 1) , \quad s \mapsto \begin{cases} e^{-1/s^{2}} , & s > 0 , \\ 0 , & s \le 0 , \end{cases}
\end{equation*}

且对 $ -\infty < \alpha < \beta < \infty $，定义 $ k : \mathbb{R} \to \mathbb{R} $，$ s \mapsto h(s - \alpha)h(\beta - s) $，以及

\[
\ell : \mathbb{R} \to \mathbb{R} , \quad t \mapsto \int_{\alpha}^{t} k(s) \, ds / \int_{\alpha}^{\beta} k(s) \, ds .
\]

然后证明

\begin{items}
\item

$ \ell \in C^{\infty}(\mathbb{R}, \mathbb{R}) $；

\item

$ \ell $ 是递增的；

\item

对 $ t \le \alpha $ 有 $ \ell(t) = 0 $，且对 $ t \ge \beta $ 有 $ \ell(t) = 1 $。

\end{items}

\item

设 $ -\infty < \alpha_{j} < \beta_{j} < \infty $，其中 $ j = 1, \ldots, m $，且 $ A := \prod_{j=1}^{m}(\alpha_{j}, \beta_{j}) \subset \mathbb{R}^{m} $。证明存在 $ g \in C^{\infty}(\mathbb{R}^{m}, \mathbb{R}) $，使得对 $ x \in A $ 有 $ g(x) > 0 $，且对 $ x \in A^{c} $ 有 $ g(x) = 0 $。

（提示：考虑

\[
g(x_{1}, \ldots, x_{m}) := k_{1}(x_{1}) \cdots k_{m}(x_{m}) \quad \text{对 } (x_{1}, \ldots, x_{m}) \in \mathbb{R}^{m} ,
\]

其中 $ k_{j}(t) := h(t - \alpha_{j})h(\beta_{j} - t) $，$ t \in [\alpha_{j}, \beta_{j}] $，$ j = 1, \ldots, m $。）

\item

设 $ K \subset \mathbb{R}^{m} $ 是紧致的，$ U $ 是 $ K $ 的开邻域。则证明存在 $ f \in C^{\infty}(\mathbb{R}^{m}, \mathbb{R}) $，使得

\begin{items}
\item

对 $ x \in \mathbb{R}^{m} $ 有 $ f(x) \in [0, 1] $；

\item

对 $ x \in K $ 有 $ f(x) = 1 $；

\item

对 $ x \in U^{c} $ 有 $ f(x) = 0 $。
（提示：对 $ x \in K $，存在 $ \varepsilon_{x} > 0 $ 使得 $ A_{x} := \overline{\mathbb{B}}(x, \varepsilon_{x}) \subset U $。然后选取 $ g_{x} \in C^{\infty}(\mathbb{R}^{m}, \mathbb{R}) $ 使得对 $ y \in A_{x} $ 有 $ g_{x}(y) > 0 $（见习题 6），并选取 $ x_{0}, \ldots, x_{n} \in K $ 使得 $ K \subset \bigcup_{j=0}^{n} A_{x_{j}} $。则 $ G := g_{x_{0}} + \cdots + g_{x_{n}} $ 属于 $ C^{\infty}(\mathbb{R}^{m}, \mathbb{R}) $，且数 $ \delta := \min_{x \in K} G(x) > 0 $。最后，令 $ \ell $ 如习题 5 中所定义，$ \alpha = 0 $，$ \beta = \delta $，并令 $ f := \ell \circ G $。）

\end{items}

\item

定义 $ E_{0} := \left\{ u \in C^{1}([0, 1], \mathbb{R}) ;\ u(0) = 0, u(1) = 1 \right\} $，并对 $ u \in E_{0} $ 令 $ f(u) := \int_{0}^{1} [t\dot{u}(t)]^{2} \, dt $。则证明 $ \inf\left\{ f(u) ;\ u \in E_{0} \right\} = 0 $，尽管不存在 $ u_{0} \in E_{0} $ 使得 $ f(u_{0}) = 0 $。

\item

对 $ a, b \in \mathbb{R}^{m} $，考虑具有固定边界条件的变分问题

\begin{equation}
\int_{\alpha}^{\beta} |\dot{u}(t)| \, dt \Rightarrow \mathrm{Min} \quad \text{对 } u \in C^{1}([\alpha, \beta], \mathbb{R}^{m}) , \quad u(\alpha) = a , \quad u(\beta) = b .\tag{6.24}
\end{equation}

\begin{items}
\item

求 (6.24) 的欧拉–拉格朗日方程的所有具有 $ |\dot{u}(t)| = \text{常数} $ 性质的解。

\item

当 $ m = 2 $ 且在 $ \dot{u}_{1}(t) \ne 0 $ 与 $ \dot{u}_{2}(t) \ne 0 $（$ t \in [\alpha, \beta] $）的假设下，(6.24) 的欧拉–拉格朗日方程如何写？

\end{items}

\item

求下列情形的欧拉–拉格朗日方程：

\begin{items}
\item

$ \int_{\alpha}^{\beta} u\sqrt{1 + \dot{u}^{2}} \, dt \Rightarrow \mathrm{Min} \quad \text{对 } u \in C^{1}([\alpha, \beta], \mathbb{R}) $；

\item

$ \int_{\alpha}^{\beta} \sqrt{(1 + \dot{u}^{2})/u} \, dt \Rightarrow \mathrm{Min} \quad \text{对 } u \in C^{1}([\alpha, \beta], (0, \infty)) $。

\end{items}

\item

设 $ f $ 由 (6.6) 定义，其中 $ L \in C^{0,2}([\alpha, \beta] \times (X \times \mathbb{R}^{m}), \mathbb{R}) $。证明 $ f \in C^{2}(U, \mathbb{R}) $，且

\begin{equation*}
\begin{aligned}
\partial^{2} f(u)[h, k] &= \int_{\alpha}^{\beta} \Big\{ \partial_{2}^{2} L(\cdot, u, \dot{u})[h, k] + \partial_{2}\partial_{3} L(\cdot, u, \dot{u})[h, \dot{k}] \\ &\qquad{}+ \partial_{2}\partial_{3} L(\cdot, u, \dot{u})[\dot{h}, k] + \partial_{3}^{2} L(\cdot, u, \dot{u})[\dot{h}, \dot{k}] \Big\} dt .
\end{aligned}
\end{equation*}

对 $ h, k \in E_{1} $。由此导出一个充分条件，以保证欧拉–拉格朗日方程的解 $ u $ 是 $ f $ 的局部极小值。

\item

设 $ T(\dot{q}) := m|\dot{q}|^{2}/2 $，且对 $ q \in \mathbb{R}^{3} $ 有 $ U(t, q) = U(q) $。证明\textbf{总能量} $ E := T + U $ 沿变分问题

\[
\int_{t_{0}}^{t_{1}} [T(\dot{q}) - U(q)] \, dt \Rightarrow \mathrm{Min} \quad \text{对 } q \in C^{2}([t_{0}, t_{1}], \mathbb{R}^{3}) .
\]

的欧拉–拉格朗日方程的每个解 $ q $ 保持常数。

\end{exercises}

\section{逆映射}

设 $ J $ 是 $ \mathbb{R} $ 中的开区间，且 $ f : J \to \mathbb{R} $ 可微。又设 $ a \in J $ 且 $ f'(a) \ne 0 $。则线性近似

\[
\mathbb{R} \to \mathbb{R} , \quad x \mapsto f(a) + f'(a)(x - a)
\]

在 $ a $ 处可逆。这在局部也成立，即存在 $ \varepsilon > 0 $ 使得 $ f $ 在 $ X := (a - \varepsilon, a + \varepsilon) $ 上可逆。此外，$ Y = f(X) $ 是开区间，且逆映射 $ f^{-1} : Y \to \mathbb{R} $ 可微，并有

\[
(f^{-1})'(f(x)) = (f'(x))^{-1} \qquad \text{对 } x \in X
\]

（与定理 IV.1.8 比较）。

在本节中，我们对多元函数建立这一情形的自然推广。

\needspace{3\baselineskip}\subsection*{线性映射之逆的导数}

下面，设

• $ E $ 与 $ F $ 是域 $ \mathbb{K} $ 上的巴拿赫空间。

我们要确定映射

\[
\operatorname{inv} : \mathcal{L}\mathrm{is}(E, F) \to \mathcal{L}(F, E) , \quad A \mapsto A^{-1}
\]

的可微性，并在必要时求出其导数。为使这一问题有意义，我们必须先验证 $ \mathcal{L}\mathrm{is}(E, F) $ 是巴拿赫空间 $ \mathcal{L}(E, F) $ 的开子集。为证明这一点，我们先证明一个关于巴拿赫代数 $ \mathcal{L}(E) $ 中“几何级数”的定理。这里令 $ I := 1_{E} $。

\begin{proposition}

设 $ A \in \mathcal{L}(E) $ 且 $ \|A\| < 1 $。则 $ I - A $ 属于 $ \mathcal{L}\mathrm{aut}(E) $，$ (I - A)^{-1} = \sum_{k=0}^{\infty} A^{k} $，且 $ \|(I - A)^{-1}\| \le (1 - \|A\|)^{-1} $。

\proofhead
我们在巴拿赫代数 $ \mathcal{L}(E) $ 中考虑几何级数 $ \sum A^{k} $。因为 $ \|A^{k}\| \le \|A\|^{k} $ 且
\begin{equation}
\sum_{k=0}^{\infty} \|A\|^{k} = 1/(1 - \|A\|) ,\tag{7.1}
\end{equation}

由优级数判别法（定理 II.8.3）知，$ \sum A^{k} $ 在 $ \mathcal{L}(E) $ 中绝对收敛。特别地，级数的值

\[
B := \sum_{k=0}^{\infty} A^{k}
\]

是 $ \mathcal{L}(E) $ 中良定义的元素。显然 $ AB = BA = \sum_{k=1}^{\infty} A^{k} $。因此 $ (I - A)B = B(I - A) = I $，从而 $ (I - A)^{-1} = B $。最后，由注 II.8.2(c) 与 (7.1) 给出估计

\[
\|(I - A)^{-1}\| = \|B\| \le (1 - \|A\|)^{-1} .\quad\blacksquare
\]
\end{proposition}

\qyxffntextdef{1}{由注 2.2(e)，当 $ E = F = \mathbb{K} $ 时，此公式化为 $ (1/z)' = -1/z^{2} $。}
\begin{proposition}

\begin{subitems}
\item

$ \mathcal{L}\mathrm{is}(E, F) $ 在 $ \mathcal{L}(E, F) $ 中是开集。

\item

映射

\[
\operatorname{inv} : \mathcal{L}\mathrm{is}(E, F) \to \mathcal{L}(F, E) , \quad A \mapsto A^{-1}
\]

无穷多次连续可微，且\qyxffnmark{1}

\begin{equation}
\partial\operatorname{inv}(A)B = -A^{-1}BA^{-1} \qquad \text{对 } A \in \mathcal{L}\mathrm{is}(E, F) \text{ 且 } B \in \mathcal{L}(E, F) .\tag{7.2}
\end{equation}

\proofhead
(i) 设 $ A_{0} \in \mathcal{L}\mathrm{is}(E, F) $，$ A \in \mathcal{L}(E, F) $ 满足 $ \|A - A_{0}\| < 1/\|A_{0}^{-1}\| $。因为
\begin{equation}
A = A_{0} + A - A_{0} = A_{0}\left[I + A_{0}^{-1}(A - A_{0})\right] ,\tag{7.3}
\end{equation}

只需验证 $ I + A_{0}^{-1}(A - A_{0}) $ 属于 $ \mathcal{L}\mathrm{aut}(E) $。因为

\[
\|-A_{0}^{-1}(A - A_{0})\| \le \|A_{0}^{-1}\|\|A - A_{0}\| < 1 ,
\]

这由命题 7.1 得到。因此 $ \mathcal{L}(E, F) $ 中以 $ A_{0} $ 为中心、以 $ 1/\|A_{0}^{-1}\| $ 为半径的开球属于 $ \mathcal{L}\mathrm{is}(E, F) $。

\end{subitems}

\begin{subitems}
\item[（ii）]

我们有 $ 2\|A - A_{0}\| < 1/\|A_{0}^{-1}\| $。由 (7.3)，我们得到

\[
A^{-1} = \left[I + A_{0}^{-1}(A - A_{0})\right]^{-1}A_{0}^{-1} ,
\]

从而

\[
A^{-1} - A_{0}^{-1} = \left[I + A_{0}^{-1}(A - A_{0})\right]^{-1}\left(I - \left[I + A_{0}^{-1}(A - A_{0})\right]\right)A_{0}^{-1} ,
\]

即

\begin{equation}
\operatorname{inv}(A) - \operatorname{inv}(A_{0}) = -\left[I + A_{0}^{-1}(A - A_{0})\right]^{-1}A_{0}^{-1}(A - A_{0})A_{0}^{-1} .\tag{7.4}
\end{equation}

由此及命题 7.1，我们推出

\begin{equation}
\|\operatorname{inv}(A) - \operatorname{inv}(A_{0})\| \le \frac{\|A_{0}^{-1}\|^{2}\|A - A_{0}\|}{1 - \|A_{0}^{-1}(A - A_{0})\|} < 2\|A_{0}^{-1}\|^{2}\|A - A_{0}\| .\tag{7.5}
\end{equation}

因此 $ \operatorname{inv} $ 连续。

\item[（iii）]

下面证明 $ \operatorname{inv} $ 可微。设 $ A \in \mathcal{L}\mathrm{is}(E, F) $。对满足 $ \|B\| < 1/\|A^{-1}\| $ 的 $ B \in \mathcal{L}(E, F) $，由 (i) 知 $ A + B $ 属于 $ \mathcal{L}\mathrm{is}(E, F) $，且

\[
(A + B)^{-1} - A^{-1} = (A + B)^{-1}\left[A - (A + B)\right]A^{-1} = -(A + B)^{-1}BA^{-1} .
\]

这给出

\[
\operatorname{inv}(A + B) - \operatorname{inv}(A) + A^{-1}BA^{-1} = \left[\operatorname{inv}(A) - \operatorname{inv}(A + B)\right]BA^{-1} .
\]

于是，

\[
\operatorname{inv}(A + B) - \operatorname{inv}(A) + A^{-1}BA^{-1} = o(\|B\|) \quad (B \to 0)
\]

由 $ \operatorname{inv} $ 的连续性得到。因此 $ \operatorname{inv} : \mathcal{L}\mathrm{is}(E, F) \to \mathcal{L}(F, E) $ 可微，且 (7.2) 成立。

\item[（iv）]

我们表示映射

\[
g : \mathcal{L}(F, E)^{2} \to \mathcal{L}\left(\mathcal{L}(E, F), \mathcal{L}(F, E)\right)
\]

为

\[
g(T_{1}, T_{2})(S) := -T_{1}ST_{2} \qquad \text{对 } T_{1}, T_{2} \in \mathcal{L}(F, E) \text{ 且 } S \in \mathcal{L}(E, F) .
\]

于是容易看出 $ g $ 是双线性且连续的，因此

\[
g \in C^{\infty}\left(\mathcal{L}(F, E)^{2}, \mathcal{L}\left(\mathcal{L}(E, F), \mathcal{L}(F, E)\right)\right)
\]

（见习题 4.4 与 5.2）。此外，

\begin{equation}
\partial\operatorname{inv}(A) = g\left(\operatorname{inv}(A), \operatorname{inv}(A)\right) \qquad \text{对 } A \in \mathcal{L}\mathrm{is}(E, F) .\tag{7.6}
\end{equation}

因此由 (ii) 得 $ \partial\operatorname{inv} $ 的连续性。于是映射 $ \operatorname{inv} $ 属于 $ C^{1}\left(\mathcal{L}\mathrm{is}(E, F), \mathcal{L}(F, E)\right) $。最后，由 (7.6) 借助链式法则及简单的归纳论证，可知该映射是光滑的。\nobreak\hspace{0.35em}■

\end{subitems}

\end{proposition}

\needspace{3\baselineskip}\subsection*{反函数定理}

现在我们可以证明一个关于可微映射局部行为的中心结果，它——正如我们稍后将看到的——具有极其深远的影响。

\qyxffntextdef{1}{即 0 的邻域。}
\begin{theorem}

（反函数）设 $ X $ 在 $ E $ 中是开集，$ x_{0} \in X $。又设对 $ q \in \mathbb{N}^{\times} \cup \{\infty\} $ 有 $ f \in C^{q}(X, F) $。最后，设

\[
\partial f(x_{0}) \in \mathcal{L}\mathrm{is}(E, F) .
\]

则存在 $ x_{0} $ 在 $ X $ 中的开邻域 $ U $ 以及 $ y_{0} := f(x_{0}) $ 的开邻域 $ V $，它们具有如下性质：

\begin{subitems}
\item

$ f : U \to V $ 是双射。

\item

$ f^{-1} \in C^{q}(V, E) $，且对每个 $ x \in U $ 有

\[
\partial f(x) \in \mathcal{L}\mathrm{is}(E, F) \quad \text{且} \quad \partial f^{-1}(f(x)) = \left[\partial f(x)\right]^{-1} .
\]

\proofhead
(i) 令 $ A := \partial f(x_{0}) $，$ h := A^{-1}f : X \to E $。则由习题 5.1 与链式法则（定理 5.7）可知，$ h $ 属于 $ C^{q}(X, E) $ 且 $ \partial h(x_{0}) = A^{-1}\partial f(x_{0}) = I $。于是我们不失一般性地考虑 $ E = F $ 且 $ \partial f(x_{0}) = I $ 的情形（见习题 5.1）。
\end{subitems}

\begin{subitems}
\item[（ii）]

我们通过令

\[
h(x) := f(x + x_{0}) - f(x_{0}) \qquad \text{对 } x \in X_{1} := X - x_{0} .
\]

再做一次简化。则 $ X_{1} $ 在 $ E $ 中是开集，且 $ h \in C^{q}(X_{1}, E) $，$ h(0) = 0 $，$ \partial h(0) = \partial f(x_{0}) = I $。因此只需考虑

\[
E = F , \quad x_{0} = 0 , \quad f(0) = 0 , \quad \partial f(0) = I .
\]
\item[（iii）]

我们证明 $ f $ 在 0 附近局部双射。确切地说，我们证明存在零点邻域\qyxffnmark{1} $ U $ 与 $ V $，使得对每个 $ y \in V $，方程 $ f(x) = y $ 在 $ U $ 中有唯一解，且 $ f(U) \subset V $。这个问题显然等价于确定零点邻域 $ U $ 与 $ V $，使得映射

\[
g_{y} : U \to E , \quad x \mapsto x - f(x) + y
\]

对每个 $ y \in V $ 恰有一个不动点。开始，我们令 $ g := g_{0} $。因为 $ \partial f(0) = I $，有 $ \partial g(0) = 0 $，且利用 $ \partial g $ 的连续性可找到 $ r > 0 $ 使得

\begin{equation}
\|\partial g(x)\| \le 1/2 \qquad \text{对 } x \in 2r\mathbb{B} .\tag{7.7}
\end{equation}

因为 $ g(0) = 0 $，应用积分形式的中值定理（定理 3.10）得 $ g(x) = \int_{0}^{1}\partial g(tx)x\,dt $。进而，我们得到

\begin{equation}
\|g(x)\| \le \int_{0}^{1}\|\partial g(tx)\|\|x\|\,dt \le \|x\|/2 \qquad \text{对 } x \in 2r\mathbb{B} .\tag{7.8}
\end{equation}

因此对每个 $ y \in r\mathbb{B} $，我们有估计

\[
\|g_{y}(x)\| \le \|y\| + \|g(x)\| \le 2r \qquad \text{对 } x \in 2r\mathbb{B} ,
\]

即对每个 $ y \in r\mathbb{B} $，$ g_{y} $ 把 $ 2r\mathbb{B} $ 映到其自身。
对 $ x_{1}, x_{2} \in 2r\mathbb{B} $，借助中值定理可得

\[
\|g_{y}(x_{1}) - g_{y}(x_{2})\| = \left\|\int_{0}^{1}\partial g\left(x_{2} + t(x_{1} - x_{2})\right)(x_{1} - x_{2})\,dt\right\| \le \frac{1}{2}\|x_{1} - x_{2}\| .
\]

于是，对每个 $ y \in r\mathbb{B} $，$ g_{y} $ 是 $ 2r\mathbb{B} $ 上的压缩映射。由此，巴拿赫不动点定理（定理 IV.4.3）保证了存在唯一的 $ x \in 2r\mathbb{B} $ 使得 $ g_{y}(x) = x $，从而使得 $ f(x) = y $。
我们令 $ V := r\mathbb{B} $，$ U := f^{-1}(V) \cap 2r\mathbb{B} $。则 $ U $ 是开零点邻域，且 $ f|U : U \to V $ 是双射。

\item[（iv）]

现在证明 $ f^{-1} : V \to E $ 连续。考虑

\[
x = x - f(x) + f(x) = g(x) + f(x) \qquad \text{对 } x \in U .
\]

因此

\[
\|x_{1} - x_{2}\| \le \frac{1}{2}\|x_{1} - x_{2}\| + \|f(x_{1}) - f(x_{2})\| \qquad \text{对 } x_{1}, x_{2} \in U ,
\]

从而

\begin{equation}
\|f^{-1}(y_{1}) - f^{-1}(y_{2})\| \le 2\|y_{1} - y_{2}\| \qquad \text{对 } y_{1}, y_{2} \in V .\tag{7.9}
\end{equation}

于是 $ f^{-1} : V \to E $ 是利普希茨连续的。

\item[（v）]

我们证明 $ f^{-1} : V \to E $ 的可微性，并证明其导数由下式给出

\begin{equation}
\partial f^{-1}(y) = \left[\partial f(x)\right]^{-1} \qquad \text{其中 } x := f^{-1}(y) .\tag{7.10}
\end{equation}

首先证明对 $ x \in U $，$ \partial f(x) $ 属于 $ \mathcal{L}\mathrm{aut}(E) $。由 $ f(x) = x - g(x) $ 可知，对 $ x \in U $ 有 $ \partial f(x) = I - \partial g(x) $。注意到 (7.7)，可见 $ \partial f(x) \in \mathcal{L}\mathrm{aut}(E) $ 由命题 7.1 得到。
现在设 $ y, y_{0} \in V $，令 $ x := f^{-1}(y) $，$ x_{0} := f^{-1}(y_{0}) $。则

\[
f(x) - f(x_{0}) = \partial f(x_{0})(x - x_{0}) + o(\|x - x_{0}\|) \quad (x \to x_{0}) ,
\]

且当 $ x \to x_{0} $ 时，我们得到

\begin{equation*}
\begin{aligned}
&\|f^{-1}(y) - f^{-1}(y_{0}) - \left[\partial f(x_{0})\right]^{-1}(y - y_{0})\| \\ &\quad = \|x - x_{0} - \left[\partial f(x_{0})\right]^{-1}(f(x) - f(x_{0}))\| \le c\left\|o(\|x - x_{0}\|)\right\|
\end{aligned}
\end{equation*}

其中 $ c = \left\|\left[\partial f^{-1}(x_{0})\right]^{-1}\right\| $。因为 (7.9) 给出估计 $ 2\|y - y_{0}\| \ge \|x - x_{0}\| $，我们最终有

\begin{equation}
\frac{\|f^{-1}(y) - f^{-1}(y_{0}) - \left[\partial f(x_{0})\right]^{-1}(y - y_{0})\|}{\|y - y_{0}\|} \le \frac{2c\left\|o(\|x - x_{0}\|)\right\|}{\|x - x_{0}\|}\tag{7.11}
\end{equation}

当 $ x \to x_{0} $ 时。现在令 $ y \to y_{0} $ 取极限，由 (iv) 知 $ x \to x_{0} $，而 (7.11) 表明 $ f^{-1} $ 在 $ y_{0} $ 处可微，其导数为 $ \left[\partial f(x_{0})\right]^{-1} $。

\item[（vi）]

还需验证 $ f^{-1} $ 属于 $ C^{q}(V, E) $。考虑 (7.10) 所表明的，即

\begin{equation}
\partial f^{-1} = (\partial f \circ f^{-1})^{-1} = \operatorname{inv} \circ \partial f \circ f^{-1} .\tag{7.12}
\end{equation}

由 (iv) 知 $ f^{-1} $ 属于 $ C(V, E) $，且 $ f^{-1}(V) \cap \mathbb{B}(x_{0}, 2r) = U $；又由假设知 $ \partial f \in C(U, \mathcal{L}(E)) $。利用命题 7.2 可知 $ \partial f^{-1} $ 属于 $ C(V, \mathcal{L}(E)) $，这就证明了 $ f^{-1} \in C^{1}(V, E) $。对 $ q > 1 $，可由 (7.12) 利用链式法则归纳地证明 $ f^{-1} \in C^{p}(V, E) $。\nobreak\hspace{0.35em}■

\end{subitems}

\end{theorem}

\needspace{3\baselineskip}\subsection*{微分同胚}

设 $ X $ 在 $ E $ 中是开集，$ Y $ 在 $ F $ 中是开集，$ q \in \mathbb{N} \cup \{\infty\} $。我们称映射 $ f : X \to Y $ 为从 $ X $ 到 $ Y $ 的 $ C^{q} $ 微分同胚，若它是双射且

\[
f \in C^{q}(X, Y) , \quad \text{且} \quad f^{-1} \in C^{q}(Y, E) .
\]

我们可以把 $ C^{0} $ 微分同胚称为\textbf{同胚}或\textbf{拓扑映射}。我们令

\[
\operatorname{Diff}^{q}(X, Y) := \{f : X \to Y ; \ f \text{ 是 } C^{q} \text{ 微分同胚}\} .
\]

映射 $ g : X \to F $ 是局部 $ C^{q} $ 微分同胚\footnote{当 $ q = 0 $ 时，人们说局部同胚或局部拓扑映射。}，如果每个 $ x_{0} \in X $ 都有开邻域 $ U \in \mathfrak{U}_{X}(x_{0}) $ 与 $ V \in \mathfrak{U}_{F}(g(x_{0})) $，使得 $ g|U $ 属于 $ \operatorname{Diff}^{q}(U, V) $。我们用 $ \operatorname{Diff}_{\mathrm{loc}}^{q}(X, F) $ 表示从 $ X $ 到 $ F $ 的所有局部 $ C^{q} $ 微分同胚之集。

\begin{remark}

设 $ X $ 在 $ E $ 中是开集，$ Y $ 在 $ F $ 中是开集。

\begin{items}
\item

对每个 $ q \in \mathbb{N} $，有包含关系

\begin{equation*}
\begin{aligned}
&\operatorname{Diff}^{\infty}(X, Y) \subset \operatorname{Diff}^{q+1}(X, Y) \subset \operatorname{Diff}^{q}(X, Y) \subset \operatorname{Diff}_{\mathrm{loc}}^{q}(X, F) , \\ &\operatorname{Diff}_{\mathrm{loc}}^{\infty}(X, F) \subset \operatorname{Diff}_{\mathrm{loc}}^{q+1}(X, F) \subset \operatorname{Diff}_{\mathrm{loc}}^{q}(X, F) .
\end{aligned}
\end{equation*}
\item

设 $ f \in C^{q+1}(X, Y) $ 且 $ f \in \operatorname{Diff}^{q}(X, Y) $。则一般不能推出 $ f $ 属于 $ \operatorname{Diff}^{q+1}(X, Y) $。

\proofhead
令 $ E := F := \mathbb{R} $，$ X := Y := \mathbb{R} $，考虑 $ f(x) := x^{3} $。则 $ f $ 是光滑的拓扑映射，即 $ f \in C^{\infty}(X, Y) \cap \operatorname{Diff}^{0}(X, Y) $。然而 $ f^{-1} $ 在 0 处不可微。\nobreak\hspace{0.35em}■
\item

局部拓扑映射是\textbf{开映射}，即它们把开集映为开集（见习题 III.2.14）。

\proofhead
这容易由定理 I.3.8(ii) 与 III.2.4(ii) 得到。\nobreak\hspace{0.35em}■
\item

对 $ f \in \operatorname{Diff}_{\mathrm{loc}}^{1}(X, Y) $，有 $ \partial f(x) \in \mathcal{L}\mathrm{is}(E, F) $ 对 $ x \in X $ 成立。

\proofhead
这是链式法则的推论（见习题 1）。\nobreak\hspace{0.35em}■
\end{items}

\end{remark}

\begin{corollary}

设 $ X $ 在 $ E $ 中是开集，$ q \in \mathbb{N}^{\times} \cup \{\infty\} $，且 $ f \in C^{q}(X, F) $。则

\[
f \in \operatorname{Diff}_{\mathrm{loc}}^{q}(X, F) \iff \partial f(x) \in \mathcal{L}\mathrm{is}(E, F) \ \text{对 } x \in X .
\]

\proofhead
“$ \Rightarrow $” 由注 7.4(a)，归纳可得 $ f $ 属于 $ \operatorname{Diff}_{\mathrm{loc}}^{1}(X, F) $。于是结论由注 7.4(d) 得到。

“$ \Leftarrow $” 这由反函数定理得到。\nobreak\hspace{0.35em}■
\end{corollary}

\begin{remark}

在推论 7.5 的假设下，对 $ x \in X $ 有 $ \partial f(x) $ 属于 $ \mathcal{L}\mathrm{is}(E, F) $。则 $ f $ 是局部拓扑的，从而 $ Y := f(X) $ 是开集。然而一般地，$ f $ 不是从 $ X $ 到 $ Y $ 的 $ C^{q} $ 微分同胚。

\proofhead
令 $ X := E := F := \mathbb{C} $，$ f(z) := e^{z} $。则 $ f $ 光滑，且对 $ z \in \mathbb{C} $ 有 $ \partial f(z) = e^{z} \ne 0 $。由于其 $ 2\pi i $-周期性，$ f $ 不是单射。\nobreak\hspace{0.35em}■
\end{remark}

\needspace{3\baselineskip}\subsection*{非线性方程组的可解性}

我们现在对 $ E = F = \mathbb{R}^{m} $ 的情形陈述反函数定理。我们将利用 $ \mathbb{R}^{m} $ 上的每个线性映射都连续，以及线性映射当其行列式不为零时可逆这一事实。

本节其余部分，设 $ X $ 在 $ \mathbb{R}^{m} $ 中是开集，$ x_{0} \in X $。又设 $ q \in \mathbb{N}^{\times} \cup \{\infty\} $，$ f = (f^{1}, \ldots, f^{m}) \in C^{q}(X, \mathbb{R}^{m}) $。

\begin{theorem}

若 $ \det(\partial f(x_{0})) \ne 0 $，则存在 $ x_{0} $ 的开邻域 $ U $ 与 $ f(x_{0}) $ 的开邻域 $ V $，使得 $ f|U $ 属于 $ \operatorname{Diff}^{q}(U, V) $。

\proofhead
由定理 1.6 知 $ \operatorname{Hom}(\mathbb{R}^{m}, \mathbb{R}^{m}) = \mathcal{L}(\mathbb{R}^{m}) $。线性代数告诉我们
\[
\partial f(x_{0}) \in \mathcal{L}\mathrm{aut}(\mathbb{R}^{m}) \iff \det(\partial f(x_{0})) \ne 0
\]

（例如 [Gab96, A3.6(a)]）。于是结论由定理 7.3 得到。\nobreak\hspace{0.35em}■

\end{theorem}

\begin{corollary}

若 $ \det(\partial f(x_{0})) \ne 0 $，则存在 $ x_{0} $ 的开邻域 $ U $ 与 $ f(x_{0}) $ 的开邻域 $ V $，使得方程组

\begin{equation}
\begin{aligned}
f^{1}(x^{1}, \ldots, x^{m}) &= y^{1} , \\ &\vdots \\ f^{m}(x^{1}, \ldots, x^{m}) &= y^{m}
\end{aligned}\tag{7.13}
\end{equation}

对每个 $ m $-元组 $ (y^{1}, \ldots, y^{m}) \in V $ 在 $ U $ 中恰有一个解

\[
x^{1} = x^{1}(y^{1}, \ldots, y^{m}), \ldots, x^{m} = x^{m}(y^{1}, \ldots, y^{m}) .
\]

函数 $ x^{1}, \ldots, x^{m} $ 属于 $ C^{q}(V, \mathbb{R}) $。

\end{corollary}

\begin{remark}

\begin{items}
\item

根据注 1.18(b) 与推论 2.9，线性映射 $ \partial f(x) $ 的行列式可用雅可比矩阵 $ \left[\partial_{k}f^{j}(x)\right] $ 来计算，即 $ \det(\partial f(x)) = \det\left[\partial_{k}f^{j}(x)\right] $。它称为 $ f $（在点 $ x $ 处）的\textbf{雅可比行列式}。我们也可以简单地称它为\textbf{雅可比}，并把它写作

\[
\frac{\partial(f^{1}, \ldots, f^{m})}{\partial(x^{1}, \ldots, x^{m})}(x) .
\]
\item

由于反函数定理的证明是构造性的，原则上可用它来计算（近似地）$ f^{-1}(x) $。特别地，通过在 $ x_{0} $ 的邻域中考察方程组 (7.13)，我们可以逼近推论 7.8 所描述的有限维情形的解。

\proofhead
这是因为定理 7.3 的证明基于压缩定理，因而基于逐次逼近法。\nobreak\hspace{0.35em}■
\end{items}

\end{remark}

\exerciseshead

\begin{exercises}
\item

设 $ X $ 是 $ E $ 中的开集，$ Y $ 是 $ F $ 中的开集，且 $ f \in \operatorname{Diff}^{1}_{\mathrm{loc}}(X,Y) $。证明：对 $ x_0 \in X $，$ \partial f(x_0) $ 属于 $ \mathcal{L}\mathrm{is}(E,F) $，且 $ \partial f^{-1}(y_0) = [\partial f(x_0)]^{-1} $，其中 $ y_0 := f(x_0) $。

\item

设 $ m,n \in \mathbb{N}^{\times} $，且 $ X $ 是 $ \mathbb{R}^n $ 中的开集。证明：若 $ \operatorname{Diff}^{1}_{\mathrm{loc}}(X,\mathbb{R}^m) $ 非空，则 $ m = n $。

\item

对 (a)–(d) 中定义的映射，求 $ Y := f(X) $ 与 $ f^{-1} $。并判定是否有 $ f \in \operatorname{Diff}^{q}_{\mathrm{loc}}(X,\mathbb{R}^2) $ 或 $ f \in \operatorname{Diff}^{q}(X,Y) $。

\begin{items}
\item

$ X := \mathbb{R}^2 $，且对 $ (a,b) \in \mathbb{R}^2 $ 有 $ f(x,y) := (x + a, y + b) $；

\item

$ X := \mathbb{R}^2 $，且 $ f(x,y) := (x^2 - x - 2, 3y) $；

\item

$ X := \mathbb{R}^2 \setminus \{(0,0)\} $，且 $ f(x,y) := (x^2 - y^2, 2xy) $；

\item

$ X := \{(x,y) \in \mathbb{R}^2 \; ; \; 0 < y < x\} $，且 $ f(x,y) := (\log xy, 1/(x^2 + y^2)) $。
（(c) 的提示：$ \mathbb{R}^2 \leftrightarrow \mathbb{C} $。）

\end{items}

\item

设

\[
f: \mathbb{R}^2 \to \mathbb{R}^2 , \quad (x,y) \mapsto (\cosh x \cos y, \sinh x \sin y)
\]

且 $ X := \{(x,y) \in \mathbb{R}^2 \; ; \; x > 0\} $，$ Y := f(X) $。证明：

\begin{items}
\item

$ f|X \in \operatorname{Diff}^{\infty}_{\mathrm{loc}}(X,\mathbb{R}^2) $；

\item

$ f|X \notin \operatorname{Diff}^{\infty}(X,Y) $；

\item

对 $ U := \{(x,y) \in X \; ; \; 0 < y < 2\pi\} $ 与 $ V := Y \setminus ([0,\infty) \times \{0\}) $，$ f|U $ 属于 $ \operatorname{Diff}^{\infty}(U,V) $；

\item

$ Y = \mathbb{R}^2 \setminus ([-1,1] \times \{0\}) $。

\end{items}

\item

设 $ X $ 是 $ \mathbb{R}^m $ 中的开集，且 $ f \in C^1(X,\mathbb{R}^m) $。证明：

\begin{items}
\item

若对 $ x \in X $ 有 $ \partial f(x) \in \mathcal{L}\mathrm{is}(\mathbb{R}^m) $，则 $ x \mapsto |f(x)| $ 在 $ X $ 中无最大值；

\item

若对 $ x \in X $ 有 $ \partial f(x) \in \mathcal{L}\mathrm{is}(\mathbb{R}^m) $ 且 $ f(x) \ne 0 $，则 $ x \mapsto |f(x)| $ 无最小值。

\end{items}

\item

设 $ H $ 是实希尔伯特空间，且

\[
f: H \to H , \quad x \mapsto x/\sqrt{1 + |x|^2} .
\]

令 $ Y := \operatorname{im}(f) $，证明 $ f \in \operatorname{Diff}^{\infty}(H,Y) $。$ f^{-1} $ 与 $ \partial f $ 是什么？

\item

设 $ X $ 是 $ \mathbb{R}^m $ 中的开集，且 $ f \in C^1(X,\mathbb{R}^m) $。又设存在 $ \alpha > 0 $ 使得

\begin{equation}
|f(x) - f(y)| \ge \alpha |x - y| \quad \text{对 } x,y \in X .\tag{7.14}
\end{equation}

证明 $ Y := f(X) $ 是 $ \mathbb{R}^m $ 中的开集，且 $ f \in \operatorname{Diff}^1(X,Y) $。并证明当 $ X = \mathbb{R}^m $ 时有 $ Y = \mathbb{R}^m $。
（提示：由 (7.14) 可得，对 $ x \in X $ 有 $ \partial f(x) \in \mathcal{L}\mathrm{is}(X,Y) $。若 $ X = \mathbb{R}^m $，则由 (7.14) 推出 $ Y $ 在 $ \mathbb{R}^m $ 中是闭集。）

\item

设 $ f \in \operatorname{Diff}^1(\mathbb{R}^m,\mathbb{R}^m) $ 且 $ g \in C^1(\mathbb{R}^m,\mathbb{R}^m) $，并且下列任一假设成立：

\begin{items}
\item

$ f^{-1} $ 与 $ g $ 是利普希茨连续的；

\item

$ g $ 在 $ \mathbb{R}^m $ 的某个有界子集之外为零。
证明：于是存在 $ \varepsilon_0 > 0 $ 使得对 $ \varepsilon \in (-\varepsilon_0,\varepsilon_0) $ 有 $ f + \varepsilon g \in \operatorname{Diff}^1(\mathbb{R}^m,\mathbb{R}^m) $。
（提示：考虑 $ \operatorname{id}_{\mathbb{R}^m} + f^{-1} \circ (\varepsilon g) $，并应用习题 7。）

\end{items}

\end{exercises}

\section{隐函数}

在前几节中，我们（在有限维情形下）研究了非线性方程组的可解性。我们集中于方程个数等于变量个数的情形。现在，我们探讨未知量多于方程的非线性系统的可解性。

我们将证明主要结果——隐函数定理，而不去为在一般巴拿赫空间上证明它所需的本质上是辅助性的考虑所困扰。为说明这一基本定理，我们证明常微分方程的基本存在性与连续性定理。隐函数定理的有限维应用将在随后两节中处理。

作为动机，我们考虑函数 $ f:\mathbb{R}^{2}\to\mathbb{R} $，$ (x,y)\mapsto x^{2}+y^{2}-1 $。设 $ (a,b)\in\mathbb{R}^{2} $，其中 $ a\ne\pm 1 $，$ b>0 $ 且 $ f(a,b)=0 $。则存在开区间 $ A $ 与 $ B $，$ a\in A $、$ b\in B $，使得对每个 $ x\in A $ 恰存在一个 $ y\in B $ 满足 $ f(x,y)=0 $。通过对应 $ x\mapsto y $，我们定义映射 $ g:A\to B $，满足对 $ x\in A $ 有 $ f(x,g(x))=0 $。显然，我们有 $ g(x)=\sqrt{1-x^{2}} $。此外，存在开区间 $ \widetilde{B} $ 使得 $ -b\in\widetilde{B} $，以及一个 $ \widetilde{g}:A\to\widetilde{B} $ 使得对 $ x\in A $ 有 $ f(x,\widetilde{g}(x))=0 $。这里，我们有 $ \widetilde{g}(x)=-\sqrt{1-x^{2}} $。函数 $ g $ 在 $ A $ 上由 $ f $ 与 $ (a,b) $ 唯一确定，而函数 $ \widetilde{g} $ 由 $ f $ 与 $ (a,-b) $ 定义。因此我们说 $ g $ 与 $ \widetilde{g} $ 分别在 $ (a,b) $ 与 $ (a,-b) $ 附近由 $ f $ 隐式定义。函数 $ g $ 与 $ \widetilde{g} $ 分别是方程 $ f(x,y)=0 $ 在 $ (a,b) $ 与 $ (a,-b) $ 附近关于 $ y $（作为 $ x $ 的函数）的\textbf{解}。我们说 $ f(x,y)=0 $ 在 $ (a,b) $ 附近\textbf{隐式地}定义了 $ g $。

\bookfigure[0.44]{figures/4f577e681a213afd.png}

这样的解显然不能在 $ (1,0) $ 或 $ (-1,0) $ 的邻域中给出。为提示障碍所在，我们注意到，对 $ a=\pm 1 $，有 $ \partial_{2}f(a,b)=0 $；而对 $ a\ne\pm 1 $，这一量不消失，因为 $ \partial_{2}f(a,b)=2b $。

\needspace{3\baselineskip}\subsection*{积空间上的可微映射}

这里，设

- $ E_{1} $、$ E_{2} $ 与 $ F $ 是 $ \mathbb{K} $ 上的巴拿赫空间；
  $ q\in\mathbb{N}^{\times}\cup\{\infty\} $。

设 $ X_{j} $ 在 $ E_{j} $ 中开，$ j=1,2 $，且 $ f:X_{1}\times X_{2}\to F $ 在 $ (a,b) $ 处可微。则函数 $ f(\cdot,b):X_{1}\to F $ 与 $ f(a,\cdot):X_{2}\to F $ 也分别在 $ a $ 与 $ b $ 处可微。为避免与标准偏导数混淆，我们用 $ D_{1}f(a,b) $ 记 $ f(\cdot,b) $ 在 $ a $ 处的导数，用 $ D_{2}f(a,b) $ 记 $ f(a,\cdot) $ 在 $ b $ 处的导数，

\begin{remark}

\begin{items}
\item

显然 $ D_{j}f:X_{1}\times X_{2}\to\mathcal{L}(E_{j},F) $，$ j=1,2 $。

\item

陈述

\begin{subitems}
\item

$ f\in C^{q}(X_{1}\times X_{2},F) $ 且

\item

$ D_{j}f\in C^{q-1}(X_{1}\times X_{2},\mathcal{L}(E_{j},F)) $，$ j=1,2 $
是等价的。据此，我们有

\[
\partial f(a,b)(h,k)=D_{1}f(a,b)h+D_{2}f(a,b)k
\]

对 $ (a,b)\in X_{1}\times X_{2} $ 与 $ (h,k)\in E_{1}\times E_{2} $（见定理 5.4 与命题 2.8）。

\proofhead
“(i)$ \Rightarrow $(ii)” 这一蕴含是显然的。
“(ii)$ \Rightarrow $(i)” 设 $ (a,b)\in X_{1}\times X_{2} $ 且
\[
A(h,k):=D_{1}f(a,b)h+D_{2}f(a,b)k\qquad\text{对 }(h,k)\in E_{1}\times E_{2} \ .
\]

容易验证 $ A $ 属于 $ \mathcal{L}(E_{1}\times E_{2},F) $。积分形式的中值定理（定理 3.10）给出

\begin{equation*}
\begin{aligned}
f(a+h,b+k)-f(a,b)-A(h,k)&=f(a+h,b+k)-f(a,b+k)\\ &\quad{}+f(a,b+k)-f(a,b)-A(h,k)\\ &=\int_{0}^{1}\left[D_{1}f(a+th,b+k)-D_{1}f(a,b)\right]h\,dt\\ &\qquad{}+f(a,b+k)-f(a,b)-D_{2}f(a,b)k
\end{aligned}
\end{equation*}

只要 $ \max\{\|h\|,\|k\|\} $ 足够小。由此，我们得到估计

\[
\|f(a+h,b+k)-f(a,b)-A(h,k)\|\le\varphi(h,k)\max\{\|h\|,\|k\|\}
\]

其中

\[
\varphi(h,k):=\max_{0\le t\le 1}\|D_{1}f(a+th,b+k)-D_{1}f(a,b)\|+\frac{\|f(a,b+k)-f(a,b)-D_{2}f(a,b)k\|}{\|k\|}\ .
\]

现在 $ D_{1}f $ 的连续性给出，当 $ (h,k)\to(0,0) $ 时 $ \varphi(h,k)\to 0 $。于是我们看到

\[
f(a+h,b+k)-f(a,b)-A(h,k)=o(\|(h,k)\|)\qquad((h,k)\to(0,0)) \ .
\]

因此 $ f $ 在 $ (a,b) $ 处可微，且 $ \partial f(a,b)=A $。最后，对 $ D_{j}f $ 的正则性假设与 $ A $ 的定义蕴含 $ \partial f\in C^{q-1}(X_{1}\times X_{2},F) $，从而 $ f\in C^{q}(X_{1}\times X_{2},F) $。\nobreak\hspace{0.35em}■

\end{subitems}

\item

在特殊情形 $ E_{1}=\mathbb{R}^{m} $、$ E_{2}=\mathbb{R}^{n} $、$ F=\mathbb{R}^{\ell} $ 下，我们有

\begin{equation*}
[D_{1}f]=\begin{bmatrix}\partial_{1}f^{1} & \cdots & \partial_{m}f^{1}\\ \vdots & & \vdots\\ \partial_{1}f^{\ell} & \cdots & \partial_{m}f^{\ell}\end{bmatrix}\ ,\qquad[D_{2}f]=\begin{bmatrix}\partial_{m+1}f^{1} & \cdots & \partial_{m+n}f^{1}\\ \vdots & & \vdots\\ \partial_{m+1}f^{\ell} & \cdots & \partial_{m+n}f^{\ell}\end{bmatrix}
\end{equation*}

其中 $ f=(f^{1},\ldots,f^{\ell}) $。

\proofhead
这由 (b) 与推论 2.9 推出。\nobreak\hspace{0.35em}■
\end{items}

\end{remark}

\needspace{3\baselineskip}\subsection*{隐函数定理}

借助我们刚学到的内容，我们可以转而研究非线性方程的可解性。这一结果是基本的。

\qyxffnnumdef{1}
\afterpage{\qyxffntextof{1}{此处及类似情形下，我们使用自然的矩阵记号。}}
\begin{theorem}

（隐函数） 设 $ W $ 在 $ E_{1}\times E_{2} $ 中开且 $ f\in C^{q}(W,F) $。进一步，设 $ (x_{0},y_{0})\in W $ 使得

\[
f(x_{0},y_{0})=0\quad\text{且}\quad D_{2}f(x_{0},y_{0})\in\operatorname{Lis}(E_{2},F) \ .
\]

则存在开邻域 $ U\in\mathfrak{U}_{W}(x_{0},y_{0}) $ 与 $ V\in\mathfrak{U}_{E_{1}}(x_{0}) $，以及唯一的 $ g\in C^{q}(V,E_{2}) $，使得

\begin{equation}
\left((x,y)\in U\text{ 且 }f(x,y)=0\right)\iff\left(x\in V\text{ 且 }y=g(x)\right) \ .\tag{8.1}
\end{equation}

此外，

\begin{equation}
\partial g(x)=-\left[D_{2}f(x,g(x))\right]^{-1}D_{1}f(x,g(x))\qquad\text{对 }x\in V \ .\tag{8.2}
\end{equation}

证明

\begin{subitems}
\item

令 $ A:=D_{2}f(x_{0},y_{0})\in\operatorname{Lis}(E_{2},F) $ 且 $ \widetilde{f}:=A^{-1}f\in C^{q}(W,E_{2}) $。则

\[
\widetilde{f}(x_{0},y_{0})=0\quad\text{且}\quad D_{2}\widetilde{f}(x_{0},y_{0})=I_{E_{2}} \ .
\]

于是，根据习题 5.1，我们不失一般性地考虑情形 $ F=E_{2} $ 且 $ D_{2}f(x_{0},y_{0})=I_{E_{2}} $。此外，对开邻域 $ W_{1}\in\mathfrak{U}_{E_{1}}(x_{0}) $ 与 $ W_{2}\in\mathfrak{U}_{E_{2}}(y_{0}) $，我们可以假设 $ W=W_{1}\times W_{2} $。

\item

对映射

\[
\varphi:W_{1}\times W_{2}\to E_{1}\times E_{2}\ ,\quad(x,y)\mapsto(x,f(x,y))
\]

我们有 $ \varphi\in C^{q}(W_{1}\times W_{2},E_{1}\times E_{2}) $，且\qyxffnmark{1}

\begin{equation*}
\partial\varphi(x_{0},y_{0})=\begin{bmatrix}I_{E_{1}} & 0\\ D_{1}f(x_{0},y_{0}) & I_{E_{2}}\end{bmatrix}\in\mathcal{L}(E_{1}\times E_{2}) \ .
\end{equation*}

我们立即验证

\begin{equation*}
\begin{bmatrix}I_{E_{1}} & 0\\ -D_{1}f(x_{0},y_{0}) & I_{E_{2}}\end{bmatrix}\in\mathcal{L}(E_{1}\times E_{2})
\end{equation*}

是 $ \partial\varphi(x_{0},y_{0}) $ 的逆。因此，我们有 $ \partial\varphi(x_{0},y_{0})\in\operatorname{Laut}(E_{1}\times E_{2}) $，又因为 $ \varphi(x_{0},y_{0})=(x_{0},0) $，反函数定理（定理 7.3）保证存在开邻域 $ U\in\mathfrak{U}_{E_{1}\times E_{2}}(x_{0},y_{0}) $ 与 $ X\in\mathfrak{U}_{E_{1}\times E_{2}}(x_{0},0) $，使得 $ \varphi|U\in\operatorname{Diff}^{q}(U,X) $。我们令 $ \psi:=(\varphi|U)^{-1}\in\operatorname{Diff}^{q}(X,U) $，并把 $ \psi $ 写成

\[
\psi(\xi,\eta)=(\psi_{1}(\xi,\eta),\psi_{2}(\xi,\eta))\qquad\text{对 }(\xi,\eta)\in X \ .
\]

则 $ \psi_{j}\in C^{q}(X,E_{2}) $，$ j=1,2 $，而 $ \varphi $ 的定义表明

\[
(\xi,\eta)=\varphi(\psi(\xi,\eta))=(\psi_{1}(\xi,\eta),f(\psi_{1}(\xi,\eta),\psi_{2}(\xi,\eta)))\qquad\text{对 }(\xi,\eta)\in X \ .
\]

因此，我们识别出

\begin{equation}
\psi_{1}(\xi,\eta)=\xi\quad\text{且}\quad\eta=f(\xi,\psi_{2}(\xi,\eta))\qquad\text{对 }(\xi,\eta)\in X \ .\tag{8.3}
\end{equation}

此外，$ V:=\{x\in E_{1};(x,0)\in X\} $ 是 $ x_{0} $ 在 $ E_{1} $ 中的开邻域，且对 $ g(x):=\psi_{2}(x,0) $（$ x\in V $），我们有 $ g\in C^{q}(V,E_{2}) $ 且

\begin{equation*}
\begin{aligned}
(x,f(x,g(x)))&=(\psi_{1}(x,0),f(\psi_{1}(x,0),\psi_{2}(x,0)))\\ &=\varphi(\psi_{1}(x,0),\psi_{2}(x,0))=\varphi\circ\psi(x,0)=(x,0)
\end{aligned}
\end{equation*}

对 $ x\in V $。这与 (8.3) 一起蕴含 (8.1)。$ g $ 的唯一性是显然的。

\item

我们令 $ h(x):=f(x,g(x)) $，$ x\in V $。则 $ h=0 $，且链式法则与注 8.1(b) 一起蕴含

\begin{equation}
\partial h(x)=D_{1}f(x,g(x))I_{E_{1}}+D_{2}f(x,g(x))\partial g(x)=0\qquad\text{对 }x\in V \ .\tag{8.4}
\end{equation}

由 $ q\ge 1 $ 可得

\[
D_{2}f\in C^{q-1}(U,\mathcal{L}(E_{2}))\subset C(U,\mathcal{L}(E_{2})) \ .
\]

根据命题 7.2(i)，$ \operatorname{Laut}(E_{2}) $ 在 $ \mathcal{L}(E_{2}) $ 中开。因此

\[
\left((D_{2}f)^{-1}(\operatorname{Laut}(E_{2}))\right)\cap U=\{(x,y)\in U;D_{2}f(x,y)\in\operatorname{Laut}(E_{2})\}
\]

是 $ (x_{0},y_{0}) $ 在 $ U $ 中的开邻域。通过缩小 $ U $，我们因此可以假设对 $ (x,y)\in U $ 有 $ D_{2}f(x,y)\in\operatorname{Laut}(E_{2}) $。于是由 (8.4) 得到 (8.2)。\nobreak\hspace{0.35em}■

\end{subitems}

\end{theorem}

\begin{remark}

定理 8.2 说的是，在 $ (x_{0},y_{0}) $ 附近，纤维 $ f^{-1}(0) $ 是一个 $ C^{q} $ 函数的图像。\nobreak\hspace{0.35em}■

通过在特殊情形 $ E_{1}=\mathbb{R}^{m} $ 与 $ E_{2}=F=\mathbb{R}^{n} $ 下表述隐函数定理，我们得到关于“依赖于参数的非线性方程组的局部可解性”的陈述。

\end{remark}

\begin{corollary}

设 $ W $ 在 $ \mathbb{R}^{m+n} $ 中开且 $ f\in C^{q}(W,\mathbb{R}^{n}) $。又设 $ (a,b)\in W $ 使得 $ f(a,b)=0 $，即

\begin{equation*}
\begin{array}{c}f^{1}(a^{1},\ldots,a^{m},b^{1},\ldots,b^{n})=0\ ,\\ \vdots\\ f^{n}(a^{1},\ldots,a^{m},b^{1},\ldots,b^{n})=0\ .\end{array}
\end{equation*}

则若

\[
\frac{\partial(f^{1},\ldots,f^{n})}{\partial(x^{m+1},\ldots,x^{m+n})}(a,b):=\det\left[\partial_{m+k}f^{j}(a,b)\right]_{1\le j,k\le n}\ne 0\ ,
\]

则存在 $ (a,b) $ 在 $ W $ 中的开邻域 $ U $、$ a $ 在 $ \mathbb{R}^{m} $ 中的开邻域 $ V $，以及 $ g\in C^{q}(V,\mathbb{R}^{n}) $，使得

\[
\left((x,y)\in U\text{ 且 }f(x,y)=0\right)\iff\left(x\in V\text{ 且 }y=g(x)\right) \ .
\]

也就是说，存在 $ a $ 在 $ \mathbb{R}^{m} $ 中的邻域 $ V $，使得方程组

\begin{equation*}
\begin{array}{c}f^{1}(x^{1},\ldots,x^{m},y^{1},\ldots,y^{n})=0\ ,\\ \vdots\\ f^{n}(x^{1},\ldots,x^{m},y^{1},\ldots,y^{n})=0\end{array}
\end{equation*}

在 $ b=(b^{1},\ldots,b^{n}) $ 附近，对每个 $ m $ 元组 $ (x^{1},\ldots,x^{m})\in V $ 恰有一个解

\begin{equation*}
\begin{array}{c}y^{1}=g^{1}(x^{1},\ldots,x^{m})\ ,\\ \vdots\\ y^{n}=g^{n}(x^{1},\ldots,x^{m})\end{array}
\end{equation*}

此外，解 $ g^{1},\ldots,g^{n} $ 是来自 $ V $ 的“参数” $ (x^{1},\ldots,x^{m}) $ 的 $ C^{q} $ 函数，且

\begin{equation*}
[\partial g(x)]=-\begin{bmatrix}\partial_{m+1}f^{1} & \cdots & \partial_{m+n}f^{1}\\ \vdots & & \vdots\\ \partial_{m+1}f^{n} & \cdots & \partial_{m+n}f^{n}\end{bmatrix}^{-1}\begin{bmatrix}\partial_{1}f^{1} & \cdots & \partial_{m}f^{1}\\ \vdots & & \vdots\\ \partial_{1}f^{n} & \cdots & \partial_{m}f^{n}\end{bmatrix}
\end{equation*}

对 $ x\in V $，其中导数 $ \partial_{k}f^{j} $ 在 $ (x,g(x)) $ 处取值。

\proofhead
因为 $ D_{2}f(x,y) $ 属于 $ \operatorname{Laut}(\mathbb{R}^{n}) $ 当且仅当
\[
\frac{\partial(f^{1},\ldots,f^{n})}{\partial(x^{m+1},\ldots,x^{m+n})}(x,y)\ne 0\ ,
\]

所以由定理 8.2 与注 8.1(c) 验证所有断言。\nobreak\hspace{0.35em}■

\end{corollary}

\needspace{3\baselineskip}\subsection*{正则值}

设 $ X $ 在 $ \mathbb{R}^{m} $ 中开，且 $ f:X\to\mathbb{R}^{n} $ 可微。若 $ \partial f(x)\in\mathcal{L}(\mathbb{R}^{m},\mathbb{R}^{n}) $ 是满射，则称 $ x\in X $ 为 $ f $ 的\textbf{正则点}。若 $ X $ 中每一点都正则，则称该映射\textbf{正则}或为\textbf{淹没}。最后，若纤维 $ f^{-1}(y) $ 只由正则点组成，则称 $ y\in\mathbb{R}^{n} $ 为 $ f $ 的\textbf{正则值}。

\begin{remark}

\begin{items}
\item

若 $ m<n $，则 $ f $ 没有正则点。

\proofhead
这由 $ \operatorname{rank}(\partial f(x))\le m $ 推出。\nobreak\hspace{0.35em}■
\item

若 $ n\le m $，则 $ x\in X $ 是 $ f $ 的正则点当且仅当 $ \partial f(x) $ 的秩\footnote{设 $ E $ 与 $ F $ 是有限维巴拿赫空间，且 $ A\in\mathcal{L}(E,F) $。则 $ A $ 的秩定义为 $ \operatorname{rank}(A):=\dim(\operatorname{im}(A)) $。显然，$ \operatorname{rank}(A) $ 与表示矩阵 $ [A]_{\mathcal{E},\mathcal{F}} $ 在 $ E $ 的任意基 $ \mathcal{E} $ 与 $ F $ 的任意基 $ \mathcal{F} $ 下的线性代数秩相同。}为 $ n $。

\item

若 $ n=1 $，则 $ x $ 是 $ f $ 的正则点当且仅当 $ \nabla f(x)\ne 0 $。

\item

每个 $ y\in\mathbb{R}^{n}\setminus\operatorname{im}(f) $ 都是 $ f $ 的正则值。

\item

设 $ x_{0}\in X $ 是 $ f $ 的正则点，且 $ f(x_{0})=0 $。则存在 $ n $ 个变量，它们唯一地解出方程组

\begin{equation*}
\begin{array}{c}f^{1}(x^{1},\ldots,x^{m})=0\ ,\\ \vdots\\ f^{n}(x^{1},\ldots,x^{m})=0\end{array}
\end{equation*}

在 $ x_{0} $ 的邻域内作为其余 $ m-n $ 个变量的函数。当 $ f $ 属于 $ C^{q} $ 类时，这些解也是 $ C^{q} $ 函数。

\proofhead
由 (a)，我们有 $ m\ge n $。此外，我们可以通过适当地置换 $ \mathbb{R}^{m} $ 中的坐标，即通过 $ \mathbb{R}^{m} $ 的一个适当的正交变换，使得
\[
\frac{\partial(f^{1},\ldots,f^{n})}{\partial(x^{m-n+1},\ldots,x^{m})}(x_{0})\ne 0
\]

于是该注由推论 8.4 推出。\nobreak\hspace{0.35em}■

\item

设 $ 0\in\operatorname{im}(f) $ 是 $ f\in C^{q}(X,\mathbb{R}^{n}) $ 的正则值。则对每个 $ x_{0}\in f^{-1}(0) $，存在 $ \mathbb{R}^{m} $ 中的邻域 $ U $，使得 $ f^{-1}(0)\cap U $ 是 $ m-n $ 个变量的一个 $ C^{q} $ 函数的图像。

\proofhead
这由 (e) 与注 8.3 推出。\nobreak\hspace{0.35em}■
\end{items}

\end{remark}

\needspace{3\baselineskip}\subsection*{常微分方程}

在第 1 节中，我们利用指数映射来处理\emph{线性}微分方程的存在性与唯一性问题。接下来，我们处理形如 $ \dot{x}=f(t,x) $ 的微分方程，其中 $ f $ 现在可以是非线性的。

在本节其余部分，设

- $ J $ 是 $ \mathbb{R} $ 中的开区间；
  $ E $ 是巴拿赫空间；
  $ D $ 是 $ E $ 的开子集；
  $ f\in C(J\times D,E) $。

函数 $ u:J_{u}\to D $ 称为 $ E $ 中微分方程 $ \dot{x}=f(t,x) $ 的\textbf{解}，如果 $ J_{u} $ 是 $ J $ 的完全子区间，$ u\in C^{1}(J_{u},D) $，且

\[
\dot{u}(t)=f(t,u(t))\qquad\text{对 }t\in J_{u} \ .
\]

给定 $ (t_{0},x_{0})\in J\times D $，则这一对

\begin{equation}
\dot{x}=f(t,x)\text{ 且 }x(t_{0})=x_{0}\tag*{(8.5)${}_{(t_0,x_0)}$}
\end{equation}

成为 $ \dot{x}=f(t,x) $ 的一个\textbf{初值问题}。若 $ u $ 满足微分方程 $ \dot{x}=f(t,x) $ 且 $ u(t_{0})=x_{0} $，则映射 $ u:J_{u}\to D $ 是 $ (8.5)_{(t_{0},x_{0})} $ 的一个\textbf{解}。如果没有 $ (8.5)_{(t_{0},x_{0})} $ 的解 $ v:J_{v}\to D $ 使得 $ v\supset u $ 且 $ v\ne u $，则它是\textbf{不可延拓的}（或\textbf{极大的}）。此时 $ J_{u} $ 是 $ (8.5)_{(t_{0},x_{0})} $ 的一个\textbf{最大存在区间}。当 $ J_{u}=J $ 时，我们称 $ u $ 是 $ (8.5)_{(t_{0},x_{0})} $ 的一个\textbf{全局}解。

\begin{remark}

\begin{items}
\item

全局解是不可延拓的。

\item

设 $ J_{u} $ 是 $ J $ 的完全子区间。则 $ u\in C(J_{u},D) $ 是 $ (8.5)_{(t_{0},x_{0})} $ 的解当且仅当 $ u $ 满足积分方程

\[
u(t)=x_{0}+\int_{t_{0}}^{t}f(s,u(s))\,ds\qquad\text{对 }t\in J_{u} \ .
\]

\proofhead
因为 $ s\mapsto f(s,u(s)) $ 连续，该注由微积分基本定理容易推出。\nobreak\hspace{0.35em}■
\item

设 $ A\in\mathcal{L}(E) $ 且 $ g\in C(\mathbb{R},E) $。则初值问题

\[
\dot{x}=Ax+g(t)\text{ 且 }x(t_{0})=x_{0}
\]

有唯一的全局解

\[
u(t)=e^{(t-t_{0})A}x_{0}+\int_{t_{0}}^{t}e^{(t-s)A}g(s)\,ds\qquad\text{对 }t\in\mathbb{R} \ .
\]

\proofhead
这是定理 1.17 的推论。\nobreak\hspace{0.35em}■
\item

当 $ E=\mathbb{R} $ 时，称 $ \dot{x}=f(t,x) $ 为\textbf{标量}微分方程。

\item

初值问题

\begin{equation}
\dot{x}=x^{2}\text{ 且 }x(t_{0})=x_{0}\tag{8.6}
\end{equation}

对 $ (t_{0},x_{0})\in\mathbb{R}^{2} $ 至多有一个解。

\proofhead
设 $ u\in C^{1}(J_{u},\mathbb{R}) $ 与 $ v\in C^{1}(J_{v},\mathbb{R}) $ 是 (8.6) 的解。对 $ t\in J_{u}\cap J_{v} $，令 $ w(t):=u(t)-v(t) $，并设 $ I $ 是 $ J_{u}\cap J_{v} $ 的紧致子区间，且 $ t_{0}\in I $。则只需验证 $ w $ 在 $ I $ 上消失。由 (b)，我们有
\[
w(t)=u(t)-v(t)=\int_{t_{0}}^{t}\left(u^{2}(s)-v^{2}(s)\right)ds=\int_{t_{0}}^{t}(u(s)+v(s))w(s)\,ds\qquad\text{对 }t\in I \ .
\]

令 $ \alpha:=\max_{s\in I}|u(s)+v(s)| $，可得

\[
|w(t)|\le\alpha\left|\int_{t_{0}}^{t}|w(s)|\,ds\right|\qquad\text{对 }t\in I\ ,
\]

而该断言由格朗沃尔引理推出。\nobreak\hspace{0.35em}■

\item

初值问题 (8.6) 对 $ x_{0}\ne 0 $ 没有全局解。

\proofhead
设 $ u\in C^{1}(J_{u},\mathbb{R}) $ 是 (8.6) 的满足 $ x_{0}\ne 0 $ 的解。因为 0 是初值 $ (t_{0},0) $ 的唯一全局解，由 (e) 可得对 $ t\in J_{u} $ 有 $ u(t)\ne 0 $。因此 (8.6) 蕴含
\[
\frac{\dot{u}}{u^{2}}=\left(-\frac{1}{u}\right)^{\cdot}=1
\]

在 $ J_{u} $ 上。通过积分，我们得到

\[
\frac{1}{x_{0}}-\frac{1}{u(t)}=t-t_{0}\qquad\text{对 }t\in J_{u}\ ,
\]

因此

\[
u(t)=1/(t_{0}-t+1/x_{0})\qquad\text{对 }t\in J_{u} \ .
\]

于是我们得到 $ J_{u}\ne J=\mathbb{R} $。\nobreak\hspace{0.35em}■

\item

初值问题 (8.6) 对每个 $ (t_{0},x_{0})\in\mathbb{R}^{2} $ 有唯一的不可延拓解 $ u(\cdot,t_{0},x_{0})\in C^{1}(J(t_{0},x_{0}),\mathbb{R}) $，其中

\begin{equation*}
J(t_{0},x_{0})=\begin{cases}(t_{0}+1/x_{0},\infty)\ , & x_{0}<0\ ,\\ \mathbb{R}\ , & x_{0}=0\ ,\\ (-\infty,t_{0}+1/x_{0})\ , & x_{0}>0\ ,\end{cases}
\end{equation*}

且

\begin{equation*}
u(t,t_{0},x_{0})=\begin{cases}0\ , & x_{0}=0\ ,\\ 1/(t_{0}-t+1/x_{0})\ , & x_{0}\ne 0\ .\end{cases}
\end{equation*}

\proofhead
这由 (e) 与 (f) 的证明推出。\nobreak\hspace{0.35em}■
\item

标量初值问题

\begin{equation}
\dot{x}=2\sqrt{|x|}\ ,\quad x(0)=0\tag{8.7}
\end{equation}

有不可数多个全局解。

\proofhead
显然 $ u_{0}\equiv 0 $ 是一个全局解。对 $ \alpha<0<\beta $，令
\begin{equation*}
u_{\alpha,\beta}(t):=\begin{cases}-(t-\alpha)^{2}\ , & t\in(-\infty,\alpha]\ ,\\ 0\ , & t\in(\alpha,\beta)\ ,\\ (t-\beta)^{2}\ , & t\in[\beta,\infty)\ .\end{cases}
\end{equation*}

容易验证 $ u_{\alpha,\beta} $ 是 (8.7) 的一个全局解。\nobreak\hspace{0.35em}■

\end{items}

\bookfigure[0.30]{figures/ab018178b98553d7.png}
\end{remark}

\needspace{3\baselineskip}\subsection*{分离变量法}

一般不可能“显式地”解出给定的微分方程。不过，对某些标量微分方程，可以借助隐函数定理给出初值问题的唯一局部解。

函数 $ \Phi:J\times D\to\mathbb{R} $ 称为 $ \dot{x}=f(t,x) $ 的\textbf{首次积分}，如果对 $ \dot{x}=f(t,x) $ 的每个解 $ u:J_{u}\to D $，映射

\[
J_{u}\to\mathbb{R}\ ,\quad t\mapsto\Phi(t,u(t))
\]

是常值。若此外

\[
\Phi\in C^{1}(J\times D,\mathbb{R})\quad\text{且}\quad\partial_{2}\Phi(t,x)\ne 0\qquad\text{对 }(t,x)\in J\times D\ ,
\]

则 $ \Phi $ 是\textbf{正则首次积分}。

\begin{proposition}

设 $ \Phi $ 是标量微分方程 $ \dot{x}=f(t,x) $ 的正则首次积分，且

\[
f(t,x)=-\partial_{1}\Phi(t,x)/\partial_{2}\Phi(t,x)\qquad\text{对 }(t,x)\in J\times D \ .
\]

则对每个 $ (t_{0},x_{0})\in J\times D $，存在 $ J $ 的开区间 $ I $ 与 $ D $ 中的 $ U $，$ (t_{0},x_{0})\in I\times U $，使得初值问题

\[
\dot{x}=f(t,x)\ ,\quad x(t_{0})=x_{0}
\]

在 $ I $ 上恰有一个满足 $ u(I)\subset U $ 的解 $ u $。这通过对 $ x $ 解方程 $ \Phi(t,x)=\Phi(t_{0},x_{0}) $ 得到。

\proofhead
这由定理 8.2 立即推出。\nobreak\hspace{0.35em}■
\end{proposition}

通过适当地选取 $ \Phi $，我们得到一类重要的具有“分离变量”的标量微分方程。如下一推论所示，我们可以用求积法解这些方程。

\begin{proposition}

设 $ g\in C(J,\mathbb{R}) $ 且 $ h\in C(D,\mathbb{R}) $，对 $ x\in D $ 有 $ h(x)\ne 0 $。又设 $ G\in C^{1}(J,\mathbb{R}) $ 与 $ H\in C^{1}(D,\mathbb{R}) $ 分别是 $ g $ 与 $ h $ 的原函数。则

\[
\Phi(t,x):=G(t)-H(x)\qquad\text{对 }(t,x)\in J\times D\ ,
\]

是 $ \dot{x}=g(t)/h(x) $ 的正则首次积分。

\proofhead
显然 $ \Phi $ 属于 $ C^{1}(J\times D,\mathbb{R}) $，且 $ \partial_{2}\Phi(t,x)=-h(x)\ne 0 $。现在设 $ u\in C^{1}(J_{u},D) $ 是 $ \dot{x}=g(t)/h(x) $ 的解。则由链式法则可得
\begin{equation*}
\begin{aligned}
\left[\Phi(t,u(t))\right]^{\cdot}&=\partial_{1}\Phi(t,u(t))+\partial_{2}\Phi(t,u(t))\dot{u}(t)\\ &=g(t)-h(u(t))\cdot\frac{g(t)}{h(u(t))}=0\ .
\end{aligned}
\end{equation*}

因此 $ t\mapsto\Phi(t,u(t)) $ 在 $ J_{u} $ 上常值。\nobreak\hspace{0.35em}■

\end{proposition}

\begin{corollary}

（分离变量法） 设 $ g\in C(J,\mathbb{R}) $ 且 $ h\in C(D,\mathbb{R}) $，对 $ x\in D $ 有 $ h(x)\ne 0 $。则初值问题

\[
\dot{x}=g(t)/h(x)\ ,\quad x(t_{0})=x_{0}
\]

对每个 $ (t_{0},x_{0})\in J\times D $ 有唯一的局部解。该解可通过解方程

\begin{equation}
\int_{x_{0}}^{x}h(\xi)\,d\xi=\int_{t_{0}}^{t}g(\tau)\,d\tau\tag{8.8}
\end{equation}

关于 $ x $ 得到。

\proofhead
这是定理 8.7 与 8.8 的直接推论。\nobreak\hspace{0.35em}■
\end{corollary}

\begin{remark}

在推论 8.9 的假设下，

\[
\frac{dx}{dt}=\frac{g(t)}{h(x)}\ ,\quad x(t_{0})=x_{0}
\]

形式上蕴含 $ h(x)\,dx=g(t)\,dt $，即“分离变量方程”。形式积分然后给出 (8.8)。\nobreak\hspace{0.35em}■

\end{remark}

\begin{example}

\begin{items}
\item

我们考虑初值问题

\begin{equation}
\dot{x}=1+x^{2}\ ,\quad x(t_{0})=x_{0}\ .\tag{8.9}
\end{equation}

令 $ g:=1 $ 与 $ h:=1/(1+X^{2}) $，由推论 8.9 可得

\[
\arctan x-\arctan x_{0}=\int_{x_{0}}^{x}\frac{d\xi}{1+\xi^{2}}=\int_{t_{0}}^{t}d\tau=t-t_{0} \ .
\]

因此

\[
x(t)=\tan(t-\alpha)\qquad\text{对 }t\in(\alpha-\pi/2,\alpha+\pi/2)
\]

且 $ \alpha:=t_{0}-\arctan x_{0} $ 是 (8.9) 的唯一极大解。特别地，(8.9) 没有全局解。

\item

设 $ \alpha>0 $，$ D:=(0,\infty) $ 且 $ x_{0}>0 $。对

\begin{equation}
\dot{x}=x^{1+\alpha}\ ,\quad x(0)=x_{0}\ ,\tag{8.10}
\end{equation}

推论 8.9 蕴含

\[
\int_{x_{0}}^{x}\xi^{-1-\alpha}\,d\xi=t \ .
\]

因此

\[
x(t)=(x_{0}^{-\alpha}-\alpha t)^{-1/\alpha}\qquad\text{对 }-\infty<t<x_{0}^{-\alpha}/\alpha
\]

是 (8.10) 的唯一极大解，因此它不是全局可解的。

\item

在 $ D:=(1,\infty) $ 上，我们考虑初值问题

\begin{equation}
\dot{x}=x\log x\ ,\quad x(0)=x_{0}\ .\tag{8.11}
\end{equation}

因为 $ \log\circ\log $ 是 $ x\mapsto 1/(x\log x) $ 在 $ D $ 上的原函数，由推论 8.9 我们得到

\[
\log(\log x)-\log(\log x_{0})=\int_{x_{0}}^{x}\frac{d\xi}{\xi\log\xi}=\int_{0}^{t}d\tau=t \ .
\]

由此我们导出

\[
x(t)=x_{0}^{(e^{t})}\qquad\text{对 }t\in\mathbb{R}\ ,
\]

是 (8.11) 的唯一全局解。

\item

设 $ -\infty<a<b<\infty $ 且 $ f\in C((a,\infty),\mathbb{R}) $，使得对 $ x\in(a,b) $ 有 $ f(x)>0 $ 且 $ f(b)=0 $。对 $ x_{0}\in(a,b) $ 与 $ t_{0}\in\mathbb{R} $，我们考虑初值问题

\[
\dot{x}=f(x)\ ,\quad x(t_{0})=x_{0} \ .
\]

根据推论 8.9，我们通过解

\[
t=t_{0}+\int_{x_{0}}^{x}\frac{d\xi}{f(\xi)}=:H(x)
\]

关于 $ x $ 得到唯一局部解 $ u:J_{u}\to(a,b) $。因为 $ H'=1/f $，$ H $ 严格递增。因此 $ u=H^{-1} $，且 $ J_{u} $ 与 $ H((a,b)) $ 一致。此外

\[
T^{*}:=\lim_{x\to b-0}H(x)=\sup J_{u}
\]

在 $ \overline{\mathbb{R}} $ 中存在，且

\[
T^{*}<\infty\iff\int_{x_{0}}^{b}\frac{d\xi}{f(\xi)}<\infty \ .
\]
\end{items}

\bookfigure[0.93]{figures/7359315e1949390f.png}

现在 $ \lim_{t\to T^{*}-0}u(t)=b $ 蕴含

\[
\lim_{t\to T^{*}-0}\dot{u}(t)=\lim_{t\to T^{*}-0}f(u(t))=0 \ .
\]

在 $ T^{*}<\infty $ 的情形下，则由命题 IV.1.16 可得

\begin{equation*}
v(t):=\begin{cases}u(t)\ , & t\in J_{u}\ ,\\ b\ , & t\in[T^{*},\infty)\ ,\end{cases}
\end{equation*}

给出 $ u $ 的一个延拓。对

\[
T_{*}:=\lim_{x\to a+0}H(x)=\inf J_{u}
\]

的类似考虑留给你。

\begin{items}
\item[（e）]

对 $ a\in C(J,\mathbb{R}) $ 与 $ (t_{0},x_{0})\in J\times\mathbb{R} $，我们考虑初值问题

\begin{equation}
\dot{x}=a(t)x\ ,\quad x(t_{0})=x_{0}\tag{8.12}
\end{equation}

对\textbf{标量线性齐次微分方程} $ \dot{x}=a(t)x $，其\textbf{时变系数}为 $ a $。因此 (8.12) 有唯一的全局解

\[
x(t)=x_{0}e^{\int_{t_{0}}^{t}a(\tau)\,d\tau}\qquad\text{对 }t\in J \ .
\]

\proofhead
由格朗沃尔引理可得，(8.12) 至多有一个解。若 $ u\in C^{1}(J_{u},\mathbb{R}) $ 是一个解，且存在 $ t_{1}\in J_{u} $ 使得 $ u(t_{1})=0 $，则初值问题 $ \dot{x}=a(t)x $，$ x(t_{1})=0 $ 的唯一性结果蕴含 $ u=0 $。
因此设 $ x_{0}\ne 0 $。则由分离变量法我们得到
\[
\log|x|-\log|x_{0}|=\int_{x_{0}}^{x}\frac{d\xi}{\xi}=\int_{t_{0}}^{t}a(\tau)\,d\tau
\]

从而

\[
|x(t)|=|x_{0}|e^{\int_{t_{0}}^{t}a(\tau)\,d\tau}\qquad\text{对 }t\in J \ .
\]

因为对 $ t\in J $ 有 $ x(t)\ne 0 $，断言得证。\nobreak\hspace{0.35em}■

\end{items}

\end{example}

\needspace{3\baselineskip}\subsection*{利普希茨连续性与唯一性}

注 8.6(h) 表明，初值问题 $ (8.5)_{(t_{0},x_{0})} $ 一般不是唯一可解的。下面我们将看到，如果要求 $ f $ 比仅仅连续稍微正则一些，就能保证唯一性。

设 $ E $ 与 $ F $ 是巴拿赫空间，$ X\subset E $ 且 $ I\subset\mathbb{R} $。则称 $ f\in C(I\times X,F) $ 在 $ x\in X $ 中\textbf{局部利普希茨连续}，如果 $ I\times X $ 中的每个 $ (t_{0},x_{0}) $ 有邻域 $ U\times V $，使得

\[
\|f(t,x)-f(t,y)\|\le L\|x-y\|\qquad\text{对 }t\in U\text{ 且 }x,y\in V
\]

对某个 $ L\ge 0 $。我们令

\[
C^{0,1-}(I\times X,F):=\{f\in C(I\times X,F);f\text{ 在 }x\in X\text{ 中局部利普希茨连续}\} \ .
\]

如果 $ I $ 是单点，从而 $ f $ 有形式 $ f:X\to F $，我们称 $ f $ \textbf{局部利普希茨连续}，且

\[
C^{1-}(X,F):=\{f:X\to F;f\text{ 局部利普希茨连续}\} \ .
\]

\begin{remark}

\begin{items}
\item

显然，我们有 $ C^{0,1-}(I\times X,F)\subset C(I\times X,F) $。

\item

设 $ X $ 在 $ E $ 中开，且 $ f\in C(I\times X,F) $。假设 $ \partial_{2}f $ 存在且属于 $ C(I\times X,\mathcal{L}(E,F)) $。则 $ f\in C^{0,1-}(I\times X,F) $。

\proofhead
设 $ (t_{0},x_{0})\in I\times X $。因为 $ \partial_{2}f $ 属于 $ C(I\times X,\mathcal{L}(E,F)) $ 且 $ X $ 在 $ E $ 中开，存在 $ \varepsilon>0 $ 使得 $ \mathbb{B}(x_{0},\varepsilon)\subset X $ 且
\begin{equation}
\|\partial_{2}f(t_{0},x_{0})-\partial_{2}f(t,x)\|\le 1\qquad\text{对 }(t,x)\in U\times V\ ,\tag{8.13}
\end{equation}

其中我们令 $ U:=(t_{0}-\varepsilon,t_{0}+\varepsilon)\cap I $ 与 $ V:=\mathbb{B}(x_{0},\varepsilon) $。令 $ L:=1+\|\partial_{2}f(t_{0},x_{0})\| $，由中值定理（定理 3.9）与 (8.13) 可得

\[
\|f(t,x)-f(t,y)\|\le\sup_{0\le s\le 1}\|\partial_{2}f(t,x+s(y-x))\|\,\|x-y\|\le L\|x-y\|
\]

对 $ (t,x),(t,y)\in U\times V $。因此 $ f $ 在 $ x $ 中局部利普希茨连续。\nobreak\hspace{0.35em}■

\item

次数 $ \ge 2 $ 的多项式是局部利普希茨连续的，但不是利普希茨连续的。

\item

设 $ I\times X $ 紧致且 $ f\in C^{0,1-}(I\times X,F) $。则 $ f $ 在 $ x\in X $ 中\textbf{一致利普希茨连续}，即存在 $ L\ge 0 $ 使得

\[
\|f(t,x)-f(t,y)\|\le L\|x-y\|\qquad\text{对 }x,y\in X\text{ 且 }t\in I \ .
\]

\proofhead
对每个 $ (t,x)\in I\times X $，存在 $ \varepsilon_{t}>0 $、$ \varepsilon_{x}>0 $ 与 $ L(t,x)\ge 0 $，使得
\[
\|f(s,y)-f(s,z)\|\le L(t,x)\|y-z\|\qquad\text{对 }(s,y),(s,z)\in\mathbb{B}(t,\varepsilon_{t})\times\mathbb{B}(x,\varepsilon_{x}) \ .
\]

因为 $ I\times X $ 紧致，存在 $ (t_{0},x_{0}),\ldots,(t_{m},x_{m})\in I\times X $，使得

\[
I\times X\subset\bigcup_{j=0}^{m}\mathbb{B}(t_{j},\varepsilon_{t_{j}})\times\mathbb{B}(x_{j},\varepsilon_{x_{j}}/2) \ .
\]

因为 $ f(I\times X) $ 紧致，存在 $ R>0 $ 使得 $ f(I\times X)\subset R\mathbb{B} $。我们令

\[
\delta:=\min\{\varepsilon_{x_{0}},\ldots,\varepsilon_{x_{m}}\}/2>0
\]

且

\[
L:=\max\{L(t_{0},x_{0}),\ldots,L(t_{m},x_{m}),2R/\delta\}>0 \ .
\]

现在设 $ (t,x),(t,y)\in I\times X $。则存在 $ k\in\{0,\ldots,m\} $ 使得

\[
(t,x)\in\mathbb{B}(t_{k},\varepsilon_{t_{k}})\times\mathbb{B}(x_{k},\varepsilon_{x_{k}}/2) \ .
\]

若 $ \|x-y\|<\delta $，我们发现 $ (t,y) $ 位于 $ \mathbb{B}(t_{k},\varepsilon_{t_{k}})\times\mathbb{B}(x_{k},\varepsilon_{x_{k}}) $ 中，且得到

\[
\|f(t,x)-f(t,y)\|\le L(t_{k},x_{k})\|x-y\|\le L\|x-y\| \ .
\]

另一方面，若 $ \|x-y\|\ge\delta $，则可得

\[
\|f(t,x)-f(t,y)\|\le\frac{2R}{\delta}\cdot\delta\le L\|x-y\| \ .
\]

因此 $ f $ 在 $ x\in X $ 中一致利普希茨连续。\nobreak\hspace{0.35em}■

\item

设 $ X $ 在 $ E $ 中紧致且 $ f\in C^{1-}(X,F) $。则 $ f $ 是利普希茨连续的。

\proofhead
这是 (d) 的特殊情形。\nobreak\hspace{0.35em}■
\end{items}

\end{remark}

\begin{theorem}

设 $ f\in C^{0,1-}(J\times D,E) $，并设 $ u:J_{u}\to D $ 与 $ v:J_{v}\to D $ 是 $ \dot{x}=f(t,x) $ 的解，且对某个 $ t_{0}\in J_{u}\cap J_{v} $ 有 $ u(t_{0})=v(t_{0}) $。则对每个 $ t\in J_{u}\cap J_{v} $ 有 $ u(t)=v(t) $。

\proofhead
设 $ I\subset J_{u}\cap J_{v} $ 是紧致区间，$ t_{0}\in I $，且 $ w:=u-v $。只需验证 $ w|I=0 $。

因为 $ K:=u(I)\cup v(I)\subset D $ 紧致，注 8.12(d) 保证存在 $ L\ge 0 $ 使得
\[
\|f(s,u(s))-f(s,v(s))\|\le L\|u(s)-v(s)\|=L\|w(s)\|\qquad\text{对 }s\in I \ .
\]

因此，由

\[
w(t)=u(t)-v(t)=\int_{t_{0}}^{t}\left(f(s,u(s))-f(s,v(s))\right)ds\qquad\text{对 }t\in J_{u}\cap J_{v}\ ,
\]

（见注 8.6(b)）不等式

\[
\|w(t)\|\le L\left|\int_{t_{0}}^{t}\|w(s)\|\,ds\right|\qquad\text{对 }t\in I
\]

成立，且格朗沃尔引理蕴含 $ w|I=0 $。\nobreak\hspace{0.35em}■

\end{theorem}

\needspace{3\baselineskip}\subsection*{皮卡–林德勒夫定理}

在下文中，设

- $ J\subset\mathbb{R} $ 是开区间；
  $ E $ 是有限维巴拿赫空间；
  $ D $ 在 $ E $ 中开；
  $ f\in C^{0,1-}(J\times D,E) $。

我们现在证明常微分方程的基本局部存在性与唯一性定理。

\begin{theorem}

（皮卡–林德勒夫） 设 $ (t_{0},x_{0})\in J\times D $。则存在 $ \alpha>0 $，使得初值问题

\begin{equation}
\dot{x}=f(t,x)\ ,\quad x(t_{0})=x_{0}\tag{8.14}
\end{equation}

在 $ I:=[t_{0}-\alpha,t_{0}+\alpha] $ 上有唯一解。

证明

\begin{subitems}
\item

因为 $ J\times D\subset\mathbb{R}\times E $ 开，存在 $ a,b>0 $ 使得

\[
R:=[t_{0}-a,t_{0}+a]\times\overline{\mathbb{B}}(x_{0},b)\subset J\times D \ .
\]

由 $ f $ 在 $ x $ 中的局部利普希茨连续性，我们找到 $ L>0 $ 使得

\[
\|f(t,x)-f(t,y)\|\le L\|x-y\|\qquad\text{对 }(t,x),(t,y)\in R \ .
\]

最后，由 $ R $ 的紧致性与注 8.12(a) 可得，存在 $ M>0 $ 使得

\[
\|f(t,x)\|\le M\qquad\text{对 }(t,x)\in R \ .
\]
\item

设 $ \alpha:=\min\{a,b/M,1/(2L)\}>0 $ 且 $ I:=[t_{0}-\alpha,t_{0}+\alpha] $。我们令

\[
T(y)(t):=x_{0}+\int_{t_{0}}^{t}f(\tau,y(\tau))\,d\tau\qquad\text{对 }y\in C(I,D)\ ,\quad t\in I \ .
\]

如果我们能证明映射 $ T:C(I,D)\to C(I,E) $ 恰有一个不动点 $ u $，则由注 8.6(b) 可得 $ u $ 是 (8.14) 的唯一解。

\item

令

\[
X:=\{y\in C(I,E);y(t_{0})=x_{0},\max_{t\in I}\|y(t)-x_{0}\|\le b\} \ .
\]

则 $ X $ 是巴拿赫空间 $ C(I,E) $ 的闭子集，因此是完备度量空间（见习题 II.6.4）。
对 $ y\in X $，对 $ t\in I $ 有 $ y(t)\in\overline{\mathbb{B}}(x_{0},b)\subset D $。因此 $ T(y) $ 有定义，且 $ T(y)(t_{0})=x_{0} $。此外

\[
\|T(y)(t)-x_{0}\|=\left\|\int_{t_{0}}^{t}f(\tau,y(\tau))\,d\tau\right\|\le\alpha\sup_{(t,\xi)\in R}\|f(t,\xi)\|\le\frac{b}{M}\cdot M=b \ .
\]

这表明 $ T $ 把空间 $ X $ 映入自身。

\item

对 $ y,z\in X $，我们得到

\begin{equation*}
\begin{aligned}
\|T(y)(t)-T(z)(t)\|&=\left\|\int_{t_{0}}^{t}\left(f(\tau,y(\tau))-f(\tau,z(\tau))\right)d\tau\right\|\\ &\le\alpha\max_{t\in I}\|f(t,y(t))-f(t,z(t))\|\le\alpha L\max_{t\in I}\|y(t)-z(t)\|
\end{aligned}
\end{equation*}

对 $ t\in I $。利用 $ \alpha $ 的定义可得

\[
\|T(y)-T(z)\|_{C(I,E)}\le\frac{1}{2}\|y-z\|_{C(I,E)} \ .
\]

因此 $ T:X\to X $ 是压缩映射。于是压缩定理（定理 IV.4.3）给出 $ X $ 中的唯一不动点 $ u $。\nobreak\hspace{0.35em}■

\end{subitems}

\end{theorem}

\begin{remark}

\begin{items}
\item

(8.14) 在 $ I $ 上的解可以用逐次逼近法（或迭代地）通过令

\[
u_{m+1}(t):=x_{0}+\int_{t_{0}}^{t}f(\tau,u_{m}(\tau))\,d\tau\qquad\text{对 }m\in\mathbb{N}\text{ 且 }t\in I\ ,
\]

来计算，其中对 $ t\in I $ 有 $ u_{0}(t)=x_{0} $。序列 $ (u_{m}) $ 在 $ I $ 上一致收敛到 $ u $，且我们有误差估计

\begin{equation}
\|u_{m}-u\|_{C(I,E)}\le\alpha M/2^{m-1}\qquad\text{对 }m\in\mathbb{N}^{\times} \ .\tag{8.15}
\end{equation}

\begin{subitems}
\item

\proofhead
第一个断言由上一证明与定理 IV.4.3(ii) 立即推出。该定理的断言 (iii) 还给出误差估计
\[
\|u_{m}-u\|_{C(I,E)}\le 2^{-m+1}\|u_{1}-u_{0}\|_{C(I,E)} \ .
\]

现在该断言成立，因为对 $ t\in I $ 我们有

\[
\|u_{1}(t)-u_{0}(t)\|=\left\|\int_{t_{0}}^{t}f(\tau,u_{0}(\tau))\,d\tau\right\|\le\alpha M \ .\quad\blacksquare
\]
\end{subitems}

\item

误差估计 (8.15) 可改进为

\[
\|u_{m}-u\|_{C(I,E)}<\frac{M\alpha\sqrt{e}}{2^{m}(m+1)!}\qquad\text{对 }m\in\mathbb{N}\ ,
\]

（见习题 11）。

\item

$ E $ 有限维的假设只用于证明 $ f(R) $ 上界 $ M $ 的存在性。因此，若把“$ E $ 有限维”的假设换成“$ f $ 在有界集上有界”，则定理 8.14 及其证明仍然正确。

\item

虽然 $ f $ 的连续性不足以保证解 (8.14) 的唯一性，但它足以证明解的存在性（见 [Ama95, 定理 II.7.3]）。\nobreak\hspace{0.35em}■

\end{items}

最后我们证明，定理 8.14 的局部解可以延拓为唯一的不可延拓解。

\end{remark}

\begin{theorem}

对每个 $ (t_{0},x_{0})\in J\times D $，初值问题

\[
\dot{x}=f(t,x)\ ,\quad x(t_{0})=x_{0} \ .
\]

恰有一个不可延拓解 $ u(\cdot,t_{0},x_{0}):J(t_{0},x_{0})\to D $。最大存在区间是开的，即 $ J(t_{0},x_{0})=(t^{-}(t_{0},x_{0}),t^{+}(t_{0},x_{0})) $。

\proofhead
设 $ (t_{0},x_{0})\in J\times D $。根据定理 8.14，存在 $ \alpha_{0} $ 与 $ (8.5)_{(t_{0},x_{0})} $ 在 $ I_{0}:=[t_{0}-\alpha_{0},t_{0}+\alpha_{0}] $ 上的唯一解 $ u $。我们令 $ x_{1}:=u(t_{0}+\alpha_{0}) $ 与 $ t_{1}:=t_{0}+\alpha_{0} $，并把定理 8.14 应用于初值问题 $ (8.5)_{(t_{1},x_{1})} $。则存在 $ \alpha_{1}>0 $ 与 $ (8.5)_{(t_{1},x_{1})} $ 在 $ I_{1}:=[t_{1}-\alpha_{1},t_{1}+\alpha_{1}] $ 上的唯一解 $ v $。定理 8.13 表明对 $ t\in I_{0}\cap I_{1} $ 有 $ u(t)=v(t) $。因此
\begin{equation*}
u_{1}(t):=\begin{cases}u(t)\ , & t\in I_{0}\ ,\\ v(t)\ , & t\in I_{1}\ ,\end{cases}
\end{equation*}

是 $ (8.5)_{(t_{0},x_{0})} $ 在 $ I_{0}\cup I_{1} $ 上的解。类似的论证表明 $ u $ 可以向左延拓越过 $ t_{0}-\alpha_{0} $。现在设

\[
t^{+}:=t^{+}(t_{0},x_{0}):=\sup\{\beta\in\mathbb{R};(8.5)_{(t_{0},x_{0})}\text{ 在 }[t_{0},\beta]\text{ 上有解}\}
\]

且

\[
t^{-}:=t^{-}(t_{0},x_{0}):=\inf\{\gamma\in\mathbb{R};(8.5)_{(t_{0},x_{0})}\text{ 在 }[\gamma,t_{0}]\text{ 上有解}\} \ .
\]

上面的考虑表明 $ t^{+}\in(t_{0},\infty] $ 与 $ t^{-}\in[-\infty,t_{0}) $ 有良好定义，且 $ (8.5)_{(t_{0},x_{0})} $ 在 $ (t^{-},t^{+}) $ 上有不可延拓解 $ u $。现在由定理 8.13 可得 $ u $ 是唯一的。\nobreak\hspace{0.35em}■

\end{theorem}

\begin{example}

\begin{items}
\item

（$ m $ 阶常微分方程） 设 $ E:=\mathbb{R}^{m} $ 且 $ g\in C^{0,1-}(J\times D,\mathbb{R}) $。则

\begin{equation}
x^{(m)}=g(t,x,\dot{x},\ldots,x^{(m-1)})\tag{8.16}
\end{equation}

是 \textbf{$ m $ 阶常微分方程}。函数 $ u:J_{u}\to\mathbb{R} $ 是 (8.14) 的\textbf{解}，如果 $ J_{u}\subset J $ 是完全区间，$ u $ 属于 $ C^{m}(J_{u},\mathbb{R}) $，

\[
\left(u(t),\dot{u}(t),\ldots,u^{(m-1)}(t)\right)\in D\qquad\text{对 }t\in J_{u}\ ,
\]

且

\[
u^{(m)}(t)=g\left(t,u(t),\dot{u}(t),\ldots,u^{(m-1)}(t)\right)\qquad\text{对 }t\in J_{u} \ .
\]

对 $ t_{0}\in J $ 与 $ x_{0}:=(x_{0}^{0},\ldots,x_{0}^{m-1})\in D $，对

\begin{equation}
\begin{array}{l}x^{(m)}=g\left(t,x,\dot{x},\ldots,x^{(m-1)}\right)\\ x(t_{0})=x_{0}^{0},\ldots,x^{(m-1)}(t_{0})=x_{0}^{m-1}\end{array}\tag*{(8.17)${}_{(t_0,x_0)}$}
\end{equation}

称为 (8.16) 的具有初值 $ x_{0} $ 与初时刻 $ t_{0} $ 的\textbf{初值问题}。
在这种语言下我们有：对每个 $ (t_{0},x_{0})\in J\times D $，存在 $ (8.17)_{(t_{0},x_{0})} $ 的唯一极大解 $ u:J(t_{0},x_{0})\to\mathbb{R} $。最大存在区间 $ J(t_{0},x_{0}) $ 是开的。

\proofhead
我们通过
\begin{equation}
f(t,y):=\left(y^{2},y^{3},\ldots,y^{m},g(t,y)\right)\qquad\text{对 }t\in J\text{ 且 }y=(y^{1},\ldots,y^{m})\in D \ .\tag{8.18}
\end{equation}

来定义 $ f:J\times D\to\mathbb{R}^{m} $。则 $ f $ 属于 $ C^{0,1-}(J\times D,\mathbb{R}^{m}) $。利用定理 8.16，存在 $ (8.5)_{(t_{0},x_{0})} $ 的唯一不可延拓解 $ z:J(t_{0},x_{0})\to\mathbb{R}^{m} $。然后可以容易验证 $ u:=\operatorname{pr}_{1}\circ z:J(t_{0},x_{0})\to\mathbb{R} $ 是 $ (8.17)_{(t_{0},x_{0})} $ 的解。
反之，若 $ v:J(t_{0},x_{0})\to\mathbb{R} $ 是 $ (8.17)_{(t_{0},x_{0})} $ 的解，则向量函数 $ (v,\dot{v},\ldots,v^{(m-1)}) $ 在 $ J(t_{0},x_{0}) $ 上解出初值问题 $ (8.5)_{(t_{0},x_{0})} $，其中 $ f $ 由 (8.18) 定义。由此可得 $ v=u $。类似的论证表明 $ u $ 不可延拓。\nobreak\hspace{0.35em}■

\item

（一维牛顿运动方程） 设 $ V $ 在 $ \mathbb{R} $ 中开且 $ f\in C^{1}(V,\mathbb{R}) $。对某个 $ x_{0}\in V $，设

\[
U(x):=\int_{x_{0}}^{x}f(\xi)\,d\xi\text{ 对 }x\in V\qquad\text{且}\qquad T(y):=y^{2}/2\text{ 对 }y\in\mathbb{R} \ .
\]

最后令 $ D:=V\times\mathbb{R} $ 与 $ L:=T-U $。根据例 6.14(a)，微分方程

\begin{equation}
-\ddot{x}=f(x)\tag{8.19}
\end{equation}

是受保守力 $ -\nabla U=f $ 作用的质点的（一维）运动的牛顿运动方程。由 (a)，我们知道 (8.19) 等价于系统

\begin{equation}
\dot{x}=y\ ,\quad\dot{y}=-f(x)\tag{8.20}
\end{equation}

因此等价于一阶微分方程

\begin{equation}
\dot{u}=F(u)\ ,\tag{8.21}
\end{equation}

其中对 $ u=(x,y)\in D $ 有 $ F(u):=(y,-f(x)) $。
函数 $ E:=T+U:D\to\mathbb{R} $，称为\textbf{总能量}，是 (8.21) 的首次积分。这意味着质点的运动\textbf{能量守恒}\footnote{见习题 6.12。}，即 (8.21) 的每个解 $ x\in C^{2}(J,V) $ 对每个 $ t\in J $ 与 $ t_{0}\in J $ 都满足 $ E(x(t),\dot{x}(t))=E(x(t_{0}),\dot{x}(t_{0})) $。

\proofhead
显然 $ E $ 属于 $ C^{1}(D,\mathbb{R}) $。(8.20) 的每个解 $ u:J_{u}\to D $ 都有
\[
\left[E(u(t))\right]^{\cdot}=\left[y^{2}(t)/2+U(x(t))\right]^{\cdot}=y(t)\dot{y}(t)+f(x(t))\dot{x}(t)=0\qquad\text{对 }t\in J_{u} \ .
\]

因此 $ E(u) $ 是常值。\nobreak\hspace{0.35em}■

\item

由 (b)，(8.20) 的每个解都位于总能量的一个水平集 $ E^{-1}(c) $ 中。
反之，存在性断言 (a) 意味着对 $ E $ 的某个水平集上的每一点 $ (x_{0},y_{0}) $，存在 (8.20) 的满足 $ u(0)=(x_{0},y_{0}) $ 的解。因此，$ E $ 的水平集的全体与 (8.20) 的所有（极大）解曲线之集相同；这一集合称为 (8.19) 的\textbf{相图}。

\end{items}

\bookfigure[0.43]{figures/92bc15fb8492707f.png}

相图具有以下性质：

\begin{subitems}
\item

$ E $ 的临界点恰是满足 $ f(x)=0 $ 的点 $ (x,0)\in D $。这意味着 (8.20) 的驻点位于 $ x $ 轴上，且恰是势 $ U $ 的临界点。

\item

每个水平集关于 $ x $ 轴对称。

\item

若 $ c $ 是 $ E $ 的正则值，则 $ E^{-1}(c) $ 可局部表示为某个 $ C^{2} $ 函数的图像。更精确地，对每个 $ u_{0}\in E^{-1}(c) $，存在正数 $ b $ 与 $ \varepsilon $ 以及 $ \varphi\in C^{2}((-\varepsilon,\varepsilon),\mathbb{R}) $，$ \varphi(0)=0 $，使得

\begin{equation}
E^{-1}(c)\cap\mathbb{B}(u_{0},b)=\operatorname{tr}(g)\ ,\tag{8.22}
\end{equation}

其中 $ g\in C^{2}((-\varepsilon,\varepsilon),\mathbb{R}^{2}) $ 定义为

\begin{equation*}
g(s):=\begin{cases}(s,\varphi(s))+u_{0}\ , & \partial_{2}E(u_{0})\ne 0\ ,\\ (\varphi(s),s)+u_{0}\ , & \partial_{1}E(u_{0})\ne 0\ .\end{cases}
\end{equation*}
\item

设 $ (x_{0},0) $ 是 $ E $ 的正则点且 $ c:=U(x_{0}) $。则水平集 $ E^{-1}(c) $ 与 $ x $ 轴正交相交。

\proofhead
(i) 这由 $ \nabla E(x,y)=(y,f(x)) $ 推出。
\end{subitems}

\begin{subitems}
\item[（ii）]

该陈述是 $ E(x,-y)=E(x,y) $（$ (x,y)\in D $）的推论。

\item[（iii）]

这由注 8.4(e) 与 (f) 推出。

\item[（iv）]

我们沿用 (iii) 的记号。因为

\[
\nabla E(x_{0},0)=(f(x_{0}),0)\ne(0,0)\ ,
\]

我们看到 (8.22) 对 $ u_{0}:=(x_{0},0) $ 与 $ g(s):=(\varphi(s)+x_{0},s) $（$ s\in(-\varepsilon,\varepsilon) $）成立。此外，由例 3.6(c) 可得

\[
0=\left(\nabla E(x_{0},0)\,\middle|\,\dot{g}(0)\right)=\left((f(x_{0}),0)\,\middle|\,(\dot{\varphi}(0),1)\right)=f(x_{0})\dot{\varphi}(0) \ .
\]

因此 $ \dot{\varphi}(0)=0 $，即 $ E^{-1}(c) $ 在 $ (x_{0},0) $ 处的切线与 $ y $ 轴平行（见注 IV.1.4(a)）。\nobreak\hspace{0.35em}■

\end{subitems}

\end{example}

\exerciseshead

\begin{exercises}
\item

证明方程组

\begin{equation*}
\begin{aligned}
x^{2} + y^{2} - u^{2} - v &= 0 , \\ x^{2} + 2y^{2} + 3u^{2} + 4v^{2} &= 1
\end{aligned}
\end{equation*}

在 $ (1/2,0,1/2,0) $ 附近可解出 $ (u,v) $。$ u $ 与 $ v $ 关于 $ (x,y) $ 的一阶导数是什么？

\item

设 $ E $、$ F $、$ G $ 与 $ H $ 为巴拿赫空间，且 $ X $ 是 $ H $ 中的开集。进一步设

\[
A \in C^{1}(X,\mathcal{L}(E,F)) \quad \text{且} \quad B \in C^{1}(X,\mathcal{L}(E,G))
\]

且

\[
(A(x),B(x)) \in \mathcal{L}\mathrm{is}(E,F\times G) \quad \text{对 } x \in X .
\]

最后，设 $ (f,g) \in F\times G $ 且

\[
\varphi : X \to E , \quad x \mapsto (A(x),B(x))^{-1}(f,g) .
\]

证明

\[
\partial\varphi(x)h = -S(x)\partial A(x)[h,\varphi(x)] - T(x)\partial B(x)[h,\varphi(x)] \quad \text{对 } (x,h) \in X\times H ,
\]

其中，对 $ x \in X $，

\[
S(x) := (A(x),B(x))^{-1}\big|(F\times\{0\}) \quad \text{且} \quad T(x) := (A(x),B(x))^{-1}\big|(\{0\}\times G) .
\]
\item

求具有\textbf{时变系数} $ a,b \in C(J,\mathbb{R}) $ 的\textbf{标量线性非齐次微分方程}

\[
\dot{x} = a(t)x + b(t)
\]

的通解。

\item

设 $ D $ 是 $ \mathbb{R} $ 中的开集，且 $ f \in C(D,\mathbb{R}) $。证明\textbf{相似微分方程}

\[
\dot{x} = f(x/t)
\]

等价于可分离变量微分方程

\[
\dot{y} = (f(y)-y)/t
\]

利用变换 $ y := x/t $。

\item

$ \dot{x} = (x^{2}+tx+t^{2})/t^{2} $ 的通解是什么？（提示：习题 4。）

\item

求初值问题 $ \dot{x}+tx^{2} = 0 $，$ x(0) = 2 $ 的解。

\item

设 $ a,b \in C(J,\mathbb{R}) $ 且 $ \alpha \ne 1 $。证明\textbf{伯努利微分方程}

\[
\dot{x} = a(t)x + b(t)x^{\alpha}
\]

在变换 $ y := x^{1-\alpha} $ 下化为线性微分方程

\[
\dot{y} = (1-\alpha)(a(t)y+b(t))
\]
\item

计算\textbf{逻辑斯蒂微分方程}

\[
\dot{x} = (\alpha-\beta x)x \quad \text{对 } \alpha,\beta > 0 .
\]

的通解 $ u $。并求 $ \lim_{t\to t^{+}} u(t) $ 以及 $ u $ 的拐点。

\item

利用代换 $ y = x/t $，求 $ \dot{x} = (t+x)/(t-x) $ 的通解。

\item

设 $ f \in C^{1-}(D,E) $，并令 $ u : J(\xi) \to \mathbb{R} $ 表示如下初值问题的不可延拓解：

\[
\dot{x} = f(x) , \quad x(0) = \xi .
\]

若 $ f $ 由

\begin{equation*}
\begin{aligned}
& \mathbb{R} \to \mathbb{R} , \quad x \mapsto x^{1+\alpha} , \quad \alpha \ge 0 ; \\ & E \to E , \quad x \mapsto Ax , \quad A \in \mathcal{L}(E) ; \\ & \mathbb{R} \to \mathbb{R} , \quad x \mapsto 1 + x^{2} .
\end{aligned}
\end{equation*}

给出，求 $ \mathcal{D}(f) := \{(t,\xi) \in \mathbb{R}\times D \; ; \; t \in J(\xi)\} $。

\item

证明注 8.14(b) 的误差估计。（提示：采用定理 8.13 证明中的记号，由归纳法可得

\[
\|u_{m+1}(t)-u_{m}(t)\| \le ML^{m}\frac{|t-t_{0}|^{m+1}}{(m+1)!} \quad \text{对 } m \in \mathbb{N} \text{ 且 } t \in I ,
\]

从而 $ \|u_{m+1}-u_{m}\|_{C(I,E)} < M\alpha 2^{-m}/(m+1)! $ 对 $ m \in \mathbb{N} $。）

\item

设 $ A \in C(J,\mathcal{L}(E)) $ 且 $ b \in C(J,E) $。证明初值问题

\[
\dot{x} = A(t)x + b(t) , \quad x(t_{0}) = x_{0}
\]

对每个 $ (t_{0},x_{0}) \in J\times E $ 都有唯一全局解。（提示：迭代一个积分方程，并利用习题 IV.4.1 与 $ \int_{t_{0}}^{t}\int_{t_{0}}^{s} d\sigma\, ds = (t-t_{0})^{2}/2 $。）

\end{exercises}

\section{流形}

从注 8.5(e) 我们知道，非线性方程组在正则点附近的解集可以用图像来描述。在本节中，我们将更精确地研究 $ \mathbb{R}^{n} $ 中可局部地通过图像表示的子集；即，我们研究 $ \mathbb{R}^{n} $ 的子流形。我们引入正则参数化与局部坐标卡这两个重要概念，它们使我们能够用函数局部地表示子流形。

• 在本节全节中，$ q $ 属于 $ \mathbb{N}^{\times} \cup \{\infty\} $。

\needspace{3\baselineskip}\subsection*{$ \mathbb{R}^{n} $ 的子流形}

$ \mathbb{R}^{n} $ 的子集 $ M $ 称为 $ \mathbb{R}^{n} $ 的 \textbf{$ m $ 维 $ C^{q} $ 子流形}，如果对每个 $ x_{0} \in M $，在 $ \mathbb{R}^{n} $ 中存在 $ x_{0} $ 的开邻域 $ U $、$ \mathbb{R}^{n} $ 中的开集 $ V $ 以及 $ \varphi \in \operatorname{Diff}^{q}(U, V) $，使得 $ \varphi(U \cap M) = V \cap (\mathbb{R}^{m} \times \{0\}) $。

\bookfigure[0.89]{figures/3291f3a227d1d7a5.png}

$ \mathbb{R}^{n} $ 的一维与二维子流形分别称为 $ \mathbb{R}^{n} $ 中的（\textbf{嵌入}）\textbf{曲线}与（\textbf{嵌入}）\textbf{曲面}。$ \mathbb{R}^{n} $ 的 $ n - 1 $ 维（或\textbf{余维数} 1 的）子流形称为（嵌入）\textbf{超曲面}（在 $ \mathbb{R}^{n} $ 中）。代替“$ \mathbb{R}^{n} $ 的 $ C^{q} $ 子流形”这一说法，我们常——尤其当“环绕空间” $ \mathbb{R}^{n} $ 不重要时——简单地说 $ C^{q} $ 流形。\footnote{注意到空集是每个维数 $ \le n $ 的流形。非空流形的维数是唯一确定的，正如注 9.16(a) 将表明的那样。}

$ \mathbb{R}^{n} $ 的 $ m $ 维子流形是在 $ \mathbb{R}^{n} $ 中局部定义的，即用开邻域来定义。至多相差一个小的形变，每个这样的邻域在 $ \mathbb{R}^{n} $ 中的位置，与 $ \mathbb{R}^{m} $ 作为向量子空间含于 $ \mathbb{R}^{n} $ 中时相同。

\begin{example}

\begin{items}
\item

$ \mathbb{R}^{n} $ 的子集 $ X $ 是 $ \mathbb{R}^{n} $ 的 $ n $ 维 $ C^{\infty} $ 子流形当且仅当 $ X $ 在 $ \mathbb{R}^{n} $ 中开。

\proofhead
如果 $ X $ 是 $ \mathbb{R}^{n} $ 的 $ n $ 维 $ C^{\infty} $ 子流形且 $ x_{0} \in X $，则存在 $ x_{0} $ 的开邻域 $ U $、$ \mathbb{R}^{n} $ 中的开集 $ V $ 以及 $ \varphi \in \operatorname{Diff}^{\infty}(U, V) $，使得 $ \varphi(U \cap X) = V $。因此 $ U \cap X = \varphi^{-1}(V) = U $。于是 $ U $ 属于 $ X $，这表明 $ X $ 是开集。
现在设 $ X $ 在 $ \mathbb{R}^{n} $ 中开。我们令 $ U := X $、$ V := X $ 且 $ \varphi := \operatorname{id}_{X} $。于是我们看到 $ X $ 是 $ \mathbb{R}^{n} $ 的 $ n $ 维 $ C^{\infty} $ 子流形。 \nobreak\hspace{0.35em}■
\item

设 $ M := \{x_{0}, \ldots, x_{k}\} \subset \mathbb{R}^{n} $。则 $ M $ 是 $ \mathbb{R}^{n} $ 的 0 维 $ C^{\infty} $ 子流形。

\proofhead
我们令 $ \alpha := \min\{|x_{i} - x_{j}| ; \ 0 \le i, j \le k, \ i \ne j\} $ 并选取 $ y \in M $。则 $ \mathbb{B}(y, \alpha) $ 是 $ y $ 在 $ \mathbb{R}^{n} $ 中的开邻域，且对满足 $ x \in \mathbb{B}(y, \alpha) $ 的 $ \varphi(x) := x - y $，我们有 $ \varphi \in \operatorname{Diff}^{\infty}(\mathbb{B}(y, \alpha), \alpha\mathbb{B}) $ 且 $ \varphi(\mathbb{B}(y, \alpha) \cap M) = \{0\} $。 \nobreak\hspace{0.35em}■
\item

设 $ \psi \in \operatorname{Diff}^{q}(\mathbb{R}^{n}, \mathbb{R}^{n}) $ 且 $ M $ 是 $ \mathbb{R}^{n} $ 的 $ m $ 维 $ C^{q} $ 子流形。则 $ \psi(M) $ 是 $ \mathbb{R}^{n} $ 的 $ m $ 维 $ C^{q} $ 子流形。

\proofhead
我们把这留给你作为习题。 \nobreak\hspace{0.35em}■
\item

$ \mathbb{R}^{n} $ 的每个 $ C^{q} $ 子流形对 $ 1 \le r \le q $ 也是 $ \mathbb{R}^{n} $ 的 $ C^{r} $ 子流形。 \nobreak\hspace{0.35em}■

\item

上述定义中的微分同胚 $ \varphi $ 与 $ \mathbb{R}^{n} $ 的开集 $ V $ 可以这样选取，使得 $ \varphi(x_{0}) = 0 $。

\proofhead
只需把给定的 $ \varphi $ 与 $ C^{\infty} $ 微分同胚 $ y \mapsto y - \varphi(x_{0}) $ 复合即可。 \nobreak\hspace{0.35em}■
\end{items}

\end{example}

\needspace{3\baselineskip}\subsection*{图像}

下一个定理表明图像是流形，从而给出了一大类例子。

\begin{proposition}

设 $ X $ 在 $ \mathbb{R}^{m} $ 中开且 $ f \in C^{q}(X, \mathbb{R}^{n}) $。则 $ \operatorname{graph}(f) $ 是 $ \mathbb{R}^{m+n} $ 的 $ m $ 维 $ C^{q} $ 子流形。

\proofhead
我们令 $ U := X \times \mathbb{R}^{n} $ 并考虑
\[
\varphi : U \to \mathbb{R}^{m+n} = \mathbb{R}^{m} \times \mathbb{R}^{n} , \quad (x, y) \mapsto (x, y - f(x)) .
\]

于是我们有 $ \varphi \in C^{q}(U, \mathbb{R}^{m} \times \mathbb{R}^{n}) $，且 $ \operatorname{im}(\varphi) = U $。此外，$ \varphi : U \to U $ 是双射，且 $ \varphi^{-1}(x, z) = (x, z + f(x)) $。因此 $ \varphi $ 是 $ U $ 到自身的 $ C^{q} $ 微分同胚，且 $ \varphi(U \cap \operatorname{graph}(f)) = X \times \{0\} = U \cap (\mathbb{R}^{m} \times \{0\}) $。 \nobreak\hspace{0.35em}■

\end{proposition}

\needspace{3\baselineskip}\subsection*{正则值定理}

下一个定理给出注 8.5(f) 的新解释。它提供了理解子流形的最自然的方式之一。

\begin{theorem}

（正则值）设 $ X $ 在 $ \mathbb{R}^{m} $ 中开且 $ c $ 是 $ f \in C^{q}(X, \mathbb{R}^{n}) $ 的正则值。则 $ f^{-1}(c) $ 是 $ \mathbb{R}^{m} $ 的 $ (m - n) $ 维 $ C^{q} $ 子流形。

\proofhead
这由注 8.5(f) 与命题 9.2 立即得到。 \nobreak\hspace{0.35em}■
\end{theorem}

\begin{corollary}

设 $ X $ 在 $ \mathbb{R}^{n} $ 中开且 $ f \in C^{q}(X, \mathbb{R}) $。如果对 $ x \in f^{-1}(c) $ 有 $ \nabla f(x) \ne 0 $，则 $ f $ 的\textbf{水平集} $ f^{-1}(c) $ 是 $ \mathbb{R}^{n} $ 的 $ C^{q} $ 超平面。

\proofhead
见注 8.5(c)。 \nobreak\hspace{0.35em}■
\end{corollary}

\begin{example}

\begin{items}
\item

函数

\[
f : \mathbb{R}^{2} \to \mathbb{R} , \quad (x, y) \mapsto x^{2} - y^{2}
\]

以 $ (0, 0) $ 为其唯一的临界点。其水平集是双曲线。

\end{items}

\bookfigure[0.20]{figures/834ee3fa2ae44d17.png}

\begin{items}
\item[（b）]

（欧几里得）$ n $ 球面 $ S^{n} := \{x \in \mathbb{R}^{n+1} ; \ |x| = 1\} $ 是 $ \mathbb{R}^{n+1} $ 中的 $ n $ 维 $ C^{\infty} $ 超曲面。

\proofhead
函数 $ f : \mathbb{R}^{n+1} \to \mathbb{R} $，$ x \mapsto |x|^{2} $ 是光滑的，且 $ S^{n} = f^{-1}(1) $。因为 $ \nabla f(x) = 2x $，我们知道 1 是 $ f $ 的正则值。于是结论由推论 9.4 得到 \nobreak\hspace{0.35em}■
\item[（c）]

\textbf{正交群}\footnote{这里 $ 1_{n} $ 表示 $ \mathbb{R}^{n \times n} $ 中的单位矩阵。关于 $ O(n) $ 的更多内容，另见习题 1 与 2。} $ O(n) := \{A \in \mathbb{R}^{n \times n} ; \ A^{\top}A = 1_{n}\} $ 是 $ \mathbb{R}^{n \times n} $ 的 $ n(n - 1)/2 $ 维 $ C^{\infty} $ 子流形。

\end{items}

证明

\begin{subitems}
\item

由习题 1.5 我们知道（因为 $ A^{\top} = A^{*} $）对每个 $ A \in \mathbb{R}^{n \times n} $，$ A^{\top}A $ 是对称的。容易验证 $ \mathbb{R}^{n \times n}_{\mathrm{sym}} $ 的维数为 $ n(n + 1)/2 $（即 $ (n \times n) $ 矩阵对角线及其上方的元素个数）。对映射

\[
f : \mathbb{R}^{n \times n} \to \mathbb{R}^{n \times n}_{\mathrm{sym}} , \quad A \mapsto A^{\top}A ,
\]

我们有 $ O(n) = f^{-1}(1_{n}) $。首先，我们观察到

\[
g : \mathbb{R}^{n \times n} \times \mathbb{R}^{n \times n} \to \mathbb{R}^{n \times n} , \quad (A, B) \mapsto A^{\top}B
\]

是双线性的，从而光滑。于是，因为 $ f(A) = g(A, A) $，映射 $ f $ 也光滑。此外，我们有（见命题 4.6）

\[
\partial f(A)B = A^{\top}B + B^{\top}A \qquad \text{对 } A, B \in \mathbb{R}^{n \times n} .
\]
\item

设 $ A \in f^{-1}(1_{n}) = O(n) $ 且 $ S \in \mathbb{R}^{n \times n}_{\mathrm{sym}} $。对 $ B := AS/2 $，我们则有

\[
\partial f(A)B = \frac{1}{2}A^{\top}AS + \frac{1}{2}SA^{\top}A = S
\]

因为 $ A^{\top}A = 1_{n} $。因此对每个 $ A \in f^{-1}(1_{n}) $，$ \partial f(A) $ 是满射。于是 $ 1_{n} $ 是 $ f $ 的正则值。因为 $ \dim(\mathbb{R}^{n \times n}) = n^{2} $ 且 $ n^{2} - n(n + 1)/2 = n(n - 1)/2 $，结论由定理 9.3 得到。 \nobreak\hspace{0.35em}■

\end{subitems}

\end{example}

\needspace{3\baselineskip}\subsection*{浸入定理}

设 $ X $ 在 $ \mathbb{R}^{m} $ 中开。如果对每个 $ x \in X $，$ \partial f(x) \in \mathcal{L}(\mathbb{R}^{m}, \mathbb{R}^{n}) $ 是单射，则称映射 $ f \in C^{1}(X, \mathbb{R}^{n}) $ 为（$ X $ 在 $ \mathbb{R}^{n} $ 中的）\textbf{浸入}。于是 $ f $ 是 $ F := f(X) $ 的\textbf{正则参数化}。最后，$ F $ 是 $ m $ 维（正则）参数化超曲面，而 $ X $ 是其\textbf{参数域}。一维或二维的参数化超曲面分别是（\textbf{正则}）\textbf{参数化曲线}或（\textbf{正则}）\textbf{参数化曲面}。

\begin{examplenote}

\begin{items}
\item

如果 $ f \in C^{1}(X, \mathbb{R}^{n}) $ 是浸入，则 $ m \le n $。

\proofhead
对 $ n < m $，不存在单射 $ A \in \mathcal{L}(\mathbb{R}^{m}, \mathbb{R}^{n}) $ 这一点由秩公式 (2.4) 立即得到。 \nobreak\hspace{0.35em}■
\item

对每个 $ \ell \in \mathbb{N}^{\times} $，$ \mathbb{R} \to \mathbb{R}^{2} $，$ t \mapsto (\cos(\ell t), \sin(\ell t)) $ 在 $ (0, 2\pi) $ 上的限制是 $ C^{\infty} $ 浸入。$ [0, 2\pi) $ 的像是单位圆 $ S^{1} $，它被遍历 $ \ell $ 次。

\item

映射 $ (-\pi, \pi) \to \mathbb{R}^{2} $，$ t \mapsto (1 + 2\cos t)(\cos t, \sin t) $ 是光滑浸入。其像的闭包称为\textbf{帕斯卡蜗线}。

\item

容易看出 $ (-\pi/4, \pi/2) \to \mathbb{R}^{2} $，$ t \mapsto \sin 2t(-\sin t, \cos t) $ 是单射的 $ C^{\infty} $ 浸入。 \nobreak\hspace{0.35em}■

\end{items}

\bookfigure[0.64]{figures/3f5fe048423cb300.png}

下一个定理表明 $ m $ 维参数化超曲面可局部地用流形表示。

\end{examplenote}

\begin{theorem}

（浸入）设 $ X $ 在 $ \mathbb{R}^{m} $ 中开且 $ f \in C^{q}(X, \mathbb{R}^{n}) $ 是浸入。则对每个 $ x_{0} \in X $，存在 $ X $ 中的开邻域 $ X_{0} $，使得 $ f(X_{0}) $ 是 $ \mathbb{R}^{n} $ 的 $ m $ 维 $ C^{q} $ 子流形。

证明

\begin{subitems}
\item

通过置换坐标，我们可以不失一般性地假设雅可比矩阵 $ [\partial f(x_{0})] $ 的前 $ m $ 行线性无关。因此

\[
\frac{\partial(f^{1}, \ldots, f^{m})}{\partial(x^{1}, \ldots, x^{m})}(x_{0}) \ne 0 .
\]
\item

我们考虑 $ \mathbb{R}^{n} $ 中的开集 $ X \times \mathbb{R}^{n-m} $ 以及映射

\[
\psi : X \times \mathbb{R}^{n-m} \to \mathbb{R}^{n} , \quad (x, y) \mapsto f(x) + (0, y) .
\]

显然，$ \psi $ 属于 $ C^{q} $ 类，且 $ \partial\psi(x_{0}, 0) $ 有矩阵

\begin{equation*}
[\partial\psi(x_{0}, 0)] = \begin{bmatrix} A & 0 \\ B & 1_{n-m} \end{bmatrix}
\end{equation*}

其中

\begin{equation*}
A := \begin{bmatrix} \partial_{1}f^{1} & \cdots & \partial_{m}f^{1} \\ \vdots & & \vdots \\ \partial_{1}f^{m} & \cdots & \partial_{m}f^{m} \end{bmatrix}(x_{0}) \quad \text{且} \quad B := \begin{bmatrix} \partial_{1}f^{m+1} & \cdots & \partial_{m}f^{m+1} \\ \vdots & & \vdots \\ \partial_{1}f^{n} & \cdots & \partial_{m}f^{n} \end{bmatrix}(x_{0}) .
\end{equation*}

因此 $ \partial\psi(x_{0}, 0) $ 属于 $ \mathcal{L}\mathrm{aut}(\mathbb{R}^{n}) $，从而

\[
\det[\partial\psi(x_{0}, 0)] = \det A = \frac{\partial(f^{1}, \ldots, f^{m})}{\partial(x^{1}, \ldots, x^{m})}(x_{0}) \ne 0 .
\]

于是反函数定理（定理 7.3）表明存在开邻域 $ V \in \mathfrak{U}_{\mathbb{R}^{n}}(x_{0}, 0) $ 与 $ U \in \mathfrak{U}_{\mathbb{R}^{n}}(\psi(x_{0}, 0)) $，满足 $ \psi|V \in \operatorname{Diff}^{q}(V, U) $。

\end{subitems}

\bookfigure[0.95]{figures/7c264da1286f5185.png}

我们现在令 $ \Phi := (\psi|V)^{-1} \in \operatorname{Diff}^{q}(U, V) $ 且 $ X_{0} := \{x \in \mathbb{R}^{m} ; \ (x, 0) \in V\} $。则 $ X_{0} $ 是 $ x_{0} $ 在 $ \mathbb{R}^{m} $ 中的开邻域，且

\[
\Phi(U \cap f(X_{0})) = \Phi(\psi(X_{0} \times \{0\})) = X_{0} \times \{0\} = V \cap (\mathbb{R}^{m} \times \{0\}) .\quad\blacksquare
\]
\end{theorem}

\begin{corollary}

设 $ I $ 是开区间且 $ \gamma \in C^{q}(I, \mathbb{R}^{n}) $。又设 $ t_{0} \in I $ 且 $ \dot{\gamma}(t_{0}) \ne 0 $。则存在 $ I $ 的含 $ t_{0} $ 的开子区间 $ I_{0} $，使得 $ \gamma|I_{0} $ 的迹是 $ \mathbb{R}^{n} $ 中的嵌入 $ C^{q} $ 曲线。

\proofhead
这立即由定理 9.7 得到。 \nobreak\hspace{0.35em}■
\end{corollary}

\begin{remark}

\begin{items}
\item

光滑道路

\[
\gamma : \mathbb{R} \to \mathbb{R}^{2} , \quad t \mapsto (t^{3}, t^{2})
\]

对 $ t \ne 0 $ 满足 $ \dot{\gamma}(t) \ne 0 $。在 $ t = 0 $ 处，$ \gamma $ 的导数消失。$ \gamma $ 的迹，即\textbf{尼尔抛物线}，在那里“不光滑”，而是缩成一点。

\end{items}

\bookfigure[0.23]{figures/9297b746d70c70df.png}

\begin{items}
\item[（b）]

设 $ f \in C^{q}(X, \mathbb{R}^{n}) $ 是浸入。则 $ f(X) $ 一般不是 $ \mathbb{R}^{n} $ 的 $ C^{q} $ 子流形，因为 $ f(X) $ 可以有“自交”；见例 9.6(c)。

\item[（c）]

例 9.6(c) 中给出的浸入不是单射。而且，单射浸入的像一般也不是子流形；见例 9.6(d) 与习题 16。另一方面，例 9.6(b) 表明非单射浸入的像确实可以是子流形。

\item[（d）]

设 $ X $ 在 $ \mathbb{R}^{m} = \mathbb{R}^{m} \times \{0\} \subset \mathbb{R}^{n} $ 中开，且 $ f \in C^{q}(X, \mathbb{R}^{n}) $ 是浸入。则对每个 $ x_{0} \in X $，存在 $ (x_{0}, 0) $ 在
$ \mathbb{R}^{n} $ 中的开邻域 $ V $、$ \mathbb{R}^{n} $ 中的开集 $ U $ 以及 $ \psi \in \operatorname{Diff}^{q}(V, U) $，使得对满足 $ (x, 0) \in V $ 的 $ x \in X $ 有 $ \psi(x, 0) = f(x) $。

\proofhead
这立即由定理 9.7 的证明得到。 \nobreak\hspace{0.35em}■
\end{items}

\end{remark}

\needspace{3\baselineskip}\subsection*{嵌入}

设 $ g : I \to \mathbb{R}^{2} $ 是例 9.6(d) 的单射 $ C^{\infty} $ 浸入。我们已经确定 $ S = \operatorname{im}(g) $ 表示 $ \mathbb{R}^{2} $ 中的嵌入曲线。这就提出了问题：单射浸入还必须具有哪些其他性质，其像才是子流形？如果分析上面的例子，容易看出 $ g^{-1} : S \to I $ 不连续。因此映射 $ g : I \to S $ 不是拓扑映射。事实上，下一个定理表明，如果单射浸入有连续逆映射，则其像是子流形。如果 $ f : X \to f(X) $ 是拓扑映射，我们称（$ C^{q} $）浸入 $ f : X \to \mathbb{R}^{n} $ 为 $ X $ 在 $ \mathbb{R}^{n} $ 中的（$ C^{q} $）\textbf{嵌入}。\footnote{自然，$ f(X) $ 带有从 $ \mathbb{R}^{n} $ 诱导的拓扑。如果上下文使 $ X $ 与 $ \mathbb{R}^{n} $ 的含义无歧义，我们就简称为嵌入。}

\begin{proposition}

设 $ X $ 在 $ \mathbb{R}^{m} $ 中开且 $ f \in C^{q}(X, \mathbb{R}^{n}) $ 是嵌入。则 $ f(X) $ 是 $ \mathbb{R}^{n} $ 的 $ m $ 维 $ C^{q} $ 子流形。

\proofhead
我们令 $ M := f(X) $ 并选取 $ y_{0} \in M $。根据定理 9.7，$ x_{0} := f^{-1}(y_{0}) $ 在 $ X $ 中有开邻域 $ X_{0} $，使得 $ M_{0} := f(X_{0}) $ 是 $ \mathbb{R}^{n} $ 的 $ m $ 维子流形。
\bookfigure[0.86]{figures/d76adad9b642af21.png}

因此存在 $ y_{0} $ 的开邻域 $ U_{1} $ 与 0 在 $ \mathbb{R}^{n} $ 中的开邻域 $ V_{1} $，以及从 $ U_{1} $ 到 $ V_{1} $ 的 $ C^{q} $ 微分同胚 $ \Phi $，使得 $ \Phi(M_{0} \cap U_{1}) = V_{1} \cap (\mathbb{R}^{m} \times \{0\}) $。因为 $ f $ 是拓扑映射，$ M_{0} $ 在 $ M $ 中开。因此存在 $ \mathbb{R}^{n} $ 中的开集 $ U_{2} $，使得 $ M_{0} = M \cap U_{2} $（见命题 III.2.26）。因此 $ U := U_{1} \cap U_{2} $ 是 $ y_{0} $ 在 $ \mathbb{R}^{n} $ 中的开邻域，且 $ V := \Phi(U) $ 是 0 在 $ \mathbb{R}^{n} $ 中的开邻域，满足 $ \Phi(M \cap U) = V \cap (\mathbb{R}^{m} \times \{0\}) $。因为这对每个 $ y_{0} \in M $ 都成立，结论得证。 \nobreak\hspace{0.35em}■

\end{proposition}

\begin{example}

\begin{items}
\item

（球坐标）
设

\[
f_{3} : \mathbb{R}^{3} \to \mathbb{R}^{3} , \quad (r, \varphi, \vartheta) \mapsto (x, y, z)
\]

由

\begin{equation}
\begin{aligned}
x &= r\cos\varphi\sin\vartheta , \\ y &= r\sin\varphi\sin\vartheta , \\ z &= r\cos\vartheta ,
\end{aligned}\tag{9.1}
\end{equation}

定义，并设 $ V_{3} := (0, \infty) \times (0, 2\pi) \times (0, \pi) $。则 $ g_{3} := f_{3}|V_{3} $ 是 $ V_{3} $ 在 $ \mathbb{R}^{3} $ 中的 $ C^{\infty} $ 嵌入，且 $ F_{3} = g_{3}(V_{3}) = \mathbb{R}^{3} \setminus H_{3} $，其中 $ H_{3} $ 表示闭半平面 $ \mathbb{R}^{+} \times \{0\} \times \mathbb{R} $。

\end{items}

\bookfigure[0.27]{figures/85da6de7078b0c25.png}

如果把 $ (r, \varphi, \vartheta) $ 限制到 $ V_{3} $ 的如下形式的子集

\[
(r_{0}, r_{1}) \times (\varphi_{0}, \varphi_{1}) \times (\vartheta_{0}, \vartheta_{1}) ,
\]

就得到“球壳”的正则参数化。

\bookfigure[0.25]{figures/fcd552a34aa7dc85.png}

考虑 $ f_{3}(\overline{V_{3}}) = \mathbb{R}^{3} $。我们还有，对 $ W_{3} := (0, \infty) \times [0, 2\pi) \times (0, \pi) $ 有 $ f_{3}(W_{3}) = \mathbb{R}^{3} \setminus (\{0\} \times \{0\} \times \mathbb{R}) $，且 $ f_{3} $ 把 $ W_{3} $ 双射地映到 $ \mathbb{R}^{3} \setminus (\{0\} \times \{0\} \times \mathbb{R}) $ 上。于是，凭借 (9.1)，每个点 $ (x, y, z) \in \mathbb{R}^{3} \setminus (\{0\} \times \{0\} \times \mathbb{R}) $ 都可以用\textbf{球坐标} $ (r, \varphi, \vartheta) \in W_{3} $ 唯一地描述。换句话说，$ f_{3}|W_{3} $ 是 $ \mathbb{R}^{3} $ 除去 $ z $ 轴的参数化。然而，$ W_{3} $ 在 $ \mathbb{R}^{3} $ 中不开。

\proofhead
显然 $ f_{3} \in C^{\infty}(\mathbb{R}^{3}, \mathbb{R}^{3}) $，且 $ g_{3} : V_{3} \to F_{3} $ 是拓扑映射（见习题 11）。我们也求出
\begin{equation*}
[\partial f_{3}(r, \varphi, \vartheta)] = \begin{bmatrix} \cos\varphi\sin\vartheta & -r\sin\varphi\sin\vartheta & r\cos\varphi\cos\vartheta \\ \sin\varphi\sin\vartheta & r\cos\varphi\sin\vartheta & r\sin\varphi\cos\vartheta \\ \cos\vartheta & 0 & -r\sin\vartheta \end{bmatrix} .
\end{equation*}

$ \partial f_{3} $ 的行列式可沿最后一行展开来计算，对 $ (r, \varphi, \vartheta) \in V_{3} $ 给出值 $ -r^{2}\sin\vartheta \ne 0 $。因此 $ g_{3} $ 是浸入，从而是嵌入。其余论断是显然的。 \nobreak\hspace{0.35em}■

\begin{items}
\item[（b）]

（球坐标\footnote{注意此处对该术语的再次使用。}）
我们定义

\[
f_{2} : \mathbb{R}^{2} \to \mathbb{R}^{3} , \quad (\varphi, \vartheta) \mapsto (x, y, z)
\]

由

\begin{equation}
\begin{aligned}
x &= \cos\varphi\sin\vartheta , \\ y &= \sin\varphi\sin\vartheta , \\ z &= \cos\vartheta
\end{aligned}\tag{9.2}
\end{equation}

并设 $ V_{2} := (0, 2\pi) \times (0, \pi) $。则限制 $ g_{2} := f_{2}|V_{2} $ 是 $ V_{2} $ 在 $ \mathbb{R}^{3} $ 中的 $ C^{\infty} $ 嵌入，且 $ F_{2} := g_{2}(V_{2}) = S^{2} \setminus H_{3} $。换句话说，$ F_{2} $ 是从 $ S^{2} $ 中去掉半平面 $ H_{3} $ 与 $ S^{2} $ 相交所得的半圆得到的。把 $ (\varphi, \vartheta) $ 限制到 $ V_{2} $ 的如下形式的子集 $ (\varphi_{0}, \varphi_{1}) \times (\vartheta_{0}, \vartheta_{1}) $，就得到 $ S^{2} $ 的一个截段的正则参数化，这意味着矩形 $ \overline{V_{2}} $ 被 $ f_{2} $ “弯成（部分）球面”。还注意到，对 $ W_{2} := [0, 2\pi) \times (0, \pi) $，我们有 $ f_{2}(W_{2}) = S^{2}\setminus\{\pm e_{3}\} $，其中 $ e_{3} $ 是\textbf{北极}，$ -e_{3} $ 是\textbf{南极}。$ f_{2} $ 也把参数域 $ W_{2} $ 双射地映到 $ S^{2}\setminus\{\pm e_{3}\} $ 上。利用 (9.2)，$ S^{2}\setminus\{\pm e_{3}\} $ 因此可以通过\textbf{球坐标} $ (\varphi, \vartheta) \in W_{2} $ 来描述，尽管 $ f_{2}|W_{2} $ 不是 $ S^{2} \setminus \{\pm e_{3}\} $ 的正则参数化，因为 $ W_{2} $ 在 $ \mathbb{R}^{2} $ 中不开。

\end{items}

\bookfigure[0.28]{figures/fdd13147712973f2.png}

\proofhead
显然 $ f_{2} = f_{3}(1, \cdot, \cdot) \in C^{\infty}(\mathbb{R}^{2}, \mathbb{R}^{3}) $ 且 $ F_{2} = S^{2} \cap F_{3} $。因为 $ F_{3} $ 在 $ \mathbb{R}^{3} $ 中开，由此可得 $ F_{2} $ 在 $ S^{2} $ 中开。此外，$ g_{2} = g_{3}(1, \cdot, \cdot) $ 把 $ V_{2} $ 双射地映到 $ F_{2} $ 上，且 $ g_{2}^{-1} = g_{3}^{-1}|F_{2} $。因此 $ g_{2}^{-1} : F_{2} \to V_{2} $ 连续。因为
\begin{equation}
[\partial f_{2}(\varphi, \vartheta)] = \begin{bmatrix} -\sin\varphi\sin\vartheta & \cos\varphi\cos\vartheta \\ \cos\varphi\sin\vartheta & \sin\varphi\cos\vartheta \\ 0 & -\sin\vartheta \end{bmatrix}\tag{9.3}
\end{equation}

由正则矩阵 $ [\partial f_{3}(1, \varphi, \vartheta)] $ 的最后两列组成，我们看到对 $ (\varphi, \vartheta) \in V_{2} $，$ \partial g_{2}(\varphi, \vartheta) $ 是单射。于是 $ g_{2} $ 是正则 $ C^{\infty} $ 参数化。 \nobreak\hspace{0.35em}■

\begin{items}
\item[（c）]

（柱坐标）定义

\[
f : \mathbb{R}^{3} \to \mathbb{R}^{3} , \quad (r, \varphi, z) \mapsto (x, y, z)
\]

由

\begin{equation}
\begin{aligned}
x &= r\cos\varphi , \\ y &= r\sin\varphi , \\ z &= z ,
\end{aligned}\tag{9.4}
\end{equation}

并设 $ V := (0, \infty) \times (0, 2\pi) \times \mathbb{R} $。则 $ g := f|V $ 是 $ V $ 在 $ \mathbb{R}^{3} $ 中的 $ C^{\infty} $ 嵌入，且 $ g(V) = F_{3} = \mathbb{R}^{3} \setminus H_{3} $。进一步把 $ g $ 限制到 $ V $ 的如下形式的子集
$ R := (r_{0}, r_{1}) \times (\varphi_{0}, \varphi_{1}) \times (z_{0}, z_{1}) $ 上，就得到“柱壳片段”的正则参数化。换句话说，$ f $ 把矩形体 $ R $ “弯成”这个形状。

\end{items}

\bookfigure[0.30]{figures/928c0ee47c7fb290.png}
\bookfigure[0.23]{figures/e9f734979b22dfaf.png}

我们还有 $ f(\overline{V}) = \mathbb{R}^{3} $，且对 $ W := (0, \infty) \times [0, 2\pi) \times \mathbb{R} $ 有 $ f(W) = \mathbb{R}^{3} \setminus (\{0\} \times \{0\} \times \mathbb{R}) =: Z $，并且 $ f|W $ 把 $ W $ 双射地映到 $ Z $ 上，即 $ \mathbb{R}^{3} $ 除去 $ z $ 轴。因此 (9.4) 通过\textbf{柱坐标}描述 $ Z $，尽管 $ f|W $ 不是 $ Z $ 的正则参数化，因为 $ W $ 不开。

\proofhead
显然 $ g $ 双射地映到 $ F_{3} $ 上，且容易看出 $ g^{-1} $ 是光滑的（见习题 12）。我们也有
\begin{equation*}
[\partial f(r, \varphi, z)] = \begin{bmatrix} \cos\varphi & -r\sin\varphi & 0 \\ \sin\varphi & r\cos\varphi & 0 \\ 0 & 0 & 1 \end{bmatrix} ,
\end{equation*}

这由 $ \det\partial f(r, \varphi, z) = r $ 得到。因此 $ f|(0, \infty) \times \mathbb{R} \times \mathbb{R} $ 是开半空间 $ (0, \infty) \times \mathbb{R} \times \mathbb{R} $ 在 $ \mathbb{R}^{3} $ 中的 $ C^{\infty} $ 浸入，且 $ g $ 是 $ V $ 在 $ \mathbb{R}^{3} $ 中的 $ C^{\infty} $ 嵌入。其余论断也是直截了当的。 \nobreak\hspace{0.35em}■

\begin{items}
\item[（d）]

（柱坐标）
我们在 (9.4) 中令 $ r = 1 $。则 $ g(1, \cdot, \cdot) $ 是 $ (0, 2\pi) \times \mathbb{R} $ 在 $ \mathbb{R}^{3} $ 中的 $ C^{\infty} $ 嵌入，对它有 (c) 的显然类比。

\end{items}

\bookfigure[0.20]{figures/79cfba3d09b20632.png}

\begin{items}
\item[（e）]

（旋转曲面）设 $ J $ 是 $ \mathbb{R} $ 中的开区间，$ \rho, \sigma \in C^{q}(J, \mathbb{R}) $ 满足对 $ t \in J $ 有 $ \rho(t) > 0 $ 且 $ (\dot{\rho}(t), \dot{\sigma}(t)) \ne (0, 0) $。则

\[
r : J \times \mathbb{R} \to \mathbb{R}^{3} , \quad (t, \varphi) \mapsto (\rho(t)\cos\varphi, \rho(t)\sin\varphi, \sigma(t))
\]

是 $ J \times \mathbb{R} $ 在 $ \mathbb{R}^{3} $ 中的 $ C^{q} $ 浸入。
映射

\[
\gamma : J \to \mathbb{R}^{3} , \quad t \mapsto (\rho(t), 0, \sigma(t))
\]

是 $ J $ 在 $ \mathbb{R}^{3} $ 中的 $ C^{q} $ 浸入。$ \gamma $ 的像 $ \Gamma $ 是位于 $ x $-$ z $ 平面中的正则参数化曲线，且 $ R := r(J \times \mathbb{R}) $ 由 $ \Gamma $ 绕 $ z $ 轴旋转生成。于是 $ R $ 称为\textbf{旋转曲面}，而 $ \Gamma $ 是 $ R $ 的\textbf{子午线}。如果 $ \gamma $ 不是单射，则 $ R $ 含有平行于 $ x $-$ y $ 轴的平面中的圆，在这些圆处 $ R $ 自交。

\end{items}

\bookfigure[0.25]{figures/0577e7909d086e0b.png}
\bookfigure[0.75]{figures/9b28de638277ee4a.png}

设 $ I $ 是 $ J $ 的开子区间，使得 $ \gamma|I $ 是嵌入，且 $ K $ 是 $ [0, 2\pi] $ 的开子区间。则 $ r|I \times K $ 也是嵌入。

\proofhead
$ r $ 的雅可比矩阵是
\begin{equation}
\begin{bmatrix} \dot{\rho}(t)\cos\varphi & -\rho(t)\sin\varphi \\ \dot{\rho}(t)\sin\varphi & \rho(t)\cos\varphi \\ \dot{\sigma}(t) & 0 \end{bmatrix} .\tag{9.5}
\end{equation}

划去最后一行所得的 $ (2 \times 2) $ 矩阵的行列式为 $ \rho(t)\dot{\rho}(t) $。因此，如果 $ \dot{\rho}(t) \ne 0 $，矩阵 (9.5) 的秩为 2。如果 $ \dot{\rho}(t) = 0 $，则 $ \dot{\sigma}(t) \ne 0 $，且由 (9.5) 划去第一行或第二行所得的两个行列式中至少有一个不等于 0。这证明了 $ r $ 是浸入。我们把其余论断留给你。 \nobreak\hspace{0.35em}■

\begin{items}
\item[（f）]

（环面）设 $ 0 < r < a $。方程 $ (x - a)^{2} + z^{2} = r^{2} $ 定义 $ x $-$ z $ 平面中的一个圆，并且，通过把它绕 $ z $ 轴旋转，我们得到一个 (2-)环面 $ \mathrm{T}^{2}_{a,r} $。

\end{items}

\bookfigure[0.74]{figures/4b541a7d955090fc.png}

对由

\[
\tau_{2}(t, \varphi) := ((a + r\cos t)\cos\varphi, (a + r\cos t)\sin\varphi, r\sin t) ,
\]

定义的 $ \tau_{2} \in C^{\infty}(\mathbb{R}^{2}, \mathbb{R}^{3}) $，我们求出 $ \tau_{2}([0, 2\pi]^{2}) = \mathrm{T}^{2}_{a,r} $，且 $ \tau_{2}|(0, 2\pi)^{2} $ 是嵌入。开正方形 $ (0, 2\pi)^{2} $ 在 $ \tau_{2} $ 下的像是从 2-环面上“剪去”两个圆所得的曲面：第一个是 $ x $-$ y $ 平面中的 $ x^{2} + y^{2} = (r + a)^{2} $，第二个是 $ x $-$ z $ 平面中的 $ (x - a)^{2} + z^{2} = r^{2} $。

$ \tau_{2} $ 的像也可以描述如下：正方形 $ [0, 2\pi]^{2} $ 先通过“粘合”其两条对边做成一个管子。然后，在把这个管子扭转成一个环之后，其两个圆形端面也被粘合。

\bookfigure[0.33]{figures/943aab07d55bc572.png}

\proofhead
映射 $ t \mapsto (a + r\cos t, 0, r\sin t) $ 是 $ \mathbb{R} $ 在 $ \mathbb{R}^{3} $ 中的 $ C^{\infty} $ 浸入，其像是由 $ (x - a)^{2} + z^{2} = r^{2} $ 描述的圆。因此结论由 (e) 取 $ \rho(t) := a + r\cos t $ 与 $ \sigma(t) := r\sin t $ 容易得到。 \nobreak\hspace{0.35em}■

命题 9.10 的如下逆命题表明每个子流形局部地都是某个嵌入的像。
\end{example}

\begin{theorem}

设 $ M $ 是 $ \mathbb{R}^{n} $ 的 $ m $ 维 $ C^{q} $ 子流形。则 $ M $ 的每个点 $ p $ 在 $ M $ 中有邻域 $ U $，使得 $ U $ 是 $ \mathbb{R}^{m} $ 中开集在 $ C^{q} $ 嵌入下的像。

\proofhead
对每个 $ p \in M $，存在 $ p $ 与 0 在 $ \mathbb{R}^{n} $ 中的开邻域 $ \tilde{U} $ 与 $ \tilde{V} $，以及 $ C^{q} $ 微分同胚 $ \Phi : \tilde{U} \to \tilde{V} $，使得 $ \Phi(M \cap \tilde{U}) = \tilde{V} \cap (\mathbb{R}^{m} \times \{0\}) $。我们令 $ U := M \cap \tilde{U} $，$ V := \{x \in \mathbb{R}^{m} ; \ (x, 0) \in \tilde{V}\} $，以及
\[
g : V \to \mathbb{R}^{n} , \quad x \mapsto \Phi^{-1}((x, 0)) .
\]

则 $ U $ 在 $ M $ 中开，$ V $ 在 $ \mathbb{R}^{m} $ 中开；且 $ g $ 属于 $ C^{q}(V, \mathbb{R}^{n}) $。此外，$ g $ 把集合 $ V $ 双射地映到 $ U $ 上，且对 $ x \in V $ 有 $ \operatorname{rank}\partial g(x) = m $，因为 $ [\partial g(x)] $ 由正则矩阵 $ [\partial\Phi^{-1}((x, 0))] $ 的前 $ m $ 列组成。因为 $ g^{-1} = \Phi|U $，我们清楚地看到，把 $ g $ 解释为 $ \mathbb{R}^{n} $ 的拓扑子空间 $ U $ 中的映射，它是从 $ V $ 到 $ U $ 的拓扑映射。 \nobreak\hspace{0.35em}■

\end{theorem}

\needspace{3\baselineskip}\subsection*{局部坐标卡与参数化}

我们现在利用已建立的内容，用 $ M $ 与 $ \mathbb{R}^{m} $ 的开子集之间的映射来局部地描述 $ \mathbb{R}^{n} $ 的 $ m $ 维子流形 $ M $。在第 III 卷中，我们将学会——借助局部坐标卡——如何描述不（先验地）嵌入欧几里得空间的“抽象”流形。这种描述在很大程度上与“环绕空间”无关。

设 $ M $ 是 $ \mathbb{R}^{n} $ 的子集且 $ p \in M $。我们用

\[
i_{M} : M \to \mathbb{R}^{n} , \quad x \mapsto x
\]

表示 $ M $ 到 $ \mathbb{R}^{n} $ 的典范单射。映射 $ \varphi $ 称为 $ M $ 在 $ p $ 附近的 $ m $ 维（局部）$ C^{q} $ 坐标卡，如果

• $ U := \operatorname{dom}(\varphi) $ 是 $ p $ 在 $ M $ 中的开邻域；

• $ \varphi $ 是 $ U $ 到 $ \mathbb{R}^{m} $ 的开集 $ V := \varphi(U) $ 上的同胚；

• $ g := i_{M} \circ \varphi^{-1} $ 是 $ C^{q} $ 浸入。

集合 $ U $ 是 $ \varphi $ 的\textbf{坐标卡覆盖的区域}，$ V $ 是\textbf{参数值域}，而 $ g $ 是 $ U $ 在 $ \varphi $ 中的\textbf{参数化}。偶尔，我们用 $ (\varphi, U) $ 记 $ \varphi $，用 $ (g, V) $ 记 $ g $。$ M $ 的 $ m $ 维 $ C^{q} $ 图册是 $ M $ 的 $ m $ 维 $ C^{q} $ 坐标卡族 $ \{\varphi_{\alpha} ; \ \alpha \in A\} $，其坐标卡覆盖的区域 $ U_{\alpha} := \operatorname{dom}(\varphi_{\alpha}) $ 覆盖集合 $ M $，即 $ M = \bigcup_{\alpha} U_{\alpha} $。最后，$ (x^{1}, \ldots, x^{m}) := \varphi(p) $ 是 $ p \in U $ 在坐标卡 $ \varphi $ 中的\textbf{局部坐标}。

下一个定理表明流形可以通过坐标卡来描述。这特别意味着 $ \mathbb{R}^{n} $ 的 $ m $ 维子流形局部地“看起来像 $ \mathbb{R}^{m} $”，即它局部同胚于 $ \mathbb{R}^{m} $ 的开子集。

\bookfigure[0.83]{figures/8b17973a60e469c8.png}

\begin{theorem}

设 $ M $ 是 $ \mathbb{R}^{n} $ 的 $ m $ 维 $ C^{q} $ 子流形。则对每个 $ p \in M $，存在 $ M $ 在 $ p $ 附近的 $ m $ 维 $ C^{q} $ 坐标卡。因此 $ M $ 有 $ m $ 维 $ C^{q} $ 图册。如果 $ M $ 紧致，则 $ M $ 有有限图册。

\proofhead
定理 9.12 保证每个 $ p \in M $ 在其附近有 $ m $ 维 $ C^{q} $ 坐标卡 $ \varphi_{p} $。因此 $ \{\varphi_{p} ; \ p \in M\} $ 是图册。因为这个图册的坐标卡覆盖的区域形成 $ M $ 的开覆盖，最后的论断由紧致性的定义得到。 \nobreak\hspace{0.35em}■
\end{theorem}

\begin{example}

\begin{items}
\item

设 $ X $ 在 $ \mathbb{R}^{n} $ 中开。则 $ X $ 有只含一个坐标卡的 $ C^{\infty} $ 图册，即\textbf{平凡坐标卡} $ \operatorname{id}_{X} $。

\item

设 $ X $ 在 $ \mathbb{R}^{m} $ 中开且 $ f \in C^{q}(X, \mathbb{R}^{n}) $。则根据命题 9.2，$ \operatorname{graph}(f) $ 是 $ \mathbb{R}^{m+n} $ 的 $ m $ 维 $ C^{q} $ 子流形，且它有只由坐标卡 $ \varphi $ 组成的图册，其中对 $ x \in X $ 有 $ \varphi((x, f(x))) = x $。

\item

球面 $ S^{2} $ 有恰含两个坐标卡的 $ C^{\infty} $ 图册。

\proofhead
在例 9.11(b) 的记号下，令 $ U_{2} := g_{2}(V_{2}) $ 与 $ \varphi_{2} : U_{2} \to V_{2} $，并令 $ (x, y, z) \mapsto (\varphi, \vartheta) $ 为 $ g_{2} $ 的逆映射。则 $ \varphi_{2} $ 是 $ S^{2} $ 的 $ C^{\infty} $ 坐标卡。
现在我们通过 $ \tilde{V}_{2} := (\pi, 3\pi) \times (0, \pi) $ 与
\[
\tilde{g}_{2}(\varphi, \vartheta) := (\cos\varphi\sin\vartheta, \cos\vartheta, \sin\varphi\sin\vartheta) .
\]

定义 $ \tilde{g}_{2} : \tilde{V}_{2} \to S^{2} $。则 $ \tilde{g}_{2}(\tilde{V}_{2}) $ 是从 $ S^{2} $ 中去掉它与坐标半平面 $ (-\mathbb{R}^{+}) \times \mathbb{R} \times \{0\} $ 相交所得的半圆得到的。显然 $ U_{2} \cup \tilde{U}_{2} = S^{2} $，其中 $ \tilde{U}_{2} := \tilde{g}_{2}(\tilde{V}_{2}) $，且例 9.11(b) 的证明表明 $ \tilde{\varphi}_{2} := g_{2}^{-1} $ 是 $ S^{2} $ 的 $ C^{\infty} $ 坐标卡。 \nobreak\hspace{0.35em}■

\item

对每个 $ n \in \mathbb{N} $，$ n $ 球面 $ S^{n} $ 有恰含两个光滑坐标卡的图册。对 $ n \ge 1 $，这些坐标卡是\textbf{球极投影} $ \varphi_{\pm} $，其中 $ \varphi_{+} $ [$ \varphi_{-} $] 把每个 $ p \in S^{n}\setminus\{e_{n+1}\} $ [$ p \in S^{n}\setminus\{-e_{n+1}\} $] 映到“穿刺点”，即连接\textbf{北极} $ e_{n+1} $ [\textbf{南极} $ -e_{n+1} $] 与 $ p $ 的直线与\textbf{赤道超平面} $ \mathbb{R}^{n} \times \{0\} $ 相交的点。

\end{items}

\bookfigure[0.47]{figures/22954b5a5c9811e1.png}

\proofhead
由例 9.5(b) 我们知道 $ S^{n} $ 是 $ \mathbb{R}^{n} $ 的 $ n $ 维 $ C^{\infty} $ 子流形。当 $ n = 0 $ 时，$ S^{n} $ 只由 $ \mathbb{R} $ 中的两个点 $ \pm 1 $ 组成，命题是平凡的。

因此设 $ n \in \mathbb{N}^{\times} $。经过 $ p $ 与 $ \pm e_{n+1} \in S^{n} $ 的直线 $ t \mapsto tp \pm (1 - t)e_{n+1} $ 在 $ t = 1/(1 \mp p^{n+1}) $ 时与超平面 $ x^{n+1} = 0 $ 相交。因此穿刺点的坐标为 $ x = p'/(1 \mp p^{n+1}) \in \mathbb{R}^{n} $，其中 $ p = (p', p^{n+1}) \in \mathbb{R}^{n} \times \mathbb{R} $。这些定义了映射 $ \varphi_{\pm} : S^{n}\setminus\{\pm e_{n+1}\} \to \mathbb{R}^{n} $，$ p \mapsto x $，它们显然连续。

为了计算其逆映射，我们考虑直线

\[
t \mapsto t(x, 0) \pm (1 - t)e_{n+1}
\]

经过 $ (x, 0) \in \mathbb{R}^{n} \times \mathbb{R} $ 与 $ \pm e_{n+1} \in S^{n} $。它们当 $ t > 0 $ 且 $ t^{2}|x|^{2} + (1 - t)^{2} = 1 $ 时与 $ S^{n}\setminus\{\pm e_{n+1}\} $ 相交，因此当 $ t = 2/(1 + |x|^{2}) $ 时相交。由此我们得到

\[
\varphi_{\pm}^{-1}(x) = \frac{(2x, \pm(|x|^{2} - 1))}{1 + |x|^{2}} \qquad \text{对 } x \in \mathbb{R}^{n} .
\]

于是 $ g_{\pm} := i_{S^{n}} \circ \varphi_{\pm}^{-1} $ 属于 $ C^{\infty}(\mathbb{R}^{n}, \mathbb{R}^{n+1}) $。对 $ x \in \mathbb{R}^{n} $ 有 $ \operatorname{rank}\partial g_{\pm}(x) = n $ 这一点容易验证。 \nobreak\hspace{0.35em}■

\begin{items}
\item[（e）]

环面 $ \mathrm{T}^{2}_{a,r} $ 是 $ \mathbb{R}^{3} $ 中的 $ C^{\infty} $ 超曲面，且它有含三个坐标卡的图册。

\proofhead
设 $ X := \mathbb{R}^{3}\setminus\{0\} \times \{0\} \times \mathbb{R} $ 且
\[
f : X \to \mathbb{R} , \quad (x, y, z) \mapsto \left(\sqrt{x^{2} + y^{2}} - a\right)^{2} + z^{2} - r^{2} .
\]

则 $ f \in C^{\infty}(X, \mathbb{R}) $，且 0 是 $ f $ 的正则值。容易验证 $ f^{-1}(0) = \mathrm{T}^{2}_{a,r} $。因此，根据定理 9.3，$ \mathrm{T}^{2}_{a,r} $ 是 $ \mathbb{R}^{3} $ 中的光滑曲面。把
$ \tau_{2}|(0, 2\pi)^{2} $ 求逆的映射 $ \varphi $ 是 $ \mathrm{T}^{2}_{a,r} $ 的二维 $ C^{\infty} $ 坐标卡。第二个这样的坐标卡 $ \tilde{\varphi} $ 作为 $ \tau_{2}|(\pi, 3\pi)^{2} $ 的逆映射。最后，我们定义第三个坐标卡 $ \hat{\varphi} $ 为 $ \tau_{2}|(\pi/2, 5\pi/2)^{2} $ 的逆映射，于是 $ \{\varphi, \tilde{\varphi}, \hat{\varphi}\} $ 是 $ \mathrm{T}^{2}_{a,r} $ 的图册。 \nobreak\hspace{0.35em}■

\end{items}

\end{example}

\needspace{3\baselineskip}\subsection*{坐标卡的更换}

曲线或曲面（一般地，$ \mathbb{R}^{n} $ 的子流形）的局部几何意义与它通过局部坐标卡的描述无关。对于具体计算，必须使用特定坐标卡。当我们在流形上移动时，计算可能把我们带到所用坐标卡的边界，迫使“更换坐标卡”。因此，我们应该理解这是如何进行的。除了这一实践上的理由之外，理解坐标卡的更换将帮助我们理解流形的局部描述是如何“拼接”成整体描述的。

设 $ \{(\varphi_{\alpha}, U_{\alpha}) ; \ \alpha \in A\} $ 是 $ M \subset \mathbb{R}^{n} $ 的 $ m $ 维 $ C^{q} $ 图册。我们称映射\footnote{此处及邻近各处，当我们使用过渡函数 $ \varphi_{\beta} \circ \varphi_{\alpha}^{-1} $ 时，总假设 $ U_{\alpha} \cap U_{\beta} \ne \varnothing $。}

\[
\varphi_{\beta} \circ \varphi_{\alpha}^{-1} : \varphi_{\alpha}(U_{\alpha} \cap U_{\beta}) \to \varphi_{\beta}(U_{\alpha} \cap U_{\beta}) \quad \text{对 } \alpha, \beta \in A ,
\]

为\textbf{过渡函数}。它们规定图册 $ \{\varphi_{\alpha} ; \ \alpha \in A\} $ 中的坐标卡如何“缝合在一起”。

\bookfigure[0.80]{figures/3af02cb400dadd54.png}

\qyxffnnumdef{1}
\begin{proposition}

如果 $ (\varphi_{\alpha}, U_{\alpha}) $ 与 $ (\varphi_{\beta}, U_{\beta}) $ 是 $ m $ 维 $ C^{q} $ 流形的 $ m $ 维 $ C^{q} $ 坐标卡，则

\[
\varphi_{\beta} \circ \varphi_{\alpha}^{-1} \in \operatorname{Diff}^{q}\left(\varphi_{\alpha}(U_{\alpha} \cap U_{\beta}), \varphi_{\beta}(U_{\alpha} \cap U_{\beta})\right) ,
\]

其中 $ (\varphi_{\beta} \circ \varphi_{\alpha}^{-1})^{-1} = \varphi_{\alpha} \circ \varphi_{\beta}^{-1} $。

证明

\begin{subitems}
\item

显然 $ \varphi_{\beta} \circ \varphi_{\alpha}^{-1} $ 是双射，其逆映射是 $ \varphi_{\alpha} \circ \varphi_{\beta}^{-1} $。因此只需验证 $ \varphi_{\beta} \circ \varphi_{\alpha}^{-1} $ 属于 $ C^{q}(\varphi_{\alpha}(U_{\alpha} \cap U_{\beta}), \mathbb{R}^{m}) $。

\item

我们令 $ V_{\gamma} := \varphi_{\gamma}(U_{\gamma}) $，且 $ g_{\gamma} $ 是 $ \gamma \in \{\alpha, \beta\} $ 时属于 $ (\varphi_{\gamma}, U_{\gamma}) $ 的参数化。又设 $ x_{\gamma} \in \varphi_{\gamma}(U_{\alpha} \cap U_{\beta}) $ 满足 $ g_{\alpha}(x_{\alpha}) = g_{\beta}(x_{\beta}) =: p $。因为 $ g_{\gamma} $ 是单射 $ C^{q} $ 浸入，注 9.9(d) 给出 $ p $ 的开邻域
$ \tilde{U}_{\gamma} $ 与 $ (x_{\gamma}, 0) $ 在 $ \mathbb{R}^{n} $ 中的开邻域 $ \tilde{V}_{\gamma} $，以及 $ \psi_{\gamma} \in \operatorname{Diff}^{q}(\tilde{V}_{\gamma}, \tilde{U}_{\gamma}) $，使得对所有满足 $ (y, 0) \in \tilde{V}_{\gamma} $ 的 $ y \in V_{\gamma} $ 有 $ \psi_{\gamma}(y, 0) = g_{\gamma}(y) $。我们现在令

\[
V := \{x \in V_{\alpha} ; \ (x, 0) \in \tilde{V}_{\alpha}\} \cap \varphi_{\alpha}(U_{\alpha} \cap U_{\beta}) .
\]

显然 $ V $ 是 $ x_{\alpha} $ 在 $ \mathbb{R}^{m} $ 中的开邻域。利用 $ g_{\alpha} $ 的连续性，我们可以（必要时通过缩小 $ V $）假设 $ g_{\alpha}(V) $ 含于 $ \tilde{U}_{\beta} $。因此

\[
\varphi_{\beta} \circ \varphi_{\alpha}^{-1}(x) = \psi_{\beta}^{-1} \circ g_{\alpha}(x) \qquad \text{对 } x \in V .
\]

因为 $ g_{\alpha} \in C^{q}(V_{\alpha}, \mathbb{R}^{n}) $，链式法则表明 $ \psi_{\beta}^{-1} \circ g_{\alpha} $ 属于 $ C^{q}(V, \mathbb{R}^{n}) $。\qyxffnmark{1} 因此由此可得 $ \varphi_{\beta} \circ \varphi_{\alpha}^{-1} \in C^{q}(V, \mathbb{R}^{m}) $，因为 $ \psi_{\beta}^{-1} \circ g_{\alpha} $ 的像与 $ \varphi_{\beta} \circ \varphi_{\alpha}^{-1} $ 的像重合，从而位于 $ \mathbb{R}^{m} $ 中。因为这对每个 $ x_{\alpha} \in \varphi_{\alpha}(U_{\alpha} \cap U_{\beta}) $ 都成立，且属于 $ C^{q} $ 类是局部性质，命题得证。 \nobreak\hspace{0.35em}■

\end{subitems}

\end{proposition}
\qyxffntextof{1}{此处及其他各处，我们常对映射及其到定义域子集上的限制使用同一个记号。}

\begin{remark}

\begin{items}
\item

$ \mathbb{R}^{n} $ 的子流形的维数是唯一的。因此对给定的 $ m $ 维流形，简单地谈论其“坐标卡”而不是“$ m $ 维坐标卡”是有意义的。

\proofhead
设 $ M $ 是 $ \mathbb{R}^{n} $ 的 $ m $ 维 $ C^{q} $ 子流形且 $ p \in M $。则根据定理 9.13，存在 $ p $ 附近的 $ m $ 维 $ C^{q} $ 坐标卡 $ (\varphi, U) $。设 $ (\psi, V) $ 是 $ p $ 附近的 $ m' $ 维 $ C^{q} $ 坐标卡。则命题 9.15 的证明表明
\[
\psi \circ \varphi^{-1} \in \operatorname{Diff}^{q}\left(\varphi(U \cap V), \psi(U \cap V)\right) ,
\]

其中 $ \varphi(U \cap V) $ 在 $ \mathbb{R}^{m} $ 中开，$ \psi(U \cap V) $ 在 $ \mathbb{R}^{m'} $ 中开。这蕴含着 $ m = m' $（见 (7.2)），正如所求。 \nobreak\hspace{0.35em}■

\item

通过坐标卡 $ (\varphi_{1}, U_{1}) $，坐标卡覆盖的区域 $ U_{1} $ 可以用局部坐标 $ (x^{1}, \ldots, x^{m}) = \varphi_{1}(q) \in \mathbb{R}^{m} $（对 $ q \in U $）来描述。如果 $ (\varphi_{2}, U_{2}) $ 是第二个坐标卡，则 $ U_{2} $ 有其自己的局部坐标 $ (y^{1}, \ldots, y^{m}) = \varphi_{2}(q) \in \mathbb{R}^{m} $。于是 $ U_{1} \cap U_{2} $ 在两个坐标系 $ (x^{1}, \ldots, x^{m}) $ 与 $ (y^{1}, \ldots, y^{m}) $ 中有描述。过渡函数 $ \varphi_{2} \circ \varphi_{1}^{-1} $ 正是把坐标 $ x $ 转换成坐标 $ y $ 的\textbf{坐标变换} $ (x^{1}, \ldots, x^{m}) \mapsto (y^{1}, \ldots, y^{m}) $。\footnote{坐标变换将在第 III 卷中详细讨论。} \nobreak\hspace{0.35em}■

\end{items}

\end{remark}

\exerciseshead

\begin{exercises}
\item

设 $ H $ 是有限维实希尔伯特空间，且设 $ \varphi : H \to H $ 是满足 $ \varphi(0) = 0 $ 的等距。验证：

\begin{items}
\item

$ \varphi $ 是线性的；

\item

$ \varphi^{*}\varphi = \mathrm{id}_{H} $；

\item

$ \varphi \in \mathcal{L}\mathrm{aut}(H) $，且 $ \varphi^{-1} $ 是等距。

\end{items}

\item

\begin{items}
\item

对 $ A \in \mathbb{R}^{n\times n} $，证明以下陈述等价：

\begin{subitems}
\item

$ A $ 是等距；

\item

$ A \in O(n) $；

\item

$ \left|\det A\right| = 1 $；

\item

列向量 $ a^{\bullet}_{k} = (a^{1}_{k}, \ldots, a^{n}_{k}) $（$ k = 1, \ldots, n $）构成 $ \mathbb{R}^{n} $ 的一个标准正交基；

\item

行向量 $ a^{j}_{\bullet} = (a^{j}_{1}, \ldots, a^{j}_{n}) $（$ k = 1, \ldots, n $）构成 $ \mathbb{R}^{n} $ 的一个 标准正交基。

\end{subitems}

\item

\proofhead
$ O(n) $ 是矩阵群。
\end{items}

\item

证明例 9.1(c) 的陈述。

\item

设 $ M $ 与 $ N $ 分别是 $ \mathbb{R}^{k} $ 与 $ \mathbb{R}^{\ell} $ 的 $ m $ 维与 $ n $ 维 $ C^{q} $ 子流形。证明 $ M \times N $ 是 $ \mathbb{R}^{k+\ell} $ 的一个 $ (m + n) $ 维 $ C^{q} $ 子流形。

\item

判断 $ \mathcal{L}\mathrm{aut}(\mathbb{R}^{n}) $ 是否为 $ \mathcal{L}(\mathbb{R}^{n}) $ 的子流形。

\item

通过下式定义 $ f : \mathbb{R}^{3} \to \mathbb{R}^{2} $：

\[
f(x, y, z) := (x^{2} + xy - y - z, 2x^{2} + 3xy - 2y - 3z) .
\]

证明 $ f^{-1}(0) $ 是 $ \mathbb{R}^{3} $ 中的嵌入曲线。

\item

设 $ f, g : \mathbb{R}^{4} \to \mathbb{R}^{3} $ 由下式给出：

\[
f(x, y, z, u) := (xz - y^{2}, yu - z^{2}, xu - yz) ,
\]
\[
g(x, y, z, u) := (2(xz + yu), 2(xu - yz), z^{2} + u^{2} - x^{2} - y^{2}) .
\]

验证：

\begin{items}
\item

$ f^{-1}(0)\setminus\{0\} $ 是 $ \mathbb{R}^{4} $ 中的嵌入曲面；

\item

对每个 $ a \in \mathbb{R}^{3}\setminus\{0\} $，$ g^{-1}(a) $ 是 $ \mathbb{R}^{4} $ 中的嵌入曲线。

\end{items}

\item

下列集合中哪些

\begin{items}
\item

$ K := \{ (x, y) \in \mathbb{R}^{n} \times \mathbb{R} \;;\; \left|x\right|^{2} = y^{2} \} $，

\item

$ \{ (x, y) \in K \;;\; y > 0 \} $，

\item

$ K \setminus \{(0, 0)\} $
是 $ \mathbb{R}^{n+1} $ 的子流形？

\end{items}

\item

对 $ A \in \mathbb{R}^{n\times n}_{\mathrm{sym}} $，证明 $ \{ x \in \mathbb{R}^{n} \;;\; (x|Ax) = 1 \} $ 是 $ \mathbb{R}^{n} $ 的光滑超曲面。对 $ n = 2 $ 的情形画出曲线，对 $ n = 3 $ 的情形画出曲面。（提示：习题 4.5。）

\item

对\textbf{特殊正交群} $ SO(n) := \{ A \in O(n) \;;\; \det A = 1 \} $，证明：

\begin{items}
\item

$ SO(n) $ 是 $ O(n) $ 的子群；

\item

$ SO(n) $ 是 $ \mathbb{R}^{n\times n} $ 的光滑子流形。它的维数是多少？

\end{items}

\item

通过下式递归地定义 $ f_{n} \in C^{\infty}(\mathbb{R}^{n}, \mathbb{R}^{n}) $：

\[
f_{2}(y) := (y^{1}\cos y^{2}, y^{1}\sin y^{2}) \qquad \text{对 } y = (y^{1}, y^{2}) \in \mathbb{R}^{2} ,
\]

且

\[
f_{n+1}(y) := (f_{n}(y') \sin y^{n+1}, y^{1}\cos y^{n+1}) \qquad \text{对 } y = (y', y^{n+1}) \in \mathbb{R}^{n} \times \mathbb{R} .
\]

又设 $ V_{2} := (0, \infty) \times (0, 2\pi) $，且对 $ n \ge 3 $ 有 $ V_{n} := V_{2} \times (0, \pi)^{n-2} $。

\begin{items}
\item

显式求出 $ f_{n} $。

\item

证明：

\end{items}

\begin{subitems}
\item

$ \left|f_{n}(1, y^{2}, \ldots, y^{n})\right| = 1 $ 且 $ \left|f_{n}(y)\right| = \left|y^{1}\right| $；

\item

$ f_{n}(V_{n}) = \mathbb{R}^{n}\setminus H_{n} $，其中 $ H_{n} := \mathbb{R}^{+} \times \{0\} \times \mathbb{R}^{n-2} $；

\item

$ f_{n}(\overline{V_{n}}) = \mathbb{R}^{n} $；

\item

$ g_{n} := f_{n}|V_{n} : V_{n} \to \mathbb{R}^{n}\setminus H_{n} $ 是拓扑的；

\item

对 $ y \in V_{n} $ 且 $ n \ge 3 $，有 $ \det\partial f_{n}(y) = (-1)^{n}(y^{1})^{n-1}\sin^{n-2}(y^{n}) \cdots \sin(y^{3}) $。
因此 $ g_{n} $ 是 $ V_{n} $ 到 $ \mathbb{R}^{n} $ 的 $ C^{\infty} $ 嵌入。由 $ f_{n} $ 诱导的坐标称为 \textbf{$ n $ 维极坐标（或球坐标）}（参见例 9.11(a)）。

\end{subitems}

\item

用 $ f : \mathbb{R}^{3} \to \mathbb{R}^{3} $ 表示 $ \mathbb{R}^{3} $ 中的柱坐标（参见例 9.11(c)）。又设 $ V := (0, \infty) \times (0, 2\pi) \times \mathbb{R} $ 且 $ g := f|V $。证明：

\begin{items}
\item

$ g $ 是 $ V $ 到 $ \mathbb{R}^{3} $ 的 $ C^{\infty} $ 嵌入；

\item

$ g(1, \cdot, \cdot) $ 是 $ (0, 2\pi) \times \mathbb{R} $ 到 $ \mathbb{R}^{3} $ 的 $ C^{\infty} $ 嵌入；

\item

$ f(V) = \mathbb{R}^{3}\setminus H_{3} $（参见习题 11）；

\item

$ f(\overline{V}) = \mathbb{R}^{3} $。

\end{items}

\item

\begin{items}
\item

证明\textbf{椭圆柱面}\footnote{$ M_{a,a} $ 是通常的柱面。}

\[
M_{a,b} := \{ (x, y, z) \in \mathbb{R}^{3} \;;\; x^{2}/a^{2} + y^{2}/b^{2} = 1 \} \qquad \text{对 } a, b \in (0, \infty) ,
\]

是 $ \mathbb{R}^{3} $ 中的 $ C^{\infty} $ 超曲面。

\item

设 $ W := (0, 2\pi) \times \mathbb{R} $，且

\[
f_{1} : W \to \mathbb{R}^{3} \quad \text{对 } (\varphi, z) \mapsto (a\cos\varphi, b\sin\varphi, z) ,
\]
\[
f_{2} : W \to \mathbb{R}^{3} \quad \text{对 } (\varphi, z) \mapsto (-a\sin\varphi, b\cos\varphi, z) .
\]

又设对 $ j = 1, 2 $ 有 $ U_{j} := f_{j}|W $ 与 $ \varphi_{j} := (f_{j}|U_{j})^{-1} $。证明 $ \{(\varphi_{1}, \varphi_{2})\} $ 是 $ M $ 的一个图册，并计算过渡函数 $ \varphi_{1} \circ \varphi_{2}^{-1} $。

\end{items}

\item

设 $ M $ 是 $ \mathbb{R}^{n} $ 的非空紧致 $ m $ 维 $ C^{1} $ 子流形，且 $ m \ge 1 $。证明 $ M $ 不存在只含一个坐标卡的图册。

\item

证明由例 9.11(e) 中最后一个映射给出的曲面不是 $ \mathbb{R}^{3} $ 的子流形。

\item

对 $ g : (-\pi/4, \pi/2) \to \mathbb{R}^{2} $，$ t \mapsto \sin(2t)(-\sin t, \cos t) $，验证：

\begin{items}
\item

$ g $ 是单射的 $ C^{\infty} $ 浸入；

\item

$ \operatorname{im}(g) $ 不是 $ \mathbb{R}^{2} $ 中的嵌入曲线。

\end{items}

\item

对图册 $ \{\varphi_{-}, \varphi_{+}\} $ 计算过渡函数 $ \varphi_{-} \circ \varphi_{+}^{-1} $，其中 $ \varphi_{\pm} $ 是 $ S^{n} $ 的球极投影。

\item

设 $ M $ 与 $ N $ 分别是 $ \mathbb{R}^{m} $ 与 $ \mathbb{R}^{n} $ 的子流形，且设 $ \{ (\varphi_{\alpha}, U_{\alpha}) \;;\; \alpha \in \mathbf{A} \} $ 与 $ \{ (\psi_{\beta}, V_{\beta}) \;;\; \beta \in \mathbf{B} \} $ 分别是 $ M $ 与 $ N $ 的图册。验证：

\[
\{ (\varphi_{\alpha} \times \psi_{\beta}, U_{\alpha} \times V_{\beta}) \;;\; (\alpha, \beta) \in \mathbf{A} \times \mathbf{B} \} ,
\]

其中 $ \varphi_{\alpha} \times \psi_{\beta}(p, q) := (\varphi_{\alpha}(p), \psi_{\beta}(q)) $，是 $ M \times N $ 的一个图册。

\end{exercises}

\section{切线与法向}

我们现在引入线性结构，这些结构使我们能够给子流形之间的映射赋予导数。这些结构将用局部坐标来描述，而局部坐标在具体计算中当然非常有用。

为了说明为什么需要这样的结构，我们考虑 $ \mathbb{R}^{3} $ 中单位球面 $ S^{2} $ 上的实值函数 $ f : S^{2} \to \mathbb{R} $。通过差商极限来定义 $ f $ 的导数的朴素尝试注定立即失败：对 $ p \in S^{2} $ 以及满足 $ h \ne 0 $ 的 $ h \in \mathbb{R}^{3} $，$ p + h $ 一般并不落在 $ S^{2} $ 中，因此 $ f $ 的“增量” $ f(p + h) - f(p) $ 甚至在点 $ p $ 处根本没有定义。

\needspace{3\baselineskip}\subsection*{$ \mathbb{R}^{n} $ 中的切向}

我们从 $ \mathbb{R}^{n} $ 的 $ n $ 维子流形这一简单情形开始（见例 9.1(a)）。设 $ X $ 在 $ \mathbb{R}^{n} $ 中开，且 $ p \in X $。$ X $ 在点 $ p $ 处的\textbf{切向空间} $ T_{p}X $ 是集合 $ \{p\}\times\mathbb{R}^{n} $，并赋以 $ \mathbb{R}^{n} = (\mathbb{R}^{n}, (\cdot|\cdot)) $ 的诱导欧几里得向量空间结构，即

\[
(p,v) + \lambda(p,w) := (p, v + \lambda w) \quad \text{且} \quad ((p,v)|(p,w))_{p} := (v|w)
\]

对 $ (p,v), (p,w) \in T_{p}X $ 和 $ \lambda \in \mathbb{R} $。$ T_{p}X $ 中的元素 $ (p,v) $ 称为 $ X $ 在 $ p $ 处的一个\textbf{切向向量}，也记作 $ (v)_{p} $。我们称 $ v $ 为 $ (v)_{p} $ 的\textbf{切部分}。\footnote{我们区分“切向的”与“切的”：切向向量含有“基点” $ p $，而其切部分不含。}

\begin{remark}

切向空间 $ T_{p}X $ 与 $ \mathbb{R}^{n} $ 是等距同构的希尔伯特空间。把 $ \mathbb{R}^{n} $“附着”到点 $ p \in X $ 上，可以清楚地表达这个等距同构。\nobreak\hspace{0.35em}■

设 $ Y $ 在 $ \mathbb{R}^{\ell} $ 中开，且 $ f \in C^{1}(X,Y) $。则线性映射

\[
T_{p}f : T_{p}X \to T_{f(p)}Y , \quad (p,v) \mapsto (f(p), \partial f(p)v)
\]

称为 $ f $ 在点 $ p $ 处的\textbf{切向}。

\end{remark}

\begin{remark}

\begin{items}
\item

显然

\[
T_{p}f \in \mathcal{L}(T_{p}X, T_{f(p)}Y) .
\]

因为当 $ v $ 在 $ \mathbb{R}^{n} $ 的零邻域中时，映射 $ v \mapsto f(p) + \partial f(p)v $ 在点 $ p $ 处（直到一阶）逼近映射 $ f $，所以我们知道 $ \operatorname{im}(T_{p}f) $ 是 $ T_{f(p)}Y $ 的向量子空间，它在点 $ f(p) $ 处（直到一阶）逼近 $ f(X) $。

\end{items}

\bookfigure[0.88]{figures/0f9f65d21a26b217.png}

\begin{items}
\item[（b）]

若 $ Z $ 在 $ \mathbb{R}^{s} $ 中开且 $ g \in C^{1}(Y,Z) $，则\textbf{链式法则}为

\[
T_{p}(g\circ f) = T_{f(p)}g\circ T_{p}f .
\]

换言之，下面是交换图。

\begin{equation*}
\begin{array}{ccc} X & \xrightarrow{f} & Y \\ g\circ f \searrow & & \swarrow g \\ & Z & \end{array} \qquad \begin{array}{ccc} T_{p}X & \xrightarrow{T_{p}f} & T_{f(p)}Y \\ T_{p}(g\circ f) \searrow & & \swarrow T_{f(p)}g \\ & T_{g(f(p))} & \end{array}
\end{equation*}

\proofhead
这由 $ C^{1} $ 映射的链式法则（见定理 3.3）推出。\nobreak\hspace{0.35em}■
\item[（c）]

对 $ f \in \operatorname{Diff}^{1}(X,Y) $，有

\[
T_{p}f \in \mathcal{L}\mathrm{is}(T_{p}X, T_{f(p)}Y) \quad \text{且} \quad (T_{p}f)^{-1} = T_{f(p)}f^{-1} \quad \text{对 } p \in X .
\]

\proofhead
这是 (b) 的推论。\nobreak\hspace{0.35em}■
\end{items}

\end{remark}

\needspace{3\baselineskip}\subsection*{切向空间}

在本节余下部分，设

• $ M $ 是 $ \mathbb{R}^{n} $ 的 $ m $ 维 $ C^{q} $ 子流形；
$ q \in \mathbb{N}^{\times} \cup \{\infty\} $；
$ p \in M $，且 $ (\varphi,U) $ 是 $ M $ 在 $ p $ 附近的一个坐标卡；
$ (g,V) $ 是属于 $ (\varphi,U) $ 的参数化。

$ M $ 在点 $ p $ 处的\textbf{切向空间} $ T_{p}M $ 是 $ T_{\varphi(p)}V $ 在 $ T_{\varphi(p)}g $ 下的像，因此 $ T_{p}M = \operatorname{im}(T_{\varphi(p)}g) $。$ T_{p}M $ 的元素称为 $ M $ 在 $ p $ 处的\textbf{切向向量}，而 $ TM := \bigcup_{p\in M}T_{p}M $ 是 $ M $ 的\textbf{切向丛}\footnote{对 $ p \in M $，有包含关系 $ T_{p}M \subset M\times\mathbb{R}^{n} $。因此 $ TM $ 是 $ M\times\mathbb{R}^{n} $ 的子集。这里我们用更简单的词“切”，因为它更标准，且无需“切部分之丛”。}。

\begin{remark}

\begin{items}
\item

$ T_{p}M $ 是良定义的，也就是说，它与所选的坐标卡 $ (\varphi,U) $ 无关。

\proofhead
设 $ (\tilde{\varphi},\tilde{U}) $ 是 $ M $ 在 $ p $ 附近的另一个坐标卡，其相应的参数化为 $ (\tilde{g},\tilde{V}) $。不失一般性，我们可以假设坐标卡覆盖的区域 $ U $ 与 $ \tilde{U} $ 重合。否则考虑 $ U\cap\tilde{U} $。由链式法则（注 10.2(b)），图
\begin{equation*}
\begin{array}{ccc} & T_{p}\mathbb{R}^{n} & \\ T_{\varphi(p)}g \nearrow & & \nwarrow T_{\tilde{\varphi}(p)}\tilde{g} \\ T_{\varphi(p)}V & \xrightarrow{\ T_{\varphi(p)}(\tilde{\varphi}\circ\varphi^{-1})\ } & T_{\tilde{\varphi}(p)}\tilde{V} \end{array}
\end{equation*}

交换。由命题 9.15 与注 10.2(c)，$ T_{\varphi(p)}(\tilde{\varphi}\circ\varphi^{-1}) $ 是同构，由此立即得到断言。\nobreak\hspace{0.35em}■

\item

若 $ M $ 是 $ \mathbb{R}^{n} $ 的子集，则上述 $ T_{p}M $ 的定义与注 10.1 中的定义一致。特别地，有 $ TM = M\times\mathbb{R}^{n} $。

\item

切向空间 $ T_{p}M $ 是 $ T_{p}\mathbb{R}^{n} $ 的 $ m $ 维向量子空间，因而是一个 $ m $ 维希尔伯特空间，其内积 $ (\cdot|\cdot)_{p} $ 由 $ T_{p}\mathbb{R}^{n} $ 诱导。

\proofhead
令 $ x_{0} := \varphi(p) $。对 $ (v)_{x_{0}} \in T_{x_{0}}V $，有 $ (T_{x_{0}}g)(v)_{x_{0}} = (p, \partial g(x_{0})v) $。因此可得
\[
\operatorname{rank}T_{x_{0}}g = \operatorname{rank}\partial g(x_{0}) = m ,
\]

从而证得断言。\nobreak\hspace{0.35em}■

\item

因为 $ T_{\varphi(p)}g : T_{\varphi(p)}V \to T_{p}\mathbb{R}^{n} $ 是单射且 $ T_{p}M = \operatorname{im}(T_{\varphi(p)}g) $，所以恰有一个 $ A \in \mathcal{L}\mathrm{is}(T_{p}M, T_{\varphi(p)}V) $ 使得 $ (T_{\varphi(p)}g)(A) = i_{T_{p}M} $，其中 $ i_{T_{p}M} $ 是 $ T_{p}M $ 到 $ T_{p}\mathbb{R}^{n} $ 的典范单射。换言之，当把 $ T_{\varphi(p)}g $ 理解为 $ T_{\varphi(p)}V $ 到其像 $ T_{p}M $ 上的映射时，$ A $ 是 $ T_{\varphi(p)}g $ 的逆。我们称 $ T_{p}\varphi := A $ 为坐标卡 $ \varphi $ 在点 $ p $ 处的\textbf{切向}。此外，$ (T_{p}\varphi)v \in T_{\varphi(p)}V $ 是切向向量 $ v \in T_{p}M $ 在 $ \varphi $ 诱导的\textbf{局部坐标}下的\textbf{表示}。若 $ (\tilde{\varphi},\tilde{U}) $ 是 $ M $ 在 $ p $ 附近的另一个坐标卡，则图

\begin{equation*}
\begin{array}{ccc} & T_{p}M & \\ T_{p}\varphi \swarrow & \cong & \searrow T_{p}\tilde{\varphi} \\ T_{\varphi(p)}\varphi(U) & \xrightarrow[\cong]{\ T_{\varphi(p)}(\tilde{\varphi}\circ\varphi^{-1})\ } & T_{\tilde{\varphi}(p)}\tilde{\varphi}(\tilde{U}) \end{array}
\end{equation*}

交换，其中 $ \cong $ 表示“同构”。

\proofhead
不失一般性，假设 $ U = \tilde{U} $。对 $ \tilde{g} := i_{M}\circ\tilde{\varphi}^{-1} $，由 $ \tilde{g} = (i_{M}\circ\varphi^{-1})\circ(\varphi\circ\tilde{\varphi}^{-1}) = g\circ(\varphi\circ\tilde{\varphi}^{-1}) $ 与注 10.2(b) 的链式法则可得
\[
T_{\tilde{\varphi}(p)}\tilde{g} = T_{\varphi(p)}g\,T_{\tilde{\varphi}(p)}(\varphi\circ\tilde{\varphi}^{-1}) .
\]

由此，$ T_{p}\varphi $ 与 $ T_{p}\tilde{\varphi} $ 的定义、注 10.2(c) 以及命题 9.15 给出关系

\[
T_{p}\tilde{\varphi} = \left(T_{\tilde{\varphi}(p)}(\varphi\circ\tilde{\varphi}^{-1})\right)^{-1}T_{p}\varphi = T_{\varphi(p)}(\tilde{\varphi}\circ\varphi^{-1})\,T_{p}\varphi ,
\]

从而证得断言。\nobreak\hspace{0.35em}■

\end{items}

\bookfigure[0.84]{figures/016b85d65b6bd6ab.png}

\begin{items}
\item[（e）]

（局部坐标中的内积）设 $ x_{0} := \varphi(p) $，且

\[
g_{jk}(x_{0}) := \left(\partial_{j}g(x_{0})\big|\partial_{k}g(x_{0})\right) \quad \text{对 } 1 \le j,k \le m .
\]

则 $ [g_{jk}] \in \mathbb{R}^{m\times m}_{\mathrm{sym}} $ 称为 $ M $ 关于坐标卡 $ \varphi $（或参数化 $ g $）在 $ p $ 处（即 $ x_{0} $ 处）的\textbf{（第一）基本矩阵}。它是正定的。
对 $ v,w \in T_{p}M $，有\footnote{另见注 2.17(c)。}

\begin{equation}
(v|w)_{p} = \sum_{j,k=1}^{m}g_{jk}(x_{0})v^{j}w^{k} ,\tag{10.1}
\end{equation}

其中 $ v^{j} $ 与 $ w^{k} $ 分别是 $ (T_{p}\varphi)v $ 与 $ (T_{p}\varphi)w $ 的切部分关于标准基的局部表示的分量。

\proofhead
对 $ \xi,\eta \in \mathbb{R}^{m} $，有
\begin{equation}
\sum_{j,k=1}^{m}g_{jk}(x_{0})\xi^{j}\eta^{k} = \left(\partial g(x_{0})\xi\big|\partial g(x_{0})\eta\right) .\tag{10.2}
\end{equation}

特别地，因为 $ \partial g(x_{0}) $ 是单射，有

\[
\left([g_{jk}](x_{0})\xi\big|\xi\right) = |\partial g(x_{0})\xi|^{2} > 0 \quad \text{对 } \xi \in \mathbb{R}^{m}\setminus\{0\} .
\]

因此基本矩阵是正定的。
因为 $ v = T_{x_{0}}g\,(T_{p}\varphi)v $ 且 $ (T_{p}\varphi)v = \sum_{j=1}^{m}v^{j}e_{j} $，由 $ (\cdot|\cdot)_{p} $ 的定义可得

\[
(v|w)_{p} = \left((T_{x_{0}}g)(T_{p}\varphi)v\big|(T_{x_{0}}g)(T_{p}\varphi)w\right)_{p} = \left(\partial g(x_{0})\sum_{j=1}^{m}v^{j}e_{j}\Big|\partial g(x_{0})\sum_{k=1}^{m}w^{k}e_{k}\right) .
\]

从而 (10.1) 蕴含 (10.2)。\nobreak\hspace{0.35em}■

\end{items}

\end{remark}

\begin{example}

（图像的切向空间）设 $ X $ 在 $ \mathbb{R}^{m} $ 中开，并令 $ f \in C^{q}(X,\mathbb{R}^{\ell}) $。由命题 9.2，$ M := \operatorname{graph}(f) $ 是 $ \mathbb{R}^{m}\times\mathbb{R}^{\ell} = \mathbb{R}^{m+\ell} $ 的 $ m $ 维 $ C^{q} $ 子流形。则对 $ x \in X $ 令 $ g(x) := (x, f(x)) $，$ g $ 是 $ M $ 的一个 $ C^{q} $ 参数化，且对 $ p = (x_{0}, f(x_{0})) \in M $，有

\[
(T_{x_{0}}g)(v)_{x_{0}} = \left(p, (v, \partial f(x_{0})v)\right) \quad \text{对 } (v)_{x_{0}} \in T_{x_{0}}X .
\]

因此，我们得到

\[
T_{p}M = \left\{ \left(p, (v, \partial f(x_{0})v)\right) ;\ v \in \mathbb{R}^{m} \right\} ,
\]

也就是说，$ T_{p}M $ 是“附着于点 $ p = (x_{0}, f(x_{0})) $”的 $ \partial f(x_{0}) $ 的图像。

\bookfigure[0.61]{figures/8c98067fadf4c971.png}

对 $ n = m + \ell $，$ T_{p}M $ 在 $ \mathbb{R}^{n} $ 中的表示是通过把 $ (\eta)_{p} = (p,\eta) \in T_{p}\mathbb{R}^{n} $ 等同于 $ p + \eta \in \mathbb{R}^{n} $ 而得到的。由此可得

\begin{equation*}
\begin{aligned}
T_{p}M \ni &\left((x_{0}, f(x_{0})), (v, \partial f(x_{0})v)\right) \\ &= \left(x_{0} + v, f(x_{0}) + \partial f(x_{0})v\right) \in \mathbb{R}^{m}\times\mathbb{R}^{\ell} = \mathbb{R}^{n} ,
\end{aligned}
\end{equation*}

并且我们得到

\begin{equation*}
\begin{aligned}
T_{p}M &= \left\{ \left(x, f(x_{0}) + \partial f(x_{0})(x - x_{0})\right) ;\ x \in \mathbb{R}^{m} \right\} \\ &= \operatorname{graph}\left(x \mapsto f(x_{0}) + \partial f(x_{0})(x - x_{0})\right) .
\end{aligned}
\end{equation*}

这一表示再次表明 $ T_{p}M $ 是曲线切线概念的高维推广（见注 IV.1.4）。\nobreak\hspace{0.35em}■

设 $ \varepsilon > 0 $，使得对 $ t \in (-\varepsilon,\varepsilon) $ 和 $ j \in \{1,\ldots,m\} $ 有 $ \varphi(p) + te_{j} \in V $。则

\[
\gamma_{j}(t) := g\left(\varphi(p) + te_{j}\right) \quad \text{对 } t \in (-\varepsilon,\varepsilon) ,
\]

称为过 $ p $ 的第 $ j $ 条\textbf{坐标道路}。

\end{example}

\begin{remark}

对 $ x_{0} := \varphi(p) $，有

\[
T_{p}M = \operatorname{span}\left\{ \left(\partial_{1}g(x_{0})\right)_{p}, \ldots, \left(\partial_{m}g(x_{0})\right)_{p} \right\} ,
\]

也就是说，$ p $ 处在坐标道路 $ \gamma_{j} $ 上的切向向量构成 $ T_{p}M $ 的一组基。

\proofhead
对 $ [\partial g(x_{0})] $ 的第 $ j $ 列，有
\[
\partial_{j}g(x_{0}) = \partial g(x_{0})e_{j} = \dot{\gamma}_{j}(0) .
\]

该注成立，因为 $ g $ 是浸入，从而 $ [\partial g(x_{0})] $ 的列向量线性无关，且 $ \dim T_{p}M = m $。\nobreak\hspace{0.35em}■

\bookfigure[0.27]{figures/6387f089ada6bfc3.png}
\end{remark}

\needspace{3\baselineskip}\subsection*{切向空间的刻画}

$ M $ 在 $ p $ 处的切向向量可以描述为 $ M $ 中正则道路的切向向量。下一个定理给出切向空间的几何解释。

\begin{theorem}

对每个 $ p \in M $，有

\begin{equation*}
\begin{aligned}
T_{p}M = \left\{ (v)_{p} \in T_{p}\mathbb{R}^{n} ;\ \exists\,\varepsilon > 0,\ \exists\,\gamma \in C^{1}\left((-\varepsilon,\varepsilon),\mathbb{R}^{n}\right) \text{ 使得 } \right. \\ \left. \operatorname{im}(\gamma) \subset M,\ \gamma(0) = p,\ \dot{\gamma}(0) = v \right\} .
\end{aligned}
\end{equation*}

换言之，对每个 $ (v)_{p} \in T_{p}M \subset T_{p}\mathbb{R}^{n} $，存在 $ \mathbb{R}^{n} $ 中过 $ p $ 的一条 $ C^{1} $ 道路，它整个包含在 $ M $ 中，且以 $ (v)_{p} $ 作为它在 $ p $ 处的切向向量。这种道路的每个切向向量都属于 $ T_{p}M $。

证明

\begin{subitems}
\item

设 $ (v)_{p} \in T_{p}M $ 且 $ x_{0} := \varphi(p) $。则存在 $ \xi \in \mathbb{R}^{m} $ 使得 $ v = \partial g(x_{0})\xi $。因为 $ V = \varphi(U) $ 在 $ \mathbb{R}^{m} $ 中开，所以存在 $ \varepsilon > 0 $ 使得对 $ t \in (-\varepsilon,\varepsilon) $ 有 $ x_{0} + t\xi \in V $。现在对 $ t \in (-\varepsilon,\varepsilon) $ 令 $ \gamma(t) := g(x_{0} + t\xi) $，于是 $ \gamma $ 是 $ M $ 中的一条 $ C^{1} $ 道路，满足 $ \gamma(0) = p $ 且 $ \dot{\gamma}(0) = \partial g(x_{0})\xi = v $。

\item

设 $ \gamma \in C^{1}\left((-\varepsilon,\varepsilon),\mathbb{R}^{n}\right) $，其中 $ \operatorname{im}(\gamma) \subset M $ 且 $ \gamma(0) = p $。由注 9.9(d)，存在 $ (x_{0},0) $ 在 $ \mathbb{R}^{n} $ 中的开邻域 $ \tilde{V} $、$ \mathbb{R}^{n} $ 中的开邻域 $ \tilde{U} $，以及一个 $ \psi \in \operatorname{Diff}^{q}(\tilde{V},\tilde{U}) $ 使得对 $ x \in V $ 有 $ \psi(x,0) = g(x) $。通过缩小 $ \varepsilon $，可以假设 $ \operatorname{im}(\gamma) \subset U\cap\tilde{U} $。由此可得

\[
\gamma(t) = (g\circ\varphi\circ\gamma)(t) = (g\circ\operatorname{pr}_{\mathbb{R}^{m}}\circ\psi^{-1}\circ\gamma)(t) ,
\]

并由链式法则得到

\[
\dot{\gamma}(0) = \partial g(x_{0})\left(\operatorname{pr}_{\mathbb{R}^{m}}\circ\psi^{-1}\circ\gamma\right)^{\cdot}(0) .
\]

对 $ \xi := \left(\operatorname{pr}_{\mathbb{R}^{m}}\circ\psi^{-1}\circ\gamma\right)^{\cdot}(0) \in \mathbb{R}^{m} $ 和 $ v := \partial g(x_{0})\xi \in \mathbb{R}^{n} $，有 $ (v)_{p} \in T_{p}M $，证毕。\nobreak\hspace{0.35em}■

\end{subitems}

\end{theorem}

若 $ X $ 在 $ \mathbb{R}^{n} $ 中开，且点 $ c \in \mathbb{R}^{\ell} $ 是 $ f \in C^{q}(X,\mathbb{R}^{\ell}) $ 的正则值，则由定理 9.3 我们知道 $ M = f^{-1}(c) $ 是 $ \mathbb{R}^{n} $ 的 $ (n-\ell) $ 维 $ C^{q} $ 子流形。在下一个定理中，我们证明“线性化”也有类似的结论，即

\[
T_{p}M = (T_{p}f)^{-1}(c) = \ker(T_{p}f) .
\]

\begin{theorem}

（正则值）设 $ X $ 在 $ \mathbb{R}^{n} $ 中开，且 $ c \in \mathbb{R}^{\ell} $ 是 $ f \in C^{q}(X,\mathbb{R}^{\ell}) $ 的正则值。对 $ \mathbb{R}^{n} $ 的 $ (n-\ell) $ 维 $ C^{q} $ 子流形 $ M := f^{-1}(c) $，则对 $ p \in M $ 有 $ T_{p}M = \ker(T_{p}f) $。

\proofhead
设 $ (v)_{p} \in T_{p}M \subset T_{p}\mathbb{R}^{n} $。由定理 10.6，存在 $ \varepsilon > 0 $ 和一条道路 $ \gamma \in C^{1}\left((-\varepsilon,\varepsilon),\mathbb{R}^{n}\right) $ 使得 $ \operatorname{im}(\gamma) \subset M $，且 $ \gamma(0) = p $、$ \dot{\gamma}(0) = v $。特别地，对每个 $ t \in (-\varepsilon,\varepsilon) $ 有 $ f\left(\gamma(t)\right) = c $，对此关系求导可得
\[
\partial f\left(\gamma(0)\right)\dot{\gamma}(0) = \partial f(p)v = 0 .
\]

由此可得 $ T_{p}M \subset \ker(T_{p}f) $。

因为 $ p $ 是 $ f $ 的正则点，有

\[
\dim\left(\operatorname{im}(T_{p}f)\right) = \dim\left(\operatorname{im}(\partial f(p))\right) = \ell .
\]

因此，秩公式 (2.4) 给出

\[
\dim\left(\ker(T_{p}f)\right) = n - \ell = \dim(T_{p}M) .
\]

所以 $ T_{p}M $ 不是 $ \ker(T_{p}f) $ 的真向量子空间。\nobreak\hspace{0.35em}■

\end{theorem}

\needspace{3\baselineskip}\subsection*{可微映射}

设 $ N $ 是 $ \mathbb{R}^{\ell} $ 的 $ C^{r} $ 子流形，且 $ 1 \le s \le \min\{q,r\} $。又设 $ f \in C(M,N) $，$ (\psi,W) $ 是 $ N $ 在 $ f(p) $ 附近的一个坐标卡。则 $ U\cap f^{-1}(W) $ 是 $ p $ 在 $ M $ 中的开邻域。因此通过缩小 $ U $，可以不失一般性地假设 $ f(U) \subset W $。若映射

\[
f_{\varphi,\psi} := \psi\circ f\circ\varphi^{-1} : \varphi(U) \to \psi(W)
\]

在点 $ \varphi(p) $ 处是（$ s $ 次）[连续]可微的，则称函数 $ f $ 在 $ p $ 处（$ s $ 次）[连续]可微。

若 $ f $ 在 $ M $ 的每一点处都（$ s $ 次）[连续]可微，则称映射 $ f \in C(M,N) $（$ s $ 次）[连续]可微。我们用 $ C^{s}(M,N) $ 表示从 $ M $ 到 $ N $ 的所有 $ s $ 次连续可微函数之集，而

\[
\operatorname{Diff}^{s}(M,N) := \left\{ f \in C^{s}(M,N) ;\ f \text{ 是双射},\ f^{-1} \in C^{s}(N,M) \right\}
\]

是从 $ M $ 到 $ N $ 的所有 $ C^{s} $ 微分同胚之集。最后，若 $ \operatorname{Diff}^{s}(M,N) $ 非空，则称 $ M $ 与 $ N $ $ C^{s} $-微分同胚。

\begin{remark}

\begin{items}
\item

上述定义与坐标卡的选择无关。

\proofhead
若 $ (\tilde{\varphi},\tilde{U}) $ 与 $ (\tilde{\psi},\tilde{W}) $ 分别是 $ M $ 在 $ p $ 附近和 $ N $ 在 $ f(p) $ 附近的另一些坐标卡，使得 $ f(\tilde{U}) \subset \tilde{W} $，则
\begin{equation}
f_{\tilde{\varphi},\tilde{\psi}} = \tilde{\psi}\circ f\circ\tilde{\varphi}^{-1} = (\tilde{\psi}\circ\psi^{-1})\circ f_{\varphi,\psi}\circ(\varphi\circ\tilde{\varphi}^{-1}) .\tag{10.3}
\end{equation}

由此，结合命题 9.15 与链式法则，即得本定理。\nobreak\hspace{0.35em}■

\item

若 $ M $ 与 $ N $ 的维数分别为 $ n $ 与 $ \ell $，即 $ M $ 在 $ \mathbb{R}^{n} $ 中开、$ N $ 在 $ \mathbb{R}^{\ell} $ 中开，则上述定义与第 5 节中的定义一致。

\item

函数 $ f_{\varphi,\psi} $ 是 $ f $ 在坐标卡 $ \varphi $ 与 $ \psi $ 下的\textbf{局部表示}，或\textbf{局部坐标表示}。与映射于“弯曲集合” $ M $ 与 $ N $ 之间的函数 $ f $ 不同，$ f_{\varphi,\psi} $ 是欧几里得空间开子集之间的映射。

\item

若 $ s > \min\{q,r\} $，则一般不可能合理地定义一个与坐标无关的 $ M $ 到 $ N $ 的 $ C^{s} $ 映射概念。

\proofhead
这由 (10.3) 推出，因为过渡函数只分别属于 $ C^{q} $ 或 $ C^{r} $。\nobreak\hspace{0.35em}■
\end{items}

设 $ f : M \to N $ 在 $ p $ 处可微，且 $ (\psi,W) $ 是 $ N $ 在 $ f(p) $ 附近使得 $ f(U) \subset W $ 的坐标卡。则图

\begin{equation}
\begin{array}{ccc} M \supset U & \xrightarrow{f} & W \subset N \\ \varphi \downarrow \cong & & \cong \downarrow \psi \\ \mathbb{R}^{m} \supset \varphi(U) & \xrightarrow{f_{\varphi,\psi}} & \psi(W) \subset \mathbb{R}^{\bar{n}} \end{array}\tag{10.4}
\end{equation}

交换，其中 $ \cong $ 表示 $ C^{1} $-微分同胚\footnote{注意局部表示 $ \varphi_{\varphi,\operatorname{id}} = \operatorname{id}_{\varphi(U)} $。}，$ \bar{n} $ 是 $ N $ 的维数。我们

现在通过要求图

\begin{equation}
\begin{array}{ccc} T_{p}M & \xrightarrow{T_{p}f} & T_{f(p)}N \\ T_{p}\varphi \downarrow \cong & & \cong \downarrow T_{f(p)}\psi \\ T_{\varphi(p)}\varphi(U) & \xrightarrow{T_{\varphi(p)}f_{\varphi,\psi}} & T_{\psi(f(p))}\psi(W) \end{array}\tag{10.5}
\end{equation}

交换来定义 $ f $ 在 $ p $ 处的\textbf{切向} $ T_{p}f $，这里 $ \cong $ 表示“同构”。

\end{remark}

\begin{remark}

\begin{items}
\item

切向 $ T_{p}f $ 与坐标无关，且 $ T_{p}f \in \mathcal{L}(T_{p}M, T_{f(p)}N) $。

\proofhead
设 $ (\tilde{\varphi},\tilde{U}) $ 是 $ M $ 在 $ p $ 附近的坐标卡，$ (\tilde{\psi},\tilde{W}) $ 是 $ f(p) $ 附近使得 $ f(\tilde{U}) \subset \tilde{W} $ 的坐标卡。则 (10.3) 与注 10.2(b) 的链式法则给出
\begin{equation*}
\begin{aligned}
T_{\tilde{\varphi}(p)}f_{\tilde{\varphi},\tilde{\psi}} &= T_{\tilde{\varphi}(p)}\left((\tilde{\psi}\circ\psi^{-1})\circ f_{\varphi,\psi}\circ(\varphi\circ\tilde{\varphi}^{-1})\right) \\ &= T_{\psi(f(p))}(\tilde{\psi}\circ\psi^{-1})\circ T_{\varphi(p)}f_{\varphi,\psi}\circ T_{\tilde{\varphi}(p)}(\varphi\circ\tilde{\varphi}^{-1}) .
\end{aligned}
\end{equation*}

于是该注由注 10.3(d) 推出。\nobreak\hspace{0.35em}■

\item

设 $ O $ 是另一个流形，且 $ g : N \to O $ 在 $ f(p) $ 处可微。则有\textbf{链式法则}

\begin{equation}
T_{p}(g\circ f) = T_{f(p)}g\,T_{p}f\tag{10.6}
\end{equation}

以及

\begin{equation}
T_{p}\operatorname{id}_{M} = \operatorname{id}_{T_{p}M} .\tag{10.7}
\end{equation}

若 $ f $ 属于 $ \operatorname{Diff}^{1}(M,N) $，则

\[
T_{p}f \in \mathcal{L}\mathrm{is}(T_{p}M, T_{f(p)}N) \quad \text{且} \quad (T_{p}f)^{-1} = T_{f(p)}f^{-1} \quad \text{对 } p \in M .
\]

\proofhead
陈述 (10.6) 与 (10.7) 由 (10.5) 图的交换性以及注 10.2(b) 的链式法则易得。其余断言是 (10.6) 与 (10.7) 的直接推论。\nobreak\hspace{0.35em}■
\end{items}

\end{remark}

\begin{example}

\begin{items}
\item

$ M $ 到 $ \mathbb{R}^{n} $ 的典范单射 $ i_{M} $ 属于 $ C^{q}(M,\mathbb{R}^{n}) $，且对 $ \psi := \operatorname{id}_{\mathbb{R}^{n}} $，有

\[
T_{\varphi(p)}(i_{M})_{\varphi,\psi} = T_{\varphi(p)}g .
\]

因此 $ T_{p}i_{M} $ 是 $ T_{p}M $ 到 $ T_{p}\mathbb{R}^{n} $ 的典范单射。

\proofhead
这显然由 $ (i_{M})_{\varphi,\psi} = i_{M}\circ\varphi^{-1} = g $ 推出。\nobreak\hspace{0.35em}■
\item

设 $ X $ 是 $ M $ 的开邻域，且 $ \tilde{f} \in C^{s}(X,\mathbb{R}^{\ell}) $。又设 $ \tilde{f}(M) \subset N $。则 $ f := \tilde{f}|M $ 属于 $ C^{s}(M,N) $，且

\[
T_{f(p)}i_{N}\,T_{p}f = T_{p}\tilde{f}\,T_{p}i_{M} \quad \text{对 } p \in M ,
\]

也就是说，图

\begin{equation*}
\begin{array}{ccc} M & \xrightarrow{i_{M}} & X \\ f \downarrow & & \downarrow \tilde{f} \\ N & \xrightarrow{i_{N}} & \mathbb{R}^{\ell} \end{array} \qquad \begin{array}{ccc} T_{p}M & \xrightarrow{T_{p}i_{M}} & T_{p}X \\ T_{p}f \downarrow & & \downarrow T_{p}\tilde{f} \\ T_{f(p)}N & \xrightarrow{T_{f(p)}i_{N}} & T_{f(p)}\mathbb{R}^{\ell} \end{array}
\end{equation*}

交换。

\proofhead
因为 $ N $ 是 $ \mathbb{R}^{\ell} $ 的 $ C^{r} $ 子流形，我们（可能通过缩小 $ W $）可以假设存在 $ W $ 在 $ \mathbb{R}^{\ell} $ 中的开邻域 $ \bar{W} $，以及 $ \bar{W} $ 到 $ \mathbb{R}^{\ell} $ 的某个开子集上的 $ C^{r} $ 微分同胚 $ \Psi $，使得 $ \Psi \supset \psi $。于是我们得到
\[
f_{\varphi,\psi} = \psi\circ f\circ\varphi^{-1} = \Psi\circ\tilde{f}\circ g \in C^{s}(V,\mathbb{R}^{\ell}) .
\]

因为这对 $ M $ 的每个坐标卡 $ (\varphi,U) $ 与 $ N $ 的每个使得 $ f(U) \subset W $ 的坐标卡 $ (\psi,W) $ 都成立，所以可以断定 $ f $ 属于 $ C^{s}(M,N) $。因为 $ i_{N}\circ f = \tilde{f}\circ i_{M} $，断言的最后部分是链式法则的推论。\nobreak\hspace{0.35em}■

\item

设 $ X $ 在 $ \mathbb{R}^{n} $ 中开，$ Y $ 在 $ \mathbb{R}^{\ell} $ 中开，且 $ f \in C^{q}(X,\mathbb{R}^{\ell}) $ 满足 $ f(X) \subset Y $。则 $ f $ 属于 $ C^{q}(X,Y) $，其中 $ X $ 与 $ Y $ 分别理解为 $ \mathbb{R}^{n} $ 与 $ \mathbb{R}^{\ell} $ 的 $ n $ 维与 $ \ell $ 维子流形。此外，有

\[
T_{p}f = (f(p), \partial f(p)) \quad \text{对 } p \in X .
\]

\proofhead
这是 (b) 的特例。\nobreak\hspace{0.35em}■
\end{items}

\end{example}

\needspace{3\baselineskip}\subsection*{微分与梯度}

若 $ f : M \to \mathbb{R} $ 在 $ p $ 处可微，则

\[
T_{p}f : T_{p}M \to T_{f(p)}\mathbb{R} = \{f(p)\}\times\mathbb{R} \subset \mathbb{R}\times\mathbb{R} .
\]

利用到第二个因子的典范投影 $ \operatorname{pr}_{2} : \mathbb{R}\times\mathbb{R} \to \mathbb{R} $，令

\[
d_{p}f := \operatorname{pr}_{2}\circ T_{p}f \in \mathcal{L}(T_{p}M,\mathbb{R}) = (T_{p}M)' ,
\]

并称 $ d_{p}f $ 为 $ f $ 在点 $ p $ 处的\textbf{微分}。因此 $ d_{p}f $ 是 $ T_{p}M $ 上的连续线性形式。因为切向空间是 $ m $ 维希尔伯特空间，由里斯表示定理，存在唯一的向量 $ \nabla_{p}f := \nabla_{p}^{M}f \in T_{p}M $ 使得

\[
(d_{p}f)v = (\nabla_{p}f|v)_{p} \quad \text{对 } v \in T_{p}M ,
\]

我们称之为 $ f $ 在点 $ p $ 处的\textbf{梯度}。

\begin{remark}

\begin{items}
\item

设 $ X $ 在 $ \mathbb{R}^{n} $ 中开且 $ f \in C^{1}(X,\mathbb{R}) $。则

\[
\nabla_{p}f = (p, \nabla f(p)) \quad \text{对 } p \in X ,
\]

其中 $ \nabla f(p) $ 是 $ f $ 在点 $ p $ 处的梯度，如第 2 节所定义。因此 $ \nabla_{p}f $ 与 $ \nabla f(p) $ 这两个定义是一致的。

\proofhead
我们用平凡坐标卡 $ (\operatorname{id}_{X}, X) $ 来描述流形 $ X $。则 $ d_{p}f $ 的局部表示与 $ \partial f(p) $ 一致。于是断言由 $ (\cdot|\cdot)_{p} $ 与梯度的定义推出。\nobreak\hspace{0.35em}■
\item

（局部坐标表示）设 $ f \in C^{1}(M,\mathbb{R}) $，且 $ f_{\varphi} := f\circ\varphi^{-1} $ 是 $ f $ 在 $ M $ 的坐标卡 $ \varphi $ 与 $ \mathbb{R} $ 的坐标卡 $ \operatorname{id}_{\mathbb{R}} $ 下的局部表示，即 $ f_{\varphi} = f_{\varphi,\operatorname{id}_{\mathbb{R}}} $。又设 $ [g^{jk}] $ 是基本矩阵 $ [g_{jk}] $ 关于 $ \varphi $ 在 $ p $ 处的逆。则在 $ \varphi $ 诱导的局部坐标下，梯度 $ \nabla_{p}f \in T_{p}M $\footnote{将此与公式 (2.6) 比较。} 的局部表示 $ (T_{p}\varphi)\nabla_{p}f $ 有切部分

\[
\left(\sum_{k=1}^{m}g^{1k}(x_{0})\partial_{k}f_{\varphi}(x_{0}), \ldots, \sum_{k=1}^{m}g^{mk}(x_{0})\partial_{k}f_{\varphi}(x_{0})\right)
\]

其中 $ x_{0} := \varphi(p) $。

\proofhead
利用 $ \nabla_{p}f $ 与 $ d_{p}f $ 的定义，由命题 2.8(i) 可得
\begin{equation}
(\nabla_{p}f|v)_{p} = (d_{p}f)v = \partial(f\circ\varphi^{-1})(x_{0})(T_{p}\varphi)v = \sum_{j=1}^{m}\partial_{j}f_{\varphi}(x_{0})v^{j}\tag{10.8}
\end{equation}

对 $ v \in T_{p}M $ 与 $ (T_{p}\varphi)v = \sum_{j=1}^{m}v^{j}e_{j} $。对 $ (T_{p}\varphi)\nabla_{p}f = \sum_{j=1}^{m}w^{j}e_{j} $，利用注 10.3(e) 得到

\begin{equation}
(\nabla_{p}f|v)_{p} = \sum_{j,k=1}^{m}g_{jk}(x_{0})w^{j}v^{k} \quad \text{对 } v \in T_{p}M .\tag{10.9}
\end{equation}

现在 (10.8)、(10.9) 与 $ [g_{jk}] $ 的对称性蕴含

\[
\sum_{k=1}^{m}g_{jk}(x_{0})w^{k} = \partial_{j}f_{\varphi}(x_{0}) \quad \text{对 } 1 \le j \le m ,
\]

因此乘以 $ [g_{jk}]^{-1} = [g^{jk}] $ 后即得结论。\nobreak\hspace{0.35em}■

\end{items}

下一个定理给出 $ p \in M $ 是 $ M $ 上可微实值函数极值点的必要条件。它推广了定理 3.13。

\end{remark}

\begin{theorem}

若 $ p \in M $ 是 $ f \in C^{1}(M,\mathbb{R}) $ 的局部极值点，则 $ \nabla_{p}f = 0 $。

\proofhead
因为 $ f_{\varphi} \in C^{1}(V,\mathbb{R}) $ 在 $ x_{0} = \varphi(p) $ 处有局部极值，所以 $ \partial f_{\varphi}(x_{0}) = 0 $（见命题 2.5 与定理 3.13）。因此，对 $ v \in T_{p}M $ 以及 $ (T_{p}\varphi)v $ 的切部分 $ \xi $，有
\[
(\nabla_{p}f|v)_{p} = (d_{p}f)v = \partial f_{\varphi}(x_{0})\xi = 0 .\quad\blacksquare
\]
\end{theorem}

\needspace{3\baselineskip}\subsection*{法向}

$ T_{p}M $ 在 $ T_{p}\mathbb{R}^{n} $ 中的正交补称为 $ M $ 在 $ p $ 处的\textbf{法空间}，记作 $ T_{p}^{\perp}M $。$ T_{p}^{\perp}M $ 中的向量是 $ M $ 在 $ p $ 处的\textbf{法向量}，而 $ T^{\perp}M := \bigcup_{p\in M}T_{p}^{\perp}M $ 是 $ M $ 的\textbf{法丛}。

\begin{proposition}

设 $ X $ 在 $ \mathbb{R}^{n} $ 中开，且 $ c $ 是 $ f \in C^{q}(X,\mathbb{R}^{\ell}) $ 的正则值。若 $ M := f^{-1}(c) $ 非空，则 $ \{\nabla_{p}f^{1}, \ldots, \nabla_{p}f^{\ell}\} $ 是 $ T_{p}^{\perp}M $ 的一组基。

证明

\begin{subitems}
\item

由注 10.11(a)，$ \nabla_{p}f^{j} $ 有切部分 $ \nabla f^{j}(p) $。由 $ \partial f(p) $ 的满射性可得，$ \mathbb{R}^{n} $ 中的向量 $ \nabla f^{1}(p), \ldots, \nabla f^{\ell}(p) $——从而 $ T_{p}\mathbb{R}^{n} $ 中的 $ \nabla_{p}f^{1}, \ldots, \nabla_{p}f^{\ell} $——线性无关。

\item

设 $ v \in T_{p}M $。由定理 10.6，我们知道存在 $ \varepsilon > 0 $ 和一条 $ \gamma \in C^{1}\left((-\varepsilon,\varepsilon),\mathbb{R}^{n}\right) $ 使得 $ \operatorname{im}(\gamma) \subset M $，$ \gamma(0) = p $，且 $ \gamma(0) = v $。因为对 $ t \in (-\varepsilon,\varepsilon) $ 有 $ f^{j}\left(\gamma(t)\right) = c^{j} $，可得

\[
0 = (f^{j}\circ\gamma)'(0) = \left(\nabla f^{j}(\gamma(0))\big|\dot{\gamma}(0)\right) = (\nabla_{p}f^{j}|v)_{p} \quad \text{对 } 1 \le j \le \ell .
\]

因为这对每个 $ v \in T_{p}M $ 都成立，我们看到 $ \nabla_{p}f^{1}, \ldots, \nabla_{p}f^{\ell} $ 属于 $ T_{p}^{\perp}M $。于是 $ \dim(T_{p}^{\perp}M) = n - \dim(T_{p}M) = \ell $。\nobreak\hspace{0.35em}■

\end{subitems}

\end{proposition}

\begin{example}

\begin{items}
\item

对 $ \mathbb{R}^{n} $ 中的球面 $ S^{n-1} $，当

\[
f : \mathbb{R}^{n} \to \mathbb{R} , \quad x \mapsto |x|^{2} .
\]

时有 $ S^{n-1} = f^{-1}(1) $。因为 $ \nabla_{p}f = (p, 2p) $，有

\[
T_{p}^{\perp}S^{n-1} = (p, \mathbb{R}p) .
\]
\end{items}

\bookfigure[0.20]{figures/3a1a9a94dcb7e97b.png}

\begin{items}
\item[（b）]

对 $ X := \mathbb{R}^{3}\setminus\left(\{0\}\times\{0\}\times\mathbb{R}\right) $ 和

\[
f(x_{1},x_{2},x_{3}) := \left(\sqrt{x_{1}^{2} + x_{2}^{2}} - 2\right)^{2} + x_{3}^{2} ,
\]

有 $ f \in C^{\infty}(X,\mathbb{R}) $，且\footnote{见例 9.11(f)。} $ f^{-1}(1) = \mathbb{T}_{2,1}^{2} $。则在点 $ p = (p_{1},p_{2},p_{3}) $ 处的法向量由 $ \nabla_{p}f = (p, 2(p - k)) $ 给出，其中

\[
k := \left(\frac{2p_{1}}{\sqrt{p_{1}^{2} + p_{2}^{2}}}, \frac{2p_{2}}{\sqrt{p_{1}^{2} + p_{2}^{2}}}, 0\right) .
\]
\end{items}

\bookfigure[0.31]{figures/7bd1c69296de4fe4.png}

\proofhead
这由例 9.14(e) 与重算得到。\nobreak\hspace{0.35em}■

\begin{items}
\item[（c）]

设 $ X $ 在 $ \mathbb{R}^{n} $ 中开且 $ f \in C^{q}(X,\mathbb{R}) $。则一个\textbf{单位法向量}，即长度为 1 的法向量，在 $ M := \operatorname{graph}(f) $ 上的 $ p := (x, f(x)) $ 处由 $ \nu_{p} := (p, \nu(x)) \in T_{p}\mathbb{R}^{n+1} $ 给出，其中

\[
\nu(x) := \frac{(-\nabla f(x), 1)}{\sqrt{1 + |\nabla f(x)|^{2}}} \in \mathbb{R}^{n+1} .
\]

\proofhead
对由 $ g(x) := (x, f(x)) $（$ x \in X $）定义的参数化 $ g $，有
\[
\partial_{j}g(x) := (e_{j}, \partial_{j}f(x)) \quad \text{对 } x \in X , \quad 1 \le j \le n .
\]

显然 $ \nu(x) $ 是 $ \mathbb{R}^{n+1} $ 中长度为 1 的向量，且与每个向量 $ \partial_{j}g(x) $ 正交。因此断言由注 10.5 推出。\nobreak\hspace{0.35em}■

\end{items}

\end{example}

\needspace{3\baselineskip}\subsection*{约束极值}

在许多应用中，我们寻求在某种约束下函数 $ F : \mathbb{R}^{n} \to \mathbb{R} $ 的极值点。换言之，约束意味着并非 $ \mathbb{R}^{n} $ 的每个点都允许使 $ F $ 取极值，相反所求的点必须属于某个子集 $ M $。这些约束常常由形如 $ h^{1}(p) = 0, \ldots, h^{\ell}(p) = 0 $ 的方程来描述，而这些方程的解集是一个子流形，即集合 $ M $。若 $ F|M $ 在 $ p \in M $ 处有局部极值，则称 $ p $ 是 $ F $ \textbf{在约束} $ h^{1}(p) = 0, \ldots, h^{\ell}(p) = 0 $ \textbf{下的极值点}。

下一个在实践中很重要的定理给出了一个点是约束极值点的必要条件。

\begin{theorem}

（拉格朗日乘子）设 $ X $ 在 $ \mathbb{R}^{n} $ 中开，并给定函数 $ F, h^{1}, \ldots, h^{\ell} \in C^{1}(X,\mathbb{R}) $，其中 $ \ell < n $。又设 0 是映射 $ h := (h^{1}, \ldots, h^{\ell}) $ 的正则值，且设 $ M := h^{-1}(0) $ 非空。若 $ p \in M $ 是 $ F $ 在约束 $ h^{1}(p) = 0, \ldots, h^{\ell}(p) = 0 $ 下的极值点，则存在唯一的实数 $ \lambda_{1}, \ldots, \lambda_{\ell} $，即\textbf{拉格朗日乘子}，使得 $ p $ 是下述函数的临界点：

\[
F - \sum_{j=1}^{\ell}\lambda_{j}h^{j} \in C^{1}(X,\mathbb{R})
\]

\proofhead
由正则值定理，我们知道 $ M $ 是 $ \mathbb{R}^{n} $ 的 $ (n-\ell) $ 维 $ C^{1} $ 子流形。由例 10.10(b)，$ f := F|M $ 属于 $ C^{1}(M,\mathbb{R}) $，且因为 $ i_{\mathbb{R}} = \operatorname{id}_{\mathbb{R}} $，有 $ T_{p}f = T_{p}F\,T_{p}i_{M} $。由此可得 $ d_{p}f = d_{p}F\,T_{p}i_{M} $，从而
\begin{equation}
\left(\nabla_{p}f\big|(v)_{p}\right)_{p} = d_{p}F\,T_{p}i_{M}(v)_{p} = \left(\nabla_{p}F\big|T_{p}i_{M}(v)_{p}\right)_{p} = \left(\nabla F(p)\big|v\right)\tag{10.10}
\end{equation}

对 $ (v)_{p} \in T_{p}M \subset T_{p}\mathbb{R}^{n} $。

若 $ p $ 是 $ f $ 的临界点，定理 10.12 断言 $ \nabla_{p}f = 0 $。现在由 (10.10) 可得 $ \nabla F(p) \in T_{p}^{\perp}M $。因此命题 10.13 表明存在唯一的实数 $ \lambda_{1}, \ldots, \lambda_{\ell} $ 使得

\[
\nabla F(p) = \sum_{j=1}^{\ell}\lambda_{j}\nabla h^{j}(p) ,
\]

由于注 3.14(a)，这就完成了证明。\nobreak\hspace{0.35em}■

\end{theorem}

\begin{remark}

拉格朗日乘子法把在约束 $ h^{1}(p) = 0, \ldots, h^{\ell}(p) = 0 $ 下求 $ F $ 极值的问题转化为求函数

\[
F - \sum_{j=1}^{\ell}\lambda_{j}h^{j} \in C^{1}(X,\mathbb{R})
\]

（无约束）临界点的问题。我们通过求解关于 $ n + \ell $ 个未知量 $ p_{1}, \ldots, p_{n}, \lambda_{1}, \ldots, \lambda_{\ell} $ 的 $ \ell + n $ 个方程

\begin{equation*}
\begin{aligned}
h^{j}(p) &= 0 && \text{对 } 1 \le j \le \ell , \\ \partial_{k}\left(F - \sum_{j=1}^{\ell}\lambda_{j}h^{j}\right)(p) &= 0 && \text{对 } 1 \le k \le n ,
\end{aligned}
\end{equation*}

来确定临界点与拉格朗日乘子，其中 $ p = (p^{1}, \ldots, p^{n}) $。随后，我们必须确定这些临界点中哪些实际上是极值点。\nobreak\hspace{0.35em}■

\end{remark}

\needspace{3\baselineskip}\subsection*{拉格朗日乘子的应用}

在下面这些具有独立意义的例子中，我们演示拉格朗日乘子的一些非平凡应用。特别地，我们给出主轴变换定理的一个简短证明，该定理在线性代数中是用其他方法证明的。

\begin{example}

\begin{items}
\item

给定任意向量 $ a_{j} \in \mathbb{R}^{n} $（$ 1 \le j \le n $），\textbf{阿达马不等式}断言

\[
\left|\det[a_{1}, \ldots, a_{n}]\right| \le \prod_{j=1}^{n}|a_{j}| .
\]

证明

\begin{subitems}
\item

由线性代数可知，行列式是列向量的 $ n $ 重线性函数。因此只需验证

\[
-1 \le \det[a_{1}, \ldots, a_{n}] \le 1 \quad \text{对 } a_{j} \in S^{n-1} \text{ 且 } 1 \le j \le n .
\]
\item

令

\[
F(x) := \det[x_{1}, \ldots, x_{n}] \quad \text{对 } h^{j}(x) := |x_{j}|^{2} - 1 , \quad h = (h^{1}, \ldots, h^{n}) ,
\]

且 $ x = (x_{1}, \ldots, x_{n}) \in \mathbb{R}^{n}\times\cdots\times\mathbb{R}^{n} = \mathbb{R}^{n^{2}} $。则 $ F $ 属于 $ C^{\infty}(\mathbb{R}^{n^{2}},\mathbb{R}) $，$ h $ 属于 $ C^{\infty}(\mathbb{R}^{n^{2}},\mathbb{R}^{n}) $，且\footnote{$ x_{j}^{\top} $ 表示将 $ x_{j} $ 理解为行向量。}

\begin{equation*}
[\partial h(x)] = 2\begin{bmatrix} x_{1}^{\top} & & & & 0 \\ & x_{2}^{\top} & & & \\ & & \ddots & & \\ & & & x_{n-1}^{\top} & \\ 0 & & & & x_{n}^{\top} \end{bmatrix} \in \mathbb{R}^{n\times n^{2}} .
\end{equation*}

显然对每个 $ x \in h^{-1}(0) $，$ \partial h(x) $ 的秩都是极大的。因此 0 是 $ h $ 的正则值，且 $ M := h^{-1}(0) $ 是 $ \mathbb{R}^{n^{2}} $ 的 $ n(n-1) $ 维 $ C^{\infty} $ 子流形。此外，因为 $ M = S^{n-1}\times\cdots\times S^{n-1} $，$ M $ 是紧致的。因此 $ f := F|M \in C^{\infty}(M,\mathbb{R}) $ 取到最小值与最大值。

\item

设 $ p = (p_{1}, \ldots, p_{n}) \in M $ 是 $ f $ 的极值点。利用拉格朗日乘子，存在 $ \lambda_{1}, \ldots, \lambda_{n} \in \mathbb{R} $ 使得

\begin{equation}
\begin{aligned}
\nabla F(p) &= \sum_{j=1}^{n}\lambda_{j}\nabla h^{j}(p) = 2\sum_{j=1}^{n}\lambda_{j}(0, \ldots, 0, p_{j}, 0, \ldots, 0) \\ &= 2(\lambda_{1}p_{1}, \ldots, \lambda_{n}p_{n}) \in \mathbb{R}^{n^{2}} .
\end{aligned}\tag{10.11}
\end{equation}

此外，例 4.8(a) 蕴含

\begin{equation}
\frac{\partial F}{\partial x_{j}^{k}}(p) = \det[p_{1}, \ldots, p_{j-1}, e_{k}, p_{j+1}, \ldots, p_{n}] \quad \text{对 } 1 \le j,k \le n .\tag{10.12}
\end{equation}

令 $ B := [p_{1}, \ldots, p_{n}] $，并用 $ B^{\sharp} := [b_{jk}^{\sharp}]_{1\le j,k\le n} $ 表示 $ B $ 的\textbf{余因子矩阵}，其元素为 $ b_{jk}^{\sharp} := (-1)^{j+k}\det B_{jk} $，其中 $ B_{jk} $ 是从 $ B $ 中删去第 $ k $ 行与第 $ j $ 列得到的（见 [Gab96, § A.3.7]）。则 (10.11) 与 (10.12) 给出

\[
\frac{\partial F}{\partial x_{j}^{k}}(p) = b_{jk}^{\sharp} = 2\lambda_{j}p_{j}^{k} .
\]

因为 $ B^{\sharp}B = (\det B)1_{n} $，有

\begin{equation}
\delta_{ij}\det B = \sum_{k=1}^{n}b_{ik}^{\sharp}p_{j}^{k} = 2\lambda_{i}(p_{i}|p_{j}) \quad \text{对 } 1 \le i,j \le n .\tag{10.13}
\end{equation}

在 (10.13) 中取 $ i = j $，得到

\begin{equation}
2\lambda_{1} = \cdots = 2\lambda_{n} = \det B .\tag{10.14}
\end{equation}

在 $ \det B = 0 $ 的情形，断言显然成立。若 $ \det B \ne 0 $，则 (10.13) 与 (10.14) 表明对 $ i \ne j $，$ p_{i} $ 与 $ p_{j} $ 正交。因此 $ B $ 属于 $ O(n) $，且得到 $ |\det B| = 1 $（见习题 9.2）。现在断言成立，因为 $ F(p) = \det B $。\nobreak\hspace{0.35em}■

\end{subitems}

\item

（主轴变换）设 $ A \in \mathcal{L}_{\mathrm{sym}}(\mathbb{R}^{n}) $。则存在实数 $ \lambda_{1} \ge \lambda_{2} \ge \cdots \ge \lambda_{n} $ 和 $ x_{1}, \ldots, x_{n} \in S^{n-1} $ 使得对 $ 1 \le k \le n $ 有 $ Ax_{k} = \lambda_{k}x_{k} $，也就是说，$ x_{k} $ 是以 $ \lambda_{k} $ 为特征值的特征向量。$ x_{1}, \ldots, x_{n} $ 构成 $ \mathbb{R}^{n} $ 的一个标准正交基。在这个基下，$ A $ 有矩阵 $ [A] = \operatorname{diag}(\lambda_{1}, \ldots, \lambda_{n}) $。我们还有

\[
\lambda_{k} = \max\left\{ (Ax|x) ;\ x \in S^{n-1}\cap E_{k} \right\} \quad \text{对 } k = 1, \ldots, n ,
\]

其中 $ E_{1} := \mathbb{R}^{n} $，且对 $ k = 2, \ldots, n $ 有 $ E_{k} := \left(\operatorname{span}\{x_{1}, \ldots, x_{k-1}\}\right)^{\perp} $。

\proofhead
(i) 对 $ x \in \mathbb{R}^{n} $，令 $ h^{0}(x) := |x|^{2} - 1 $ 与 $ F(x) := (Ax|x) $。则 0 是 $ h^{0} \in C^{\infty}(\mathbb{R}^{n},\mathbb{R}) $ 的正则值，且 $ F $ 在 $ S^{n-1} = h^{-1}(0) $ 上取到最大值。设 $ x_{1} \in S^{n-1} $ 是 $ f := F|S^{n-1} $ 的最大点。利用拉格朗日乘子，存在 $ \lambda_{1} \in \mathbb{R} $ 使得 $ \nabla F(x_{1}) = 2Ax_{1} = 2\lambda_{1}x_{1} $（见习题 4.5）。因此 $ x_{1} $ 是 $ A $ 的以 $ \lambda_{1} $ 为特征值的特征向量。因为 $ x_{1} \in S^{n-1} $，我们还有
\[
\lambda_{1} = \lambda_{1}(x_{1}|x_{1}) = (Ax_{1}|x_{1}) = f(x_{1}) .
\]
\end{items}

(ii) 现在我们递归地构造 $ x_{2}, \ldots, x_{n} $。假设对 $ k \ge 2 $ 已知道 $ x_{1}, \ldots, x_{k-1} $，令 $ h := (h^{0}, h^{1}, \ldots, h^{k-1}) $，其中对 $ 1 \le j \le k-1 $ 有 $ h^{j}(x) := 2(x_{j}|x) $。则 $ h^{-1}(0) = S^{n-1}\cap E_{k} $ 是紧致的，且存在 $ x_{k} \in S^{n-1}\cap E_{k} $ 使得对 $ x \in S^{n-1}\cap E_{k} $ 有 $ f(x) \le f(x_{k}) $。此外，（因为对 $ B \in \mathbb{R}^{k\times n} $ 有 $ \operatorname{rank} B = \operatorname{rank} B^{\top} $）可以验证

\[
\operatorname{rank}\partial h(x) = \operatorname{rank}[x, x_{1}, \ldots, x_{k-1}] = k \quad \text{对 } x \in S^{n-1}\cap E_{k} .
\]

因此 0 是 $ h $ 的正则值，我们利用定理 10.15 找到实数 $ \mu_{0}, \ldots, \mu_{k-1} $ 使得

\begin{equation}
2Ax_{k} = \nabla F(x_{k}) = \sum_{j=0}^{k-1}\mu_{j}\nabla h^{j}(x_{k}) = 2\mu_{0}x_{k} + 2\sum_{j=1}^{k-1}\mu_{j}x_{j} .\tag{10.15}
\end{equation}

因为对 $ 1 \le j \le k-1 $ 有 $ (x_{k}|x_{j}) = 0 $，有

\[
(Ax_{k}|x_{j}) = (x_{k}|Ax_{j}) = \lambda_{j}(x_{k}|x_{j}) = 0 \quad \text{对 } 1 \le j \le k-1 .
\]

因此由 (10.15) 可得

\[
0 = (Ax_{k}|x_{j}) = \mu_{j} \quad \text{对 } j = 1, \ldots, k-1 .
\]

由 (10.15) 还看到 $ x_{k} $ 是 $ A $ 的以 $ \mu_{0} $ 为特征值的特征向量。最后得到

\[
\mu_{0} = \mu_{0}(x_{k}|x_{k}) = (Ax_{k}|x_{k}) = f(x_{k}) .
\]

于是我们完成了证明（见注 5.13(b)）。\nobreak\hspace{0.35em}■

\end{example}

\begin{remark}

设 $ A \in \mathcal{L}_{\mathrm{sym}}(\mathbb{R}^{n}) $ 满足 $ n \ge 2 $，并令 $ \lambda_{1} \ge \lambda_{2} \ge \cdots \ge \lambda_{n} $ 为 $ A $ 的特征值。又设 $ x_{1}, \ldots, x_{n} $ 是 $ \mathbb{R}^{n} $ 的一个标准正交基，其中 $ x_{k} $ 是 $ A $ 的以 $ \lambda_{k} $ 为特征值的特征向量。则对 $ x = \sum_{j=1}^{n}\xi^{j}x_{j} \in \mathbb{R}^{n} $，有

\begin{equation}
(Ax|x) = \sum_{j=1}^{n}\lambda_{j}(\xi^{j})^{2} .\tag{10.16}
\end{equation}

现在假设对某个 $ k \in \{0, \ldots, m\} $ 有

\[
\lambda_{1} \ge \cdots \ge \lambda_{k} > 0 > \lambda_{k+1} \ge \cdots \ge \lambda_{m}
\]

则 $ \gamma \in \{-1,1\} $ 是 $ C^{\infty} $ 映射

\[
a : \mathbb{R}^{n} \to \mathbb{R} , \quad x \mapsto (Ax|x)
\]

的正则值（见习题 4.5）。因此，由正则值定理，

\[
a^{-1}(\gamma) = \left\{ x \in \mathbb{R}^{n} ;\ (Ax|x) = \gamma \right\}
\]

是 $ \mathbb{R}^{n} $ 中的 $ C^{\infty} $ 超曲面。令 $ \alpha_{j} := 1/\sqrt{|\lambda_{j}|} $，由 (10.16) 可得，对 $ x = \sum_{j=1}^{n}\xi^{j}x_{j} \in a^{-1}(\gamma) $ 有

\begin{equation}
\sum_{j=1}^{k}\left(\frac{\xi^{j}}{\alpha_{j}}\right)^{2} - \sum_{j=k+1}^{n}\left(\frac{\xi^{j}}{\alpha_{j}}\right)^{2} = \gamma\tag{10.17}
\end{equation}

若 $ A $ 正定，则由注 5.13(b) 可得 $ \lambda_{1} \ge \cdots \ge \lambda_{n} > 0 $。此时，由 (10.17) 读出 $ a^{-1}(1) $ 是一个 $ (n-1) $ 维\textbf{椭球}，其\textbf{主轴}为 $ \alpha_{1}x_{1}, \ldots, \alpha_{n}x_{n} $。若 $ A $ 不定，则 (10.17) 表明 $ a^{-1}(\pm 1) $（一般地）是一个\textbf{双曲面}，其\textbf{主轴}为 $ \alpha_{1}x_{1}, \ldots, \alpha_{n}x_{n} $。

\bookfigure[0.74]{figures/d2c299065ec8da8f.png}

这就把已知的偶维情形推广到了更高维。

\bookfigure[0.74]{figures/ac7187efb0296046.png}

这些考虑阐明了“主轴变换”这一名称。若 $ A $ 的一个或多个特征值为零，则 $ a^{-1}(\gamma) $ 是一个具有椭球形或双曲面形截面的柱面。\nobreak\hspace{0.35em}■

\end{remark}

\exerciseshead

\begin{exercises}
\item

用 $ (g_{n},V_{n}) $ 记 $ F_{n} := \mathbb{R}^{n}\setminus H_{n} $ 的由 $ n $ 维极坐标给出的正则参数化（见习题 9.11）。

\begin{items}
\item

\proofhead
$ F_{n} $ 关于 $ g_{n} $ 的第一基本矩阵 $ [(g_{n})_{jk}] $ 由
\begin{equation*}
\begin{bmatrix} 1 & 0 \\ 0 & r^{2} \end{bmatrix} \quad \text{对 } n = 2 ,
\end{equation*}

给出，或当 $ n \ge 3 $ 时由

\[
\operatorname{diag}\left( 1, r^{2}\sin^{2}(y^{3}) \cdot \cdots \cdot \sin^{2}(y^{n}), r^{2}\sin^{2}(y^{3}) \cdot \cdots \cdot r^{2}\sin^{2}(y^{n-1}), \ldots, r^{2}\sin^{2}(y^{3}), r^{2} \right)
\]

给出，其中 $ (r,y^{2},\ldots,y^{n}) \in V_{n} $。

\item

设 $ f \in C^{1}(F_{n},\mathbb{R}) $，并令 $ \varphi_{n} $ 记属于 $ g_{n} $ 的坐标卡。计算 $ \nabla_{p}f $ 在 $ n $ 维极坐标下，即在由 $ \varphi_{n} $ 诱导的局部坐标下的表示。

\end{items}

\item

设 $ (g,V) $ 是 $ \mathbb{R}^{n} $ 的 $ m $ 维 $ C^{1} $ 子流形 $ M $ 的一个参数化。又用

\[
\sqrt{g}(y) := \sqrt{\det\left[ g_{jk}(y) \right]} \quad \text{对 } y \in V
\]

记 $ M $ 关于 $ g $ 的\textbf{格拉姆行列式}。验证下列陈述：

\begin{items}
\item

若 $ m = 1 $，则 $ \sqrt{g} = |\dot{g}| $。

\item

若 $ m = 2 $ 且 $ n = 3 $，则

\[
\sqrt{g} = \sqrt{\left( \frac{\partial(g^{2},g^{3})}{\partial(x,y)} \right)^{2} + \left( \frac{\partial(g^{3},g^{1})}{\partial(x,y)} \right)^{2} + \left( \frac{\partial(g^{1},g^{2})}{\partial(x,y)} \right)^{2}} = |\partial_{1}g \times \partial_{2}g| ,
\]

其中 $ \times $ 表示叉积（或向量积）（另见第 VIII.2 节）。

\item

设 $ V $ 在 $ \mathbb{R}^{n} $ 中开且 $ f \in C^{1}(V,\mathbb{R}) $。对 $ f $ 的图像的参数化 $ g : V \to \mathbb{R}^{n+1} $，$ x \mapsto (x,f(x)) $，证明 $ \sqrt{g} = \sqrt{1 + |\nabla f|^{2}} $。

\end{items}

\item

求 $ T_{p}S^{2} $ 在 $ p = (0,1,0) $ 处，分别用

\begin{items}
\item

球坐标（见例 9.11(b)）；

\item

由球极投影给出的坐标。

\end{items}

\item

求 $ T_{p}\mathrm{T}_{2,1}^{2} $ 在 $ p = (\sqrt{2},\sqrt{2},1) $ 处。

\item

设 $ M $ 是 $ \mathbb{R}^{n} $ 的 $ m $ 维 $ C^{q} $ 子流形，且 $ (\varphi,U) $ 是 $ p \in M $ 附近的一个坐标卡。证明

\begin{items}
\item

$ M $ 的每个开子集都是 $ \mathbb{R}^{n} $ 的 $ m $ 维 $ C^{q} $ 子流形；

\item

若把 $ U $ 理解为流形，则 $ \varphi $ 属于 $ \operatorname{Diff}^{q}(U,\varphi(U)) $，且坐标卡的切向 $ T_{p}\varphi $ 与 $ \varphi \in C^{q}(U,\varphi(U)) $ 在点 $ p $ 处的切向一致。

\end{items}

\item

对 $ A \in GL(n) := \mathcal{L}\mathrm{aut}(\mathbb{R}^{n}) $，求 $ T_{A}GL(n) $。

\item

证明正交群 $ O(n) $ 在 $ 1_{n} $ 处的切空间是斜对称 $ (n \times n) $ 矩阵所成的向量空间，即

\[
T_{1_{n}}O(n) = (1_{n}, \{ A \in \mathbb{R}^{n\times n} \; ; \; A + A^{\top} = 0 \}) .
\]

（提示：利用例 9.5(c) 与定理 10.7。）

\item

证明特殊正交群 $ SO(n) $ 在 $ 1_{n} $ 处的切空间为

\[
T_{1_{n}}SO(n) = (1_{n}, \{ A \in \mathbb{R}^{n\times n} \; ; \; \operatorname{tr}(A) = 0 \}) .
\]

（提示：对满足 $ \operatorname{tr}(A) = 0 $ 的 $ A \in \mathbb{R}^{n\times n} $，注意到 $ \gamma(t) = e^{tA} $（$ t \in \mathbb{R} $），并见定理 10.6。）

\item

证明

\begin{items}
\item

对 $ \psi \in \operatorname{Diff}^{q}(M,N) $ 与 $ p \in M $，有 $ T_{p}\psi \in \mathcal{L}\mathrm{is}(T_{p}M,T_{\psi(p)}N) $；

\item

若 $ M $ 与 $ N $ 是 $ \mathbb{R}^{n} $ 的微分同胚的 $ C^{q} $ 子流形，则它们的维数相同。

\end{items}

\item

证明

\begin{items}
\item

$ S^{1} \times \mathbb{R} $（见习题 9.4）与 $ \{ (x,y,z) \in \mathbb{R}^{3} \; ; \; x^{2} + y^{2} = 1 \} $ 微分同胚；

\item

$ S^{1} \times S^{1} $ 与 $ \mathrm{T}_{2,1}^{2} $ 微分同胚；

\item

$ S^{n} $ 与 $ rS^{n} $（$ r > 0 $）微分同胚。

\end{items}

\item

设 $ M := \{ (x,y,z) \in \mathbb{R}^{3} \; ; \; x^{2} + y^{2} = 1 \} $，且

\[
\nu : M \to S^{2} , \quad (x,y,z) \mapsto (x,y,0) .
\]

则 $ \nu \in C^{\infty}(M,S^{2}) $。又 $ T_{p}\nu $ 对 $ p \in M $ 对称，且其特征值为 0 和 1。

\item

设 $ X $ 在 $ \mathbb{R}^{n} $ 中开且 $ f \in C^{1}(X,\mathbb{R}) $。又设 $ \nu : M \to S^{n} $ 是 $ M := \operatorname{graph}(f) $ 上的单位法向量。证明 $ \nu $ 属于 $ C^{\infty}(M,S^{n}) $，且 $ T_{p}\nu $ 对 $ p \in M $ 对称。

\item

设 $ M $ 与 $ \nu $ 同习题 12。又设

\[
\varphi_{\alpha} : M \to \mathbb{R}^{n+1} , \quad p \mapsto p + \alpha\nu(p)
\]

其中 $ \alpha \in \mathbb{R} $。证明存在 $ \alpha_{0} > 0 $，使得对每个 $ \alpha \in (-\alpha_{0},\alpha_{0}) $，$ \varphi_{\alpha}(M) $ 是与 $ M $ 微分同胚的光滑超曲面。

\item

求含于椭球面 $ \{ (x,y,z) \in \mathbb{R}^{3} \; ; \; (x/a)^{2} + (y/b)^{2} + (z/c)^{2} = 1 \} $（$ a,b,c > 0 $）的最大长方体的体积。

\end{exercises}
